% ============================================================
%  H. Brøchner, Problemet om Tro og Viden
%  En historisk-kritisk Afhandling
%  Kjøbenhavn: P. G. Philipsens Forlag, 1868.
%  TRANSCRIPTION (Danish)
%  Source: KB scan, ~/bibliotek/Brøchner, Hans/1868-problemet-tro-viden.pdf,
%  sha256 dfb663931ed9f23fcd7f97d598437ac7b614bf4487e8b0d44ac8b87ca622040c,
%  241 PDF pp.  Page map: PDF = printed + 9, uniform (202 folios in the text
%  layer, and every folio re-read at the image during collation).
%  [Until 2026-09-21 this line read "cross-checked where possible". That was
%  never true of the words: no page had been compared against the scan.]
%
% ---------------------------------------------------------------------
%  ACCURACY: FULL COLLATION  [2026-09-21]
% ---------------------------------------------------------------------
%  Method (REPAIR-PLAYBOOK.md 5C): all 226 printed pages read at the image
%  against this file, word by word, body and footnotes, by 23 agents of 10
%  pp. each in 4 waves.  Each agent also adjudicated an ocrdiff.py hint sheet
%  against the embedded ABBYY layer (812 candidates, 21 of them real, 791
%  the witness's fault, 0 unresolved).  Every printer's-error claim, every
%  doubtful reading, every sense-changing correction and a random sample of
%  the rest (217 questions in all) went to a BLIND second reader, asked only
%  what the page prints; disagreements went to a third reading.  2
%  collator claims and 1 printer's-error claim were refuted that way and
%  NOT applied (p.50 and p.55 "Conseqv-": the page reads Consequ-; p.199
%  "kages": the page reads tages).
%
%  Transcriber's errors corrected: 199 in words, spelling and punctuation
%  (45 of them in the Forord, pp.1-7, which had been silently modernised:
%  "skarp"/klar, "gode Grunde"/klare Grunde, "man maa therefore"/maa
%  derfor, Konseqvents/Conseqvents, andetsteds/andensteds, forvanker/
%  forvansker ...), incl. lost run-in heads "S. Kierkegaard. ---" (p.118) and
%  "R. Nielsen. ---" (p.126), dropped short words (e.g. pp.15, 16, 17, 19,
%  20, 22, 34, 43, 75, 78, 136, 175, 187), ſk misread as st (uadskillelig,
%  skee, fremskridende, skilles ...), and Greek that had been transliterated
%  or set in math mode (pp.4, 12, 13): the page prints Greek type
%  throughout, never transliteration.  Five ø restored where both readers
%  saw the slash (Kløften, Kløft p.214; Religiøsitet- pp.222, 223).
%  Emphasis: ~1,100 corrections to letterspacing (\emph{}) and antiqua
%  (\textit{}); pp.1-120 had almost none marked.  Proper names ARE
%  letterspaced in this print.  Paragraphing: 9 breaks the print does not
%  have removed, 2 missing ones restored (pp.182, 226); one rule added
%  (p.160).  Page markers: all 226 now \opage, read at the image (the
%  former \apage were machine-placed, most a word or a line off, and
%  pp.53, 57, 112, 137, 164, 178, 205 had none).
%
%  Printer's errors kept as printed, each with a "% PRINTED AS IS" comment
%  at the site (26): pp.5, 16, 18, 51, 52, 57, 65, 88, 90, 109 (x2),
%  110 (x2), 119, 136, 143, 156, 162, 171, 172, 175, 182, 189, 190, 198,
%  199.  All but p.65 (liberium) had been silently corrected before.
%
%  Post-repair re-sample (seed "|post", pp.48, 75, 103, 106, 149, 217):
%  0 word errors in 6 pages (bounds the residual word rate at <=0.5/pg,
%  95%; the collation itself read every page).  8 emphasis-extent
%  disagreements on pp.103 and 106, not applied: EMPHASIS EXTENT IS THE
%  SOFT CLASS OF THIS EDITION.  Short words at the edge of a spaced run are
%  often undecidable at the scan's 250 ppi, and readers split on them.
%
%  STILL OPEN (PRN / for a second eye; the file is unchanged at each):
%    - o/ø: not swept systematically (single readers find what they
%      seek).  E.g. "Kloft(en)" pp.198-199 left plain; p.80 "Oeconomie"
%      where one reader saw "Øconomie".
%    - fat-face: single observations at pp.29 ("og"), 45 ("i"), 221 ("i");
%      nothing applied.
%    - emphasis extents flagged doubtful: pp.79, 80, 102, 103, 106, 150,
%      185, 200, 206.
%    - p.11 paragraph break before "Til det samme Resultat" (no visible
%      indent; left); p.131 "Uklarhed" (faint l, left).
%    - Decorative rules not recorded: Forord end (p.7), double rule and
%      tail-piece pp.224, 226, rule under "Slutning." p.225.  Where the
%      print has a short centred rule and this file has an editorial
%      \section/\subsection, the heading stands for the rule (standing
%      convention; pp.49, 56, 66, 71, 126, 130, 171, 195, 199, 224).
%    - The book's Indhold (PDF 8-9) is not transcribed.
%  Headings are editorial.  The print's own heads are numbered run-in or
%  display heads ("1. Problemets Forberedelse i Oldtiden.", "a. Det Nye
%  Testamentes Udsagn ...", "2. Det Ethiskes og det Religiøses Forhold
%  ..."); many editorial \subsection titles have no printed head at all.
% ---------------------------------------------------------------------
% ============================================================
\documentclass[12pt,a4paper]{book}
\usepackage[utf8]{inputenc}
\usepackage[T1]{fontenc}
\usepackage{libertinus}
\usepackage{libertinust1math}
\usepackage{textalpha}   % Greek words (μυστήριον, γνῶσις, etc.) typed directly
\usepackage{graphicx}    % for \reflectbox (the old Danish »ɔ:« id-est mark)
\DeclareUnicodeCharacter{0254}{\reflectbox{c}}  % ɔ (U+0254) = reversed c
\usepackage[danish]{babel}
\usepackage[protrusion=true,expansion=true]{microtype}
\usepackage{setspace}
\usepackage[margin=1.2in]{geometry}
\usepackage{enumitem}
\usepackage{fancyhdr}
\setlength{\headheight}{28pt}
\addtolength{\topmargin}{-13.5pt}
\usepackage{hyperref}

\hypersetup{
  pdftitle={Problemet om Tro og Viden},
  pdfauthor={H. Brøchner},
  hidelinks
}

\pagestyle{fancy}
\fancyhf{}
\fancyhead[LE,RO]{\thepage}
\fancyhead[RE]{\textit{Problemet om Tro og Viden}}
\fancyhead[LO]{\textit{\leftmark}}

\setstretch{1.3}

\title{\textbf{Problemet om Tro og Viden}\\[6pt]
  \large En historisk-kritisk Afhandling}
\author{H.\ Brøchner\\[4pt]
  \small Professor i Philosophie}
\date{Kjøbenhavn: P.\ G.\ Philipsens Forlag, 1868}

% ---- original-page markers (paginate.py) ----
\usepackage{marginnote}
\newcommand{\opage}[1]{\marginnote{\footnotesize\textit{[#1]}}}
\newcommand{\apage}[1]{\marginnote{\footnotesize\textit{[#1]}}}
\begin{document}
\maketitle
\thispagestyle{empty}
\newpage

\tableofcontents
\newpage

%% ============================================================
%%  FORORD (pp. 1--7)
%% ============================================================
\chapter*{Forord}
\addcontentsline{toc}{chapter}{Forord}
\label{ch:forord}

\noindent \opage{1}Dette Skrift gjengiver, i alt Væsentligt uforandret,
Indholdet af de Forelæsninger, jeg i Efteraarshalvaaret 1867 har holdt
ved Kjøbenhavns Universitet over \emph{Forholdet mellem Tro og Viden, mellem Religion, Philosophie og Ethik}. Disse Forelæsninger omfattede
Undersøgelsens historisk-kritiske Deel og dannede Forberedelsen til et
Forsøg paa en selvstændig Løsning af det behandlede Problem, der i
Løbet af dette Aar vil blive offentliggjort.

Den \emph{historisk-kritiske} Deel af Undersøgelsen giver, foruden en kort
Skizze af Problemets Tilbliven, en udførlig Kritik af de
Løsningsforsøg, der hos os have faaet en særegen Betydning, og af
hvilke eet: det, der hidrører fra S.\ Kierkegaard, er gjort med
en Originalitet, en Alvor og Dybde, der til alle Tider vil bevare det
denne Betydning. Gjennem Kritiken af disse Løsningsforsøg, som under
deres Udvikling og Begrundelse have staaet i et vedvarende polemisk
eller mæglende Forhold til de tidligere og andensteds fremkomne, og
deels kunne betragtes som Resultater af tidligere udprægede
Hovedretninger, deels væsentligen have fuldbyrdet Kritiken af
saadanne, ere de betydeligste af de ældre og samtidige Løsningsforsøg
belyste og tildeels de fuldstændige Forudsætninger givne for en endelig
Dom om dem, og denne historiske Kritik vil saa finde sit Supplement i
det Skrift, der knytter sig til dette.

Problemets nærværende Stilling er for vor Literaturs Vedkommende
betegnet ved Kampen mellem en Form af \emph{Theologien}, der væsentligen
laaner de Vaaben, med hvilke den vil hævde Foreningen af Tro og Viden,
paa \opage{2}eklektisk Maade fra nyere philosophiske Blandingsformer, og det \emph{antitheologisk-religiøse} Standpunkt, der ved den principielle
Modsætning mellem Tro og Viden vil sikre Religionens og Videnskabens
Gyldighed i deres gjensidige Uafhængighed og Muligheden til deres
fredelige Samværen i Bevidstheden, medens den i Kraft af samme
principielle Modsætning frakjender Theologien al Ret til at bære
Videnskabens Navn, hvilket den tilegner sig idet den gjør Fordring paa
at være en „Troesvidenskab“.

Den \emph{philosopherende} Theologie, der fastholder Muligheden af
Troesindeholdets Optagen og Udprægen i Videns Form, og derfor ikke
blot vil hævde Theologien Videnskabens Navn, men tillægger den, som
Aabenbaringsvidenskabens afsluttende Form, Principatet blandt
Videnskaberne og Dommermyndigheden over dem, har til sit historiske
Udgangspunkt havt det speculative System, der gjennem en tvetydig
Udvidelse af „Religionens“ Begreb bragte Religionen i dens specifiske
Betydning i et irenisk Forhold til Philosophien, fra hvilken ikke en
Indholdsforskjel, men kun en Formdifferents skulde sondre den.
Efterat denne tvetydige Enhed har maattet opgives under de Kampe, der
udviklede sig efter den Hegelske Religionsphilosophies Fremkomst, har
den af Philosophien paavirkede Theologie, idet den ombyttede
Speculationens Betegnelse med det Christeliges, dog bevaret
Mæglingstendentsen, men den støtter sig nu nærmest paa en anden
philosophisk Retning, og særlig paa det System, der ved at bringe en
Fordobbling ind i Philosophien skaffer Plads for en „fri Tænkning“,
der fra sin nære Slægtning: Phantasien, laaner den „Fjederham“, i hvilken den kan overflyve Mysteriets Afgrunde; og den anbringer
forsigtigt et Forbehold, der afskjærer ogsaa den „høiere Videnskabs“
Conseqventser, naar Tanken, trods „Friheden“, uvilkaarlig kommer til
at følge sine væsensnødvendige Love.

Den \emph{antitheologiske} Theorie, der ligeledes har den Hegelske Philosophie
til Forudsætning, men hvis Forhold til Forudsætningen oprindelig er
negativt, polemisk, har \opage{3}til historisk Incitament ogsaa den Protest,
der fra en Philosophie, som paa een Gang optræder som Conseqvents af
og som Modsætning til hiin speculative Mæglingslære, nedlagdes mod
Enheden af Religion og Philosophie. Som den betydeligste Repræsentant
for denne Retning staaer \emph{Feuerbach} --- en Mand, hvis Skrifter kun
meget ufuldstændigt ere kjendte og kun meget overfladisk ere bedømte
af dem, der hos os have kritiseret ham. Idet han opfattede Religionens
Standpunkt som praktisk i Modsætning til Videnskabens theoretiske, og
idet han nærmere bestemte hint ved en bestandig skarpere Udhæven af
det alogiske Moment deri, fik han den uforenelige Modsætning til
Synspunkt for Forholdet og fordrede i den sande Erkjendelses og den
sande ethiske Handlings Interesse Religionens Absorption i Videnskab
og Ethik. Religionen som saadan skulde, som et Alogisk, som et
ufuldkomment Gjennemgangstrin i Udviklingen af Menneskeaandens
Selvbevidsthed, forsvinde; kun dens Navn blev tilbage i uegentlig og
overført Betydning.

Enig med Feuerbach i Modsætningen, men uenig i Begrundelsen og
Anvendelsen af denne, gjør saa S.\ \emph{Kierkegaard} Sondringen mellem
Religion og Philosophie i Religionens Interesse. Menneskeaanden, der i
Existentsens Uro ikke kan gribe det Evige i Erkjendelsens objective
Vished, og selv erkjender sin Erkjendelses Skranke, faaer kun i det
religiøse Forhold, i Troen, en subjectiv existentiel Vished om det,
der er dens dybeste Interesse, den finder først der sin sande
Forsoning. Videns Gyldighed benægtes ikke, men den objective Viden
begrændses til et Omraade, paa hvilket Aandens dybeste Interesser ikke
finde deres Tilfredsstillelse. Tro og Viden faae saaledes hver sit
Gebeet, og derved og ved Videns Subordination skal Foreneligheden af
dem i samme Bevidsthed blive mulig.

Denne Kierkegaardske Opfattelse af Problemet om Forholdet mellem Tro
og Viden er bleven optagen af Professor R.\ \emph{Nielsen} og fremsat paa en
saadan Maade, at man har villet stille ham i Originalitet lige med
Theoriens Ophavsmand og har tillagt ham Løsningen af en Opgave \opage{4}af
samme Oprindelighed som den, hiin stræbte at løse. Hvor lidet
berettiget denne Anskuelse er, vil i dette Skrift blive paaviist. Det
vil vise sig, at Problemet ved Prof.\ Nielsen ikke er blevet ført
videre, at der ved den Form, han har givet Løsningen, kun er bragt
Usikkerhed og Vaklen ind i de afgjørende Bestemmelser, at Forsøget paa
en psychologisk og metaphysisk Begrundelse af Troesprincipets Sondring
fra Vidensprincipet som absolut uensartet med dette er mislykket og
fører til uløselige Modsigelser, og at det Samme gjælder om hans
Begrundelse af den Sætning, at de absolut uensartede Principer skulle
kunne forenes i samme Bevidsthed.

Naar der da, ved Siden af S.\ Kierkegaard, den egentlige Repræsentant
for det antitheologisk-religiøse Standpunkt, bliver taget saa udførligt
Hensyn til Prof.\ Nielsen, saa er det ikke fordi Theorien hos ham har
faaet en dybere Begrundelse eller en rigere Udvikling, men fordi den
ved ham har ombyttet den esoteriske Læres Charakteer med den exoteriskes,
er trængt ind i „de Dannedes“ vide Kreds, og i denne fra en
Videnssag er ifærd med at blive til en Troessag og at fastslaaes som
Dogme i Kraft af et inappellabelt
% [collation 2026-09-21] Greek re-encoded from math mode; the page prints αὐτὸς ἔφα (Doric ἔφα), in Greek type.
αὐτὸς ἔφα.

Som Repræsentant for den philosopherende Theologie fremtræder i dette
Skrift Biskop \emph{Martensen}; men idet Udviklingen nøie knytter sig til det
i hans Skrifter Givne, er ingensteds det Almene, det Typiske i
Standpunktet tabt af Sigte, eller glemt over det Individuelle. Til andre
theologiske Forsøg hos os har det ikke været nødvendigt directe at
tage Hensyn.

Den endelige Opgave, hvis Løsning forberedes gjennem Kritiken, der
umiddelbart kun kan give et negativt Resultat, er Hævdelsen af et
Standpunkt, der i Forhandlingerne hos os ikke er kommet til Orde, og
som kun i Former, der ikke kunne tilfredsstille, er gjort gjældende
andensteds. Dette Standpunkt fastholder en saadan Modsætning mellem
Videnskaben, særlig Philosophien, og de positive Religionsformer,
mellem Viden og Tro i Betydning af Aabenbaringstro, at \opage{5}Viden og denne
Tro ikke kunne bringes til Enhed i samme Bevidsthed; --- det fastholder
Umuligheden af, at en sand og indholdsfyldig, det hele, virkelige
Menneskeliv omfattende Ethik kan udfolde sig paa Basis af den positive
Religion med dens supranaturale Handlingslov; --- og det fastholder den
indre, uopløselige Enhed mellem den sande Erkjendelse og den sande
ethiske Handlen, dennes nødvendige Betingethed ved hiin. Men medens
det fastholder hiin Modsætning, hævder det paa den anden Side ligesaa
afgjørende, idet det sondrer det Religiøses \emph{Væsens}bestemmelse fra de
positive Religioners phænomenale Former, i hvilke det kun seer et
partielt og ufuldkomment, endeliggjørende Udtryk for hiin, idet det
altsaa sondrer det i Religionerne Religiøse fra Religionerne i deres
historiske Former, den indre Enhed af det Religiøse, det Philosophiske
og det Ethiske, og gjennem denne indre Enhed Menneskelivets indre
Enhed, Betingelsen for dets sande og fulde Forsoning, Forsoningen
indenfor Virkeligheden, den ene og udelte Menneskenaturs Forsoning med
sig i sin Væsensgrund. Idet det hævder Væsentligheden og
Nødvendigheden af det religiøse Forhold for Mennesket, hvis
Selvbevidsthed først bliver fuldstændig virkeliggjort i dette, gjennem
Gudsbevidstheden, og som ifølge sit Væsens Grundforhold aldrig har
været og aldrig vil være uden Religion; --- idet det tillige seer den
ethiske Handlen og den sande Erkjendelse som vindende en høiere
Betydning gjennem Enheden med det religiøse Grundforhold; --- og idet
det lader den sande Viden og den sande ethiske Handlen blive en
adæquat Virkelighedsform for det Religiøse, opfatter det denne Viden
og denne Handlen som i sig indesluttende et Høieste og Ubetinget, i
hvilket det i Religionerne Væsentlige er bevaret. Dette Standpunkt
% PRINTED AS IS, p.5 [collation 2026-09-21]: the page prints "Bevidstheds." (genitive -s, letterspaced) where sense wants "Bevidsthed"; confirmed by a second blind reader.
betegne vi som \emph{den religiøs forsonede humane Bevidstheds}.

Det adskiller sig fra det af \emph{Feuerbach} indtagne ikke blot derved, at
dets metaphysiske Grundbegreb er et andet: et saadant, paa hvilket
en Erkjendelse og en Handlen med Idealitetens og det Almeengyldiges
Præg kan bygges, \opage{6}og ud fra hvilket Begrebet om den menneskelige Frihed
kan hævdes og et religiøst Forhold begrundes til det Absolute, med
hvilket Mennesket er i Væsensenhed; men ogsaa derved, at det
kæmper mod Theologien og Dogmet ikke blot i Ethikens og Videnskabens,
men først og fremmest i selve Religionens, i det religiøse Forholds
Uendeligheds Interesse.

Det adskiller sig fra den \emph{gamle Hegelske Skoles} Opfattelse af
Religionens Forhold til Erkjendelsen, idet det ikke seer det Religiøse
blot som en Videns-Form, men som Udtrykket for et det \emph{hele} menneskelige
Væsen i dets Concretion omfattende Grundforhold, og idet det ikke
lader Aandslivets Medium være et illusorisk abstract Evighedselement,
men Virkeligheden med dens reale Betingelser og Skranker.

Med den \emph{antitheologisk-religiøse} Theorie mødes det i Opfattelsen af
Theologiens Forhold til Videnskaben. Idet det hviler paa en
videnskabelig Erkjenden og i det udelte Menneskevæsen giver Erkjendelsen
Principatet, maa det være dets Opgave at hævde Videnskabens klare
Begreb og dens i dens eget Væsen begrundede Autonomie mod
Tvetydigheder og paa Tvetydigheder støttede Forsøg paa at underlægge
Videnskaben under en fremmed, uensartet Autoritets Herredømme. Men
medens der saaledes i et væsentligt Punkt er en Overensstemmelse
mellem det og hiin Theorie, træder det i en afgjort Opposition til
denne, naar Forholdet mellem Tro og Viden, mellem Religion, Philosophie
og Ethik skal bestemmes. Denne Opposition gjælder \emph{Kierkegaards}
Opfattelse, med Hensyn til hvilken det vil blive viist, at den rokker
al Erkjendelses Vished og at den lader det Ethiske forsvinde i et
abstract og virkelighedsløst Religiøst. Men den gjælder endnu
bestemtere den Form, i hvilken Theorien er bleven udført og fremstillet
af \emph{Nielsen} --- en Form, i hvilken Kierkegaard sikkerlig ikke vilde
vedkjende sig den. Den Form, i hvilken Nielsen udpræger Theorien,
fører, trods alle de mere eller mindre behændige Vendinger, der gjøres,
og ved hvilke denne Consequents søges undgaaet, til en factisk Dualisme
i Menneske\opage{7}væsenet, til en Deling af Aandens Liv fra Grunden af ved de
„absolut uensartede“ Principer, til Bevidsthedens uundgaaelige
Selvopløsning; den forvansker det Ethiske i dets Grund og aabner
Adgangen til det Religiøse for alt det Superstitiøse og Absurde. Med
den Overbevisning, at den af Nielsen statuerede \emph{absolute} Uensartethed
mellem Videns Princip paa den ene Side, Troens og Handlingens paa den
anden er „en for Religiøsiteten, Sædeligheden og al sand
Charakteerudvikling fordærvelig Misforstaaelse“, har jeg underkastet
hans Theorie en gjennemgaaende Kritik, og søgt at godtgjøre dens
Uholdbarhed, dens mangelfulde Begrundelse, dens Selvmodsigelser, dens
fordærvelige Consequentser.

Idet jeg offentliggjør dette Bidrag til Undersøgelsen af et Problem,
der gjennem Aarhundreder har havt og bestandig vil have den meest
afgjørende Betydning for Menneskeaanden, er det mit Haab, at det maa
give Anledning til en fornyet, alvorlig og samvittighedsfuld Drøftelse,
til en klar og aaben Kamp mellem Principer. En saadan Kamp, selv om
den føres med skarpe Vaaben, kan kun være i Sagens Interesse, og maa derfor være kjær for den Kæmpende, for hvem Sagen, ikke et
egoistisk Selvhensyn, er det Høieste. En fast og dyb personlig
Overbevisning skal jeg altid agte, ved klare Grunde skal jeg altid være
ligesaa paavirkelig, som jeg vil være upaavirkelig ved Andet end
disse.


%% ============================================================
%%  HISTORISKE MOMENTER (pp. 8--31)
%% ============================================================
\chapter{Historiske Momenter til Belysning af Problemets Fremkomst} \opage{8}
\label{ch:historiske}

\section{Problemets Forberedelse i Oldtiden}
\label{sec:oldtiden}

Først hos \emph{Grækerne} er det, at en bestemt Tilnærmelse til en
Tydeliggjørelse af Problemet om Troens og Videns Forhold viser sig. I \emph{Orienten}, f.~Ex.\ i Indien, bliver man sig vel, hvor en af Religionen
relativ uafhængig Tænkning udvikler sig, Uovereensstemmelsen mellem
dennes Resultater og Religionens Lærdomme bevidst, og der danner sig
en Modsætning mellem en heterodox og en orthodox Lærebygning; men den
orientalske Philosopherens Eiendommelighed, dens Mangel paa
gjennemført Uafhængighed af Religionslæren, paa Evne til at udvikle
sit Indhold i Begrebets Form, og paa klar Bevidsthed om Philosophiens
Princip, bevirker, at der ikke naaes til en Reflexion over
Differentsens \emph{principielle} Grund, altsaa ikke til en Bevidsthed om
Troens og Videns \emph{Væsensbestemmelse}. Den bevidste Differents er kun
\emph{Indholdets} factiske Differents.

I \emph{Grækenland}, hvor den vidende Aand bliver sig bevidst i Philosophien,
og, idet den udfolder sin Videns Indhold, sammenknytter det ved eet
Princip, kan den med dette Videns Indhold, hvis Princip og hvis
Sammenhæng med Principet den er sig bevidst, træde i Forhold til
Religionen saaledes som den er given i Folkereligionernes \opage{9}Skikkelse, og
Problemet antydes, ja det synes endog, ved den klassiske Oldtids
Slutning, at komme frem i sin afgjørende Form. Det vil dog vise sig,
at der, selv hvor Spørgsmaalet om Forholdet mellem Philosophiens og
Religionens Lærdomme og om Religionens Betydning for Aandens
Erkjenden stærkest accentueres, ikke naaes ud over Antydninger, og at
det, der forhindrer dette, er den græske Aands specifiske Eiendommelighed.

Allerede i den græske Philosophies ældste Tid see vi Philosophien staae
i et kritisk-polemisk Forhold til Folkereligionen og dens Myther.
Saaledes bekæmper den eleatiske Enhedslæres Begrunder \emph{Xenophanes}
Mythernes anthropomorphistiske Forestillinger om Guderne, og hans
Polemik gjentager sig hos andre af hiin fjerne Tids Tænkere, men uden
at vi hos Nogen af dem finde en egentlig Reflexion over
Folkereligionens Betydning og Væsen. I den græske Philosophies
Culminationstid, hos \emph{Plato} og \emph{Aristoteles}, træder det kritiske Forhold
til det Mythiske frem i Forbindelse med en klarere Reflexion over
Philosophiens Stilling til Folkereligionens Myther og Cultus. Her
bliver navnlig \emph{Aristoteles's} Opfattelse af Betydning. Det
Eiendommelige for ham, som i det Hele for den græske Philosophie
indtil hans Tid, er det, at den philosophiske Tanke, hvilende og
tilfredsstillet i sig selv, ud fra sig selv fælder Dommen over det
positivt Religiøse uden at nære Tvivl om Tankens Myndighed. I
Religionens Lærdomme anerkjender Aristoteles kun det som sandt, der
stadfæstes ved Philosophien, den mythiske Indklædning derimod
frakjender han uden Betænkning Sandhed, idet han fører den tilbage til
bevidst Digtning eller ubevidst Anthropomorphiseren, og Cultusformerne
faae kun Betydning for ham som det i Staten Lovgyldige.

Anderledes viser Forholdet sig i den \emph{efteraristoteliske} Tid og i denne
gjennemløber det to væsentlig forskjellige Stadier. Idet den græske
Philosophie overskrider sit Culminationspunkt og bliver sig den indre
Splid bevidst, \opage{10}som den ikke kan forsone, saa viser der sig \emph{først} en
større Hensyntagen til det Religiøse, idet Forholdet til dette, det
anerkjendende og det polemiske, faaer en væsentligere Betydning for den
ikke mere trygt i sig selv hvilende Tanke. Vi see saaledes hos \emph{Stoikerne} et Forsøg paa at rationalisere Mythen, og vi see dette
Forsøg medføre det Mythiskes Tilbageslag paa den philosophiske
Theorie, et Tilbageslag, der bliver tydeligt i Stoikernes Theodicee og
i deres Lære om Mantiken. Paa den modsatte Side see vi hos \emph{Epikuræerne}
en rationel Modsætning til Mythen, en Kamp for at befrie Aanden fra
Gudefrygtens forstyrrende Indvirkning, i hvilken \emph{Epikur} seer
Religionens Særkjende; og hos de \emph{nyere Akademikere}, f.~Ex.\ \emph{Karneades},
ligesom hos de sildigere \emph{Skeptikere} og hos Fornyerne af \emph{Kynismen},
finde vi en Opposition, deels mod de theologiske Forsøg, deels mod
Folkereligionen.

Et \emph{nyt} Stadium i Forholdets Udvikling, og et saadant, der i denne
Sammenhæng har væsentlig Interesse, begynder henimod den Tid, da
Christendommen fremtræder, historisk forberedt ved den græske og
orientalske Aands Udvikling, i \emph{Nypythagoræismen} og i \emph{Slutningsformerne af de ældre philosophiske Skoler}. Jo mere Tanken udtømtes, jo mere
Aanden theoretisk og praktisk tabte sit Hvilepunkt og ved Skepsis
bragtes til at mistvivle om sin egen og om Virkelighedens Sandhed,
desto mere søgtes Sandheden, Erkjendelsens og Handlingens Sandhed,
udenfor Menneskeaanden i det transscendente Guddommelige, der gjennem
Mellemvæsener aabenbarer sig for Aanden, og idet saadanne
guddommelige Aabenbaringer saaes i Religionerne, søgtes i dem et
Videns Udgangspunkt. Og Forholdet til Aabenbaringen --- til den
gjennem Mellemvæsener stedfindende eller i Religionerne allerede givne
--- sattes saa i \emph{Nyplatonismen} som betinget ved religiøse Luttringer
og ved religiøse Skikke, overhovedet ved et religiøst Liv og et
religiøst \emph{Sindelag}. Philosophien, der støtter sig paa dette
Aabenbaringens Indhold og af det vil øse Sandhed, \opage{11}forvandles saa til
Theologie, og i Nyplatonismen, der lader Philosophien faae den Opgave
at retfærdiggjøre Mythens Indhold, og, idet den assimilerer sig den
idealiserede Mythe, at indtage Religionernes Plads som universel
Religion, see vi indenfor Hedenskabet Parallelen til den
christelig-middelalderlige Scholastik.

Paa dette Punkt er der naaet saa nær hen til i den afgjørende Form at
stille Problemet om Troens og Videns Forhold, som der overhovedet kan
naaes hos Grækerne. Her er reflekteret over den til sig selv overladte
Aands Mangel paa Evne til at gribe det i praktisk og theoretisk
Henseende i sidste Instants Afgjørende; Philosophien, den autonome
Tankes Udfoldelse, er stillet sammen med den guddommelige Aabenbaring,
og det er udtalt som Resultat af denne Sammenstillen, at Forholdet til
det Absolute, den absolute Sandhed, ikke kan naaes ved Viden, men kun
ved det religiøse Liv og det religiøse Sindelag. Men dog kommer
Problemet i sin Skarphed selv her ikke frem; der naaes endnu kun til
Antydningen, og den afgjørende Grund er den, at den græske Aand
udelukkende opfatter sig under \emph{Videns} Bestemmelse, at den lader Viden
constituere Menneskets \emph{Væsen}, hvorved oprindelig Muligheden afskjæres
af et i sin Eiendommelighed Viden sideordnet Princip og endnu mere af
en gyldig Modsætning til Viden, men kun gjennem Sondringen og
Modsætningen kan Problemet bringes frem og den forsonende Enhed
erkjendes.

At den \emph{græske} Aand kun bliver sig bevidst i \emph{Videns Bestemthed} kan
gjøres indlysende fra et dobbelt Synspunkt. Fæste vi først Blikket paa
den \emph{græske Ethik}, saa er det strax iøinefaldende, at i denne \emph{Dydens} Begreb bliver det centrale Begreb, medens Pligtbegrebet træder
tilbage. Men undersøge vi saa, hvorledes Dyden bestemmes, saa finde vi
den bestemt \emph{som Viden}, og da Dyden for Grækerne udtrykker Væsenets
Virksomhed efter dets uforstyrrede Natur, saa bliver det, at \emph{Dyd} er Viden, enstydigt med, at Viden er Menneskets \emph{Væsensbestemmelse}.

Til det samme Resultat føres vi, naar vi betragte den \opage{12}Maade, paa
hvilken de største græske Tænkere bestemme \emph{det Guddommelige}. For \emph{Plato}
er det Høieste Ideen, nærmere det Godes Idee, men Ideen er det
objektiverede og hypostaserede Tænkte, der af Plato faaer Bestemmelsen
Νοῦς og φρόνησις.
Hos \emph{Aristoteles} bestemmes det Guddommelige som Νοῦς og nærmere som νόησις τῆς νοήσεως,
der i evig, uafbrudt Energie forholder sig tænkende
til sig selv og fra hvilken al Bestemmelse af Villie er udelukket. Men
det Guddommeliges Væsensbestemmelse er Menneskeaandens Ideal, og den
bliver ogsaa for disse Tænkere det høieste Udtryk for Menneskeaandens
Væsensbestemmelse. For \emph{Plato} er Menneskesjælen nærbeslægtet med Ideen,
uopløseligt sammenknyttet med den, og for \emph{Aristoteles} er
Menneskeaanden i sin høieste Form som
νοῦς ποιητικός
i Enhed med selve den guddommelige Aand, hvis
Væsen udtømmes i Tanke.

Hvorledes den græske Bevidsthed, trods de Forsøg, der gjøres derpaa,
ikke definitivt naaer ud over Videns Bestemmelse, viser sig, naar vi
nærmere betragte Forholdet mellem \emph{Viden} og \emph{Villie} og mellem \emph{Viden} og \emph{Tro}.

Hvad det første Forhold angaaer, saa bliver, selv hvor Villiessiden i
den menneskelige Aand stærkere accentueres --- hvilket allerede skeer
hos Aristoteles, og endnu mere i den græske Philosophies sidste
Periode --- Villiens eiendommelige Væsen dog aldrig ret erkjendt, og
Forsøgene paa at paavise denne Ejendommeligheds Kilde mislykkes, saa
at Villien dog bestandig igjen absorberes af Viden. Ja selv hvor det,
ved den græske Philosophies Slutning, hos Nyplatonikerne, udtales, at
den sande Erkjendelse er betinget ved det sande Sindelag, hvor altsaa
formelt et Villiens Principat i Forhold til Tanken er udtalt, selv der
bliver Bevægelsen dog kun en Kredsbevægelse fra Viden tilbage til
Viden, og Viden viser sig i sidste Instants som det ene Oprindelige.
Det rette Sindelag nemlig: det Sindelag, der vender sig bort fra det
Endelige og som saadant Usande og Forsvindende, bliver \emph{umiddelbart} bestemt ved en Viden af dette det Endeliges Grundbestemmelse, saaledes,
at altsaa Erkjendelsen, \opage{13}idet den bestemmes ved dette Sindelag, kun
forholder sig til sig selv gjennem Villien, og i sin Bestemthed som
concret Erkjenden til sin egentlige Betingelse og Forudsætning har den abstracte \emph{Erkjendelse} af Grundforholdet. Villiens Forhold til
Erkjendelsen bliver ogsaa tydelig gjennem den Abstraction, i hvilken
Villien forbliver: Villien som Villie faaer intet eiendommeligt
positivt Indhold; den bevæger sig kun i det Negatives Almindelighed,
dens høieste Maal bliver Aandens ekstatiske Hensynken i det
bestemmelsesløse Uendelige.

Hvad der gjælder om Erkjendelsens Forhold til Villien, gjælder ogsaa
om dens Forhold til \emph{Troen}. Naar vi tale om Tro: πίστις, maae vi gjøre en
bestemt Forskjel mellem Begrebets Anvendelse i den tidligere Tid og i
en sildigere. Hos \emph{Plato} stilles πίστις sammen med εἰκασία og staaer som et
Lavere i Modsætning til Videns sande Form, hos \emph{Aristoteles} faaer
πίστεις (i Rhetoriken) Betydningen af
Overtalelsesmidler. I Betydning af et \emph{Høiere} i Forhold til Viden
bruges Begrebet Tro først i den græske Philosophies Afslutningstid.
Hos \emph{Philo}, der idetmindste ligesaameget er jødisk Troende som græsk
Tænker, kommer Begrebet frem i Betydning af den tillidsfulde Optagelse
og Anerkjendelse af Aabenbaringens Indhold, og efter ham bevares
Begrebet som Begrebet om det, der er høiere end Viden. For \emph{Jamblichus} bliver den første Betingelse for den sande Gudserkjendelse den Tro, at
Intet er umuligt for Gud. For \emph{Proklos} betegner Troen det \emph{umiddelbare} Forhold til den guddommelige Enhed, i hvilken Foreningen med dette
gjennem Anskuelsen af dets rene Væren er given. Den danner her
Modsætningen til den altid ved Reflexionens Flerhed bestemte Viden, der
ikke skal kunne naae til det Høieste\footnote{\emph{Proklos}: „Erkjendelsen viser os
den oversandselige Verden, men den høieste Indvielse giver kun Troen; thi det
er den, der overhovedet sammenknytter det Lavere med det Høiere. Ikke gjennem
Tænkning og Reflexion trænge vi ind i de dybeste Mysterier, men gjennem
Gemyttets Stilhed, i hvilken vi gaae ind i Grundvæsenets Enhed, gjennem den
Hensætten af hele Sjælen i det Ikke-Erkjendelige, hvilken vi kun skylde Troen.“}. Og dog gaaer, trods disse
\opage{14}Bestemmelser, trods det, at Troen, som det Høiere, sondres fra Viden,
Troen ind under Viden, Bevægelsen bliver ogsaa her kun en
Kredsbevægelse fra Viden til Viden. Hos \emph{Nyplatonikerne} er den
bestemmelsesløse, uendelige Enhed, til hvilken Troen forholder sig,
netop ikke Andet end den af en om sig selv mistvivlende Viden gjennem
den yderste Abstraction sat Negation af den bestemte Viden, en
Negation, der kun er ved Tanken, ved Viden og som, i sin Negativitet,
bestemt ved den. Og hos \emph{Philo} bliver Hengivelsesforholdet til det sig
aabenbarende Guddommelige i sin høieste og sande Form den i Ekstasen
stedfindende hengivende Hensynken i den abstracte Enhed; men denne
bestemmelsesløse Enhed er den abstraherende Tænknings sidste Residuum,
en negativ Videns Gjenstand.

Den græske Bevidsthed kan altsaa naae til at see Differentsen mellem
det i Philosophiens System udfoldede Videns Indhold og Religionens i
Mythe og Cultus udfoldede Bestemmelser, og til at blive sig
Videns-Principet bevidst; men den kan, i sin Videns Monisme, ikke
erkjende Religionens Princip i dets Eiendommelighed, gjennem hvilken
det kan vise sig som et fra Philosophiens forskjelligt eller dette
modsat, og dog gyldigt, ligesom den ikke heller, ifølge sin
mangelfulde Opfattelse af Begrebets Bevægelse, af Aands-Udviklingens
Væsen, medens den gjør Viden til eneste Synspunkt, kan see Religionens
Princip som Princip for et underordnet Stadium i Aandens teleologiske
Udvikling henimod den philosophiske Erkjendelse som Aandens sande og
adæquate Form.

\section{Problemets Fremkomst i Christendommen}
\label{sec:christendommen}

\noindent Medens der indenfor den græske Philosophie --- og kun indenfor
Philosophien kunde Problemet antydes hos \opage{15}Grækerne --- kun naaes til en
Forberedelse af det, men den skarpt udprægede Form for det ikke kan findes
selv hvor alle Betingelser synes at være tilstede og hvor det Afgjørende
ligesom berøres, saa tilveiebringer den christelige Tid Betingelserne for dets
Fremkomst. Dog udvikle de sig ogsaa her kun successivt, og vi kunne i
Udviklingen adskille flere Stadier.

\subsection{Det Nye Testamentes Udsagn om Forholdet}
\label{ssec:nt}

\noindent Hvad der, idet den nye Lære forkyndes, strax maae træde frem, er
Modsætningen mellem denne og det, som den skal afløse. Denne Modsætning er en
Modsætning mellem \emph{Indholdet} af det Evangelium om Gudsrigets Komme, der
forkyndes og troende tilegnes af de Christne, og af den ikke-christne Verdens
Overbeviisning, der har sit Fællesgebet i en Tænkning, i hvilken den græske og
den østerlandske Aand mødtes. Idet hiin Overbeviisning i dets Modsætning til
Evangeliet nærmest repræsenteres af de Dannede, de Vidende, bliver
Modsætningen en Modsætning mellem Verdens Viisdom og Christi for de Enfoldige
aabenbarede Viisdom (Guds Viisdom); men ogsaa denne Modsætningens Form viser
væsentligen, idet „Viisdom“ tages i objectiv Betydning, tilbage til Modsætningen i \emph{Indhold}. Et nyt Indhold, fordum skjult for Menneskeheden, er
aabenbaret, en Guds Raadslutning, for hvilken Tidens Fylde er kommen: et
μυστήριον. Mysteriets Begreb viser hen til det nye Indhold, der omfatter ogsaa
det i Fremtiden Forborgne (saaledes Jødefolkets Fremtidsskjæbne, Opstandelsen
og Forvandlingen); men det viser tillige hen til det, der kun bliver aabenbart
for dem, i hvem Aanden virker. Idet Begrebet omfatter denne sidste Betydning,
\emph{antydes} derved Modsætningens Grund, dens \emph{Princip}, men idet der ikke kan være
Tale om, at den troende, for Viden fremmede, Bevidsthed kan søge Grunden i et
Videns Princip, saa søges Grunden til Modsætningen mellem den ikke-christelige
Bevidsthed og den christelige i en \emph{Sindelagets}, en \emph{Villiens} Mangel, der
forsaavidt hiin \opage{16}Bevidsthed har en Kjendskab til dennes Troesindhold, bliver en
Mangel paa Villie til at antage Aabenbaringens Indhold, og denne Mangel paa
Villie sees atter som begrundet i det „Kjødelige“ og det „Psychiske“, der
danner Modsætningen til det „Pneumatiske“. Bagved denne Bestemmelse af
Modsætningen ligger den principielle Modsætning, der i sin Sammenhæng med
Problemet først kommer klart og skarpt frem i Scholasticismen: Modsætningen
mellem det Naturlige og det spiritualistisk Overnaturlige. Det Naturliges
Begreb omfatter den af Synden formørkede og fordunklede Tanke, og for denne er
derfor den høiere, den overnaturlige Sandhed utilgængelig, der kun gjennem en
guddommelig Aabenbaring bliver tilgængelig for Troen.

% PRINTED AS IS, p.16 [collation 2026-09-21]: the page prints "Gjenstaud" (turned n); both readers saw it at 600 dpi.
Den christelige Bevidstheds Forhold til sin Gjenstaud betegnes deels som
πίστις og deels som γνῶσις, og γνῶσις opfattes som Troens Fuldendelse. Til
Bestemmelsen af γνῶσις har den Anskuelse støttet sig, der vil sætte Tro og
Viden i Enhed og fastholder Muligheden af en Troesvidenskab. Men Betegnelsen
γνῶσις har i det N.\ T.\ endnu kun en ubestemtere Betydning; den betyder her en
anerkjendende Kjenden, ikke en Erkjenden gjennem Tanken. Γνῶσις stiller sig fra
πίστις som det Concretere fra det Abstractere, som det Udfoldede fra det endnu
kun Begyndende, men den falder indenfor samme Sphære, bliver ikke en
Erkjendelsens, en Videns Form, thi dens Indhold er det, der overstiger Viden.
Men i Betegnelsen γνῶσις ligger en Mulighed for hiin Opfattelse, og denne
bliver greben i Gnosticismen og i den ældre Kirke, i hvilken Aanden, i sin
Trang til Erkjenden, vil udvide Erkjendelsens Gebet til det, der efter sit
Væsen ikke kan erkjendes, og derfor lægger Betydningen af en begribende
Erkjenden ind i den forjættede γνῶσις.

\subsection{Den patristiske Tids Forhold til Problemet}
\label{ssec:patristisk}

\noindent Ved den patristiske Tid forstaae vi den christelige Oldtid, der var
Dogmefrembringelsens Tid, og vi skille den som saadan fra Scholasticismens, den
christelige \opage{17}Middelalders Tid, Dogmebearbeidelsens, den formelle Dogmetilegnelses
Periode.

Patristikens Tid udfolder Dogmet i væsentlig Fuldstændighed til en
sammenhængende, christelig Lærebygning, og den gjør dette idet den benytter den
antike Philosophies Former, fra Begyndelsen af især knyttende sig til den philoniske og stoiske Philosophie, senere fortrinsviis til den restaurerede
Platonisme. Ved denne Udfolden af en christelig Lærebygning, der som Totalitet
kan stilles imod de hedenske Systemer, er en væsentlig Betingelse given for
Problemets skarpere Opfattelse. Hos Kirkefædrene sondres der mellem den troende
Erkjendelse, en christelig γνῶσις, hvis Forhold til πίστις nu udførligere
bestemmes, og Philosophiens Erkjendelse, den hedenske γνῶσις; men Reflexionen
over Forholdet fører ikke til et definitivt Resultat, idet det ikke bliver klart, om
der ved Overgangen fra πίστις til γνῶσις kun skal tænkes paa en ordnet
Forbindelse af Troesindholdet til en sammenhængende Enhed, hvis Udgangspunkt
forbliver uomsat i Erkjendelse, eller paa en ἐπιστήμη i strængere Forstand\footnote{Man lægge Mærke til den Vaklen i Anvendelsen af Begrebet, der
viser sig hos \emph{Clemens Alexandrinus}. Hans Opfattelse af πίστις viser tilbage til
Philos og har Lighed med Nyplatonikernes.}. Modsætningen bliver derfor ogsaa i
Patristikens Tid trods den mere udviklede Tænkning endnu væsentlig en Modsætning
med Hensyn til Indholdet, og selv denne Modsætning kan, fordi Bevidstheden om det
Principielle er uklar, ikke faae sin fulde Skarphed.

Dette er begrundet i den patristiske Tids Forhold til det Antike. De Mænd, der i den
østerlandske og vestlandske Kirke repræsenterede den christelige Gnosis, vare ved
Sprogets og Nationalitetens Baand, ved hele den nedarvede Cultur, saa nøie
sammenvoxede med den classiske Oldtid, at den ubevidste Paavirkning af dennes
Tænkning og Livsanskuelse ikke kunde afværges. Medens man derfor ud fra den
umiddelbare religiøse Bevidsthed satte Modsætningen mellem den christelige Troes
og den hedenske Videns \opage{18}Indhold, forhindrede Eftervirkningen af den græske
Opfattelse af Viden Opnaaelsen af den fulde Klarhed med Hensyn til Modsætningens
Princip, og under Udfoldelsen af Dogmeindholdet trængte, skjult og
uimodstaaeligt, en Tilsætning ind fra den antike Philosophie og den antike
Verdens- og Livs-Anskuelse, og den trængte ind selv i det for Christendommen
Centrale, som i Læren om Gudmennesket og om Synden, og Grændsebestemmelsen mellem
Modsætningerne blev usikker for Reflexionen.

\subsection{Problemets Udvikling i Scholasticismen}
\label{ssec:scholasticismen}

\noindent Scholasticismen har Dogmets Udfoldelse til en væsentlig afsluttet
Heelhed til Forudsætning; dens Opgave er, ved en reflecterende Tænknings Hjælp at
tilegne Bevidstheden Dogmet ved at finde, om end ei Dogmets \textit{ratio}, saa dog Dogmernes \textit{rationes}. Naar vi derfor skulle betegne dens Tid efter dens Stræben,
betegne vi den som \emph{Dogmetilegnelsens} Tid. Ligesom den ved denne Bestemmelse
adskiller sig fra Patristikens Tid, saaledes adskiller den sig fra denne ved de
forandrede Tidsbetingelser. Den Afhængighed af Oldtiden, der var begrundet i den
umiddelbare Nationalitetssammenhæng, var ophørt, og Dogmefrembringelsen var
saaledes afsluttet, at et System, omfattende \emph{Troens og Videns Totalitet}, nu kunde
forsøges. I Kirkens Tjeneste og bøiende sig under dens Autoritet stræber Tanken at
opføre en saadan christelig Lærebygning, i hvilken Troesindholdet og
Tankeindholdet skal sammensmeltes til en Heelhed. Sammenføiningens Art bliver
tydelig naar vi betragte de store Systemer, der frembragtes i Scholasticismens
Blomstringsperiode. I disse see vi den antike Verdensanskuelse, der samlede sig i
% PRINTED AS IS, p.18 [collation 2026-09-21]: the page sets final-sigma ς inside the word, "Κόςμος"; both readers saw it.
Forestillingen om den harmonisk ordnede guddommelige Κόςμος, ligesom omfattet af
det Christeliges Ramme. Anskuelsen af Verdensvirkeligheden, der maatte tages fra
den som historisk Forudsætning givne Philosophie, for hvilken Verden var en
Virkelighed, og som for Oldtiden var det Hele, blev nu, forbunden med det
christelige Troesindhold, til en Deel, indføiet i en større \opage{19}Heelhed. Det fra
Oldtiden optagne Tankeindhold udfyldte kun den midterste Deel af det
Verdensdrama, af hvilket første og sidste Act tilhørte Christendommen. Forud for
det naturlige Universum sattes den naturfrie, almægtige, skabende Gud, efter det
den christelige Himmel, Naadens overnaturlige Rige. Men i denne Sammensætning
blev, efter Scholasticismens egen Indrømmelse, Begyndelsen og Slutningen det for Tanken Utilgængelige, altsaa i Væsen uensartet med det, der udfyldte
Mellemrummet; det specifisk christelige Indhold skiller sig altsaa, trods den
forsøgte Sammenføining, som et Ubegribeligt fra Tankeindholdet. Og idet man nu søger \emph{Grunden} til Ubegribeligheden, saa finder man den i Incongruentsen imellem \emph{naturlig} Tænkning og \emph{overnaturlig} Aabenbaring. Troens „høiere Erkjendelse“, der
skylder en ved en guddommelig Undervirkning tilveiebragt Potentsation af de
naturlige Evner sin Tilværelse, har kun Navnet tilfælleds med Videns Erkjendelse,
der hviler paa disse naturlige Evner, som ikke staae i et continuerligt
Udviklingsforhold til det ved den miraculøse Potentsation Satte. Den naturlige
Erkjendelses Betydning for Troesindholdet bliver kun den: i negativ Form at hævde
det, som det Ikke-Umulige, ikke den: positivt at begrunde det. Men trods den
Væsensdifferents, der kun lader Navnets Lighed tilbage, bestemmes dog Forholdet
mellem Tro og Viden som et Forhold mellem to \emph{Erkjendelsesformer}. Alligevel ligger
der i Modsætningernes Uensartethed en Tilskyndelse til at føre dem tilbage til
forskjellige Aandsfunctioner, og en saadan Tilbageføren finde vi f.\ Ex.\ hos
Duns Scotus, naar han sætter Troen som kun blivende til ved en Villiesact og
derfor lader den paa Troen hvilende Theologie, hvis Maal ikke nærmest er
Erkjendelsen, men Saligheden, Nydelsen af Gud, staae i nærmere Forhold til Villien
end til Forstanden, idet den ikke bliver en speculativ, men en praktisk Lære, der
skal gribe vor Villie.

Naar vi følge Problemet gjennem Scholasticismens Tid, finde vi en Udvikling, der
væsentligen er betinget ved hint Foreningsforsøg, ved hvilket nødvendigviis
Bevidstheden om \opage{20}Forholdet maa klares, og Udviklingen bliver deels en Udvikling
indenfor samme Synspunkt: Videns, deels en Udvikling, ved hvilken et nyt
Synspunkt for Modsætningen antydes.

Hos \emph{Scotus Erigena}, der danner ligesom et Overgangsled mellem Patristiken og
Scholasticismen, hvis Eiendommelighed han antyder, medens han, ligesom hiin, er bunden ved Forholdet til Oldtiden, besvares Problemet gjennem et Postulat af
Harmonien mellem de af samme Kilde udsprungne Aandsvirksomheder, Troen og Fornuften: en Harmonie, i hvilken Fornuften bliver det Overgribende og Primitive.
Men i den egentlige Scholasticismes Tid gaaer Udviklingen i den Retning, at
Fornufterkjendelsen underordnes Troeserkjendelsen, og idet den i Underordningen
satte Differents skærpes mere og mere, naaes der tilsidst til en abstract
Modsætning mellem Tro og Viden, mellem Religionens og Videnskabens Sandhed, saa at
den først blot quantitative Differents bliver til en qualitativ, og endelig til en
uforsonlig Modsætning, og idet denne Udvikling nærmer sig til sin Afslutning, er
det, at der vises hen til et Princip, der, som Princip for Troen, er et andet end
Erkjendelsens og Videns Princip overhovedet.

I Scholasticismens første Periode sætter \emph{Anselm} Troen som Udgangspunkt og Norm,
men fordrer en Viden, der optager Troesindholdet i sig, og han tillægger den
begrebne Tro et høiere Værd end den ubegrebne. Hos de store Repræsentanter for
Scholasticismen i dens Culminationstid, hos \emph{Thomas Aquinas} og \emph{Duns Scotus}, sondres
en høiere Troens Erkjendelse fra den af Naaden ikke oplyste Fornufts Erkjendelse,
men \emph{Thomas} fastholder, ved Siden af Troens Nødvendighed, dens Overensstemmelse med
Fornuften, idet ogsaa denne er en Aabenbaring af Gud, og \emph{Duns Scotus} fastholder,
ved sin Lære om Proportionaliteten mellem Modsætningerne, Sammenhængen mellem den
overnaturlige og den naturlige Erkjendelse, selv om han maa indrømme, at vi ikke
kunne naae til at begribe Troen, hvilken han, som ovenfor angivet, lader have sit
menneskelige Udspring fra en Villiesact. I Scholasticismens Slut\opage{21}ningsperiode
endelig bliver ifølge den absolute Modsætning mellem det Overnaturlige og det
Naturlige Differentsen mellem Theologiens høiere Erkjendelse, der er givet ved den
ved en aandelig Skabelse frembragte Tro, og den naturlige Viden saaledes skærpet,
at den verdslige Videnskab end ikke skal kunne give en Henvisning til den
guddommelige Sandhed, en Forberedelse for vore høiere Øiemed. Der udtales nu, f.\
Ex.\ af W.\ \emph{Occam}, at Noget kan være sandt i Theologien, der er falsk i
Philosophien, og hermed har Scholasticismen selv udtalt sin Dom over den Forening,
der var dens væsentlige Opgave, den har selv sondret de to Elementer og de to
Principer, den vilde sammentvinge.

\section{Problemets Besvarelse i den nyere Tids Philosophie og Theologie}
\label{sec:nyere-tid}

\noindent I den christelige Middelalder, i hvilken Philosophien var Theologiens
Tjenerinde, Tanken behersket af Kirkens Autoritet, maatte Problemets Besvarelse
væsentlig bære Kirkens Præg, og de enkelte Røster, der hævede sig for en
uafhængig Tænkning og hævdede Fornuftens høieste Ret, maatte overdøves og
forstumme. Anderledes viser Forholdet sig fra den Tid af, da Scholasticismen selv
opløste den middelalderlige modsigende Enhed, da Philosophien og Theologien fik
hver sin Sandhed og den emanciperede Philosophie nu autonomisk udfoldede sin
Sandhed: den i sit Væsen uforstyrrede menneskelige Fornufts Sandhedsindhold. Mod
Theologiens Svar paa Problemet stod nu Philosophiens, og jo mere denne vandt
Selvstændighed og hævdede sig uanfægtet af Kirken, desto mere bestemtes
Theologiens Opfattelse af Problemet ved dens.

Den \emph{nyere Tids Philosophie} er bleven til under en Opposition mod den af Kirkens
Autoritet bundne, ufrie, formelle og virkelighedsløse scholastiske Tænkning og
mod Kirkelæren selv, en Opposition, der allerede i Overgangsperioden udviklede
sig til en lidenskabelig Kamp, under hvilken Philosophien fik sine Marthyrer. Som
Heelhed \opage{22}taget staaer den nyere Tids Philosophie i Modsætning til den christelige
Middelalders religiøse Livsanskuelse. Opfattelsen af det Naturliges Væsentlighed
og af Tankens Autonomie og dens Gyldighed som Sandhedskilde skiller afgjørende
Philosophien fra Christendommens Verdens- og Livsanskuelse. Dog bliver der, i
Forholdet til Christendommen og, som Følge deraf, i Besvarelsen af Problemet om
Videns Forhold til Troen, for hvilken denne bestemte Religionsform væsentligen
toges som Repræsentant, en stærkt udpræget Forskjellighed mellem de to
Hovedretninger, i hvilke denne philosophiske Epoche har udviklet sig: den
realistiske og den idealistiske.

Den \emph{realistiske} Retning maatte, ifølge sin Naturopfattelse, strax bestemt sondre
sig fra det Christelige, og dette Forhold har den consequent bevaret og i Formen
bestandig skærpet idet den fra en ved de ydre Forhold paabudt afværgende formel
Accommodation, gjennem en erklæret Indifferentisme naaer til en ironisk Spot og
derfra til et skarpere polemisk Forhold, der efterhaanden antager Charakteren af
et lidenskabeligt Had. Medens Tanken, der successivt saaledes giver Naturen
Overvægten, at den selv absorberes af den som Materie opfattede Natur, uforklaret
ved denne, bliver sig selv en Gaade, dømmer den, som Tanke, i Fornuftens Navn,
Religionen som det Ufornuftige, som det for Menneskets Fremskridt og Menneskets
Lykke ubetinget Fordærvelige; Religionens Kilde bestemmes som Uvidenhed, Tro
bliver Overtro, Ufornuft, endog med en Tilsætning af det sædeligt Forkastelige,
og Dommermyndigheden tilkjendes alene den Tanke, der udtaler Naturens, den
bevægede Materies, det Ene-Værendes, uforanderlige Love og hævder dem i
Anvendelsen paa det Menneskelige, der kun bliver et uselvstændigt Led i den
altomfattende Natur.

I Sammenligning med den gjennemgaaende Ensartethed i den realistiske Philosophies
Stilling til det foreliggende Problem frembyder den \emph{idealistiske} Udviklingsrække
en langt større Mangfoldighed af Former for Forholdet. Den consequente, bevidste
Gjennemførelse af hele Epochens \opage{23}Grundtanke maatte, ogsaa indenfor denne Række,
overalt bringe Modsætningen tilsyne, men den abstracte Lighed, der finder Sted
mellem Idealismen, navnlig hvor den udvikler sig som Panlogisme, og den
christelige Verdensanskuelse, og det, at, paa en enkelt Undtagelse nær, alle de
Tænkere, der repræsentere Idealismen, personlig vare knyttede til Christendommen,
har kunnet fremkalde en Stræben efter at sætte Religion og Philosophie, Tro og
Viden, i Harmonie, og denne Stræben træder frem i forskjellige Nuancer.

\emph{Modsætningen}, der gjennem hele den nyere Tids Philosophie væsentlig opfattes fra
\emph{Videns}, Erkjendelsens Synspunkt, bestemmes enten som en Modsætning mellem den
sande, Fornuftens, Erkjendelse og en usand, Indbildningskraftens, Erkjendelse
(\emph{Spinoza}: \textit{intellectus} --- \textit{imaginatio}), en Modsætning, der væsentligen falder
sammen med Modsætningen mellem Uendelighedens og Endelighedens Bestemmelse, mellem
den evige Idee og den forsvindende historiske Form (\emph{Schellings} ældre
Philosophie); --- eller som Modsætningen mellem Erkjendelsens Maal og Sandhed og
et Udviklingstrin i Erkjendelsen, der endnu kun er ufuldkomment og relativt sandt,
og som i den teleologiske Udvikling optages i den fuldkomne Sandhed, Philosophiens
Sandhed (\emph{Hegel}, naar \emph{Logikens} Bestemmelser af Forholdet mellem Form og Indhold
consequent gjennemføres.) Til denne sidste Modsætningsform ere \emph{principielt} de
Forsøg at føre tilbage, der ere gjorte paa, ved en --- moralsk eller speculativ
--- Omfortolkning at bringe de religiøse Læresætninger ind under den philosophiske
Fornufterkjendelses Omraade; thi medens der i selve disse Forsøg ligger
Forudsætningen af den teleologiske Udviklingssammenhæng, ligger der tillige deri
en Anerkjendelse af Nødvendigheden af at hæve det i kun relativ Sandhed
Fremtrædende til den rene absolute Sandheds Form. At, individuelt, Bevidstheden
om denne Sammenhæng kan være uklarere, at Omfortolkningen kan opfattes som kun
berørende et Uvæsentligt, en ydre Formdifferents, gjør, principielt, Intet til
Sagen.

\opage{24}Den Bestræbelse for at bringe Tro og Viden i Harmonie, der fremtræder indenfor
visse Udviklingsformer af den nyere Philosophie, antager, naar vi see bort fra den
overfladiske Forestilling om en „Fornuftreligion“, en „naturlig Religion“, hvis
rette Udtryk Christendommen skulde være, enten den Form, at der indrømmes en
quantitativ Forskjel til Fordeel for Aabenbaringen (et Overfornuftigt, der dog
ikke skal være et Fornuftstridigt), eller den, at der statueres en Væsenet,
Indholdet ikke berørende Formdifferents til Fordeel for Philosophien. Den sidste
Nuance repræsenteredes fornemmelig af den saakaldte høire Side af den Hegelske
Skole, der støttede sig paa Hegels Religionsphilosophie, hvis restaurerende
Tendents, idetmindste i Udtrykket, havde ført til, at det i hans Logik udviklede
Forhold mellem Form og Indhold var blevet uklart i Anvendelsen paa Religionens
Forhold til Philosophien. Fra hiin Skole overførtes saa denne mæglende Opfattelse
paa Theologiens Gebet i den speculative Dogmatik.

De Antydninger af en, Collisionen forebyggende, Sondring mellem Troens og Videns
Gebeter, der i dette Tidsrum ere givne indenfor Philosophien, faae ingen videre
Betydning. Naar \emph{Spinoza} giver Religionen en praktisk Opgave, og ved denne
Bestemmelse vil afværge Collisionen med Philosophien, saa er det en accommoderende
Vending, der ikke kan fastholdes ligeoverfor hans Philosophies Bestemmelser af
Forholdet mellem det Theoretiske og det Praktiske. Og naar \emph{Leibnitz} etsteds
antyder den samme Sondring, saa tabe hans Yttringer deres Betydning ved Siden af
hans Læres almindelige Bestemmelser, og ved Siden af hans Bestræbelse for at
hævde Dogmernes theoretiske Sandhed.

Indenfor \emph{Theologiens} Gebet træder den skarpt udprægede Modsætning mellem Tro og
Viden væsentlig frem paa Reformationens Tid og navnlig hos \emph{Luther}. Tankens
Omraade er, efter hans Opfattelse, det Verdslige, men paa det Religiøses Omraade
mægter den Intet uden at fornægte. Vel søger han undertiden at hævde Tanken en
Be\opage{25}tydning som negativ bestemmende, idet han tilkjender den Evnen til at afgjøre
hvad det Guddommelige \emph{ikke} er; men han kan ikke consequent gjennemføre denne
Indrømmelse. Hos de gamle lutherske Dogmatikere mildnes allerede Modsætningen.
Under Theologiens sildigere Udvikling gjør, directe eller indirecte, den mæglende
Philosophies Indflydelse sig gjældende, og Differentsen bliver, hvor den ikke
ganske udjævnes (som i den rationalistiske Identificeren af den naturlige Religion
og Christendommen), en \emph{quantitativ} Differents, selv hvor Benævnelserne lyde paa en
qualitativ, der i Realiteten principielt er udelukket ved Theologiens Fordring paa
at være Videnskab. Det Forsøg, der af \emph{Schleiermacher} blev gjort paa, i
Overensstemmelse med hvad der allerede tidligere var antydet af \emph{Lessing}, under
Modsætningen mellem Religion og Philosophie, at hævde Religionens Gyldighed ved at
give hver af Modsætningerne sit eiendommelige Gebet, ved at sætte hver af dem i
Forhold til en eiendommelig Aandsevne, har kun en underordnet Betydning i
Sammenligning med det Sondringsforsøg, der sildigere er gjort af Kierkegaard. Vi
støtte denne Dom om dets Betydning deels paa Schleiermachers mangelfulde
Opfattelse og vilkaarlige Begrændsning af Viden, deels, og fornemmelig, paa det
Uklare og Modsigende i hans Stilling overhovedet, idet han, paa Grundlag af et, om
end inconsequent, pantheistisk philosophisk Standpunkt vil hævde det positivt
Religiøses, særlig det Christeliges, Gyldighed, og ud fra den ved hans
philosophiske Standpunkt betingede indholdsløse og rent negative Opfattelse af det
Absolutes Idee, bestemmer Religionens Væsen og Religionens Organ, hvorved det
netop maa blive umuligt at naae til det concrete Dogme og til det for den positive
Religion uundværlige Historiske.

Mod den ved den Hegelske Religionsphilosophie og en Deel af hans Skole
repræsenterede Stræben efter at fastholde den væsentlige Enhed af Religion og
Philosophie under den som uvæsentlig opfattede Formforskjel, reistes der en
Indsigelse saavel fra Religionens som fra Philosophiens \opage{26}Side. Fra hiin nedlagdes
der en Protest mod Forflygtigelsen af det for Religionen Eiendommelige ved dets
Overføren i Philosophiens Begreber og i Begrebsprocessens Form; fra denne
hævdedes, i Overensstemmelse med Hegels logiske Begrebsudvikling,
Formforskjellens Uadstillelighed fra en tilsvarende Indholdsforskjel. Erkjendelsen
af Mæglingsforsøgets Uigjennemførlighed førte saa til Forsøget paa at løse
Problemet gjennem en Sphæresondring. I Tydskland gjør \emph{Feuerbach} dette Forsøg i
Videns Interesse. Charakteren af hans Løsningsforsøg er i det Foregaaende angivet
i de væsentlige Grundtræk; en Kritik af det vil, deels directe, deels indirecte,
blive given i det Skrift, der vil slutte sig til dette. Hos os traadte Problemet
ind i en ny Phase ved \emph{Kierkegaard} ved det Løsningsforsøg, der, medens det ligesom
Feuerbachs hviler paa en Sphæresondring, tilsigter Principernes Forening i samme
Bevidsthed gjennem begges Anerkjendelse.

Fra dette Løsningsforsøg hidrører alt det, der i det \emph{Nielsenske} har en virkelig
Betydning, og dette Forhold imellem dem vil allerede her, ved Slutningen af denne
korte Oversigt over de væsentlige Stadier i Problemets Historie, kunne gjøres
tydeligt ved en foreløbig Sammenstilling af begge med Hensyn til de Punkter, der
paa dette Sted blive af Betydning.

Hos \emph{Kierkegaard} er det Religiøse den oprindelige og bevægende personlige
Interesse og det Maal, til hvilket der fra Begyndelsen af sigtes hen; men i
Udviklingen af Standpunktet er Viden Udgangspunktet, og det er gjennem Videns
Selvbegrændsning, gjennem Erkjendelsen af den Skranke, Existentsforholdet sætter
for den existerende Vidende, at Troens og Troesforholdets Nødvendighed
anerkjendes. Videns Gyldighed bliver da ikke tilintetgjort, men anerkjendes som
relativ, indenfor sin bestemte Grændse; Troen bliver, som den for den
Existerende høieste Vishedsform, det absolut Gyldige, det, der ene giver den
Existerende den, ligesom af Uvisheden udrevne, forsonende Vished om det, der
tilfredsstiller Personlighedens dybeste Trang. Viden \opage{27}kan, netop paa Grund af sin
Begrændsning, ikke forstyrre Troens Indhold, ikke heller begribende tilegne sig
det; det Eneste, den i Forhold til Troen kan gjøre, er at blive tjenende, idet
den værner om Paradoxet, bevarer det som saadant i dets Skarphed og
Utilgængelighed. Omvendt lader Troen Viden være uforstyrret paa dens begrændsede
Omraade, og ifølge dette Forhold kunne de to Principer med deres Sphærer faae
deres Plads i samme Bevidsthed uden forstyrrende Sammenstød.

Ogsaa hos \emph{Nielsen} bliver Viden Udviklingens Udgangspunkt, ogsaa hos ham
begrændser Viden sig selv, saaledes at den, gjennem Erkjendelsen af sin ikke
blot relative, men absolute Grændse, føres til at anerkjende det med den selv
absolut uensartede Troesprincip i dets Gyldighed. Viden opklarer Forholdet
mellem Troesprincipet og Vidensprincipet og deres Sphærer, og den er istand
dertil, idet den, afspeilende, omfatter sin Modsætning, og saaledes faaer ogsaa
den til Gjenstand, ikke blot det, den begrunder. Hver af de Sphærer, der ere
bestemte ved de med hinanden absolut uensartede Principer, udfolde sig i deres
indre Uendelighed uden paa udvortes Maade at indskrænke hinanden og derfor uden
at komme i Conflict med hinanden; og saaledes adskiller, netop paa Grund af
Uensartetheden, Sphærernes og deres Principers Forening i samme Bevidsthed ikke
Bevidsthedens Enhed.

I de afgjørende Grundtræk stemme saaledes de to Former for Theorien overens.
Skulde, i de angivne Bestemmelser, en Differents søges, vilde det være den, at
Nielsen skulde sætte begge Sphærer som uendelige og uindskrænkede, begge
Principer ifølge deres absolute Uensartethed som absolut gyldige, og at han
skulde statuere et Coordinationsforhold, medens Kierkegaard tydelig udtaler
Videns blot relative Gyldighed og dens Sphæres Endelighed, og bestemt stiller
Viden i et Subordinationsforhold til Troen. Men ved en nærmere Betragtning vilde
en saadan Differents vise sig som illusorisk.

Hvad først Sphærernes Uendelighed angaaer, der umid\opage{28}delbart maatte følge som
Consequents af Principernes absolute Gyldighed, og som maatte indeslutte, at al
Virkelighed, Virkeligheden i dens Fuldstændighed, fik sin Plads indenfor hver af
de to Sphærer, saa lader den sig ikke fastholde. Med Hensyn til Videnssphæren,
der synes at faae en saadan Uendelighed derved, at Viden ogsaa har det til
Gjenstand, som den blot afspeiler, vilde Uendeligheden ialfald blive en
Uendelighed med Hensyn til Omfang, ikke til Indhold. Ved selve Adskillelsen
mellem det, der assimileres --- det ved Tanken Begrundede og Begrebne
assimileres --- og det, der kun afspeiles, bliver der i Vidensindholdet sat en
Grændse, der viser hen til en Væsensbegrændsning i Viden som menneskelig Viden,
selv om denne Begrændsning postuleres hævet i Videns Idee. Men
Væsensbegrændsningen, der umiddelbart træder frem gjennem Vidensindholdets
Begrændsning, maa uundgaaelig atter føre tilbage til en Begrændsning med Hensyn
til Omfang, paa Grund af det dialektiske Forhold mellem Omfang og Indhold og
deres nødvendige gjensidige Bestemmen. Det er umiddelbart indlysende, at det,
der kun afspeiles i Viden, ikke blot naar Viden tages i stræng Forstand som
tilegnende (assimilerende) Viden, og kun i denne Forstand er den egentlig Viden,
kun vides efter sin abstracte Væren; men at det, selv naar Viden tages i udstrakt
Betydning, ogsaa kun vides abstract, fordi den concrete Viden kun er ved
Erkjendelsen af Sammenhængen, men for den afspeilende Viden maa det Afspeiledes
Sammenhæng unddrage sig idet dens Form --- paa Grund af den absolute
Uensartethed i Principerne --- bliver absolut incommensurabel med Videns. Og
hvad der gjælder med Hensyn til Videnssphærens Uendelighed har sit Tilsvarende
med Hensyn til Troessphærens. Ikke heller denne kan fastholdes i den angivne
Betydning. Indenfor Troen findes der intet Forhold, der svarer til
Afspeilingsforholdet i Viden og ifølge hvilket Videns Indhold kunde optages i
den. Og selv om et saadant Forhold fandt Sted, vilde Troen ikke faae al
Virkelighed til Gjenstand, af den tilsvarende Grund til den, af hvilken Viden
ikke havde den. \opage{29}Troessphæren kan saaledes ikke hævde sin Uendelighed. Kun dette
Resultat søges her; i det Efterfølgende vil det paa sit Sted blive viist, at
Sphærernes factiske Begrændsethed, hvormeget end Selvbegrændsningen accentueres,
dog vil føre til et Sammenstød ved Grændsen med den Sphære, der i denne berøres.

Hvorledes Begrændsningen i begge Sphærer bliver at opfatte, er let at see naar
vi tage en concret Gjenstand. Tage vi saaledes det Naturlige, saa existerer
Naturloven og den ved Naturloven bestemte Causalsammenhæng, der begge ere
Vidensbestemmelser, og som \emph{væsentlig} Vidensgjenstand have \emph{Virkelighedsgyldighed},
ikke for Troen, der kun seer Tingene som skabte, i deres Afhængighed af Gud.
Omvendt erkjender Viden ikke det creaturlige Afhængighedsforholds
Sammenhængsform, der, som \emph{væsentlig} Troesgjenstand, efter Formsætningen,
ligeledes maatte have \emph{Virkelighedsgyldighed}. Den totale Virkelighed, hvis Væren
i Viden og Tro skulde give dem Uendelighed, reduceres altsaa for hver til en
ucomplet Virkelighed, thi hvad der er Object for Viden og Tro, er kun de
sammenhængsløse Ting, ligesom Tingene „an sich“, men ikke Forholdet, i hvilket
de staae til deres Grund, ikke Forbindelsen med denne og med den ensartede
Virkelighed, altsaa ikke det, der gjør Tingene til det, de ere, og denne, den
concrete Virkelighed bestemmende, Sammenhæng tilhører kun Viden \emph{eller} Tro, og
falder udenfor den fælleds Virkeligheds Sphære, men som en Virkelighed; thi kun
i den ere Tingene hvad de ere. Men dette vil egentlig sige: i Realiteten
fordobbles Virkeligheden og den fordobblede Virkelighed deles mellem Tro og
Viden, saa at det tilsyneladende Fælleds kun bliver et Homonymt. Og denne
Virkelighedens Fordobbling vil vise sig som modsigende og ugjennemførlig.

Hvad det andet ovenfor udhævede Punkt angaaer: Sphærernes Coordination, saa er
denne hos Nielsen, trods Bestemmelsen af Principerne som absolut uensartede, kun
tilsyneladende; i Virkeligheden findes der hos ham, ligesom hos Kierkegaard, en
Subordination Sted, saa at Uligheden \opage{30}mellem Theoriens to Fremstillingsformer
ogsaa her forsvinder. Subordinationen af Videns Princip og Videns Sphære under
Troens er hos Nielsen ikke blot indirecte sat paa flere Maader: ved Bestemmelsen
af Villien, der bringes i særligt Forhold til Troen, som det egentlig
constitutive i det menneskelige Væsen; ved Opfattelsen af den theoretiske og den
praktiske Erkjendelses væsentlige Bestemmelser og deres Forhold til
Personligheden, og ved Opfattelsen af „den religiøse Ethiks Overlegenhed over
den rationelle“; men den er udtrykkeligt udtalt i Betegnelsen af Religionens
Sandheder som \emph{over}videnskabelige Sandheder.

I de Punkter, der have en væsentlig Betydning for Problemets Udvikling, falder
altsaa den Nielsenske Theorie sammen med Kierkegaards, og Originaliteten
tilkommer ene denne. Dog bliver der, ved Kritiken af Theorien, Anledning til at
holde dens Fremstillingsformer ude fra hinanden forat belyse de Tilføininger,
den har faaet hos Nielsen ved dennes Forsøg paa at sætte Grundbestemmelsen af
Forholdet mellem Tro og Viden i Forbindelse med hans psychologiske og
metaphysiske Theorie, og de Modsigelser og Inconsequentser, af hvilke den hos
ham lider, men som Kierkegaard har vidst at undgaae.

\begin{center}
---
\end{center}

Efter det i det Foregaaende Udviklede blive altsaa de væsentlige Standpunkter,
der ere indtagne i Besvarelsen af Problemet, følgende:

\begin{enumerate}[label=\arabic*)]
\item Det religiøse Standpunkt, der sætter Troen, som det ved Guds overnaturlige
  Virken betingede Tilegnelsesforhold til den guddommelige Aabenbaring, som
  absolut hævet over den naturlige Viden, der opfattes som den, der ikke blot
  ikke kan erkjende Aabenbaringens Indhold, men maa fornægte det.
\item Theologiens og en enkelt philosophisk Forms Standpunkt, der statuerer en
  quantitativ Differents til Fordeel for Troen\footnote{Til dette Standpunkt
  gaaer i Realiteten den scholastiske Sondring af \textit{duæ veritates}
  tilbage.}.
\item \opage{31}Den rene Rationalismes Standpunkt, paa hvilket Forskjellen udviskes ved
  Forestillingen om en Fornuftreligion, en naturlig Religion, med hvilken
  Christendommen identificeres.
\item Mæglingens Standpunkt indenfor den idealistiske Philosophie, paa hvilket
  der statueres en Formdifferents til Fordeel for Philosophien.
\item Den realistiske og den consequente idealistiske Philosophies Standpunkt,
  paa hvilket, idet Forholdet bestemmes fra Videns Synspunkt, Viden faaer en
  absolut Overvægt over Troen, og bliver det eneste Sandhedsorgan.
\item Sondringen af Sphærerne for at udskille det Religiøse.
\item Sondringen af Sphærerne for at tilveiebringe begges Forening i samme
  Bevidsthed.
\end{enumerate}

Af disse Standpunkter vil, som allerede tidligere angivet, det theologiske og
det, der gjennem Sondringen af Sphærerne tilsigter deres Forsoning, blive
Gjenstand for en særskilt, udførlig Kritik; medens de øvrige, forsaavidt det
overhovedet er nødvendigt, ville blive belyste gjennem Kritiken af hine, eller
gjennem Udviklingerne i det Skrift, der danner Fortsættelsen og Afslutningen af
dette.


%% ============================================================
%%  FORELØBIGE BEGREBSBESTEMMELSER (pp. 31--43)
%% ============================================================
\chapter{Foreløbige Begrebsbestemmelser til Belysning af Problemet}
\label{ch:begreber}

Idet vi opkaste Spørgsmaalet om Troens Forhold til Viden, tage vi Tro, ikke i
den vage, udviskede Betydning --- som \emph{umiddelbar} Viden, \emph{umiddelbar} Vished om
Væren --- i hvilken Ordet bruges af Jacobi og af Theologer, der i Tvetydighedens
Interesse adoptere hans Ubestemthed, men vi tage Ordet i specifisk religiøs
Betydning, \opage{32}i Betydning af den til Religionen, bestemtere til den positive
Religion, sig forholdende Tro. Bestemmelsen af Troens Begreb bliver derfor uadskillelig fra Religionens Begreb.

Dette sidste Begreb opfatte vi først i dets \emph{ideale} Betydning, i Betydningen af
det, der er den bevægende Magt i Udviklingen af de phænomenale Former af det
Religiøse, det Maal, til hvilket Religionernes historiske Udvikling stræber hen.
Til Begrebet af Religionen efter dens ideale Betydning naae vi, idet vi gaae ud
fra, hvad der er det religiøse Forholds Øiemed. Dette kunne vi med eet Ord
bestemme som \emph{Aandens Forsoning}, dens Forsoning i sig og med sit Princip, sit
Væsens Grund. Menneskeaanden søger i Religionen Befrielse, Fred, Hvile,
Forsoning; den søger den for sig som Totalitet, efter sit Væsens fulde
Bestemthed; kun saaledes er den i \emph{sig} forsonet. Det religiøse Forhold maa derfor
være et Forhold, i hvilket Menneskeaanden efter alle sine væsentlige
Grundvirksomheder kommer til Hvile, der indeslutter Forsoning og
Tilfredsstillelse for Erkjendelsen og Følelsen som for Villien. Det er den udelte
Menneskeaands Trang, der fører til det religiøse Forhold, og naar denne Trang,
medens Forsoningen søges, nærmest og umiddelbarest fremtræder under en enkelt
Grundvirksomheds Form, saa maa Tilfredsstillelsen af dennes Trang, der her
ligesom repræsenterer de andre Virksomheders Stræben, gaae i Enhed med
Tilfredsstillelsen ogsaa af deres Trang: hele Menneskeaandens Forsoning maa være
tilstede under den enkelte Aandsfunctions Forsoning. Men den fulde Forsoning, den
efter sin Totalitet bestemte Aands Forsoning med sig og med sit og Tilværelsens
Princip, kan kun finde Sted ved det Forhold til Principet, ved hvilket Aanden, i
Bevidsthed om sin Væsensenhed med dette, i fri Hengivelse sætter det som sin
\emph{hele} Værens Princip. Dette vil med andre Ord sige: det religiøse Forhold
maa bestemmes som \emph{det absolute} Forhold \emph{til det Absolute}. Denne Bestemmelse
indeslutter et Tredobbelt. Deels viser den tilbage til det, at det Absolute maa
være i Væsensenhed med \opage{33}vort Væsen; thi kun til det, med hvilket vi ere i
Væsensenhed, med \emph{vort} Absolute, kunne vi staae i et absolut Forhold, i et, paa
een Gang vor hele \emph{Erkjenden} og vort hele \emph{Liv} omfattende Forhold. Deels viser den
hen til, at det, til hvilket vi staae i det religiøse Forhold, maa være det i
Sandhed Uendelige; thi kun i det Uendelige kan Menneskeaanden absolut finde
Hvile. Deels ligger endelig deri den Consequents, at det religiøse Forhold, netop
som det absolute, ikke kan falde udenfor de relative og særegne Forhold, hvorved
det selv maatte blive relativt, men maa omslutte disse Forhold i sin concrete,
omfattende Enhed.

Men naar det religiøse Forholds Begreb saaledes bestemmes efter sin ideale og
universelle Betydning, saa viser der sig en Incongruents mellem dette Begreb og
de positive, historiske, som særegne Former af Aandslivet udprægede Religioner.
Gjennem alle disse gaaer som Udviklingens bevægende Magt hiin Menneskeaandens,
ved Grundforholdet til Principet fremkaldte Stræben efter det i Sandhed religiøse
ɔ: det hele Menneskevæsenet og Menneskelivet omfattende absolute Forhold; det er
denne Stræben, der under Religionsformernes historiske Udvikling fører dem fremad
fra Naturreligionernes laveste Former til Aandsreligionernes høieste; men gjennem
alle de historiske, positive Religionsformer træder Religionens ideale Begreb
frem i en Relativitets- og Endelighedsbestemthed: de positive Religioner vise sig
som \emph{endelige} Udtryk for det Religiøses Væsen. Endeligheden bliver kjendelig
gjennem Formerne for Menneskeaandens Virksomhed i Religionerne, gjennem den
derved betingede Opfattelse af det Guddommelige, og gjennem Bestemmelsen af det
Menneskeliges Vexelvirkningsforhold til dette.

Vi søge først Endelighedsbestemtheden i den religiøse Bevidstheds
\emph{theoretiske} Form. Denne er, paa de positive Religioners Trin, selv om den
religiøse Bevidsthed som menneskelig Bevidsthed gjennemvirkes af Tanke, ikke den
til Klarhed og Sammenhæng udviklede, efter sit Væsen uendelige begribende Tankes,
men væsentlig den umiddelbare \opage{34}Følelses og Forestillingens, og fra denne er
Endelighedsbestemtheden uadskillelig. Men Bevidsthedsformens Eiendommelighed maa
nødvendigen bestemme Bevidstheds-Indholdets Eiendommelighed. Det aandelige
Indhold kan kun da have en ufuldkommen, med Indholdet efter dets Væsensbestemthed
uadæquat Form uden at Indholdet for Bevidstheden forvandles, naar det allerede
har været tilstede for Bevidstheden i den adæquate Form, og fra denne med
Bevidsthed er ført over i hiin. Formen bliver da kun Indklædning, Symbol,
Allegorie. Men i Menneskehedens Udvikling gaaer bestandig Religionens Form forud
for den begribende Erkjendelse: det Absolute er tilstede i dens Former før det er
tilstede for Bevidstheden i Tankens Form. Det er altsaa ikke fra en adæquatere
Form sat over i Religionens Forestillingsform: denne er for Bevidstheden ikke en
Indklædning, men den \emph{eneste} Form for Indholdet; den gjælder altsaa som den
sande og adæquate Form. Indholdet kan saaledes ikke blive uberørt af Formens
Eiendommelighed: det er kun som det saaledes formede Indhold for Bevidstheden, og en Adskillelse af det fra denne Form bliver umulig paa det oprindelige religiøse
Standpunkt. En saadan Adskillelse finder kun Sted ved en vaagnende Reflexion, og
denne kommer paa det religiøse Standpunkt kun frem i Forhold til andre, nærmest
til underordnede, religiøse Standpunkter. I Forhold til sit eget Indhold er den
religiøse Bevidsthed, efter Sagens Natur, ikke kritisk: den gjenfinder sig selv i
det saaledes formede Indhold.

Den religiøse Bevidstheds specifiske Form faaer saaledes en væsentlig Betydning,
og navnlig bliver her Forestillingsformens Indvirkning at fremhæve, idet den
religiøse Bevidstheds Indhold, der i Følelsens Umiddelbarhed er i sammensluttet,
intensiv, men uudfoldet Enhed, først gjennem Forestillingen udfoldes i
Bestemthed, bliver concret \emph{Gjenstand} for Bevidstheden.

Forestillingens Eiendommelighed bestemmes ved dens psychologiske Stilling til
Sandsning og Tanke, mellem hvilke den danner Mellemleddet. Forestillingen har,
som Mellem\opage{35}led mellem begge, et Moment af begge: den forbinder Sandsningens
receptive Afhængighed med Tankens Frihed, dens Enkelthed med Begrebets
Almindelighed. Det, der forestilles, er løsnet fra det umiddelbart værende Ydre,
hensat fra det ydre Rum i Bevidsthedens indre Rum, det er begyndt at sondres i
relativ Almindelighed fra den ydre Enkelthed og dens Tilfældighed; men efter
Indholdet er det endnu bestemt derved. Deraf følger, at det, der er i
Forestillingen, berøres af \emph{Sandselighedens} Form, at altsaa, naar det
Absolute skal opfattes gjennem Forestillingen, dets almene Væsen maa anskues i
den ved en Sandselighedens Reflex bestemte Enkelttilværelses Form, dets evige,
Tidsudviklingen omfattende og idealiserende Livsproces som en ved Tidens
Endelighed bestemt Proces, hvis enkelte Momenter falde sondrede udenfor og efter
hverandre; at det Billedliges, det Mythiskes og det Historiskes Former maae blive uadskillelige fra Opfattelsen af det Absolute.

Medens Følelsesformen endeliggjorde Indholdet ved at lade det blive uden
Udfoldelse, bringer altsaa Forestillingens Form ved sit Forhold til
Udvortesheden en Endelighedsbestemthed ind i den religiøse Bevidstheds
Opfattelse af det Absolute. Den træder frem i den Sondringens Form, ved hvilken
det Guddommeliges Transscendents bestemmes; i den Opfattelse af Gud under
Personlighedens Form, i hvilken den sandselig bestemte Enkeltpersonligheds
(Individets) Eiendommelighed skinner igjennem; i det Historiske i Religionerne, i
hvilket det Guddommelige faaer endelig Tids- og Rumsbegrændsning, og i
Opfattelsen af Guds Virken som en Virken gjennem Underets\footnote{Underet, den
positive Religions specifiske Kjendemærke, er ikke Almagtens høieste Udtryk, men
Almagtens Endeliggjørelse, Udtrykket for en illusorisk Almagtens Uendelighed
istedenfor den sande Almagtens Uendelighed, altsaa i Virkeligheden Almagtens
Negation. Almagten er endeliggjort netop ved Negationen af Naturloven, der i
Underet sættes som Udtryk for en Virken af den almægtige Villies Andet, der
modsættes Almagtens Virken.} enkelte, sondrede Acter, Vilkaarlighedens
continuitetsløse Udtryk.

\opage{36}Til Endelighedsbestemtheden i den religiøse Bevidsthed, opfattet fra dens
theoretiske Side, svarer, og ved den betinges, en lignende
Endelighedsbestemthed i den, seet fra den \emph{praktiske} Side. Tankeforholdet
til Gjenstanden er et Frihedens og Uendelighedens Forhold. Individet er i
Forholdet frit, idet det gjennem Gjenstanden, i hvilken Tanken som begribende
gjenfinder sig selv, forholder sig til sig selv, og i dette Forhold har det en
indre Uendelighed. Og som frit lader det Gjenstanden bestaae i dens Frihed, idet
det, naar det opløser Gjenstandens Væreform i Tanke, naar det søger sig selv i
Objectet, ikke søger og gjenfinder en subjectiv, individuel Tanke, men den
objective, almeengyldige Tanke, der paa een Gang constituerer Gjenstandens og
Individets Væsen. Det er saaledes, netop fordi Forholdet som Frihedsforhold er et
Uendelighedsforhold, uinteresseret, uegoistisk i Forholdet til Gjenstanden. Af en
anden Art bliver Forholdet, hvor Forestillingen er Bevidsthedens Form. Det
forestillende Subject er i en ved Forestillingens Natur betinget
Endelighedssondring fra sin Gjenstand, det gjør sig derfor gjældende efter sin
Enkelthed, mødes endnu ikke med Gjenstanden i det Almene, staaer ikke i hint
Frihedens og Uendelighedens Forhold til den. Anvendes dette paa Forholdet til
det Absolute, saa medfører det følgende Consequents. Idet det troende Individ
forholder sig til det i Forestillingens Form opfattede Absolute, staaer det som
endelig bestemt Enkeltpersonlighed ligeoverfor en guddommelig Enkeltpersonlighed,
fra hvilken i Virkeligheden Endelighedens Bestemmelser ikke ere fjernede. Men
dette Forhold, i hvilket Endeligheden er blivende, maa netop derfor faae den
aabenbare eller skjulte Egoismes Charakteer. Selv om Forholdet betegnes som et
Kjærlighedsforhold, saa gjælder det, at Kjærlighedens Hengivelse kun er absolut
fri for Egoisme, hvor Kjærligheden forholder sig til det i sin sande Uendelighed
erkjendte Absolute\footnote{jvfr. S. Kierkegaard: Kjærlighedens Gjerninger.}, og
kun hvor Kjærlighedens \opage{37}Hengivelse er absolut fri for Egoisme, er Kjærligheden den
sande og Forsoningen ved Kjærligheden, den sande Enhed med Gud gjennem den mulig.
Hvor Endelighedsbestemtheden bliver, der kan denne Enhed ikke naaes, Forsoningen
ikke opfattes som Individets Forsoning i sin Væsensgrund, men den maa opfattes som
en Forsoning med et andet Væsen, der ved sin Enkelthedsform bevarer Sondringen
fra Mennesket. Dette transscendente Væsens Fordring, dets Lovbud bliver saa
bestemmende for det religiøse Forhold, ved hvilket Forsoningen skal
tilveiebringes; til Individets Opfylden eller Ikke-Opfylden af Fordringen knytter
sig en Belønning og Straf, der ligesom Lovbuddet er bestemt ikke ud fra et
Væsensforhold, men ved den ubetingede guddommelige Villie, og derfor bliver et
Egoismens Motiv; og selv hvor Forestillingen om Saligheden nærmest berører Tanken
om Enheden med Gud, trænger en Endelighedsbestemmelse sig ind deri, deels directe
i Salighedens Opfattelse som Løn, deels indirecte gjennem Tidsforestillingen,
gjennem Forestillingen om Salighedens Fremtidighed, gjennem Sondringen af
Handlingens Følge fra Handlingen selv. Overhovedet viser den blivende
Endelighedsbestemthed sig deri, at det religiøse Forhold bliver et kun \emph{særligt} Forhold; thi dette er Eet med, at det bliver et endeligt Forhold.

Hvis vi, istedenfor at tage Udgangspunktet fra Opfattelsen af den religiøse
Bevidstheds theoretiske Bestemthed og fra denne at drage Consequentserne for
dens praktiske Forhold, umiddelbart vilde holde os til den religiøse Bevidstheds
praktiske Bestemthed, og gjøre den Synsmaade gjældende, at Udgangspunktet maatte
tages fra denne, fordi Villien var det for det religiøse Forhold Afgjørende, saa
vilde dette dog ikke føre til noget andet Resultat. Vi vilde da nemlig, idet hiin
Fremhæven af Villien hviler paa en principiel Modsætning imellem den og
Erkjendelsen, som det Bestemmende faae Villien i dens Isolerethed, men Villien i
dens Isolerethed er, hvad der senere udførlig vil blive paaviist, kun
uigjennemklaret, egoistisk Drift, altsaa en Vil\opage{38}liens
Endelighedsform\footnote{Hvormeget end Villiens Autonomie og principielle
Selvstændighed accentueres, saa undgaaes dog ikke dens Endeliggjørelse naar den
opfattes i Isolerethed, thi netop den abstracte Sondring mellem Videns
Almindelighedsform og Villiens Enkelthedsform gjør begge endelige.}, om den end
udsmykkes med de høieste Navne. Hvor Villien, i postuleret Autonomie, sættes som
det Constitutive i det menneskelige Væsen, og som det, der i sin Sondring skal
udtrykke det høieste Aandelige, naaer den netop ikke længere end til den
ufuldkomne Villiens Form, der danner Parallelen til Forestillingens Form paa det
theoretiske Gebet. Kun i Enhed med den sande Tænken naaer Villien sin Sandhed.
Hvor den ubevidst følger sin egen Lov, er den blind Drift, og hvor den bestemmes
ved et ikke til Tanke hævet bevidst Villiesmotiv („en Villieserkjendelse“), har
den altid en Affectens Bestemthed i sig; den er bestemt ved Attraa og Lidenskab,
og denne praktiske Endeligheds-Bestemthed virker uundgaaelig tilbage paa
Bevidsthedens Bestemmen af Gjenstanden, hvad der især tydeligt viser sig i den
Bedendes Opfattelse af det Guddommelige, der skal bevæges ved Bønnen. Den
Affectens Form, der træder frem i Religionen, er væsentligen Frygten og Haabet,
og det, at Frygt og Haab i Religionen bestemmes ved Forestillingen om et
\emph{overmenneskeligt} Væsen (et Udtryk, der er anvendeligt paa alle de
forskjelligste Religionsformer), befrier dem ikke fra Uklarheden og Egoismen, og
det gjør det ikke fordi det Guddommelige, opfattet gjennem en Form, der ikke er
Begrebets, har Endelighedsbestemtheden i sig.

Til Bestemmelsen af Troens Begreb i Forhold til de phænomenale Religionsformer
ere nu Forudsætningerne givne. Idet det Guddommelige faaer Charakteren af et
Transscendent og Supranaturalt, af et fra Menneskevæsenet Sondret, saa bliver det
saaledes bestemte Guddommeliges Meddelelsesforhold til Menneskeaanden opfattet
som et Aabenbaringsforhold, saaledes forestillet, at Aabenbaringen \opage{39}ikke er en af
Væsensenheden med fri Nødvendighed fremgaaende indre Væsensmeddelelse, men en
vilkaarlig, ved ydre Midler betinget, Meddelen af et Indhold, der, som Meddelelse
om Guds Væsen, er begrændset ved Guds Vilkaarlighed, og, som Meddelelse om Guds
Bud, der skal bestemme den menneskelige Handlen, har sin Gyldighed ved
Aabenbaringssubjectets ved ingen Væsensnødvendighed bundne Villie. Aabenbaringen,
taget i denne Betydning, fordrer sine særegne Redskaber, og den begrunder en
\emph{positiv} Religion, der indeholder en Lov for Erkjendelse og Handling.

Disse Bestemmelser med Hensyn til Begreberne: Aabenbaring og transscendent,
supranaturalt Aabenbaringssubject have ikke blot deres Gyldighed for
Aandsreligionerne, men ogsaa for Naturreligionerne, om de end, efter Sagens
Natur, i disse træde mindre tydeligt frem. Ogsaa i Naturreligionerne bliver,
gjennem Naturmagtens uvilkaarlige Personificeren, en Transscendents tilveiebragt,
saaledes at Naturmagten ikke blot bliver en \emph{overmenneskelig} Magt, ɔ: en
Magt, høiere i Grad end den menneskelige, men i en vis Forstand en
\emph{overnaturlig} Magt ɔ: en Magt, der i sin Hypostaserethed formelt er løsnet
fra det reale Naturlige og ved Phantasien er befriet fra Naturens, ɔ: fra \emph{sine egne} Skranker; en Magt, der, gjennem Phantasien bestemt ved Frygten og Haabet,
transscenderer sig selv, ophæver sin egen Lovmæssighed --- det vil, paa dette
Standpunkt, sige: den Lovmæssighed, Bevidstheden her opfatter: det Sædvanliges
relative Lovmæssighed. Fra dette subjectiverede Transscendente sees ogsaa i
Naturreligionerne en Aabenbaring som udgaaende, enten gjennem det Religiøses
Rørelser i de enkelte Individer, eller gjennem særegne Individer som Medier, eller
gjennem det Guddommeliges supponerede Tegn og Virkninger.

Naar saaledes Religionen som \emph{positiv, historisk} Religion bliver liig med \emph{aabenbaret} Religion i den ovenfor angivne Betydning, og constitueres ved
Forholdet gjennem et Aabenbaringsindhold til dettes transscendent-supranaturale
Kilde, saa faaer herved \emph{Troen} som \emph{religiøs} \opage{40}\emph{Tro} sin nærmere Bestemmelse. Troen som
religiøs Tro bliver saaledes Anerkjendelsesforholdet --- det theoretiske og
praktiske --- til Aabenbaringen, og det deri indesluttede Anerkjendelsesforhold
til Aabenbaringens Subject.

I denne Bestemmelse af \emph{Troens} Væsen, der opfatter det fra de positive
Religioners Synspunkt, ligger implicite, hvorledes denne Troes Forhold til Viden
og til den paa Viden hvilende Handling maa bestemmes fra Religionernes og fra
Videnskabens Synspunkt. Idet den paa den positive Religions Trin staaende
religiøse Bevidsthed opfatter sit Indhold som \emph{guddommelig Aabenbaring}, maa
den, naar den skal bestemme Forholdet --- som vi ovenfor have seet med Hensyn til
Christendommen --- sætte Troesindholdet som incommensurabelt med den menneskelige
Videns Evne, og denne Incommensurabilitet faaer sit bestemtere Præg ved at føres
tilbage til \emph{Modsætningen} mellem Uendeligt og Endeligt, mellem Overnaturligt og
Naturligt, mellem guddommelig Visdom og menneskelig Kløgt, mellem den hellige
Aands Lære og den af Synden fordærvede Menneskeaands Paafund. Den ved den
religiøse Bevidstheds Form umiddelbart satte \emph{Sondring} mellem det Guddommelige og
det Menneskelige maa nemlig, idet Reflexionen, bestemt af Forestillingsformen,
opfatter den, faae \emph{Modsætningens} Skikkelse, og idet nu Troen faaer en guddommelig
Gjenstand, og tillige --- da Gjenstanden er Aabenbaringssubject og Aabenbaringen
det Troen Betingende --- et guddommeligt Princip og Udspring, maa den
nødvendigviis, naar Forholdet skal bestemmes, stille sig i et Modsætningsforhold
til Viden med dens endelige Indhold og naturlige Princip. Den Transscendentsens
og Enkeltpersonlighedens Form, i hvilken det Guddommelige opfattes, indeslutter
paa een Gang, at Aabenbaringsindholdet, som havende det Guddommelige til
Gjenstand, ikke kan blive Menneskene meddelt uden ifølge en fri guddommelig
skaberisk Beslutning, og at det maa være et Høiere end den menneskelige tanke.
Mellem Tro og Viden danner der sig saaledes en, ifølge Skilsmissen mellem det
Guddommelige og det Menneskelige \opage{41}uophævelig, Modsætning med Hensyn til
Erkjendelsen af det guddommelige Væsen, men denne Modsætning omfatter tillige
Modsætningen mellem den aabenbarede Religions supranaturale Handlingslov og den
Handlingens Lov, der bestemmes ved den naturlige Erkjenden af Menneskeaandens
Væsen og væsentlige Forhold.

Differentsen mellem Tro og Viden vilde ogsaa i Realiteten være bestemt som en
Modsætning naar Forholdet mellem Aabenbaringen og Fornuften bestemtes ved den
tilsyneladende quantitative Forskjel, at Aabenbaringens Indhold skulde have et
videre Omfang end Videns, skulde overstige Viden, og \emph{blivende} være et
Overfornuftigt, uden at det dog endnu udtrykkelig bestemtes som et Fornuften
Modsigende. Hvad der nemlig skiller sig fra Viden og indenfor Erkjendelsens
Sphære --- og saaledes opfattes Forholdet her --- vil begrændse Viden, bliver
Videns Modsætning og ophæver i Virkeligheden Viden. Thi Viden er kun Viden som
det efter sit Begreb i sig Afsluttede, der rummer Tilværelsens Heelhed i sig,
dens Princip er et universelt Princip. Sattes et mere Omfattende end Videns
Sphære, saa maatte denne, som gaaende ind i den mere omfattende Sammenkjædning
--- de „guddommelige Sandheders“ uendelige Sammenkjædning --- bestemmes som
organisk Led i hint; men deraf vilde følge, at den først, naar den opfattedes som
saadant, blev opfattet efter sin Sandhed, og at Videns Opfattelse af sit Indhold
vilde være falsk, fordi den opfatter det, der kun er Led, som Heelhed, og
anderledes kan den ikke erkjende det, fordi Aabenbaringens overfornuftige Indhold
ikke kan blive Videns-Indhold. Kun Troens Erkjendelse vilde da være den sande
Erkjendelse, Videns vilde ikke blot være ufuldstændig, men væsentlig usand.

Opfattet fra den troende Bevidstheds Synspunkt vil saaledes Forholdet mellem Tro
og Viden vise sig som et Modsætningsforhold, men paa samme Maade maa Forholdet
bestemmes, naar det sees fra Videns Synspunkt. Idet Videnskaben erkjender den ved
de positive Religioners Forhold til Forestillingen betingede
Endelighedsbegrænds\opage{42}ning i deres Opfattelse af det Guddommelige og dets Meddelelse
til Aanden, og vindicerer sig selv som sin specifiske Form den Uendeligheden
hævdende, begribende Erkjenden, sætter den Modsætningen som en Modsætning mellem
Troens „Erkjendelse“ og Videns Erkjendelse, idet, ogsaa naar det Religiøse
væsentligen opfattes fra Villiens Side, Villiesformen sees som staaende i en
uopløselig Forbindelse med en bestemt Erkjendelsesform, saa at Erkjendelsens
Moment altsaa ogsaa fra dette Synspunkt bliver væsentligt.

De givne Begrebsbestemmelser føre os altsaa til det \emph{foreløbige} Resultat med
Hensyn til Forholdet mellem Troen i den positive Religions Betydning og Viden, at
Differentsen, saavel fra Troens som fra Videns Synspunkt, maa bestemmes som en
væsentlig og uophævelig Modsætning, der maa føre til Nødvendigheden af et Valg
mellem dem og at Enheden af Viden og Tro kun bliver mulig, naar
Endelighedsbegrændsningen fjernes fra Troen og Religionen, og disse opfattes
efter deres ideale Væsensbestemmelse\footnote{Hvorledes Forholdet mellem
Videnskaben og Religionen efter dennes ideale Begreb, mellem Viden og den til
hint Religionsbegreb svarende Tro, er at opfatte, vil blive udviklet i det Skrift,
der skal slutte sig til dette.}.

En fra denne Bestemmelse af Forholdet afgivende gives der nu paa de to
Standpunkter, der særlig ville blive Gjenstand for en kritisk Prøvelse. Fra \emph{Theologiens} Side bliver, medens en positiv Religion fastholdes som Sandhedens
absolute Form, gjennem Fordringen paa Muligheden af Troens og Videns Forening i
en Troesvidenskab, Differentsen sat som en blot quantitativ. Fra det Standpunkt,
der sondrer Modsætningerne for gjennem Sondringen at forsone dem, indrømmes
Modsætningen mellem Tro og Viden, men imod den Paastand, at et Valg imellem dem
er nødvendigt fordi Modsætningens Art umuliggjør Foreningen af dem i samme
Bevidsthed, sættes Paastanden paa en saadan Forenings Mulighed, medens der
tillige gjøres \opage{43}Indsigelse imod, at Religionens Sandhed bedømmes fra Videns
Synspunkt eller underkastes Videnskabens Dom, idet det hævdes, at dens Sandhed
hviler paa et andet Princip end Videns og følger sin egen, høiere Lov.

Gjennem Kritiken af disse to Standpunkters Besvarelse af Problemet vil da den
ovenfor givne Bestemmelses Gyldighed blive at godtgjøre.


%% ============================================================
%%  THEOLOGIENS BESVARELSE AF PROBLEMET (pp. 43--118)
%% ============================================================
\chapter{Theologiens Besvarelse af Problemet}
\label{ch:theologi}

\section{Theologiens Opgave som Troesvidenskab}
\label{sec:troesvidenskab}

Theologiens almindelige Stilling til Problemet --- en Stilling, der her altsaa
sondres bestemt fra den umiddelbare religiøse Bevidstheds --- bliver nærmere at
bestemme paa følgende Maade. Theologien som Troesvidenskab bliver netop Udtrykket
for en væsentlig Høren sammen af Tro og Erkjendelse. Den begrændser vel ikke
Troen til Erkjendelse, men sætter Erkjendelsen som væsentlig Bestemmelse i Troen.
Den nægter ikke Differentsen mellem Tro og Viden, baade naar Viden opfattes som
naturlig Viden og naar den opfattes som den ved Troen bestemte Viden, der vil
begribe Aabenbaringen, men den accentuerer Muligheden og Nødvendigheden af at
give Troen et Udtryk i Videnskabens Form. Den faaer altsaa den eiendommelige
Opgave: idet den sætter Aabenbaringen som et transscendent Væsens frie Meddelelse af det, som ingen menneskelig Kløgt kunde udfinde og som ingen kan opløse i Begreber, og
fordrer den menneskelige Videns Opgiven af sin Autonomie til Fordeel for dens
Autoritet, at bestemme Aabenbaringen og Guds Væsen og Virken saaledes, at de
drages ud af det rent Ubegribelige, og paa den anden Side at lade Videns Begreb
saaledes blive omformet, det være sig ved en „Potentsation“ eller ved en
Udvidelse, at det kan give Plads \opage{44}for det, der dog aldrig kan eller skal i egentlig
Forstand begribes, og muliggjøre en Videnskab, der følger en anden Lov end den
autonome Videns, ogsaa omfatter det for denne Viden Utilegnelige og dog ikke
ophører at være Videnskab.

Muligheden af denne Opgaves Løsning skal Kritiken af Theologien undersøge. Denne
Kritik vil, som tidligere angivet, nøie knytte sig til et bestemt theologisk
Forsøg, der i vor Literatur har faaet en særegen Betydning og hos os kan betragtes
som den egentlige Repræsentant for den theologiske Bestræbelse; men idet Kritiken
paa det Nøieste knytter sig til dette, vil den overalt føre dets Behandling af de
afgjørende Punkter tilbage til det for Theologien som Theologie Typiske og paavise
den som udtrykkende dette, saaledes at Kritiken gjennem det enkelte theologiske
Forsøg forholder sig til Theologien overhovedet.

Vi fremstille først i Sammenhæng denne theologiske Theorie, der af Biskop \emph{Martensen} er udviklet i hans forskjellige Skrifter, i dens væsentlige Grundtræk,
for derved at angive Tilknytningspunkterne for Kritiken.

Troesvidenskabens Opgave er den: at forene Tro og Viden i organisk Enhed,
saaledes at paa den ene Side Troen hævdes i sit Principat, paa den anden Side dens
Væsenssammenhæng med Viden fremtræder. Troen har Principatet fordi den, som det
Umiddelbare, er det Første, medens Viden, som det Middelbare, er det Afledede, og
fordi Troen er langt rigere end Viden: et nyt Liv i Sjælen, et Existentsforhold,
og aldrig kan udtømmes af den begribende Viden; men dens Enhed med Viden maa
hævdes fordi Troen, der i sig omfatter Villie, Følelse og Erkjendelse, selv er en Viden: den centrale Viden, saa at Viden bliver Fælledsbestemmelse for Videnskab og
Tro. Netop dette Forhold gjør en „christelig“, „religiøs“, „troende“ Videnskab
mulig. At Troen har Videns Bestemmelse i sig, fremgaaer deraf, at Gud har \emph{aabenbaret} sig for Troen, og at hans Kjærlighed og Naade har givet sig tilkjende
under \emph{Viisdommens} Form, ligesom ogsaa deraf, at den Troende, i Kraft af den
christelige \emph{Sandhedsidee}, der er \opage{45}indeholdt i den levende Tro, maa kunne sige: jeg
veed, hvad jeg troer. Var Troen ikke allerede i sig selv en Viden, da var den ikke
\emph{Sandheds} Tro, og den store Modsætning mellem Sandhed og Løgn, der udtales
paa hver Side af den hellige Skrift, vilde være uden Betydning. Var Troen ikke
selv en Viden, da kunde den Troende ikke bekjende sin Tro og kunde ikke forkaste
og bekæmpe de vrange og vildfarende Lærdomme. Denne Viden i selve Religionen er
den, der betegnedes som den primitive, den centrale Viden, og, hvor den
fremtræder med særlig Kraftighed, kan den betegnes som den centrale Anskuelse, der
i sig har Driften til at udfolde sig til Begreb. Herfra udgaaer den
videnskabelige, den reflecterende og speculerende (i Ideespeilet skuende og
Muligheder søgende) Erkjendelse.

Ifølge dette Forhold mellem Tro og Viden maa Troesvidenskaben staae i et
Afhængighedsforhold, men i et frit Afhængighedsforhold, til Aabenbaringen.
Christendommens absolute Sandhed maa anerkjendes som given forud for og uafhængig
af al Speculation; Erkjendelsen skal være en Erkjendelse i Troen og udaf Troen
(\textit{credo ut intelligam}). Den hellige Viisdomstanke maa i den troende
Bevidsthed blive Tænkningens Princip, og den menneskelige Tænkning formaaer kun
ved den at ransage Aabenbaringens Dybheder, at granske efter Sammenhæng og Grund
i de christelige Forestillinger, at stræbe efter at frembringe et aandeligt
Gjenbillede af den evige Aabenbaringsviisdom.

Den christelige Troesvidenskab vil fremstille den christelige Anskuelse som et i
sig sammenhængende Lærebegreb. Den nøies ikke med en explicativ Begriben, men
søger, ifølge en indre Drift, en speculativ, der ikke bliver staaende ved at
fremstille Sammenhængen i det Givne, men spørger om Mulighed og Grund. Ethvert
«\textit{Ita}» indeholder et skjult «\textit{Quare}», der under den grundige
Explication ikke kan Andet end komme frem og opfordre til en høiere speculativ
Begriben. Dog kan Troesvidenskaben, da dens Gjenstand ikke er en blot
Fornuftsandhed, men en positiv Sandhed, en Troessandhed, ikke udvikle sit Indhold
eller demonstrere \opage{46}sine Sandheder med den mathematiske og logiske Nødvendighed,
der er eiendommelig for den blot rationelle Videnskab: der bliver en Grændse, vel
ikke for Visheden, men for Klarheden, men dette er ikke en Mangel ved
Troesvidenskaben; det viser netop hen til dens høiere Værd.

Dette bliver klart, baade naar vi see hen til hvad der er denne Videnskabs
egentlige Organ i Menneskeaanden, og naar vi see hen til dens objective Princip,
til det Gudsbegreb, der bestemmer den.

Hvad det første Punkt angaaer, saa bestemmes Troen som en Villiens Tilegnelse i
Kjærlighed af den guddommelige Aabenbaring. Villien er altsaa den Menneskeaandens
Function, til hvilken den har sit væsentlige Forhold, men Villien er, ligesom den
i Universet er det Første, det egentlig Constitutive i den menneskelige
Personlighed: Menneskets Grundværen er Villie. Men saaledes bliver Villien det
Høieste, det, der har Principatet i Forhold til de andre Aandsfunctioner, der ere
afhængige af den. At det, der i Troen gribes gjennem Villien, ikke er
commensurabelt med Tanken, med Begrebet, viser derfor kun, at hint er det Høiere
og Rigere.

Og til det samme Resultat føres vi, naar vi betragte Troesvidenskabens objective
Princip. Troesvidenskabens Princip er et høiere end det, til hvilket den blot
logiske Tænkning naaer; det er ikke „Fornuftens døde Afgud“, men et ethisk Princip, et virkeligt Personlighedsprincip, et overnaturligt Princip. Det ethiske
Gudsbegreb er det høiere, der ikke udelukker, men indeslutter det logiske og
physiske. Det er Begrebet om den frie, villende, personlige Gud, om den ikke blot
logiske, men hellige, absolute Sandhed. Den personlige Gud, der ikke er den
ideebundne, men den ideefrie, ligesom han er den naturfrie, aabenbarer som
skabende og undergjørende sin Frihed og Almagt, og i hans Aabenbaring gjennem
Historien gjør den samme Personlighedens og Villiens Frihed sig gjældende, i
Modsætning til Naturens og det blot Logiskes Nødvendighed, og den viser sig i sin
høieste Form, hvor Gud ikke blot aabenbarer sig gjennem \opage{47}sine Virkninger, men
aabenbarer sig selv i sin personlige Fremtræden i den historiske Christus.

Med dette Gudsbegreb indtræder et andet Videnskabsbegreb end der, hvor der kun er
Tale om en blot logisk og physisk Viden: der naaes til Begrebet om en positiv
Philosophie, til Begrebet om en Aabenbaringsvidenskab, i hvilken Theologien
bliver den centrale Videnskab. Om Aabenbaringsvidenskaben overhovedet gjælder det,
der ovenfor blev sagt om Theologien, at den ikke demonstrerer sine Sandheder med
logisk og geometrisk Nødvendighed, og at den kun kan bringe det til en relativ
Begriben af de høieste Sandheder; men dette er netop Særkjendet for den høiere og
virkelige Videnskab, for den høiere Erkjendelse af Virkeligheden.
Aabenbaringsvidenskaben og særlig Theologien maa, efter Sagens Natur, være rigere
i Anskuelsen end i Begrebet, men det er Anskuelsen, der giver Dialektiken dens
Værd; thi Anskuelsen er Organet for Heelheden og rækker længere og dybere ind i
selve Tingenes Inderste, til deres første Begyndelser og Rødder, end den blot
dianoetiske Tænkning. Troesvidenskaben er derfor Videnskab i fortrinlig Forstand,
medens Naturvidenskaberne paa ingen Maade i strængere Forstand kunne kaldes
Videnskaber, og medens den blotte Fornuftphilosophie, den negative Philosophie,
kun er propædeutisk, en Børnelærdom i Sammenligning med hiin høiere Viisdom.
Dennes eiendommelige videnskabelige Charakteer og høiere Værd ligger netop i
Forholdet til den frie, personlige Gud og hans Aabenbaring, til den hellige
Villie, der er Maalestokken for hvad Tingene i sidste Betydning ere. Ligesom
Aabenbaringen er fremgaaet af den guddommelige Villies Frihed, ikke behersket af
en abstract logisk eller physisk Nødvendighed, saaledes er ogsaa dens Opfattelse
betinget med en fri Villiesact, der bestemmer Tanken, ikke bestemmes af den. Den
guddommelige Aabenbaring kan, som en Frihedsaabenbaring, en
Personlighedsaabenbaring, ikke opfattes i den logiske Nødvendigheds Former, ikke
opløses i blotte Begrebsbestemmelser; den kan, som Livsmeddelelse, ikke optages
gjennem Tankens Former \opage{48}alene. I al Frihedsaabenbaring er der et med den rent
logiske Tænkning Incommensurabelt, derfor savner denne Tænkning uundgaaelig det
Blik for det Historiske, der charakteriserer den høiere Videnskab. Men naar der i
Erkjendelsen af Frihedsaabenbaringen gaaes ud over det blot Logiskes Gebet, saa
føres vi dog ikke ud paa det Lovløses Gebet, thi den guddommelige Frihed er ingen
grundløs Vilkaarlighed; den indeholder en høiere Nødvendighed, men en saadan, der
kun lader sig frivillig erkjende. Den frie guddommelige Almagt er, som skabende og
undergjørende, bestemt ved Kjærligheden og Viisdommen, altsaa ved en ethisk
Frihedens Nødvendighed, og Nødvendigheden i almindelig Betydning faaer atter sin
relative Gyldighed, idet et Physisk bliver Betingelse for den ethiske Villie i
Gud, og idet Gud vel er naturfri, men ikke naturløs, idet der i Gud er en evig
Physis, der staaer i væsentligt Forhold til Villien. I dette Forhold er der ogsaa
givet et Tilknytningspunkt for Tanken, ved hvilket den aandelige Guds Skaben af en
materiel Natur, og hans umiddelbare miraculøse Virken, selv om Skabelsens og
Underets Hvorledes evigt unddrager sig den menneskelige Begriben, dog hæves over
det rent Ubegribelige, i et vist Maal bliver Gjenstand for en Erkjenden.

Saaledes er altsaa Troesvidenskabens Mulighed grundet i en høiere, ved Tro og
Villie bestemt Erkjendelse, en Erkjendelse, for hvilken Aabenbaringens Sandhed
oplader sig saavidt det Guddommelige kan fattes af Menneskeaandens Skrøbelighed,
medens den Viden, der ved Menneskeaandens egen Kraft og egne Midler, ved den
blotte Fornuft, vil begribe den guddommelige Aabenbaring, maa forvanske og
fornægte dennes Sandhed. Den høiere Erkjendelse er Fuldendelsen, Forløsningen af
Menneskeslægtens almindelige Fornufterkjendelse, en høiere Potents af Fornuftens
Aabenbarelse. Aabenbaringen i Christus ophæver Fornuftlovene i deres Abstraction
forat stadfæste dem i Christus-Viisdommen, der er Fornuftlovens Fylde. Og den
høiere Erkjendelse indeslutter i sig Normen og Correctivet for den lavere; ved den
er den høiere Videnskab den herskende Videnskab, \opage{49}og saaledes tilveiebringes det
rette Forhold i Tankens Verden: οὐκ ἀγαθὸν πολυκοιρανίη, εἷς κοίρανος ἔστω.

\section{Villiens Forhold til Viden og til Væsenet i det Menneskelige}
\label{sec:villie-viden}

Hvad der i denne Theorie er det Berettigede, det, der ubetinget har Krav paa
Anerkjendelse, er Bestræbelsen for at bringe den menneskelige Aands høieste
Yttringsformer, dens væsentligste Grundforhold, i indre Harmonie. Det Maal,
hvorefter der stræbes, er Aandens høieste: Aandens sande Forsoning. Om derfor end
de Midler, ved hvilke det Tilstræbte søges naaet, ere utilstrækkelige, om Veien,
der følges, ikke kan føre til Maalet, saa vil denne Theorie dog bestandig kunne
gjøre Fordring paa en agtelsesfuld, fordomsfri og indtrængende Kritik, saameget
mere som den fremtræder i en Form, der ikke blot er fiin og aandfuld, men bærer
den personlige Overbeviisnings umiskjendelige Præg.

\begin{center}
---
\end{center}

Kritiken af Theorien maa fra Begyndelsen af rettes mod et Punkt, der er den hele
Theories videnskabelige Tyngdepunkt; mod dens Bestemmelse af \emph{Villiens Forhold til Viden og til Væsenet i det Menneskelige og i det Guddommelige}.

Paa denne Bestemmelse, ifølge hvilken Villien sættes som det Primitive, det
Væsensbegrundende i den menneskelige og i den guddommelige Personlighed, og som i
sig indesluttende Viden som behersket Moment, hviler nemlig paa den ene Side
Opfattelsen af en Videnskab, der gaaer ud over den blot logiske Videns Grændse, og
dog ikke blot, ved det Vidensmoment, Villien indeslutter, hævder sig Videnskabens
Charakteer, men netop ved sin Overskriden af hiin Grændse faaer et høiere Værd;
--- og paa den anden Side hviler derpaa Opfattelsen af Gud som i sin frie Villen
hævet over Videns abstracte Nødvendighed og over den væsensbestemte Frihed, der
kun er en \textit{libera necessitas}, og dog som i sin Væren og Virken blivende
Gjenstand for en relativ Viden. Naar Villiens Principat hævdes, saa er \opage{50}Videns
Nødvendighed behersket, dens bindende Lov underkastet Subjectets uendelige Frihed
og dets individuelle Trang, saa er Erkjendelsesforholdet afhængigt af et
Villiesforhold og Villien er berettiget til, som vælgende Sandheden, at corrigere
Erkjendelsens Consequentser. Naar Villien er sat som det Primitive i Gud, saa er
Gud den skabende og undergjørende, den frit sig aabenbarende, og med dette
Gudsbegreb er Begrebet om den fordrede høiere Videnskab, om dens Tilbliven gjennem
en fri Meddelelse fra dens Princip, umiddelbart givet.

Udgangspunktet for den kritiske Undersøgelse maa tages i det \emph{Psychologiske}, hvis
Opfattelse overhovedet bliver afgjørende for Bestemmelsen af Troens og Videns
Forhold. I Opfattelsen af det Psychologiske knytter Udviklingen hos Martensen sig
paa det Nøieste til enkelte Philosopher, der staae Theologien nær. Det er
navnlig den lidet klare Tænker, J.\ H.\ Fichte, fra hvem Bestemmelsen af de
psychologiske Grundforhold hentes, og hvad han har fremsat i Indledningen til sin
endnu ufuldendte Psychologie, gjengives tildeels med hans egne Ord. Det
siges\footnote{Martensen: Om Tro og Viden S.\ 45 f.\ 123 f.\ 129.} at være
Resultatet af den dybere psychologiske Iagttagelse, et Resultat, i hvilket J.\
Bøhme, Fr.\ Baader, Schelling og den yngre Fichte stemme overens, at ikke Tanken
eller Fornuften, men Villien er det Første, Inderste, Oprindeligste og Centraleste
i Menneskesjælen, at baade vor theoretiske og vor praktiske Bevidsthed har sin
psychologiske Rod i en forbevidst Villie, som i disse tvende Hovedformer
manifesterer sig, at altsaa Selvbevidstheden og Tanken selv er en Form af Villien,
nemlig Speilingsformen, i hvilket det villende Væsen objectiverer sig, sig selv og
sin Verden. Men idet Selvbevidsthed og Tænken i sig selv ikke ere Andet end den
blotte Speilingens Act, saa forudsætte de et realt villende Subject, som deri
speiler sig. Dette reale Subject, hvis Grundværen er Villie, er det Substantielle,
Bevidstheden er kun Egenskab, Prædicat, Attribut for det: den er det, hvorigjennem
% PRINTED AS IS, p.51 [collation 2026-09-21]: the page prints "om sig om sine Forhold" (om for og); both readers saw it.
det \opage{51}reale Væsen kommer til Klarhed om sig om sine Forhold, det, gjennem hvilket den
forbevidste, potentielle Villie actualiserer sig. --- Al Iagttagelse skal ogsaa
føre os tilbage til Begjæring, Drift og Instinct som det første Psychologiske,
altsaa til en dunkel Villie, der indeslutter Fornuft og Erkjendelse, som det,
hvoraf Bevidstheden udvikler sig. Og spørge vi om Erkjendelsens og Villiens
Grundforhold i den virkelige Bevidsthed, saa er det ikke heller Erkjendelsen, men
den ethiske Villie, der har Primatet; ikke den ethiske Villie som en blot
praktisk, erkjendelsesløs Villie, men den ethiske Villie, der har Erkjendelsen,
det Contemplative, som sit eget Moment. Og som i Mennesket Villien er det Første,
ikke Tanken, saaledes er den det ogsaa i den objective Tilværelse, i Universet, og
saaledes er, hvad netop Theismen udvikler, Villie, ikke Viden, det Første i Gud, og
den ethiske Idee, det Gode, som er Villiens Idee, er den høieste, den absolute
Idee, hvilken Intet er coordineret, men Alt subordineret.

Det Uholdbare i den refererede psychologiske Theorie viser sig saasnart
Hovedpunkterne i den nærmere belyses.

1.\ Idet Villien sættes som det Primitive, gaaes der tilbage til Begrebet om en
„forbevidst Villie“ (Fichte's „vorbewußter Wille“), der skal være det, af hvilket
Aandslivets forskjellige Former udvikle sig. Men hvad der kaldes „forbevidst
Villie“, er, netop som forbevidst, ɔ: som endnu ikke bevidst, ikke Villie; thi
Villie er --- hvad Martensen selv indrømmer ved at lade Villien actualiseres
gjennem Selvbevidstheden --- kun som bevidst Villie. Derfor skal hiin Villie ogsaa
kun være en „dunkel“ Villie. Men naar Villiesbegrebet i strængere Forstand bliver
uanvendeligt, maa det Spørgsmaal opkastes: hvad det er i Bestemmelsen Villie, der
bliver anvendeligt paa hiin primitive Virken. Svaret maa blive: det er Energiens
Bestemmelse; men Energien som et Subjects, et Selvs Energie betegne vi som
Selvheds-Energie. Energien er ikke et Prædicat ved et forudsat Subject, thi
saaledes opfattet vilde Subjectet blive det Energieløse, dets Forbindelse med
Prædicatet modsigende; \opage{52}Energien er Væsensbestemmelse, Udtryk for selve Væsenet som
ideelt-realt\footnote{Det „reale, villende Væsen“ er hos Martensen og Fichte en
modsigende Reminiscents fra Herbarts uforanderlige, uvirksomme og ubevidste Reale.
„Realitet“ bliver her for Martensen en betydningsløs Bestemmelse; den kunde
ialfald betegne det Monadiske, men derved vilde, i denne Sammenhæng, Intet være
% PRINTED AS IS, p.52 [collation 2026-09-21]: the page prints "Selvhedsenergieu" (turned n); both readers saw it at 600 dpi.
naaet.}. Men Væsensenergien, Selvhedsenergieu, træder først frem gjennem den
organiske Formning; gjennem denne udformer den Betingelserne for Villien, men
ligesaafuldt og ligesaa oprindeligt Betingelserne for Erkjendelsen:
Selvhedsenergien staaer derfor i sin første Fremtrædelsesform ikke hiin nærmere end
denne, og idet Villien og Erkjendelsen virkeliggjøres, bliver Energien
Bestemmelsen for begge, det begge omfattende Begreb, og vi tale med samme Ret om en
Erkjendelsesenergie som om en Villiesenergie; i Erkjendelse og Villie virkeliggjør
Energien sin indre Mulighed i to hinanden supplerende Grundformer. At man har
overført Villiens Bestemmelse paa Energien og betegnet den som forbevidst Villie,
ogsaa paa det Stadium, hvor Villien kun vilde være physisk (organisk) virkende,
hidrører fra en Misforstaaelse af Villiens Forhold til Energie og til real Virken.
Medens man paa den ene Side overseer Erkjendelsens Bestemmelse som Energie og
Erkjendelsesvirksomhedens ubevidste Indvirken paa det Reale i Organismen, overseer
man paa den anden Side, at idet Villien som Villie udøver en bevidst Indvirkning
paa det Reale, saa forudsætter denne dens Virksomhed Selvhedsenergiens
organdannende Virksomhed, der vedvarer ogsaa efterat Villien som Villie er kommen
tilsyne, er uafhængig af Villien, og netop herved bestemt skiller sig fra denne.
Idet man kalder den oprindelige Energie Villie, skyder man Forestillingen ind om
et ligesom bag vort Væsen virksomt Væsen, hvis Virksomhed forestilles i Analogie
med den bevidste Villies Indvirkning paa et Ydre gjennem Organismen.

2.\ Finder der saaledes en Begrebsforvexling Sted i Be\opage{53}stemmelsen af det i
\emph{Væsenet} Primitive, saa lider Bestemmelsen af det, der \emph{for Iagttagelsen} er det psychologisk Første, af en lignende Forvexling. Det for Iagttagelsen Første
bliver, naar man derved forstaaer de første psychologiske Functioner, der gaae ud
over den blot organiske Formen, ikke Begjæring og Drift. Instinct kunde man til
Nød bruge som Betegnelse for dette Første, men Instinctet vilde da, ligesom i
Dyrelivet, fra hvilket Betegnelsen er tagen, omfatte Intelligentsens og Villiens
rudimentæreste Virken, de ubevidste og halvbevidste Erkjendelsesfunctioner, og de
ubevidste og halvbevidste Villiesfunctioner. Det første Psychologiske er det,
hvortil Begjæring og Drift vise tilbage. Begjæring og Drift er en Stræben efter at
fastholde eller fjerne en bestemt Tilstand; men denne Tilstand er kun en Selvets
Tilstand som fornummet, gjennem den organiske Fornemmelse eller Sandsefornemmelsen.
Det psychologisk Første er altsaa Fornemmelsen: Erkjendelsens rudimentære Form,
ligesom Driften er Villiens rudimentære Form. Dette Forhold indrømmes forsaavidt af
Martensen, som det siges, at det i Reflexionsforholdet ɔ: i Forholdet til det
Objective, skal være rigtigt, at sætte Villien som det Andet, Erkjendelsen som det
Første. Men i Forholdet til selve Subjectet skal Villien være det Første. Denne
Distinction kan imidlertid ikke fastholdes. Selvets Forhold til sig selv er
bestandig et Forhold til sig som Led i en Tilværelse, som bestemt ved det
Objective. Selv i Fornemmelsens primitiveste Form som organisk Fornemmelse, i
hvilken Selvet forholder sig til sin egen Realitetsside, er Forholdet til det
Objective, der sættes som saadant i Sandsefornemmelsen, implicite tilstede. Men
den organiske Fornemmelse er den første Form for Selvets Forsigværen.

3.\ Den Opfattelse af Selvbevidsthed og Tanke som en blot Afspeiling, der gjøres
gjældende af Martensen, er ugjennemførlig. De skulle være det reale, villende
Væsens Speiling, i hvilken det objectiverer sig selv. Men Speiling udtrykker
Reflexion. I Bevidsthedens Speiling er det reale Væsen sat som reflecteret i sig.
Men en saadan Reflexion \opage{54}er det absolut Uforklarlige naar ikke Reflexionen
oprindelig er sat som Væsensbestemmelse. Uden dette vilde Væsenet ikke blive \emph{sig} klart i Reflexionen, ikke objectivere sig \emph{for sig}. Reflexionen vilde blot blive en
Reflexion i et Andet, opfattet af en supponeret udenforstaaende Bevidsthed.
Væsenet vilde blive det for sig Dunkle trods det tiltrædende Prædicat, der som
dette Subjects Prædicat vilde blive et rent Ubegribeligt, og den Villen, der
udtrykte Væsenet, vilde, naar Reflexionens fra den virkelige Villie uadskillelige
Bestemmelse borttoges, kun være en Virken i Lighed med det ene Materielles
mechaniske Indvirken paa et andet Materielt. Sættes derimod Reflexionen som
oprindelig Væsensbestemmelse, saa er Væsenet oprindelig bestemt ved Tænkning ɔ:
ubevidst Tænkning, ved det almene Selvs Forsigværen. Hiin Opfattelse af Tanke og
Bevidsthed som det oprindelig reflexionsløse Væsens Speiling i sig vilde ligeledes
gjøre det absolut uforklarligt, hvorledes en Bevidsthed om Andet og en Begriben af
dette Andet kan finde Sted. Kun hvor Reflexionen bliver oprindelig Bestemmelse, er
i Selvets Fordobbling, i den derved satte Andethedsbestemmelse indenfor Selvet,
Tilknytningen given til et Forhold til Andetheden udenfor Selvet.

4.\ Ved den Bestemmelse, Martensen giver af det psychologiske Grundforhold, kan
han ikke undgaae at indvikle sig i Modsigelser. Det er en Modsigelse, naar der
tales om en Fælledsrod for Villien og Erkjendelsen, naar de altsaa coordineres, og
Villien dog sættes som den høiere Enhed. Det er ligeledes en Modsigelse, naar
Villien, der skal være det Primitive, som Villie skal blive actualiseret ved
Selvbevidstheden; thi det, ved hvilket Actualiseringen skeer, maa være det Første,
ligesom det maa være et Substantielt. Det er en Modsigelse, at Villien, ikke
Fornuften, skal være det Første; thi saaledes skulde Villien som fornuftløs
begrunde Fornuften. Det er endelig en Modsigelse, at Martensen udtaler, at Viden
og Villie ikke kunne adskilles, at den ene ikke kan være uden den anden (S.\ 40),
og at han overalt urgerer en \emph{væsentlig} Betydning for det Theore\opage{55}tiske ved
Siden af det Praktiske, men dog vil berøve Tanke og Selvbevidsthed
Substantialiteten og reducere dem til blot at være Egenskab og Afspeiling.

Imod de af Martensen givne psychologiske Bestemmelser stille vi da, støttede paa
de foregaaende Udviklinger, følgende:

Det menneskelige Væsens Grundbestemmelse er en Selvheds-Energie, der aabenbarer
sig i Viden og Villie.

Viden og Villie staae begge i det Forhold til Væsenet, til det Substantielle, at
de ikke kunne bestemmes som Egenskaber ved et forudsat „realt“ Væsen, men maae
opfattes som Væreformer for det subjective \emph{Væsen}, som substantielle \emph{Væsensformer}.

I det menneskelige Sjælelivs Udvikling er det psychologisk Første Erkjendelsens
første rudimentære Form.

Ifølge det, der betegner det menneskelige Væsens Eiendommelighed i Forhold til
Dyrets: Væsenets Universalitet, bliver den ved Universaliteten betingede subjective
Forsigværens Form, der, bestemt ved det Almenes Forhold til sig, træder frem i
Selvbevidsthedens og Tankens Energie, den efter Begrebet principale Væreform for
Menneskevæsenet.

Den Betydning og Magt, der af Martensen tillægges Villien, ifølge dens statuerede
Principat, i Forhold til Erkjendelsen (som beherskende Tænkningen og som vælgende
Sandheden), kan ikke fastholdes; den beror væsentlig paa en Begrebsforvexling,
deels af Villien og Erkjendelses\emph{energien}, deels af Villien og det som Selv ved \emph{Energie} bestemte menneskelige Væsen.

Af disse psychologiske Bestemmelser vil det saa fremgaae som Consequents, at der
ikke kan tilkomme Villien, idet den bestemmer sig som ethisk, en Autonomie og et
Primat i den Betydning, i hvilken Martensen tillægger den det. Villien bliver
først ethisk Villie som bestemt ved Erkjendelsen: Erkjendelsen af det Gode er
Betingelsen for at ville det\footnote{jfr.\ Martensens Udsagn (om Tro og Viden S.\
26): det værende Gode er Betingelsen for vor Villen af det Gode, for vor Stræben
efter det.}.

\opage{56}Den ethiske Villie er en \emph{Væsensvillen}, men Væsenet er principalt det til
Selvbevidsthed anlagte, efter sit Begreb tænkende Væsen: en Væsensvillen er derfor
bestemt ved Erkjendelse, ved selvbevidst Tanke.

\section{Villie og Viden i Gud}
\label{sec:villie-gud}

Det samme Grundforhold, der statueredes med Hensyn til det menneskelige Væsen,
bliver, som allerede anført, ogsaa statueret med Hensyn til det guddommelige. Det
Gudsbegreb, til hvilket Philosophien naaer gjennem Erkjendelsen af Naturen og
Fornuften, og som udtrykker Fornuftens og den tankebestemte Virkeligheds absolute
Princip, skal ikke, selv om det bestemmer det Absolute som absolut Subjectivitet,
være det sande Gudsbegreb, ikke være Udtrykket for den virkelige Gud, og det skal
det ikke fordi Gud i det er bestemt ved Viden, og ad logisk Vei. I sin Sandhed og
Virkelighed skal Gud kun opfattes naar han opfattes som den Gud, \emph{hvis Væsen er Villie} og naar det paa denne Opfattelse hvilende, ved en af Villien bestemt
Erkjendelse naaede, ethiske Gudsbegreb, der bestemmer ham som \emph{den Gode}, som den \emph{personlige} Gud, afløser det logisk-physiske\footnote{Ved det \emph{logiske} Gudsbegreb forstaaer \emph{Martensen} (jfr.\ Om Tro og Viden, S.\ 11) „Begrebet om Gud som upersonlig
Idee, Alintelligents, Νοῦς, ontologisk Subjectivitet eller desl.“; ved det \emph{physiske} „Begrebet om Gud som en altfrembringende Natur, \textit{natura naturans},
Substants o.\ s.\ v.“; ved det \emph{ethiske}: „Begrebet om Gud som Personlighed og
Villie, som det Godes, som Kjærlighedens Gud, som Naadens, og den, ikke blot
logiske, men hellige Sandheds Gud“.}, hvilket det dog ikke skal udelukke, men
bevare i sig som Moment. Det ethiske Gudsbegreb, der, i Modsætning til det
logisk-physiske, opfatter Gud som \emph{overnaturligt} Princip, sætter Villiesbestemmelser,
som: absolut Frihed, Hellighed, Kjærlighed som dette overnaturlige Aandsprincips
Bestemmelser; det bestemmer Gud som den personlige Skabergud. I denne Bestemthed
bliver Gud Gjenstand for den \opage{57}høiere Videnskab, som Martensen giver Navn af
Aabenbaringsvidenskab eller Religionsphilosophie.

Hvad denne Opfattelse af Villiens Forhold til Guds Væsen og det derved betingede
Gudsbegreb angaaer, er det nu strax indlysende, at den tilsyneladende Begrundelse,
ved hvilken for Menneskets Vedkommende Villien skulde hævdes som det Principale, og
som uden videre synes at skulle gjælde for det Guddommelige, ikke kan faae nogen
Anvendelse paa dette, paa hvilket det Menneskeliges Forhold ikke umiddelbart kunne
% PRINTED AS IS, p.57 [collation 2026-09-21]: the page prints "sattes" (a, not æ) where sense wants "sættes"; both readers saw it.
overføres. Den dunkle, forbevidste Villie, der i Mennesket sattes som Forudsætning
for den bevidste Villie og for Tænkningen, faaer ingen Plads i det guddommelige
Væsen; thi selv om der tales om en Physis, en Naturgrund, i Gud, saa skal Gud dog
ikke tænkes som vordende ud fra denne dunkle, bevidstløse Grund, men som evigt
værende i sin frie Viden om sig, i den Viden, der har sin Virkelighed i
Treenigheden. Viden og Villie have da den \emph{samme} Evighed i Gud, og der kan ikke
heller, naar der sees hen til dette Forhold, tillægges Villien et \emph{Begrebsprimat};
thi Villien i sin Klarhed vilde netop have Guds Viden af sig til
Begrebs-Forudsætning, og hvis et høiere Tredie skulde søges, saa vilde det være at
søge, ikke i Villien, men i det guddommelige Selvs Energie, ved hvilken det evigt
sætter sin Væren og sætter den som bestemt ved Alvidenhed og Almagt.

Skal altsaa Villiens Primat i det Guddommelige hævdes, maa det skee ad en anden
Vei: det maa hævdes derved, at det paavises, at idet Villien \emph{som} Villie,
d.\ v.\ s. i sin oprindelige Autonomie, ikke som bestemt ved Erkjendelsen, er det
Ethiskes Princip, den som saadant maa være det Høieste, fordi det Ethiske, der
udtrykker Friheden, er det Høiere i Forhold til det Logiske og Physiske, der ikke
naaer ud over Nødvendigheden.

\section{Kritik af Gudsbegrebet}
\label{sec:gudsbegreb}

At Villien \emph{som} Villie var det Ethiskes Princip, vilde være godtgjort, naar det
vistes, at det Ethiske saaledes var høiere end det Logiske, at det ikke kunde
optages af dette, men dog indesluttede det, paa samme Maade som det skal \opage{58}være
høiere end det Physiske, hvilket det skal indeslutte som Moment. Men en Udvikling,
ved hvilken dette paavistes, er ikke given hos Martensen, uagtet han i mange
Vendinger udtaler, at det ethiske Gudsbegreb er det høieste og det absolut sande;
og ved de Bestemmelser, han giver, forvandler det postulerede Forhold sig uvilkaarligt: det, der postuleredes at være et Høiere, gaaer tilbage til det, i
Sammenligning med hvilket det skulde være det Høiere, ja, det gaaer endog ned under
dette. Og dette er hos ham ikke noget Tilfældigt, men er begrundet i Sagens Natur.
Dette bliver her at godtgjøre med Hensyn til det Ethiskes Forhold til det \emph{Logiske}; det vil senere blive godtgjort med Hensyn til dets Forhold til det \emph{Physiske}.

\subsection{Det ethiske Gudsbegrebs Forhold til det logiske}
\label{ssec:ethisk-logisk}

Naar det ethiske Gudsbegreb af Martensen sættes som høiere end det logiske og
physiske, er dette ikke at opfatte saaledes, at det sammenfatter det Logiske og det
Physiske i en logisk-physisk Enhed, der som saadan bliver det Høiere, men saaledes,
at det, idet det i sig indeslutter det Logisk-Physiske, er et Høiere end selve denne
Enhed. Men deraf bliver den umiddelbare Consequents, der egentlig kun er en
tautologisk Omskrivning af Forholdet, at det Ethiskes Væsen ikke kan bestemmes ved
en immanent Tankeudvikling. Bestemtes det Ethiskes Væsen og Forhold ved en saadan,
vilde det selv gaae ind under det Logiske og kun blive dettes høieste Begreb;
Logiken vilde, som Metaphysik, ogsaa være det Ethiskes Metaphysik, indeslutte
Grundbegreberne for Ethiken, naturligviis i logisk Formbestemthed.

Forat fastholde Forudsætningen, maa det Ethiske sættes som transscenderende det
Logiske \emph{efter dets hele Omfang}, og dette kunde enten opfattes saaledes, at
der statueredes et andet Forhold til Virkelighed i det Ethiske end i det Logiske,
eller saaledes, at der gaves det Ethiske et Indhold, der var rigere end det
Logiskes.

Til den første Opfattelse vises vi hen naar Martensen sætter \emph{Forholdet til Existents} som det Afgjørende. Det Ethiske skal være et Høiere end det Logiske fordi
det viser hen til Existents, medens det Logiske skal være lige\opage{59}gyldigt mod
Existents, kun have ideel, ikke factisk Væren. Men denne Modsætning og disse
Bestemmelser beroe paa en Uklarhed. Det Logiske, som subjectivt, formelt Logisk,
har kun en „Tankeværen“. Det Logiske, opfattet i Enhed med det Ontologiske,
udtrykker det, der har \emph{ideal-real} Væren. „Factisk Existents“, der, i Modsætning til
denne Væren, kun vilde blive \emph{ydre, endelig} Væren, har det ikke, men det begrunder
det, der existerer i denne Form. Og det Ethiskes Væren er netop en Væren af samme
Art som det Logisk-Metaphysiskes, det ethiske \emph{Gudsbegreb} har som Begreb samme
Sphære som det logiske. Skulde det Ethiskes Existents være af en anden Art, blev
den kun hiin endelige, sandselige Existents, og det vil senere, ved Udviklingen af
Personlighedsbegrebet, vise sig, at Martensen netop ikke kan frigjøre sig fra
Forestillingen om denne. Men er nu det Ethiskes Væren af samme Art som det Logiskes,
saa vilde der altsaa ikke, gjennem Forholdet til hiint, kunne sikkres Villien det
fordrede Primat.

Hvad den anden Opfattelse angaaer, saa vil den ikke kunne fastholdes, hverken naar
vi see hen til det Forhold, der i det evige Princip maa sættes imellem dets
Virksomhedsformer, eller naar vi see hen til den menneskelige Erkjendelses Forhold.
See vi nemlig hen til det førstnævnte Forhold, saa ligger det i Begrebet om den
evigt virkelige Gud, at der i Gud maa sættes en evig Enhed og Congruents af Viden
og Villie --- begge Begreber saaledes opfattede som de kunne overføres paa det
Absolute ---, at der Intet kan være sat ved Guds Villie, der ikke er gjennemsigtigt
i hans Viden, hvilken maa tænkes som den, der gjør hans Villie aabenbar, medens
denne har hans Viden til Indhold. Og see vi hen til det andet Forhold, saa viser
det sig, at det, at det Ethiske bestemmes som rigere med Hensyn til Indhold, maa
føre til, at det Ethiskes \emph{overlogiske} Indhold maa sættes som det, der kun kan
gribes gjennem en anden Aandsfunction end det Logiske: ved en „fri Tænkning“, som \emph{Schelling} udtrykker det; ved „en Viden, der er gaaet gjennem Villien“, som \emph{Baader}
betegner \opage{60}det; ved en villende Viden, en ved Villien valgt Viden, som \emph{Martensen} vilde sige. \emph{Men dette Forholds Mulighed kan ikke godtgjøres}. At vise hen til
Villien som til et med Viden absolut uensartet Princip, saaledes som Nielsen
opfatter den, er afskaaret ved Forudsætningerne. Skulde derfor hint Forhold gjøres
tænkeligt, maatte det skee paa den Maade, at den logiske Tænkning paa een Gang
bestemte sin egen Grændse og viste ud derover, og med det Samme paapegede det
Organ, der skulde afløse den, og angav en negativ Norm med Hensyn til det, der
faldt udenfor Grændsen, en Norm, der betryggede mod, at det Vilkaarlige og Absurde
skulde sættes som en høiere Sandhed. \emph{Men en saadan Tankens Selvbegrændsning ved Tanken selv er det Umulige}. Som Udtryk for Væsensenergien er Tanken en Realitet,
enhver af dens Virksomheder er en Stadfæstelse, en Bekræftelse af denne Realitet,
en Selvbekræftelse: en Virksomhed, ved hvilken den skulde negere sig selv, er
derfor et Selvmodsigende. Forholdet maatte da bestemmes saaledes, at ikke Tanken
som saadan, men det tænkende Subject ud fra sin, ogsaa andre Functioner omfattende,
Væsenstotalitet begrændsede Tanken. Men af de i det Foregaaende givne psychologiske
Udviklinger fremgaaer det, at dette Subject i sin \emph{Enhed} maa omslutte
Grundfunctionerne som harmonisk forenede, som bestemte ved samme Maal, og at Viden,
Erkjendelsen, som Heelhed, er en Grundfunction, og den, der i Menneskeaanden har
ideel Prioritet. Heller ikke ad denne Vei kunde altsaa hint statuerede Forhold
hævdes. Lod man det Ethiske som det Høiere komme frem ved et Brud, og fastholdt det
samtidig som det, der indesluttede Viden, saa vilde den Viden, den ethiske Villie
indesluttede, kun være homonym med den Viden, der gik forud for Bruddet, og er
Viden i stræng Forstand. Forholdet vilde da i Virkeligheden være Forholdet mellem
to absolut uensartede Principer, der som saadanne fornuftigviis ikke kunde sættes
som staaende i det Overordnings- og Underordningsforhold til hinanden, Martensen
statuerer.

\opage{61}Det viser sig saaledes, ikke at være tilfældigt, at det postulerede Forhold ikke
godtgjøres som det sande: en Godtgjøren deraf er en Umulighed. Og hvis vi vilde
supponere, at Villiens Primat i den angivne Betydning var hævdet, saa vilde det
indeslutte, at vor Videns Nødvendighed ikke blot var suspenderet for det høiere
Gebets Vedkommende, men \emph{overhovedet} var umuliggjort. Hævdedes Villien som det
Væsensbegrundende, \emph{saa vilde Væsenets Begreb absorberes i Villiens},
Væsensnødvendigheden vilde da være hævet i Villiens Frihed, der, som
transscenderende Væsensnødvendigheden, vilde blive \emph{absolut Vilkaarlighed}
(«\textit{bonum est quia deus vult}»); Viden vilde i Gud, som Udtryk for den for
Væsenet forudgaaende, Væsenet bestemmende Villie, kun udtrykke \emph{den absolute Vilkaarligheds Bevidsthed om sig selv og om det ved Vilkaarligheden vilkaarligt Satte}; og i Mennesket vilde den paa samme Maade savne enhver Nødvendighedens Grund.
Og naar saaledes Villiens Principat ophævede Videns Nødvendighed, saa vilde dette
indeslutte, at det i Virkeligheden aldrig kunde blive en Gjenstand for Viden, naar
Viden tages i \emph{egentlig} Betydning.

Naar det, af de angivne Grunde, er ugjørligt at føre et Beviis for det opstillede
Forhold, saa bliver det uundgaaeligt, at Forholdet forvandles, hvor der, nærmest i
thetisk opstillede Sætningers Form, gjøres et Forsøg paa at give nærmere
Bestemmelser med Hensyn til det ethiske Gudsbegrebs høiere Værd end det logiskes.
Væsenet med dets Nødvendighed, der bliver Videns Bestemmelse, faaer uvilkaarligt
Principatet i Forhold til Villien. Paavisningen heraf bliver forsaavidt det mindre
Væsentlige, som de af Martensen givne Udviklinger kunde lide af tilfældige Mangler;
men det er dog af Interesse, efterat det almindelige Forhold ovenfor er paaviist, at
see, hvorledes det naturlige Forhold uvilkaarligt gjør sig gjældende, trods
Udgangspunktets Strid derimod.

Vi fremhæve her de Bestemmelser af Forholdet mellem det Ethiske og det Logiske,
mellem Villien og Viden, der \opage{62}komme frem idet Forholdet mellem \emph{den} Gode og \emph{det} Gode
skal angives, og vi fremhæve dem ogsaa af den Grund, at de indeslutte Momenter til
Opfattelsen af Begrebet om den Frihed, der skulde hæve Villien over Viden.
Forholdet mellem det Gode og den Gode viser tilbage til hine Forhold, idet \emph{det} Gode, som „det evigt værende Gode“ viser hen til det Theoretiske, \emph{den} Gode, ifølge
den angivne Opfattelse af det Constitutive, til Villien. Udviklingen af dette
Forhold er hos Martensen vaklende og usikker, og fører netop til det modsatte
Resultat af det tilsigtede. Der siges\footnote{Martensen: Om Tro og Viden, S.\ 67
f.}, at \emph{det} Gode --- der med Nødvendighed fordres, --- kun kan existere som en fri
Villie, og derfor maa bestemmes som \emph{den} Gode. „Nu er“, fortsættes der, „Gud som den
Gode vistnok den med Nødvendighed Existerende; men var han kun dette, da vare vi
endnu ikke komne ud over det Physiske og Logiske; thi en Gud, der kun er god med
Nødvendighed (den gode Natur) er ikke den i Sandhed ethiske Gud. Men det Gode som
hans Væsens Nødvendighed kræver netop, at han med Frihed realiserer det, eller at
han, som den, der ikke har en Lov udenfor sig eller over sig, men som er sig selv
den hellige Lov, som den absolut Autonome, actualiserer, frembringer sig selv som
den Gode, og forsaavidt han frembringer det Gode udenfor sig, ikke gjør det efter
Naturnødvendighed, men efter Frihedens Nødvendighed“. --- Netop i de givne
Bestemmelser ligger, at, naar der overhovedet skal være Tale om en Lov,
Væsensbestemmelsen og Væsensnødvendigheden bliver det Første, Forudsætningen for
Villien; altsaa bliver Væsenet den hellige Lov: det Gode bliver Lov for den Gode.
Og naar der tales om en Actualiseren, saa ligger i Actualiseringens Begreb en
Virken, altsaa en endnu ikke som „den Godes“ Villie bestemt Virksomhed, altsaa en
ved Væsensnødvendigheden, der faaer sit Udtryk i Viden, bestemt Energie, ved
hvilken den ethiske Villie er, og som, idet den er det ideelt \opage{63}Begrundende, bliver
det Principale. Og denne væsensnødvendige Virken kan, da Gud ikke skal være
vordende, og da han er evigt vidende, af Martensen ikke bringes ind under „den
dunkle Villies“ Bestemmelse.
--- Til det samme Resultat føre ogsaa de efterfølgende Udviklinger, i hvilke det
ethiske Gudsbegreb skal vises som Sandheden af det logiske og physiske. Som den
absolut villende skal Gud tillige være den absolut vidende, der veed sig selv i et
evigt Idealsystem. Men den ethiske Gud skal ikke være den i Ideen indesluttede,
ideebundne, og saa at sige i Ideen absorberede Gud. Hans Viden skal bestemme sig
som Viisdom, hvilket netop skal betegne det absolut Ideefrie i hans Existents i
Modsætning til Hedenskabets ideebundne, blot logiske Gud, og betegne hans Viden
ikke blot som en logisk, men som en ethisk, en villende Viden. --- Naar nemlig Gud
sættes som den \emph{absolut} vidende, veed han sig altsaa \emph{ogsaa i Bestemmelsen af den Gode}, og idet der hos Gud ikke kan tales om en blot afspeilende eller om en
uadæquat Viden, saa gaaer den guddommelige Villie ind under Videns Nødvendighed,
sees som Formen for Væsenets \textit{libera necessitas}. Og naar \emph{Viden} bestemmes
som Viisdom, saa fører Viisdommens \emph{Indhold} og Viisdommens \emph{Maal}
tilbage til den samme Overgriben af Viden, hvormeget der end tales om Forskjellen
mellem evige Ideer og evige Raadslutninger. Bestemmelsen af Guds Forhold til Ideen
bliver meget uklar og Ideen hævder sig, imod hvad der tilsigtes, som det Høieste.
Ideen udtrykker et Logisk, men dog skulle ikke blot Ideerne have deres
Enhedsprincip i Gud, uagtet den absolute Idee, der dog vel maatte være Ideernes
Enhedsprincip, ikke skal være Gud; men en Idee: Villiens \emph{Idee}, skal være Udtrykket
for den \emph{ethisk} bestemte Gud, altsaa for det, der er \emph{over det Logiske}, og den
\emph{absolute Personlighed} opfattes som den høieste \emph{ontologiske} Bestemmelse.
Villiens \emph{Idee} skal være det, under hvilket \emph{Alt} er subordineret. Den skal indeslutte
\emph{Existentsens} Bestemmelse, altsaa netop den Bestemmelse, der \opage{64}skal charakterisere den \emph{virkelige} Gud til Forskjel fra det \emph{logiske} Gudsbegreb.

Resultatet af det Udviklede bliver altsaa: at Forsøgene paa at bestemme Forholdet
mellem Villie og Viden i det guddommelige Væsen gjennem Forholdet mellem det
Ethiske og det Logiske maae føre tilbage til Enheden af Villien og dens Frihed med
Viden og dens Nødvendighed, saaledes at det Logiske kommer til at omfatte ogsaa
Begrebet om den ethiske Gud, og at hvad der er fremsat om Villiens Primat og hvad
der staaer i Forbindelse dermed, ikke er godtgjort og ikke kan godtgjøres.

\subsection{Den guddommelige Villies Frihed}
\label{ssec:guds-frihed}

Hvad der er udviklet med Hensyn til Forholdet mellem Villie og Viden i Sammenhæng
med Forholdet mellem det Ethiske og det Logiske, kan faae en ny Belysning derved,
at det opstillede Begreb om Villien som \emph{Frihedsprincip} analyseres, og derved den
statuerede \emph{Friheds} Art og Forhold til Videns Nødvendighed opklares, da det derved
vil vise sig, hvorfor denne Frihed ikke kan hævde Villien den fordrede Forrang i
Forhold til det ved Nødvendigheden Bestemte. Resultatet af denne Analyse er i det
Væsentlige anticiperet i de gjennem de foregaaende Udviklinger vundne Bestemmelser.

Vi see her, ligesom ved de tidligere Udviklinger, bort fra det Spørgsmaal: hvorvidt
Overførelsen af den \emph{psychologiske} Bestemmelse \emph{Villie} paa \emph{Gud} overhovedet er
berettiget, og om ikke dette Begreb nødvendigviis i denne Overførelse paa det
Absolute maa undergaae en Transformation; et Spørgsmaal, til hvilket vi komme
tilbage ved Udviklingen af Begrebet om Guds Personlighed. Vi opkaste her, idet vi
anvende Begreberne Villie og Viden paa det Guddommelige, Spørgsmaalet: hvorledes
den Frihed vilde blive at bestemme, ved hvilken Villien skulde hævdes som det
Principale, som det Overordnede i Forhold til Viden og dens Nødvendighed? Den
Frihed, der skulde hævde Villien denne Betydning, maatte være en Frihed, løsnet fra
Nødvendighedens Bestemmelse, thi denne Bestemmelse, som Bestemmelse ved Friheden,
vilde føre den \opage{65}tilbage til Væsenets og Videns i deres Grund identiske Nødvendighed.
Men en saadan Frihed vilde, som den \emph{absolut nødvendighedsløse} Frihed, være
% PRINTED AS IS, p.65 [collation 2026-09-21]: page prints "liberium" for liberum (see also EMPH p.65 on same span); confirmed by a blind second reader.
et \emph{\textit{liberium arbitrium}}, der ingen Grund har i Væsenet, og derfor kun bliver \emph{grundløs} Vilkaarlighed. At den factisk opfattes saaledes, vil blive klart naar
Skabelseslæren og Læren om Underne analyseres. Men en Frihed af denne Art, der
tilmed, ved at tilintetgjøre al Videns Nødvendighed, vilde tilintetgjøre Viden og
blive en „ikke-intelligent“ Frihed, vilde ikke kunne hævde Villien en høiere
Gyldighed. Ligesom den selv er et Lavere end den ved Fornuftnødvendigheden bestemte
Frihed, end den Væsenets \textit{libera necessitas}, der udtrykker det for sig
gjennemsigtige Væsens Frihedslov og ligemeget aabenbarer sig i Viden og Villien,
saaledes vilde den gjøre Villien, ikke til det Høieste, men til et Lavere, den
vilde sætte et Sublimat af den menneskelige Vilkaarlighed, et blot Phantasiebillede
af Villien, i Stedet for Villiens sande Begreb: at være Væsensvillen, Udtryk for
det totale Væsens Energie, i hvilken Viden og Villie ere i Enhed, og som derfor kun
kan yttre sig lovbestemt, i fri Nødvendighed.
Gjennem hint Begreb om Frihed er Villien, der betegnes som „det eneste Princip, der
kan begynde Noget“, sat ikke blot som \emph{Videns}, men som \emph{Væsenets} \textit{prius}; i
sin Aseitet \emph{sætter den paa een Gang Væsenet og transscenderer det}\footnote{Sondringen
af Villie og Væsen bliver tydelig i Læren om Skabelsen af Intet. Idet Gud skaber
ved Villien \emph{af Intet}, skaber han \emph{ikke af Væsenet}, og Sondringen mellem Villie og
Væsen er sat. Villien sætter Væren, sætter Væsenet, men transscenderer Væren og
Væsen.}. Istedenfor at være Væsensyttring, bunden ved Væsenets Frihedslov og
Væsenets Indhold, bliver den altsaa en ubegrændset, af Intet baaren Villie, der
proteusagtigt forvandler sig og kan ville alt Muligt\footnote{Dette finder sit
Udtryk i Forestillingen om en „Al-Mulighed“, og om Gud som „Mulighedernes Herre“.}
ɔ: \emph{alle Phantasiens Mulig\opage{66}heder}; thi denne Villie er i sin Lovløshed netop
kun bestemt ved \emph{den skrankeløse Phantasies lovløse Almagt}. Sattes \emph{Villien selv} som
Væsen, saa vilde dette enten føre til, at Villien, imod Forudsætningen, fik
Nødvendighedens Bestemmelse i sig, eller --- naar Friheden i hiin Betydning
fastholdtes --- til, at Friheden, d.\ v.\ s. Vilkaarligheden, i sin Ubundethed ved
alle Betingelser, altsaa \emph{ogsaa ved det af den selv Satte}, i hver Frihedsyttring
negerede dette, altsaa negerede \emph{den med Væsenet identiske Villie}, hvis Bestemmelse den \emph{selv} er, saa at Villien altsaa, medens den skulde være et Væsentligt ɔ: et
Blivende, Uophæveligt, kun vilde være i den uophørlige Forandring, der kun er
Selvophævelse.

Gjennem hint Frihedsbegreb, der staaer i nøie Forbindelse med det „ethiske“
Gudsbegreb, bliver det saaledes umuligt at hævde Villiens Forrang i det
Guddommelige. Vi vises atter hen til Nødvendigheden af at sætte Villie og Viden i
evig Enhed og Congruents i det Guddommelige, at sætte det Ethiske og Logiske i
Enhed, og det bliver indlysende, at Bestræbelserne for at hævde Villien hint høiere
Værd, idet de maae gaae ud paa at udelukke al Nødvendighedsbestemmelse fra den,
netop maa faae det modsatte Resultat af det tilsigtede, nemlig: at drage Villien og
det, der benævnes det Ethiske og Friheden i Gud, ned under det, der skulde være det
Underordnede: under Viden, under det Logiske og under Væsensnødvendigheden. Men det
Ethiskes og Frihedens Navn anvendes ogsaa kun her ved en Misbrug; det, der benævnes
saaledes, naaer end ikke til Grændsen af det Ethiske, berører ikke Frihedens sande
Begreb, overhovedet ikke Begrebets Sphære.

\subsection{Guds Personlighed}
\label{ssec:guds-personlighed}

Hvad der er sagt angaaende Opfattelsen af de omhandlede Begreber, stadfæstes ved
den derved bestemte Opfattelse af \emph{Guds Personlighed}.

For først overhovedet at kunne anvende dette Begreb paa Gud, maa den Indvending
imødegaaes, at Personlighed, \opage{67}som forudsættende Begrændsning ved et Andet, altsaa
Endelighed, er uanvendelig som Bestemmelse for det Absolute. Denne Indvending
imødegaaes væsentligen kun gjennem Postulater. Den skal være ugyldig, fordi det, at
der foruden Gud gives skabte Væsener, ikke skal være en Skranke, der krænker
Begrebet om det fuldkomne Væsen. Selvbegrændsning skal være uadskillelig fra den
fuldkomne Natur. „Gud speiler sit Væsens Fylde i sin Selvbevidstheds indre
Uendelighed og tager den derved i Sandhed i Besiddelse. Thi et alfuldkomment Væsen,
der ikke veed om sin Fuldkommenhed, mangler en meget væsentlig Fuldkommenhed. Gud
begrændser sin Magt idet han af sit evige Livsdyb fremkalder en Verden af skabte
Væsener, hvilke han giver, i afledet Betydning at have Livet i sig selv; men just
derved, at han er Magten i en fri Verden, aabenbarer han den indre Storhed i sin
Magt. Thi ikke er den Magt den sande, der ikke taaler nogen Frihedsrørelse udenfor
sig, fordi den selv umiddelbart vil være Alt og virke Alt; men den Magt er den
sande, der skaber Frihed, og som ikke destomindre formaaer at gjøre sig til Alt i
Alle“.

Disse Bestemmelser, der ikke ere naaede gjennem en Deduction, men kun thetisk
fremsatte, sammenknytte umiddelbart hinanden ophævende Modsætninger (f.\ Ex. „en Magt, \emph{der skaber Frihed}“, „\emph{i afledet Betydning at have Livet i sig selv}“), og
overføre Begreber fra en Sphære til en anden; hvor netop \emph{det} mangler, der giver dem
deres Bestemthed. Selvbevidsthed postuleres om Gud som \textit{Conditio sine qua
non} for Fuldkommenhed. Men Selvbevidsthed, saaledes som vi kjende den i
\emph{Mennesket}, d.\ v.\ s. saaledes som vi \emph{overhovedet} kjende den, er kun mulig
ved Modsætningen til Andetheden, altsaa ved Begrændsethed. Men hos Gud skal den gaae \emph{forud for} Begrændsningen ved et Andet og alligevel være Selvbevidsthed. Den for den
menneskelige Selvbevidsthed nødvendige Begrændsning bringes saa dog i Gud ind
\emph{bagefter} derved, at den selvbevidste Gud begrændser sig ved af sit evige
Livsdyb at fremkalde hiin Verden af skabte Væsener, som han giver, i \opage{68}afledet
Betydning at have Livet i sig selv. Han begrændser sig, og denne Begrændsning
bliver ikke blot en Begrændsning ved et Aandeligt, ved et i sig centraliseret
Frihedsliv, der har Muligheden til at stille sig i et Modsætningsforhold til Guds
Villie, men en Begrændsning \emph{ved et Materielt}, ved et \emph{Guds aandelige} Væsen Modsat, og
dog skal den ikke gjøre Uendeligheden Skade: Uendeligheden skal være \emph{indre}
Uendelighed. Men hvorved skiller Guds Væsen sig saa fra hver enkelt aandelig, ja,
fra hver enkelt levende Skabning, der ogsaa har indre Uendelighed? \emph{Indre Uendelighed} er relativ \emph{abstract} Uendelighed, relativ abstract Forhold til sig selv
gjennem sit uophævede Andet, eller ogsaa \emph{Mulighedens Uendelighed}; men kan dette være
en \emph{adæqvat} Bestemmelse ved Gud? er det en Bestemmelse, der fjerner Endeligheden fra
ham? Disse Spørgsmaal kunne vi kun besvare benægtende.

Men idet det ikke lykkes at frigjøre Personlighedens Bestemmelse fra
Begrændsningen, Endeligheden, saa bliver Følgen, at Gud ved hiin Bestemmelse
bestandig drages ned i den \emph{endelige Forestillings} Sphære, at Forestillingen om den
\emph{sandselig} bestemte \emph{Individualitet} skyder sig ind under Begrebet om Guds
Personlighed. Denne Forestilling om det sandselige Enkeltindivid kan og vil
Martensen ikke holde borte fra Opfattelsen af det Guddommelige. Han
siger\footnote{\emph{Martensen}: Om Tro og Viden, S.\ 47.}: „den personlige Gud maa
uundgaaelig tænkes analog med den psychologiske Bevidsthed“, og ved „analog“
forstaaes her, som det fremgaaer af de efterfølgende Udviklinger, Mere end en blot
Analogie: der forstaaes derved Formens Identitet. \emph{Nødvendigheden} af at fastholde
hiin endelige Forestilling ligger i det Dobbelte, der tilstræbes gjennem
Bestemmelsen af Gud som personlig. Der tilstræbes \emph{for det Første} at give Gud en
Væren \emph{udenfor vor Tanke}: netop Personlighedens Bestemmelse skal holde Gud i
Enkeltexistentsens Sondring fra det Men\opage{69}neskelige, hævde
Forestillings-Transscendentsen. Men denne Sondring ved et „udenfor“, ved hvilket
der lægges Eftertryk paa, at Gud som ikke blot værende i vor Tanke, vor Aand, ikke
blot har det Ideelles Væren, gjør netop Gud \emph{endelig} ved at give hans Væren
en \emph{Sandselighedens} Bestemthed, selv om denne atter \emph{formelt} negeres. Der
tilstræbes, for det \emph{Andet}, \emph{at frigjøre Gud fra Nødvendighedsloven}, at finde
et Tilknytningspunkt for den ovenfor kriticerede Opfattelse af Villie og Frihed, og
da dette ikke kan findes \emph{over} det Almene, søges det \emph{under} dette i
Forestillingen om Enkeltvæsenet, der i den Endelighed, som Bevidsthedsforholdet
antyder, har en Mulighed til Villiens vilkaarlige Løsnen fra
Fornuftnødvendighedens Almene, --- rigtignok kun paa det Vilkaar, at Villien
ubevidst falder ind under Naturnødvendighedens Almene; og det menneskelige
Enkeltvæsens Bestemmelse overføres derfor paa det Guddommelige. Netop af den Grund,
at dette Dobbelte søges i Bestemmelsen af Guds Personlighed, tilfredsstiller ikke
Begrebet om en \emph{Subjectivitetens Idee}, om en \emph{absolut Subjectivitet}, d.\ v.\ s. om en \emph{absolut concret Subjectivitet}, der ikke, som den menneskelige, i abstract
Selvbevidsthed stiller sig fra sin Bevidsthedsgjenstand og sit Bevidsthedsindhold,
men i sin Væren for sig sætter og idealiserende sammenfatter det Reales Totalitet.
Uden at man gjør sig Rede for den Transformation af Begrebet, der maa foregaae ved
Overførelsen fra den menneskelige Personlighed, der er bunden til Bevidsthedens
Endelighedsform, til det absolut Uendelige, fordres der, i Modsætning til Begrebet
om en absolut Subjectivitet, et „virkeligt“ Subject, et „realt
Subject“;\footnote{jfr.\ \emph{Martensen}: Om Tro og Viden S.\ 38, 49, 53.} men virkeligt
og realt bliver for Martensen enstydigt med begrændset, endeligt, med en Existeren i
den menneskelige Villies og Selvbevidstheds Form, skjøndt denne sidste af ham kun er
forklaret som en Speilingsform.

\opage{70}Den samme Opfattelse gjør sig ogsaa gjældende i de nærmere Bestemmelser af denne
guddommelige Personligheds Forhold til Mennesket i Aabenbaringen, og af Organet for
Menneskets Opfattelse af den personlige Gud.

Angaaende \emph{Aabenbaringsforholdet} siges der\footnote{\emph{Martensen}: Om Tro og Viden, Side
93 f.}, at hvis Gud er den personlige, villende Gud, saa kunne vi hverken erkjende
Mere eller Mindre af ham end han selv vil aabenbare os. „Det gjælder allerede i
Forholdet mellem Menneske og Menneske, at en Personlighed \emph{kun} kan erkjendes
af dens frie \emph{Handlinger}, og at ethvert nærmere Bekjendtskab beroer derpaa, at
den selv vil komme os imøde og oplade sig for os i fri \emph{Meddelelse}; og dette gjælder
i eminenteste Betydning om den guddommelige Personlighed. \textit{A priori}
construere ham kunne vi ikke; thi kun det logisk Nødvendige, men ikke den frie Aand
lader sig \textit{a priori} construere. Selv det menneskelige Genie kan ikke
\textit{a priori} construeres, men kun erkjendes naar dets underfulde Virkelighed er
given. Først efter Virkeligheden kunne vi erkjende, at Sligt er muligt, da dets
Mulighed langt overstiger de blotte Fornuftmuligheder“. --- Det er her
iøinefaldende, at Forestillingen om den sondrede Enkeltexistents atter bringes ind,
og at det derfor glemmes, at Guds \emph{Væsen} er bestemt som \emph{Kjærlighed}, Kjærligheden
altsaa som \emph{nødvendig Væsensmeddelelse}, der, netop efter Martensen, ogsaa maa
opfattes som en Meddelelse til \emph{Erkjendelsen}. Ved den skjæve Analogie med Geniet er
det ligeledes overseet, at der, i Forhold til Gud, ikke er Tale om, at construere
\textit{a priori} et historisk Fremtidigt, men om at begribe et alment, evigt
Grundforhold, et Væsensforhold. Naar der saa tilføies: „ingen
Personlighedsaabenbarelse lader sig opløse i blotte Tankebestemmelser. Den kan kun
fattes ved \emph{Anskuelse} som alt Individuelt og Oprindeligt, og her er jo Tale om det
absolute Enkeltvæsen“, --- saa er det glemt, at det „\emph{absolute Enkelt\opage{71}væsen}“
\textit{eo ipso} er \textit{ens universalissimum}, og naar denne sidste Bestemmelse
oversees, saa bliver der slet ingen Begriben, men kun en \emph{sandselig} Anskuen, der
forholder sig til et sandselig Bestemt. Men Characteren af den „Anskuelse“, ved
hvilken det høiere \emph{Erfarings}forhold skal virkeliggjøres, og som skal faae sin
Betydning ved Begrundelsen af den \emph{høiere Videnskab}, vil ved Kritiken af dennes
Begreb nærmere blive belyst.

\subsection{Det ethiske Gudsbegreb i Forhold til det physiske}
\label{ssec:ethisk-physisk}

Sammenstillingen af det \emph{logiske} Gudsbegreb med det, der betegnes som det \emph{ethiske}, har altsaa viist det førstes Overlegenhed, har viist, at det \emph{Logiske},
opfattet i sin Enhed med det \emph{Metaphysiske}, omfatter \emph{ogsaa det Ethiskes}
Væsensbestemmelser --- Noget, Martensen selv uvilkaarlig kommer til at indrømme ved
at bruge Bestemmelserne den absolute \emph{Idee}, det Godes \emph{Idee}; --- medens det, der
benævnedes det \emph{ethiske} Gudsbegreb, paa Grund af den \emph{endelige} Forestilling om
Personlighed, Villie og Frihed, \emph{ikke} naaer til det i \emph{Sandhed Ethiske}, og,
istedenfor at blive et Høiere end det Logiske, bliver i det Alogiske, hvorfor ogsaa
Forsøgene paa, at deducere dette „Gudsbegrebs“ Bestemmelser rationelt, maae
mislykkes. Og naar vi nu fra Betragtningen af det „ethiske“ Gudsbegrebs Indhold til det \emph{logiske} gaae over til Betragtningen af dets Forhold til det \emph{physiske}, saa vil
Væsenssammenhængen mellem det Logiske og det Physiske medføre, at Resultatet bliver
et analogt. Naar Naturen sees som levende Enhed, som teleologisk stræbende hen imod
Aanden, som det, hvis Væsen er Aand, og hvis Begreb først i Aanden finder sin
Fuldbyrdelse, saa sees Naturen som det objectivt værende Fornuftige, det Logiske
som immanent i den, og Logiken som Metaphysik udtrykker paa een Gang Naturens og
Aandens fornuftbestemte Væsen, bliver paa een Gang Naturens Logik og Aandens Logik.
Og i Mennesket bliver det Naturlige Villiens Grundlag, og, som det fornuftbestemte
Reale, det Ethiskes Basis, det, der skal give den frie, bevidste \opage{72}Aands Handling
Fylde og Virkelighed, som gjennemklaret ved Erkjendelsen blive Indhold for den
ethiske Villie.

I Martensens Opfattelse af det Physiske oversees ikke heller Naturens
Fornuftbestemthed. Fornuften er i Naturen, og Naturen gjør Tankens Gjerninger. Men
Naturen er dog kun det blindt Existerende: blindt i subjectiv Betydning, forsaavidt
Naturen ikke er selvtænkende, blindt i objectiv Betydning fordi Naturen kun
existerer med Nødvendighed, ikke med Frihed, og enhver ufri Nødvendighed er en
blind, saa at sige fatalistisk, Nødvendighed. Og det physiske Gudsbegreb, der har
denne blindt \emph{existerende} Natur til Underlag, skal ingen Dialektik kunne bringe i
sand Enhed med det logiske, fordi den logiske Gud er den rent ideelle, mod Existents
ligegyldige. Existents og Idee skulle kun have deres virkelige Enhedspunkt i en
Villie: kun i den ethiske Gud skal deres fuldkomne Enhed kunne findes.

I det ethiske Gudsbegreb skal derfor det physiske være indesluttet som Moment og
som nødvendigt Moment, men saaledes, at det Ethiske, ligesom det transscenderede
det Logiske efter dets hele Omfang, ogsaa i sin absolute Frihed hæver sig over \emph{alt} det Physiske, atter \emph{negerer} det som Betingelse.

\subsection{Dobbeltheden i Opfattelsen}
\label{ssec:dobbelthed}

I disse Bestemmelser viser sig en \emph{Dobbelthed} i Bestemmelsen af Forholdet.

Naar det „ethiske Gudsbegreb“ skal \emph{begribeliggjøres}, gjøres til Gjenstand for en \emph{Erkjendelse}, saa maa den Side accentueres, at det Physiske er \emph{nødvendigt} Moment, \emph{nødvendig} Betingelse for det Ethiske, og det Physiske faaer saa \emph{væsentlig} Gyldighed,
sees som bestemmende det Ethiske.

Naar derimod Accenten falder paa den „ethiske“ Guds\emph{forestilling}, paa Forestillingen
om Friheden i den ovenfor angivne Betydning, saa \emph{ophæves} atter det Naturliges \emph{Væsentlighed}. Gud bliver \emph{naturløs} og det bliver \emph{umuligt}, fra Forestillingen om denne
Gud at naae til det Physiskes Bestemmelse.

\opage{73}Dobbeltheden bliver strax tydelig i de Bestemmelser, Martensen giver, idet han
\emph{nærmere} skal angive Forholdet mellem det Ethiske og det Physiske i Gud. Vi
kunne see bort fra saadanne, efter Formen maaskee mere tilfældige, Udsagn som f.\
Ex. det, at Guds ethiske Villie ikke kan være uden en Natur, i Forhold til hvilken
den kan bestemme sig som ethisk\footnote{\emph{Martensen}: Om Tro og Viden. S.\ 51.}, et
Udsagn, der deels vilde vise tilbage til en uigjennemførlig Analogie med det
Menneskelige, deels til den af Martensen ikke billigede Lære om en vordende Gud, der
fra en dunkel Grund befrier sig til Aand og Personlighed. Vi holde os til andre
Udsagn, der, som definitivt bestemmende Forholdet, ikke kunne have denne
Tilfældighedens Charakteer. Der siges\footnote{Sammesteds. S.\ 68.}, at Gud, som den
absolut vidende og villende tillige skal være den absolut kunnende. Guds ethiske
Villie skal ikke kunne tænkes uden en evig Natur, Physis, i Gud, uden hvilken hans
Villie vilde være uden Redskaber og Productionspotentser; den skal ikke kunne
tænkes uden en evig Magtfylde. --- I disse Bestemmelser er der tillagt det Physiske
en real Gyldighed og en væsentlig Betydning for Villien og det Ethiske, netop
derved, at Villien, uden hiin Physis, sættes som værende uden Productionsmulighed,
afmægtig til Frembringen og overhovedet til Virken\footnote{Sees Villien som det,
ved hvilket evigt det guddommelige Væsens Væren sættes, saa er denne Sætten i
høieste Forstand en Virken, og, \emph{bestemtere}, \emph{en Virken med Hensyn til det, til hvis Væsen en Physis skal staae i nødvendigt Forhold}; Villiens Virken kan altsaa heller
ikke her tænkes uden „Redskaber“, d.\ v.\ s. uden et evigt værende Physisk, naar den
ikke skal kunne tænkes uden et Physisk med Hensyn til det \emph{udenfor} Gud Værende.}. Og
der ligger i dem den videregaaende Consequents, at naar Villien, forat virke,
\emph{forudsætter} Redskaber, og disse Redskaber altsaa ikke kunne tænkes som satte ved
Villien, da den i denne Sætten netop vilde have virket uden Redskaber; saa maae
Redskaberne faae lige Oprindelighed med Villien, og, naar der gjøres Alvor af
Bestemmelsen: Phy\opage{74}sis, faae vi enten samme Dualisme, som Martensen bebreider Nielsen,
eller vi faae en \emph{naturlig} Gud, for hvis physiske Nødvendighed Villien bliver
Virksomhedsformen. Men netop det \emph{Modsatte} tilstræbes, og derfor \emph{tilbagekalde} de
følgende Bestemmelser det, hvorpaa disse Consequentser støtte sig. Der siges nemlig,
i \emph{umiddelbar} Sammenhæng med det ovenfor Anførte, at hiin evige Magtfylde er en \emph{Al-Mulighed}, der forholder sig til det Ethiske i Gud som \emph{det Tjenende},
Underordnede, til det Overordnede, Herskende og Kongelige, som Guds \emph{uegentlige Andet-Væsen} til hans \emph{egentlige inderste Væsen}. „Gud er, som \emph{Frihedens} Gud, \emph{Mulighedernes Herre}, og kan frit aabenbare sig, hvilket ikke gjælder om den logiske
og physiske Gud, hvis Muligheder blindthen gaae over i deres Følger. Gud er den absolut \emph{naturfrie} (ikke \emph{naturløse}), og som saadan er han i sig selv den fuldkomne
Aand, den fuldkomne Personlighed, den sande Gud, der har det Logiske og Physiske til
Attributer, Prædicater, Egenskaber, er \emph{levende Enhed af forskjellige Momenter}.“

Idet Magtfylden sættes som liig med en Al-Mulighed, idet den bestemmes som det
Tjenende, Underordnede, altsaa bliver det, der skal magtes af Noget, der falder \emph{udenfor} Magtfylden, idet Guds Frihed i den ovenfor udviklede Betydning paaberaabes,
og idet Gud bestemmes som den absolut naturfrie, som \emph{absolut, ubetinget raadende over det Physiske}, sit Redskab, \emph{saa bliver Bestemmelsen Physis berøvet sin Realitet}, thi idet den guddommelige Villie, hvis \emph{absolute Prioritet} altsaa \emph{ogsaa i Forhold til denne Physis}, vedbliver at være Forudsætningen, \emph{ubetinget} raader over
Redskaberne, saa \emph{bestemmes Villien ikke ved deres Natur}, den kan, \emph{medens den ved at bruge dem}, \emph{tilsyneladende tillægger dem Realitet} som \emph{Redskaber}, i \emph{deres Sted substituere alt Muligt}, og ved denne \emph{Substitution er deres Realitet negeret}. Og naar
den \emph{ethiske} Guds Begreb identificeres med Begrebet om det „\emph{overnaturlige}“ Princip,
saa bliver det endnu tydeligere, at Phy\opage{75}sis i Gud kun bliver et \emph{Naturens Skin}, en \emph{Phantasiens Sublimation af den virkelige Natur}, der kun optages til Bedste for
Analogien med det Menneskelige, og til Bedste for en illusorisk Deduction. Og dog
skal det guddommelige Væsen være Enhed, sand Enhed af de forskjellige Momenter.
\emph{Dette bliver kun et Postulat}: \emph{Naturrealiteten og Naturnødvendigheden maa altid blive ubegribelig ud fra en saadan Gudsforestilling}.

\subsection{Skabelsen}
\label{ssec:skabelsen}

Denne Consequents vil stadfæste sig naar den „ethiske“, personlige Gud sættes i
Forhold til den \emph{virkelige} Natur, og naar der forsøges paa en Deduction af Læren om
\emph{Skabelsen} og om \emph{Miraklerne}.

Analysere vi først \emph{Skabelsesforestillingen} saaledes som den religiøse Bevidsthed
opfatter den, saa ligge deri følgende Bestemmelser:

\begin{enumerate}[label=\arabic*)]
\item \emph{Skabelse er Skabelse af Intet; men det, der er skabt af Intet, d.\ v.\ s. i}
  \emph{sin Oprindelse er et Intet, det er i sin Væren vedblivende et Intet, ikke}
  \emph{\textit{natura}, men \textit{creatura}. Ved Skabelsesforestillingen er derfor}
  \emph{Naturbegrebet negeret.}
\item \emph{Idet det Skabte er sat af den i sig fuldkomne, som Aand bestemte, Gud, og sat}
  \emph{i Bestemmelsen af et Materielt ɔ: i en den rent aandelige Guds Væsen absolut modsat}
  \emph{Bestemmelse, saa er Skabelsen ikke begrundet i Guds Væsens Nødvendighed, men i en}
  \emph{af Væsenet ubunden Villie, i en absolut Vilkaarlighed, hvis Virken, idet den}
  \emph{falder rent udenfor Væsensnødvendigheden og sætter det Væsenet Modsatte, er det}
  \emph{absolut Ubegribelige.}
\end{enumerate}

Mod disse Consequentser, der skarpt angive de simple Begrebsforhold, gjør Martensen
Indsigelse. Han fastholder mod den \emph{første} deels, at en Physis i Gud er Skabelsens
Forudsætning, og at, idet denne Physis bliver liig med en Al-Mulighed, det Intet, af
hvilket Verden skabes, skal være \opage{76}Eet med Guds Villies evige Muligheder; --- deels,
at Skabelsen ikke skal sætte et Intet, men et Noget, der frembringer og udvikler sig
selv i given Selvstændighed. Mod den \emph{anden} Consequents indvender han, at Skabelsens
Grund ikke maa søges i en absolut Vilkaarlighed, men i det Ethiske, i den ethisk
bestemte Almagt.

Hvad Indvendingerne imod den \emph{første} Consequents angaaer, saa er Forestillingen om en
Physis i Gud allerede ovenfor paaviist i sin indre Modsigelse. Denne Forestilling,
saaledes som den i sin Oprindelighed er fremtraadt hos \emph{Jacob Bøhme}, er fremgaaet af
en Følelse af Utilfredsstillelse ved Forestillingen om Naturens Oprindelse fra en
naturløs Gud, af Følelsen af det Intet Forklarende i Forklaringen af Naturens
Tilblivelse ved en absolut grundløs guddommelig Villiesact. Forat forklare Naturen,
henlægger man saa Naturen i Gud; man udtaler indirecte, at Naturen kun kan forklares
ved Naturen selv, at den har Væsentlighed og Selvgyldighed, man røber altsaa ubevidst
en Tvivl med Hensyn til Skabelsesforestillingen. Men den uvilkaarlige Anerkjendelse
af Naturen tager man atter tilbage, netop ved at henlægge Naturen i \emph{Gud}, ved at
ombytte den \emph{virkelige} Natur med dens Afbillede i \emph{Phantasien}, af hvilket saa hiin skal
forklares. Det er \emph{den ved Phantasien metamorphoserede Natur}, den Natur, der, skjøndt
„wahrhaftig und eigentlich“, dog er „himmlisch und geistlich“, man henlægger i Gud,
og idet man gjør den til Grund for den reale Natur, lader man paa den ene Side
Metamorphosen skjule det fra den religiøse Opfattelse Afvigende, paa den anden
Navnets Lighed tilveiebringe Skinnet af en Forklaring af den virkelige Naturs
Tilværelse. Men den aandeliggjorte, dematerialiserede, denaturerede Physis i Gud
forklarer Intet med Hensyn til den reale Naturs Væren, og en virkelig
Naturbestemthed i Gud vilde lade Skabelsesforestillingen afløses ved Begrebet om en
evig nødvendig Væsensudfoldelse.

Naar saa \emph{Naturen} i Gud gjøres til en \emph{Al-Mulighed}, og det \emph{Intet}, af hvilket Verden
skabes, sættes som ligt med \emph{Guds Villies evige Muligheder}, saa bliver \opage{77}denne
Identificeren en Ombytten af to \emph{aldeles forskjellige} Begreber. De evige
Muligheder i den evigt virkelige Gud maae jo for Gud være evigt realiserede
Muligheder ɔ: Virkeligheder, altsaa saalangt som muligt fra Intet. Skulde de være
Muligheder, der først fik Virkelighed ved den endelige Tilværelse, saa vilde deri
ligge Consequentser med Hensyn til Guds Forhold til denne Tilværelse, som Martensen
allermindst vilde kunne acceptere.

Den Indvending, at Skabelsen ikke sætter et \emph{Intet}, men et \emph{Noget}, der
\emph{frembringer} og udvikler sig selv i \emph{given Selvstændighed}, ophæver sig
selv ved sin umiddelbare Forbinden af modsatte, hinanden negerende, Bestemmelser. Af
hvilken Betragtning af den naturlige Tilværelse denne Indvending er fremgaaet, vil
blive belyst i det Efterfølgende.

Indvendingen mod den \emph{anden} ovenfor dragne Consequents udføres hos Martensen
paa følgende Maade\footnote{\emph{Martensen}: Om Tro og Viden. S.\ 59 f.}. Det skal være
aldeles overfladisk og uberettiget at føre Skabelsen tilbage til en absolut
Vilkaarlighedsact af den blotte Almagt og ifølge deraf at sætte den som det
Ubegribelige. Skabelsen skal ubetinget forudsætte det ethiske Gudsbegreb, og uden
dette skal det være aldeles urimeligt at tale om Skabelse; der skal ikke være Tale
om den blotte Almagt --- en reen physisk Bestemmelse --- men om den ved Kjærligheden
og Viisdommen, altsaa ethisk, bestemte Almagt. I Kjærlighedens ethiske „Nødvendighed“
skal Skabelsens Grund søges; gjennem den skal Skabelsen blive begribelig, vel ikke
saaledes, at man prætenderer at udgrunde dens Hvorledes, men saaledes, at den sees i
sin nødvendige Sammenhæng med Guds Væsens Bestemmelse.

Men denne Paaberaaben af det Ethiske bliver fuldstændig unyttig, og maatte allerede
af den Grund blive det, at, som ovenfor paaviist, det Ethiske i Martensens Opfattelse
ikke i Virkeligheden faaer det Physiske til sit Moment. Her \opage{78}er Tale om en Skabelse af et \emph{materielt} Existerende: i hvilken Forbindelse staaer vel den med det Ethiskes
Bestemmelse? Martensen \emph{postulerer} Forbindelsen, men gjør end ikke et Forsøg paa at
vise, hvad der er det Forbindende, og hvad der, i det ethiske Gudsbegreb, skal kunne
forklare en \emph{saadan} Skabelses \emph{Mulighed}. Imod det Postulat, at den ethiske Bestemthed
i Gud skal være det Forklarende, gjælder det, at hvad der i Sandhed skal være ethisk,
maa være bestemt ved \emph{Guds Væsen}; udelukkedes altsaa Naturmomentet fra Guds Væsen,
vilde en Skabelse af det Reale ikke have noget positivt Forhold til det Ethiske. Og
ovenfor blev det viist, at Naturmomentet i Virkeligheden blev udelukket, og at man maatte
udelukke det naar man ikke, i Stedet for en \emph{Skabelse}, skulde faae en \emph{Naturudvikling}.
Men naar derfor Naturen i Gud reduceres til et Uvirkeligt, til et blot
Phantasiebillede, saa faaer man intet Tilknytningspunkt for Muligheden af en Skabelse
af det Materielle; men \emph{hvor Muligheden mangler}, er \emph{det meningsløst at tale om Nødvendigheden}. Selv om man vilde lade Kjærligheden, der dog, efter Martensens
Deduction, allerede skal have sit virkelige Object, altsaa sin reale
Tilfredsstillelse, i Treenighedsforholdet, være Forklaringsgrund for en Sætten af
skabte \emph{Aander}, saa vilde man intet Skridt være kommen videre med Hensyn til
Skabelsen af \emph{det Materielle}. Og en saadan Forestilling om en rent aandelig Skabelse
vilde ophæve sig selv. Det ved en Skabelse Satte maatte være \emph{et Endeligt}; men et
Endeligt er \emph{et ved Tid og Rum Bestemt}, og Tid og Rum, der ikke kunne tænkes uden
hinanden, kunne kun tænkes hvor et \emph{Materielt} er. En \emph{Skaben} af et \emph{rent Aandeligt}, en
Existeren af \emph{skabte Aander}, i virkelig Sondring, med virkelig Begyndelse, med en
virkelig Livs-Historie uden den reale Tilværelses Grundformer, vilde altsaa være
selvmodsigende. Hvorledes end Skabelsens Virkning forestilles, bliver Skabelsen lige
ubegribelig.\footnote{Det \emph{forklarer} naturligviis Intet, at der siges (i Dogmatiken)
at „Logos løslader af sit aandelige Dyb den hele Mangfoldighed af Livskræfter til
Selvbevægelse“.}

\subsection{Underet}
\label{ssec:underet}

\opage{79}Hvad der gjælder med Hensyn til \emph{Skabelsen}, det \emph{universelle Under}, maa ogsaa
have Gyldighed med Hensyn til \emph{Underet}, den \emph{partielle Skabelse}.

Analysere vi Forestillingen om Underet, ville vi altsaa faae Resultater, analoge med
dem, til hvilke Analysen af Skabelsesforestillingen førte. Underet, som Virkning af
den guddommelige, absolut frie, almægtige Villie, er en Virkning, i hvilken det
\emph{Naturliges Intethed} sættes. Den sættes idet Almagten i Underet virker \emph{uden alt Realt}, og idet i Underet det existerende Reale sættes \emph{som et Intet}. Kun tilsyneladende
bruger den undergjørende efter sit Begreb \emph{enevirkende} Almagt Midler, kun
tilsyneladende har det ved Underet Tilblevne en real Forudsætning. Det tilsyneladende
Middel er intet Middel, den tilsyneladende Forudsætning er ingen virkelig
Forudsætning, thi der er \emph{ingen Væsenssammenhæng} mellem Midlet og Maalet, mellem
Forudsætningen og det, der følger efter; \emph{hints Bestemthed er i Forhold til dettes Bestemthed} et \emph{Intet}. Underet er \emph{en ny Skabelse} og som Skabelse en Skabelse af Intet:
kun det, at det er den \emph{partielle} Skabelse, giver \emph{Skinnet} af en \emph{virkelig} Forudsætning.
For en ny \emph{Skabelse} bliver det Foregaaende et \emph{Stof}, hvis \emph{Væsen negeres i Skabelsens Øieblik} ɔ: det bliver \emph{et Intet i Forhold til det Nye}; det, der gaaer forud, bliver,
med andre Ord, ikke en positiv Forudsætning, men et Saadant, der kun negeres forat
give Plads for det Nye. --- Af Underets Forhold til Almagten følger saa, at \emph{Underet, ligesom Skabelsen, maa være det Ubegribelige, en Forklaring af det en Umulighed}.
Herover er ogsaa den religiøse Bevidsthed sig klar: for den er Underet Almagtens
umiddelbare Virkning („han talede og det skeete, han bød saa stod det der“), og det, der
overstiger \emph{al} menneskelig Tanke.

\opage{80}Mod den angivne Opfattelse af Underet gjør Martensen Indsigelse fra det samme
Synspunkt, fra hvilket han gjorde Indsigelse mod de ovenfor givne
Begrebsbestemmelser med Hensyn til Skabelsen. \emph{Underet} skal ikke være uden en \emph{Naturforudsætning}, der i selve \emph{Underet} faaer en \emph{relativ} Gyldighed, og \emph{Troesvidenskaben} skal \emph{kunne forklare}, vel ikke \emph{Underets Hvorledes}, men \emph{Underets Mulighed overhovedet}, idet Underet sees i Sammenhæng med Guds Aabenbarings
Oeconomie og den frit sig aabenbarende Gud sees som „Naturens Herre“, der „kan
hvad Naturen ikke kan“. Underets relative Begribelighed hviler altsaa, ligesom
Skabelsens, paa dets Forhold til det \emph{ethiske Gudsbegreb}, nærmere: paa dets
Sammenhæng med den \emph{ethiske Teleologie}.

Forat hævde Underet en \emph{Naturforudsætning}, fastholdes det, at den ethiske
Verdensteleologie, i Sammenhæng med hvilken Underet skal opfattes, ikke skal
ophæve de lavere Tilværelseskredses Autonomie og Selvstændighed, eller modsige
Naturvidenskabens Forudsætning om en lovbestemt Tilværelsens Orden. En ethisk
Verdensanskuelse erklæres for aldeles umulig uden en \emph{relativ selvstændig} Physis,
og selve Underet for umuligt uden et \emph{ordnet System} af \emph{Naturlove}, da det netop skal
forudsætte en Naturens Orden for deri at kunne fremtræde som Underet. Ligeledes
fordres der med Hensyn til Underet, at den høiere guddommelige Virksomhed altid maa
være \emph{logisk betinget}, idet Gud ikke kan modsige sig selv, hvori ligger, at Guds
undergjørende Virken ikke maa være tilintetgjørende i Forhold til det ved Gud satte
Naturlige; hvilket ogsaa indirecte er fordret derved, at der statueres physiske
Betingelser for den guddommelige Villies Virken. --- Men disse \emph{almindelige} Bestemmelser vise sig ikke blot som \emph{uigjennemførlige} saasnart det \emph{bestemte}
Under drages frem; men de faae \emph{nærmere Bestemmelser}, ved hvilke de atter \emph{ophæves}.
Underet sættes i Forhold til det \emph{Teleologiske}. Men det, der betegnes som \emph{Under}, bestemmes som det, \emph{der ikke} kan udledes af de \emph{foregaaende Skabelsestrin}: det \opage{81}er et
specifisk Høiere, ikke blot i den Betydning, i hvilken en teleologisk Udviklings mere
fremskredne Udviklingsstadier ere et Høiere i Forhold til de lavere, og i Virkeligheden
ere det, ved hvilket disse skulle begribes, men i den Betydning, i hvilken det \emph{Overnaturlige} er høiere end det Naturlige. \emph{Idet Underet derfor træder frem, er
Teleologien brudt, Naturens Orden negeret}; det Forudgaaende er et Intet i Forhold til
Underet og dette træder frem som det rent Umiddelbare og Vilkaarlige. Det er rigtigt, at
Underet, som dette \emph{enkelte}, i Tiden begrændsede Under, i en \emph{vis Forstand} uvilkaarligt
forudsætter, \emph{om end kun forat ophæve den}, en \emph{Naturtilværelsens Orden}, forsaavidt
det forudsætter \emph{en Continuitet}, \emph{der brydes ved Underet}. Var der ikke en \emph{relativ}
Continuitet, vilde Underet ikke kunne blive \emph{synligt} som Under; thi idet Underet sattes i \emph{alle} enkelte discrete Tidspunkter, vilde Underet blive \emph{Loven}, selve Tilværelsens Orden:
\emph{Underets uendelige Discretion vilde blive Lovens uendelige Continuitet}. Men hin
relative Continuitet betegner ikke en i Naturens substantielle Væsen begrundet Orden;
den er ikke et Udtryk for selve Substantialiteten, men er kun Udtrykket for et ved den
skabende Villie vilkaarligt Sat og \emph{vilkaarligt Gjentaget}. Det er netop et
„\emph{Skabelsestrin}“, forat bruge Martensens Udtryk. Det ophæves ved Underet og sættes atter
vilkaarligt: Guds Frihed er for den religiøse Bevidsthed kun absolut, idet den ubetinget
er ubunden ved Alt, ogsaa ved de af den selv satte Forudsætninger. \emph{Consequent} vilde denne
Bevidsthed gjøre \emph{Alt} til Under; at den, i sin Umiddelbarhed, ikke gjør det og ikke kan
gjøre det uden at lade Underet blive usynligt, viser, at den aldrig udfylder
Menneskeaanden heelt, at den har den \emph{naturlige Bevidsthed} ved Siden af sig, og kun er ved
bestandig at negere denne.
Fra denne trænger oprindelig Forestillingen om det relativt Bestandige i Naturen ind i
hiin; men idet det relativt Bestandige sees som det, der \emph{kan} gjennembrydes af Underet,
har det allerede ophørt at være Natur\emph{lov}, og naar den \opage{82}\emph{religiøse} Bevidsthed seer det \emph{i sin egen Belysning}, bliver det for den religiøse Reflexion \emph{det af Almagten vilkaarligt Satte, altsaa Under}, ligesom det specifiske, iøinefaldende Under, saa at Underet \emph{atter} omfatter \emph{hele} Tilværelsen, men med det Samme bliver \emph{usynligt} som Under, hvorfor Bevidstheden
uvilkaarligt føres tilbage til hint Udgangspunkt, men kun for paany at gjennemløbe det
samme Kredsløb.

Hvad \emph{Underets Begribelighed} angaaer, saa forsøges den hævdet, deels gjennem en Polemik
mod Forestillingen om Naturlovenes Uforanderlighed, deels gjennem Tilknytningen til det
Ethiske, gjennem Bestemmelsen af det Naturliges tjenende Forhold til dette. Forat skaffe
en Plads for Underet, søges Forestillingen om \emph{Naturlovenes Uforanderlighed} hævet. Men
idet Opfattelsen af Naturlovene som uforanderlige søges gjendreven, træder der en
iøinefaldende Uklarhed frem, baade med Hensyn til Forskjellen mellem en Forandring af
Naturlovene og de samme uforandrede Naturloves forskjellige Virken under forandrede
Betingelser, og med Hensyn til Forholdet mellem Naturlovene og de mathematiske Love, hvis
Uadskillelighed Mechanikens Love viser. Forsøget paa en Gjendrivelse af hiin Forestilling
støtter sig væsentligen paa, at den Vei, ad hvilken vi komme til Bestemmelsen af
Naturlovene, maa bestemme deres Charakteer, og da Veien er Erfaringens, saa skulle
Naturlovene ikke have anden Gyldighed end den, Erfaringen kan sikkre dem. Idet
Naturlovene bestemmes ad Erfaringens Vei, gjennem Inductionen og Analogien, kunne de,
siges der\footnote{\emph{Martensen}: Om Tro og Viden. S.\ 108 ff.}, kun faae relativ
Almindelighed og Nødvendighed, deres uforanderlige Gyldighed kan aldrig faae Mere end den
høieste Grad af Sandsynlighed. Og de vilde, selv om de indenfor denne bestemte Kreds
havde ubetinget Gyldighed, som Erfaringslove dog netop kun have Gyldighed i en given
Facticitet, indenfor denne bestemte, denne nærværende Existents, men ikke have det for \emph{al} Tilværelse. \opage{83}Naturvidenskaben kan, som blot Erfaringsvidenskab, ikke afgjøre Noget om det
i absolut Forstand Mulige eller Umulige; den har end ikke, med Hensyn til den
Virkelighed, til hvilken den forholder sig, en Erkjendelse, der omfatter Begyndelse og
Ende: disse ere for den indhyllede i et uigjennemtrængeligt Mørke; den bevæger sig kun
ligesom i en Midte, indenfor en Kreds af Betingelser, der nu engang ere factisk givne,
den veed ikke engang selv hvorledes. Den kan derfor ikke modsige, at der kan gives en
høiere Existents end Naturens, der kan gjennemtrænge denne og indvirke paa den, det være
sig forløsende eller forklarende, sætte den som Organ og Middel for et Maal, der er
høiere end Naturens eget. Da Underet ikke refererer sig til de rene Begrebers og de rene
Formers Verden, men til Existents, er det netop Spørgsmaalet, om Naturen selv er den
høieste Existents, der som saadan ikke maa kunne omskiftes, eller om hiin høiere Existents
gives. Kun ud fra et physisk Gudsbegreb som det høieste kunde Naturlovenes
Uforanderlighed godtgjøres.

I denne Argumentation er det \emph{først} overseet, at det, at Naturlovene kun \emph{ved Hjælp af Erfaringen} kunne komme til vor Kundskab, ikke ophæver Muligheden af en
Begriben, der giver dem en høiere Nødvendighed end den \emph{ad den blotte Erfarings} Vei opnaaelige; med andre Ord: ikke udelukker Muligheden af, at den \emph{empiriske} Lov kan forvandles til en \emph{rationel}, en Forvandling, der netop er Naturerkjendelsens Opgave, og
hvis Mulighed det tilkommer Underets Forsvarere at modbevise. Naturerkjendelsen er og
bliver ganske vist altid en \emph{Erfaringserkjendelse}, men ikke i den Betydning, at den
skulde hvile paa en \emph{isoleret} Erfaring; den tager sit Udgangspunkt i Erfaringen for
fra det at naae til en Begriben, hvis Mulighed er indrømmet naar der tales om Fornuft i
Naturen, om, at Naturen gjør Tankens Gjerninger. --- Naar det \emph{dernæst} sluttes, at
Naturlovene, som kun gyldige i en given Facticitet, indenfor denne bestemte Existents,
ikke kunne have ubetinget Uforanderlighed og Nødvendighed, saa bliver her \opage{84}deels den
ovenfor berørte Forvexling mellem en Forandring i Naturlovene og i Betingelserne at
fjerne, deels bliver det at urgere, at Slutningens Gyldighed afhænger af det Forhold, der
sættes imellem denne Existentsform og den supponerede høiere. Er denne --- hvad
Udtrykket: \emph{Verdensteleologie}, viser hen til --- et høiere \emph{Udviklingstrin},
saa beholde Lovene, under de nye Virksomhedsbetingelser, under de nye Complicationer,
uforandret deres Gyldighed, saaledes som vi see det Uorganiskes Love beholde den indenfor
det Organiske; de gjælde for al Existents ɔ: for al \emph{Naturexistents}, og vi faae i det
Høiere ikke et Under, idethøieste et \textit{mirabile}. \emph{Naturens primitive Love
maae, da Naturen altid maa tænkes som Heelhed, paa alle Naturudviklingstrin være de
samme; thi Lovene ere netop Udtrykket for det udfoldede, i Virksomhed tænkte Væsen}.
Forat komme til Ophævelsen af Naturlovens Gyldighed maatte man forudsætte, at det Høiere
var kommet frem \emph{ved et Brud}, ved en absolut Ophævelse af Continuiteten. Men saa vilde det
Høiere være blevet til ved en ny \emph{Skabelse}, og Forholdet vilde --- hvad der jo netop skal
undgaaes --- være \emph{ubegribeligt}. Naar det Høiere, Underet, ikke skal tænkes som et nyt
Udviklingstrin, men som et Omskabende, saa fører det til den Consequents, at Underet, den \emph{partielle} Skabelse, skal træde ind i en Virkelighed, der kun paa dette Punkt forandres,
og at altsaa denne Virkelighed, der saaledes partielt forandredes, vilde blive det
Modsigende: en defect Virkelighed, en ufuldstændig Heelhed, en usamstemmende Harmonie.
Statueres der „et ordnet System af Naturlove“, saa er Bruddet paa dette ene Punkt et Brud
paa Systemets Heelhed. Og selv om man nærmest seer hen til det enkelte Reale, der i
Underet negeres, saa er en Tilintetgjørelse af selv det ubetydeligste Reale, en Ophævelse
af det \emph{imod} Loven, et Brud paa \emph{hele} Tilværelsens Orden, en Forstyrrelse af Heelheden.
Og \emph{hiin Forandring i Existentsen griber fra Existensen ind i de rene Begrebers og de
rene Formers Verden.} Man analysere \opage{85}hvilket Under man vil, det vil bringe en Modsigelse
ind i Begrebernes Sphære, idet det er en Virken uden Midler, en Foraarsagen af hvad der
negerer sin Aarsags Væsen, en Sætten af et Realt uden alt Realt, en Sætten, der negerer
det ved den evige Virksomhed Satte; og det vil, idet al \emph{physisk} Existents er \emph{mathematisk} bestemt, bringe en Modsigelse ind i de rene Formers: det Mathematiskes,
Verden; thi det er en Sætten af et Noget, et ved Quantum Bestemt, ud fra et Intet, af et
Formbestemt ud fra det Formløse. Det bliver ikke blot \emph{det kun physiske} Gudsbegreb, ud fra
hvilket Underet negeres, og \emph{Naturlovenes Uforanderlighed}, der kun er Udtrykket for Naturens \emph{Væsentlighed}, hævdes; saasnart der sættes \emph{en Naturgrund, en Physis} i Gud, er
Underet sat som umuligt. Og negeres Naturmomentet i Gud, aabenbart eller skjult, saa er
Underet atter, om end af en anden Grund, det Umulige, det rent Ubegribelige, ligesom
Skabelsen. Fordi Underet nægtes, sættes ikke Naturen \emph{som} Natur som det Høieste, men den
høiere Væren end Naturen: Aandens Væren, sees som Naturens Fuldbyrdelse, ikke som Naturens
vilkaarlige Negation. De Love, der, som gjældende paa et \emph{bestemt}
Naturudviklingstrin, gjælde for Naturen \emph{overhovedet}, gjælde for Naturen som den af Aanden
besjælede Natur, som den Natur, der har sit Maal i Aanden. Naturen er Aandens blivende,
væsentlige Grundlag.

Fra det \emph{negative} Beviis for Underets Begribelighed vende vi os til det \emph{positive}, det støtter sig paa det \emph{ethiske Gudsbegreb}. Ved dette er der, siges
der\footnote{\emph{Martensen}: Om Tro og Viden. S.\ 111 f.}, ikke blot sat et levende
Vexelforhold mellem Gud og Skabningen, men Verdensteleologien er tillige bestemt som
ethisk Teleologie; d.\ v.\ s. det Ethiske er, som Grundformaalet, saaledes ogsaa det
inderste Grundbestemmende i Tilværelsen, det, for hvis Skyld alt Andet er til, og som alt
Andet maa tjene. Dette skal dog ikke betyde, at den ethiske Teleologie paa \emph{alle} Punkter
af Tilværelsen aldeles vilkaar\opage{86}ligt og rent umiddelbart skulde træde frem, og Alt uden
videre kunne stemples med det ethiske og religiøse Afhængighedsforhold. Som ovenfor
angivet skal den ethiske Verdensteleologie netop forudsætte og stadfæste en relativ
selvstændig Physis, en lovbestemt Tilværelsens Orden. Meningen skal være, at den ethiske
Verdensanskuelse og det ethiske Gudsbegreb nødvendigen medføre, at Naturen i dens
Indretning ikke kan umuliggjøre den guddommelige Verdensplan, eller at Naturbegrebet ikke
kan fornægte Skabelsesbegrebet, at Naturen ikke kan fornægte sin Afhængighed af Skaberen.
For den ethiske Verdensanskuelse skal Naturen være et Mellemvæsen, der har den Betydning,
ikke blot at være sig selv og at fuldbyrde sine egne Formaal, men tillige at være Organ,
Middel for Friheden, og at tjene Formaal, der ere langt høiere end dens egne. Den skal
være modtagelig for Frihedens Indvirkninger, om hvilke det siges, at det, naar Frihedens
Gud forudsættes, vilde være høist besynderligt at indskrænke dem til den menneskelige
Friheds. Og naar nu den guddommelige Verdensplan førte det med sig, at der i
Menneskeslægtens Historie indtraadte Perioder, bestemte Tider, i hvilke Gud i sit Forhold
til Menneskeslægten paa prægnant og aldeles utvetydig Maade vilde aabenbare sig som den,
han altid er, som Naturens Herre, for gjennem Virkninger, der ikke kunne forklares af det
Skabte, at bringe \emph{sig selv} den menneskelige Bevidsthed nærmere, fordi Indholdet af hans
Verdensplan var en Huusholdning til Menneskets Frelse, da vilde saadanne Tider i særegen
og strængeste Betydning være \emph{Underets Tider eller de store Aabenbaringstider}. Og
Aabenbarings- eller Religionsvidenskaben maatte søge en nærmere Forstaaelse, søge at see
disse Underes Betydning i og med den store Totalsammenhæng, naturligviis uden at ville
begribe Underets Hvorledes.

Ved Henvisningen til det ethiske Gudsbegreb er dog Intet vundet. Dette ligger allerede i
hvad der blev udviklet med Hensyn til Skabelsen; det bliver, naar der tages \opage{87}Hensyn til
den Form, i hvilken Argumentationen fremtræder hos Martensen, tydeligt paa følgende
Maade.

Aandens Principat i Forhold til Naturen; det, at Naturen tjener Aanden og er
modtagelig for Frihedens Indvirkninger, gjør ikke Underet begribeligt. Naturen tjener
Aanden \emph{som} Natur, ɔ: i \emph{Overensstemmelse med sine Love}, og Frihedens Indvirkninger
ere ikke, som Underets, vilkaarlige og rent umiddelbare, men \emph{gjennem Naturmidler}, ɔ:
atter i Overensstemmelse med de erkjendte fornuftbestemte Naturlove. Idet Naturen tjener
den menneskelige fornuftige Frihed, bøier den frie Aand sig under den uendelige Fornuft,
der paa een Gang er dens egen og Naturens Væsen; kun den phantastiske Frihed drømmer om
et magisk, umiddelbart Herredømme over den underkuede Natur. Naturen, der i det Levende
umiddelbart viser det Ideelles Magt, fornægter ikke sin Bestemthed ved det Ideelle, der
netop er dens Væsen; men Naturbegrebet maa fornægte Skabelsesbegrebet, altsaa ogsaa
Underbegrebet, fordi Skabelsesbegrebet fornægter Naturbegrebet, negerer al Naturens
Realitet. Den fornuftbestemte Friheds Virkeliggjørelse gjennem den fornuftbestemte Natur
indeslutter ingen Vanskelighed for Tanken, men en Guds Aabenbarelse som \emph{Naturens Herre} i den angivne Betydning, som den, der \emph{kan} hvad \emph{Naturen ikke} kan, og, vel at mærke,
\emph{kan det indenfor Naturen, i Forhold til et Naturligt, uden alt Naturligt}, er det
Ubegribelige.

Underet kan dernæst ikke blive forklarligt ved nogen historisk Sammenhæng, thi
\emph{ligesom Underet ikke har nogen Natur, saaledes har det ikke heller nogen Historie}.
Det er, som Udtryk for den ubetingede, almægtige Villie, det absolut Pludselige, Discrete,
Sammenhængsløse, dets Begriben ud fra en Totalsammenhæng er umulig, er en Modsigelse.

Underet bliver endelig ikke heller begribeliggjort ved at skulle være Betingelse for Guds
utvetydige Aabenbaring. Gud skal i Underet aabenbare \emph{sig} paa \emph{aldeles utvetydig} Maade, ved
Virkninger, der ikke kunne forklares af det \opage{88}Skabte: Underet skulde saaledes \emph{skabe
Tro}. Men, efter Martensens egne Præmisser, kan det ikke \emph{vides}, at denne Virkning, der
foregaaer i Naturen, ikke kan have Naturaarsager. Erfaringen, der, efter Martensens
Bestemmelse, er Naturkjendelsens \emph{eneste} Kilde, kan \emph{ikke} sige det: da den jo Intet
skal vide om Mulighed eller Umulighed. Men paa hvilken Maade bliver saa Guds Aabenbaring
utvetydig? Og hvorledes er det at forstaae, at Gud i denne bestemte Virkning skal
% PRINTED AS IS, p.88 [collation 2026-09-21]: page prints "Bevidstbed" (closed-bowl b, no h tail, clear at 3x; wrong sort); confirmed by a blind second reader.
aabenbare \emph{sig selv} klarere end i sine Naturvirkninger, bringe \emph{sig selv} den menneskelige Bevidstbed nærmere end gjennem sin Væsensaabenbaring i Naturen og Fornuften? Det er kun
at forstaae saaledes, \emph{at Underet utvetydigt aabenbarer Gud naar Mennesket ved et
Under er ført til at see det som Under}. Det vil med andre Ord sige: \emph{Underet anerkjendes kun som Under gjennem Troen: gjennem et Under}; det forudsætter Troen: det \emph{primitive} Under, Underet i os, ved hvilket vi først see Underet \emph{udenfor os}; ikke det
ydre Under, men den underfulde Virkning i Aanden bliver altsaa det Afgjørende med Hensyn
til den utvetydige Aabenbaring. At Underet kun opfattes ved et Under, ikke forstaaes ved
Tanken, ligger ogsaa i hvad Martensen siger, at den søgte Forstaaelse af Underet er
utænkelig uden det for Aabenbaringen modtagelige Sind; thi Aabenbaringen og dens
Modtagelse er jo netop et Under.

Idet Underet saaledes forudsætter sig selv, og er ligesom den eneste Vei til sig selv,
maa ethvert Forsøg paa at begribeliggjøre det, ogsaa i hiin ubestemtere Betydning,
mislykkes, og Forsøget paa, i denne Betydning at begribeliggjøre det, fører uundgaaeligt,
idet Underet skal føies ind i en teleologisk Sammenhæng, der tillige bliver en
Natursammenhæng, til uvilkaarlige Forsøg paa, ogsaa at bestemme dets \emph{Hvorledes}, hvad der
var udelukket ved Forudsætningerne.

\begin{center}---\end{center}

\subsection{Den almindelige Anskuelse af Tilværelsen}
\label{ssec:alm-anskuelse}

\emph{Den almindelige Anskuelse af Tilværelsen}, til hvilken Opfattelsen af Skabelsen og
Underet viser tilbage, er \opage{89}udtalt i Dogmatiken. Det sættes som Opgaven, at tilveiebringe en \emph{christelig} Opfattelse af Tilværelsen, der dog ikke maa komme i Strid med \emph{Naturen og Fornuften efter deres Sandhed}. Synspunktet for denne Opfattelse af Tilværelsen skal
tages fra det, der er Tilværelsens Høieste, dens τέλος, fra Christi Person, der er den
hellige Histories Midtpunkt og Christendommens egentlige Indhold. Den hellige Historie
skal forklare Historien, Historien Naturen. Men ved Forestillingen om en hellig Historie
og særlig om Christus vises der hen til Underet, til en absolut ny Begyndelse, en aandelig
Skabelse i Menneskeslægten, der ikke blot skal have ethisk, men ogsaa kosmisk Betydning,
den dybeste Naturbetydning. Incarnationens Under, og det dertil knyttede Inspirationens
Under, ved hvilket en ny Sands oprinder i Menneskeslægten, en ny Begyndelse sættes i
Menneskeslægtens Bevidsthed, er Christendommens Grundundere, som den supranaturalistiske
Verdensbetragtning skal hævde mod Naturalismen og Rationalismen, men saaledes, at Naturen
og Fornuften skee deres Ret, at ingen forstyrrende Dualisme bringes ind.

Med Hensyn til Naturen sættes det da som den christelige Anskuelse, hvad der ovenfor
udvikledes: at den ikke er et i sig evigt og afsluttet System af Love og Kræfter, men et
System, der befinder sig i en fortsat teleologisk Udvikling, \emph{en fortsat Skabelse}, saa at
der maa kunne tænkes at indtræde nye Potentser, nye Love og Kræfter, hvis Indtræden er \emph{forberedt} og \emph{præfigureret} ved de foregaaende Skabelsestrin, men som \emph{ikke kunne afledes} af disse. Christi Aabenbarelse, som den sidste Potents i Skabelsesværket, skal derfor
tænkes at indvirke \emph{teleologisk forandrende} og bestemmende paa de lavere Potentser,
idet disse ikke have en evig og afsluttet Betydning i sig, men kun have en timelig
Mellembetydning. Enhedspunktet af det Naturlige og Overnaturlige skal saaledes ligge i
Naturens teleologiske Bestemmelse for Guds Rige og i den derved givne Modtagelighed og
Dannelighed for den overnaturlige Skabervirksomhed. \emph{Naturen skal ikke for\opage{90}nægte
Skabelsesbegrebet, men ligesaalidt skal Skabelsen fornægte Naturbegrebet}. Skjøndt nemlig
den nye Skabelse i Christus er en Ophævelse af \emph{denne} Naturs Love, saa skal den dog
ingenlunde være en Ophævelse af \emph{selve Naturens} Begreb. Thi det er netop \emph{Naturens Begreb} at
være, ikke den hæmmende Skranke for Friheden, men \emph{Frihedens Organ}.

Ved Friheden vises der hen til \emph{Historiens} Rige, men indenfor dette er allerede
Gjentagelsen antydet af det samme Forhold, der fandt Sted i Naturens Rige. For \emph{Rationalismen} er Historien Udfoldelsen af Menneskenaturen, af Menneskevæsenets
% PRINTED AS IS, p.90 [collation 2026-09-21]: page prints grave on the omicron of λὸγος (both words carry grave); acute expected; confirmed by a blind second reader.
naturbestemte Anlæg; den almindelige Fornuft (κοινὸς λὸγος), er for den den eneste
Erkjendelseskilde; den anerkjender ikke andre Sandheder end dem, der ere Explicationer af
Menneskeslægtens medfødte Fornuft. Christendommen derimod maa fordre, at der til den nye
Livskilde, der er opladt i Christus, svarer en ny Erkjendelseskilde, et Rige af
guddommelige, hidtil forborgne Raadslutninger, et Rige af nye Erkjendelser, der \emph{ikke} lade sig aflede af Begrebet om Menneskeslægtens fremskridende Fornuftudvikling. Men ligesom
Underet ikke stred imod Naturen, saaledes skulle disse nye, overnaturligt meddelte
Erkjendelser ingenlunde \emph{stride} mod Menneskeslægtens almindelige Fornufterkjendelser,
skjøndt de paa mange Maader \emph{nærmere bestemme} disse; thi deels forholde de sig
\emph{fuldkommende og afsluttende} til de egentlige Fornufterkjendelser, deels
\emph{forløsende}, idet de befrie dem fra det Mørke, hvormed de, paa Grund af den
altomsluttende Syndighed, ere behæftede. Der skal ved dem ligesaalidt være statueret en
Dualisme som ved Læren om de tvende Skabelser: der er kun \emph{eet Fornuftrige}, skjøndt der i dette er \emph{tvende Potentser} af Fornuftens Aabenbarelse. Objectivt skal Enheden ligge
deri, at det er \emph{den samme Logos}, der aabenbarer sig i de tvende Skabelser,
saaledes, at Logosaabenbarelsen i Christus er Aabenbarelsen i høiere Potents, derved
forskjellig fra den universelle Logosaabenbarelse, at den er den verdensfuldendende og
verdensforløsende \opage{91}Aabenbarelse, medens hiin universelle blot er den verdensskabende og
opholdende. \emph{Subjectivt} skal Enheden ligge deri, at den menneskelige Fornuft er receptiv
for Christi Aand som Verdensfuldendelsens og Verdensforløsningens Aand: en Receptivitet,
gjennem hvilken Fornuften skal ophøies til et høiere Trin af Productivitet. Det, at
Aabenbaringen i Christus skulde stride mod Fornuftens Love, skal kun indrømmes i den
Betydning, at den, paa samme Maade som de moralske Love ophæves i deres abstracte
Selvstændighed forat stadfæstes i Kjærligheden, ophæver Fornuftlovene i deres Abstraction
forat stadfæste dem i Christus-Viisdommen, der er Fornuftlovens Fylde.

Det Forhold, denne „christelige Anskuelse“ af Tilværelsen statuerer, kan altsaa i Korthed
bestemmes saaledes: \emph{Underet er Naturens høiere Potents, Aabenbaringen Fornuftens.
Underet er præfigureret og forberedt i Naturen, Aabenbaringen i Fornuften, men de første
kunne ikke udledes af de sidste.} Forholdet har saaledes en \emph{Dobbelthed} i sig: paa den ene Side vises der hen til den \emph{høiere Enhed}, paa den anden Side ophæves atter Enheden og \emph{Bruddet} indtræder mellem det ved en \emph{ny Skabelse} satte Høiere og det forudgaaende Lavere. Uigjennemførligheden af denne Dobbelthed er allerede paaviist ved Læren om
Skabelsen og Underet; hvad der der er udviklet med Hensyn til Skabelsens Begreb, der
omfatter Begrebet: „ny Skabelse“, „fortsat Skabelse“, og med Hensyn til Teleologiens
Forhold til den absolut nye Begyndelse, finder, ligesom det, der blev udviklet om Mangelen
paa en Tilknytning for det Physiske i det „ethiske Gudsbegreb“, om Naturlovenes
Uforanderlighed, om Frihedens Forhold til Naturen o.\ s.\ v., ogsaa her sin Anvendelse. De
for den givne Udvikling eiendommelige Vendinger føre ikke nærmere til det tilstræbte Maal.
Naar der saaledes tales om en \emph{Naturens Modtagelighed og Dannelighed for det
Overnaturlige}, og i Analogie hermed om den \emph{almindelige Menneskefornufts
Modtagelighed og Dannelighed for \opage{92}den overnaturlige Aabenbaring}, saa forudsætter
Modtagelighed og Dannelighed \emph{Væsensenhed}: kun hvad der stemmer med mit Væsen, kan
jeg modtage; thi Modtagelighed er uadskillelig fra en til det Modtagne svarende
Selvvirksomhed og Selvvirksomheden udtrykker Væsenet. Uden Selvvirksomhed vilde
Modtagelsen være en Tilintetgjørelse af det Modtagende. Naar der derfor tales om det
Naturliges Modtagelighed for det Overnaturlige, Fornuftens for Aabenbaringen, saa maae vi
enten sætte Forskjellen som en blot quantitativ, eller, naar vi ville fastholde den som en
afgjørende Qualitetsforskjel, lade Modtageligheden blive et tomt Ord uden Betydning, lade
den „teleologiske Tilbagevirkning“ ɔ: den bagefter satte Teleologie være en virkelig
Omskabelse, Aabenbaringens Optagelse i Jeget et Brud i Jeget, en Omskabelse, der gjør
Selvbevidsthedens Identitet ubegribelig.

Det er overhovedet charakteristisk for den givne Udvikling, at Bestemmelserne lide af
Tvetydighed og Vaghed. Naar der saaledes siges om det Høiere, der \emph{ikke} kan udledes af det
Lavere, at det er „\emph{forberedt og præfigureret}“ ved dette, saa ligger der i „Forberedelsen“
en \emph{objectiv Sammenhæng}, i „Præfigurationen“ en \emph{blot billedlig, blot subjectiv
gyldig}. Ved Forbindelsen af disse Udtryk holdes det altsaa svævende, af hvad Art
Forholdet er; det bliver paa een Gang et realt og et ikke realt Forhold, paa een Gang
indre Sammenhæng og Brud paa Sammenhængen. Og naar der tales om det Laveres
\emph{Potentsation og Stadfæstelse i Potentsationen}, og om de høiere Potentsers \emph{Tilbagevirkning}, saa faae disse Udtryk kun en Betydning for Phantasien saalænge de holdes
i det Almene, ikke forsøges bragte i Forhold til det Enkelte og Concrete. Forandringen og
Bestemmelsen af de lavere Potentser skal falde sammen med Realisationen af deres
Modtagelighed for den overnaturlige Skabervirksomhed. Men denne Realisation, der maatte
tilveiebringes ved at fjerne det, der hæmmede Muligheden, maatte, da det Hæmmende kun
kunde være et Realt, have reale Følger og \opage{93}være synligt i Naturen --- hvis ikke netop
Naturen tilintetgjordes ved denne Indvirkning, hvad der jo ikke er Forudsætningen. Men man
forsøge blot i det Concrete at paavise den Forandring i Naturens Love, i de logiske og
mathematiske Love, der skulde være foregaaet ved Incarnationen, der jo skal have den
dybeste kosmiske Betydning, eller Naturkræfternes Potentsation ved det enkelte Under, og
man vil snart opdage det Haabløse i Forsøget. Og selv uden dette: naar man forsøger paa at
tænke Tanken om en defect Natur og om en defect Fornuft, der dog var Natur og Fornuft ---
en Tanke, som hiin Forestilling om høiere Potentser forudsætter, saa vil Forestillingens
indre Modsigelse vise sig.

\begin{center}---\end{center}

\subsection{Resultatet af Kritiken af Gudsbegrebet}
\label{ssec:gudsbegreb-resultat}

Sammenfatte vi \emph{Resultaterne} af Undersøgelsen af \emph{det ethiske Gudsbegrebs Forhold til det physiske}, saa blive de følgende:

\begin{enumerate}
\item Det ethiske Gudsbegreb faaer \emph{ikke} i Virkeligheden det physiske i sig som Moment; det
  bliver Begrebet om en naturløs Gud, og derfor bliver det Naturliges Tilværelse, og guddommelige Virkninger i det Naturlige, uforklarlige ud fra det.
\item Den Frihed, ved hvilken det „ethiske“ Gudsbegreb bestemmes, giver det \emph{ikke} et Værd,
  der er høiere end det, der tilkommer den efter sit sande Begreb opfattede fornuftbestemte
  Natur; denne Frihed, der kun bliver væsenløs Vilkaarlighed, bliver tværtimod et
  Lavere end den i Naturen aabenbare Fornuftnødvendighed, og den bliver et Selvmodsigende
  og sig selv Ophævende.
\item Den Opfattelse af Tilværelsen, der hviler paa dette ethiske Gudsbegreb, bliver ikke
  den høiere. Den negerer Fornuftloven i den naturlige Tilværelse; den ophæver, med Naturens \emph{Love}, Naturens \emph{Begreb}, og Historien bliver for den, hvad enten den sees
  fra Guds eller Menneskets Synspunkt, kun en nødvendighedsløs ɔ: vilkaarlig Friheds Virken. (Netop fordi denne Op\opage{94}fattelse i Menneskevæsenet seer den fra Naturens
  Nødvendighed befriede nødvendighedsløse Villie som det Første, og derfor i
  Menneskeslægtens Historie kun seer hiin Friheds Rige, idet Aanden, som „den frie“, bliver
  den af Lov ubundne, accentuerer den „\emph{det Historiske}“ fordi det, som det, der er
  unddraget Loven, danner Analogien til den guddommelige „Frihedsvirkning“, der træder frem
  i den frie Aabenbaring og i Underet. Dog faaer dette Historiske ogsaa af en anden Grund
  saa stor en Betydning for den, nemlig fordi det, ved at være det Medium, i hvilket Gud \emph{umiddelbart} træder frem og aabenbarer „Sig Selv“, tilfredsstiller den religiøse
  Bevidstheds Trang til at faae det Evige \emph{for Anskuelsen i den Sandselighedens Form}\footnote{jfr.\ O.\ \emph{Waage}: J.\ P.\ Mynster og de philosophiske Bevægelser paa hans
  Tid i Danmark, S.\ 122, jfr.\ 116, 144, 163.}, der er uadskillelig fra denne Bevidsthed
  paa den \emph{positive} Religions Standpunkt).
\end{enumerate}

\begin{center}---\end{center}

Saaledes føres vi, ved Undersøgelsen af de af Martensen givne Bestemmelser af Guds Væsen,
fra det som ethisk benævnede, men ikke i Virkeligheden ethiske, Gudsbegreb tilbage til det
logiske, opfattet efter det Logiskes sande og fuldstændige Betydning, efter hvilken det
udtrykker Naturens og Aandens Væsensbestemmelser. Og vi føres fra Forestillingen om det
Naturliges Tilbliven ved en vilkaarlig Skaberact og om den af Villiens Vilkaarlighed
fremgaaede ydre Aabenbaring tilbage til Begrebet om en i evigt Væsensforhold til det
Guddommelige staaende, paa uforanderlige Love hvilende Natur, og om den ved det
guddommelige Væsens frie Nødvendighed bestemte evige Væsensmeddelelse til den i Gud
grundende Menneskeaand.

\begin{center}---\end{center}

\subsection{Det for Theologien Typiske i den kriticerede Opfattelse}
\label{ssec:typiske}

Hvad der i det Foregaaende er udviklet i nøie Tilknytning til det i Martensens Skrifter
foreliggende theologiske \opage{95}Forsøg, bliver nu at see som gjældende \emph{for Theologien overhovedet}, idet Martensens Opfattelse, i dens afgjørende Grundtræk, paavises som \emph{typisk} for Theologien.

Forat godtgjøre hans Opfattelses typiske Charakteer, vil det ikke være tilstrækkeligt at
paavise tilsvarende Grundtræk hos betydelige Repræsentanter for Theologien til
forskjellige Tider. Ad denne Vei, paa hvilken man ved Siden af det Overensstemmende altid
vilde kunne finde det mere eller mindre Afvigende, vilde man kun kunne naae til et \emph{Approximativt}. Opfattelsens typiske Charakteer kan kun hævdes uimodsigeligt idet der gaaes
tilbage til det, der bliver Theologiens Gjenstand: den religiøse Bevidstheds Indhold, og
naar den Form søges, som dette Indhold consequent \emph{maa} antage, naar det skal blive Indhold
for en „Troesvidenskab“.

Det vil her, hvor der handles om Opfattelsen af \emph{Gudsbegrebet}, være nok at fremdrage
eet Phænomen, der ligesom er \emph{den religiøse Bevidstheds Urphænomen} og gaaer gjennem alle, selv de laveste, Religionsformer, nemlig: \emph{Bønnen}. I Bønnens Act ligger det
positivt Religiøses Eiendommelighed concentreret; Analysen af Bønnens Begreb giver de
væsentlige Momenter til Bestemmelsen af de positive Religioners Opfattelse af det
Guddommelige.

I Bønnens Begreb ligge væsentlig \emph{to} Momenter: \emph{1) Bønnen er Bøn til den personlige
Gud, hvis Villie bestemmes ved Bønnen; 2) Bønnens Gjenstand er Underet.}

1.\ Idet den Bedende beder til Gud, i Bønnen „siger Du til sin Gud“, forholder han sig til
Gud som til en \emph{Enkeltpersonlighed} udenfor sig, og i selve Bønnens Form ligger
udtrykt, at denne Enkeltpersonlighed kan \emph{bevæges ved Bønnen}, paavirkes fra Menneskets Side. Som den, til hvem der bedes, er Gud den, der ikke er bunden ved nogen
evig Lov, men har Muligheden til bestandig nye Raadslutninger, og til Raadslutninger, der
særlig forholde sig til den Enkelte, der beder, til hans Vel, hans Ønsker, hans
Lyksalighed. „Beder, saa skal Eder gives“, \opage{96}er Udtrykket for den umiddelbare religiøse
Bevidstheds Opfattelse af Bønnen. „Gud, der er god, vil give Eder gode Gaver \emph{naar I bede ham derom}“ udtrykker Bønnens betingende Magt. „Giv os idag \emph{vort daglige Brød}“
udtrykker, at Bønnen forholder sig til en bestemt Gjenstand. Selv naar der bedes: „skee \emph{din} Villie“, saa er den Villie, om hvis Opfyldelse der bedes, den Almagtens Villie, der
vil den \emph{Troendes og Bedendes Bedste}, Almagten er \emph{Godhedens Almagt}. Bedes der: „skee \emph{ikke min}, men \emph{din} Villie“, saa gjælder Bønnen atter den samme og paa samme Maade opfattede
Almagtens Villie. I denne Bøn er en begyndende Reflexion aabenbar, der dog er bunden af
den umiddelbare religiøse Følelse. Idet Gud er den religiøse Følelses Gjenstand som den
absolut Ophøiede, Vidende, faaer Bønnen denne Resignationens Form, men Bønnens Bestemmen af Guds Villie er derfor \emph{ikke} negeret. \emph{Til en uforanderlig Villie kan der ikke bedes}. Bønnen overlader til Gud, \emph{paa hvilken} Maade han i sin Godhed bedst vil opfylde den
Bedendes Ønsker, men denne betvivler ikke Bønnens Betydning for Opfyldelsen af disse
Ønsker. Netop det, at der til den \emph{bestemte} Bøn føies \emph{Bønnen}: skee ikke min, men din Villie,
indeslutter, at den Bedende, der forestiller sig Muligheden af en guddommelig Villen, der
slet ikke refererer sig til den første Bøns Gjenstand, ved Bønnen dog vil bestemme Gud i
Forhold til denne. Opfattes Bønnen kun som det livfulde Udtryk for \emph{Contemplationen} af Gud, saa er det for Bønnen Specifiske borttaget, et \emph{andet} Begreb er skudt ind under
Bønnens Navn. Skal Bønnens Virken kun være den indirecte \emph{subjective Tilbagevirkning}
paa den Bedende, saa er Bønnens Væsen undergravet ved Reflexionen og Opfattelsen kommer i
aabenbar Strid med den umiddelbare religiøse Bevidstheds Opfattelse\footnote{Man tage blot
Hensyn til en saadan Bøn som Bønnen til Helgener eller Madonna om Intercession:
\textit{Ora pro nobis}, eller erindre f.\ Ex.\ \emph{Luthers} Udsagn om sin Bøn om den dødssyge
Melanchthons Helbredelse.}. \opage{97}En saadan Bøn uden objectivt Maal bevarer vel Bønnens Form, er \emph{formelt} en Bøn, men saaledes, at, gjennem Bevidstheden om dens blot subjective Virkning,
\emph{Bønnen som Bøn tilbagekaldes i selve Bønnen}: den er i Virkeligheden \emph{ikke Troens, men Vantroens Bøn}. Naar Bønnens Begreb tages, uden Forvanskning, i den Betydning, i hvilken
den religiøse Bevidsthed tager den, saa opfattes den som den, der \emph{forandrer Guds
Villie}, og dette forstaaes nærmere saaledes, at Bønnen bestemmer Gud, men
bestemmer Gud til en \emph{Almagtens} Yttring. Den heri liggende Modsigelse seer den religiøse
Bevidsthed som saadan ikke: \emph{Bønnens Almagt skjules af Guds Almagt}.

2.\ Det, hvorom der bedes, er \emph{Underet og kun} Underet. Begrebet om Bønnen i den
almindelige religiøse Bevidstheds Opfattelse indeslutter, at der bedes, ikke om hvad der
som ethisk Opgave hviler i vor Villies Beslutning, eller om hvad der følger af
Tilværelsens almene Lov eller af en uforanderlig guddommelig Beslutning fra Evighed, men
at der bedes om, hvad der kun har sin Grund i en absolut fri guddommelig Villie. Det er den \emph{almægtige} Villies Virken, Bønnen paakalder, og den paakalder Villien som den, der ikke
er bunden af nogen Forudsætning, hverken af en udenfor Gud værende eller af en i Guds
Væsensbestemmelse given, end ikke af Villiens forudgaaende Bestemmelser. Men den saaledes
opfattede almægtige Villie er den \emph{undergjørende} Villie. Og det Under, der forventes af
Villien, er \emph{som Under} altid et Under „\emph{i strængeste} Forstand“\footnote{jfr.\ \emph{Martensens} Yttringer: Om Tro og Viden, S.\ 112.}. \emph{Der gives ligesaalidt Grader i Underet} som der gives \emph{Grader i Almagt}; thi \emph{Underet er Almagtens Udtryk}. Der kan være en
Forskjel mellem hvad der er mere eller mindre \emph{\textit{mirabile}}; men idet dette
\textit{mirabile} sees som \emph{\textit{miraculum}}, og saaledes sees det ved Troen og ved Troen
alene, er Forskjellen hævet: \opage{98}saa sees Almagten, der ingen Quantiteren er underkastet, som
enevirkende Aarsag.

Men den Gud, til hvilken Bønnen viser hen: den ved Enkeltpersonlighedens Form bestemte
Gud, hvis almægtige Villie transcenderer hans Væsen og derved alt Væsen, al Lov; denne
Gud, i hvilken den umiddelbare religiøse Bevidsthed finder Tilfredsstillelse for sin
Villen og for sin Forestillen idet begge ere bestemte ved Phantasien\footnote{For \emph{Tanken} er her ingen Forsoning, fordi Aandens Udvikling endnu bliver i det Sporadiske,
Ugjennemarbeidede.}, --- denne Gud er netop den Gud, der, naar Forestillingen om ham
gjøres til Gjenstand for en Troes\emph{videnskab}, maa give et Gudsbegreb som det af Martensen
fremsatte. I Udførelsen kunne Nuancer finde Sted: de væsentlige Grundtræk maae være de
samme hvor Theologien med Bevidsthed fastholder sin Opgave.

Dennes Opgave som Troes\emph{videnskab} blev paa et tidligere Punkt bestemt saaledes, at den
maatte opfatte Aabenbaringen og Guds Væsen og Virken paa en saadan Maade, at de droges ud
af det rent Ubegribeliges Sphære henimod det Rationelle, at der altsaa trængte sig et
Skjær af Videns og Begrundelsens Nødvendighed ind, medens dog Indholdets suprarationelle
Charakteer ikke opgaves, og Viden, der underordnedes det Aabenbaredes Norm, ikke med sine
Former bragte sin Væsensbestemthed ind i Gjenstanden.

Undersøgelsen af hvorledes denne Opgave søges nærmere bestemt og gjennemført, og hvilken
Indvirkning overhovedet dens Gjennemførelse maa have paa Theologiens Behandling af sin
Gjenstand, bliver nu det \emph{andet} Hovedpunkt, med hvilket Kritiken af Theologiens Standpunkt
har at beskæftige sig. Kritiken af \emph{Gudsbegrebet} faaer til sit Supplement Kritiken af det \emph{Videnskabsbegreb}, der ligger til Grund for Opfattelsen af en Troesvidenskabs Mulighed og
Eiendommelighed.

\begin{center}---\end{center}

\section{Kritik af Videnskabsbegrebet}
\label{sec:videnskabsbegreb}

\opage{99}Ved Begyndelsen af dette Afsnit (S.\ 44 ff.) er det i Korthed angivet, hvorledes Martensen
søger at godtgjøre Muligheden af en Troesvidenskab som en organisk Enhed af Tro og Viden;
hvorledes han bestemmer dens Opgave, dens Form og dens videnskabelige Charakteer. Det er
tillige viist, hvorledes hans Opfattelse af det for den „høiere Videnskab“ Eiendommelige
og af denne Videnskabs Ret til Principatet hvilede paa hans Opfattelse af Grundforholdet i
det menneskelige Væsen og af Gudsbegrebet.

Gjennem den kritiske Undersøgelse af Theoriens psychologiske Grundlag og af det „ethiske
Gudsbegreb“ ere allerede i alt Væsentligt Præmisserne vundne til Dommen over en saadan
„høiere“ Videnskab. Naar der ikke kan hævdes Villien det tilstræbte Primat i det
menneskelige Væsen, og naar i det Guddommelige den Villiens Transscenderen af Væsenets og
Videns Nødvendighed, ved hvilken det ethiske Gudsbegreb skulde faae sin Eiendommelighed og
sin høiere Rang i Sammenligning med det logiske, ikke kan gjennemføres, saa er det derved
Betingede: en paa den frie guddommelige Personligheds overnaturlige Aabenbaring hvilende
Videnskab, i hvilken Videns Nødvendighed skulde begrændses, dens Love corrigeres gjennem
Villien, umulig: det Logiske hævder sin Ret og det Guddommeliges Meddelelse til Mennesket
bliver en ved Væsenets Nødvendighed bestemt, gjennem Tankens Former opfattelig. Af
Uholdbarheden af den postulerede „høiere“ Videnskabs Princip følger som nødvendig
Consequents denne Videnskabs Uholdbarhed. Dette vil stadfæste sig, naar vi, efterat have
suppleret de tidligere gjorte Bestemmelser af den høiere Videnskabs Væsen ved en
nøiagtigere Bestemmelse af dens Forhold til Videnskaben i almindelig Betydning, nærmere
belyse de givne Bestemmelser og ved nogle Exempler i Korthed paavise de Resultater, til
hvilken den høiere Videnskab fører og maa føre. Disse Exempler ville saa vise Rigtigheden
af den almindelige Formel for den høiere Videnskabs Methode, som vi opstille, og som paa
det Klareste aabenbarer de Modsigelser og Tvetydigheder, der ere uadskillelige fra
Theologien som \opage{100}Theologie og som berøver den Charakteren af Videnskab i egentlig Betydning.

\subsection{Den høiere Videnskabs Forhold til den blot rationelle}
\label{ssec:hoiere-videnskab}

Forholdet mellem den „høiere“ Videnskab og Videnskaben i almindelig Betydning
bestemmes\footnote{\emph{Martensen}: Om Tro og Viden, S.\ 94 f.} ved den fra \emph{Schelling}
hentede Modsætning mellem den \emph{positive} Videnskab, den positive Philosophie, og den
blot \emph{negative}. Denne sidste bestemmes som den \emph{søgende} Philosophie, der, gjennem en
fortsat Negation af de blot relative Principer, søger \emph{det sande Absolute, høiere end
hvilket der Intet kan tænkes eller være}, medens den \emph{positive} Philosophie tager sit
Udgangspunkt \emph{fra det saaledes fundne Princip}, udviklende dets positive Rigdom
forsaavidt denne gjennem dets factiske Aabenbarelser lader sig erkjende. Den \emph{negative} Philosophie, hvis Princip er et \emph{blot logisk}, skal være den abstracte Nødvendigheds, det \emph{Upersonliges}, Philosophie, hvis Endemaal er, fra Nødvendigheden at komme til
Friheden, og hvis Kategorier kun indeholde den negative Betingelse, \textit{conditio sine
qua non}, for Sandhedens Erkjendelse, men ikke give os den virkelige levende Sandhed; den
\emph{positive} Philosophie, hvis Princip er et \emph{Mere end et logisk}, skal være
Frihedens, det \emph{Personliges} Philosophie, der nærmere bestemmer sig som
\emph{Historiens og Aabenbaringens Philosophie}, idet den i Historien forsker efter den
personlige Guds Vidnesbyrd, og fra dette Standpunkt erkjender den tillige Naturen, der nu
vil vise sig i et nyt Lys. \emph{Den negative Philosophie kan ikke give en Forklaring af
Virkeligheden}, da den kun kan søge det Princip, der er det altforklarende; dette kan
derimod den \emph{positive} Philosophie, Philosophien for Livet, Virkelighedens Philosophie,
gjøre.

Den \emph{positive} Philosophie bestemmes nemlig som \emph{Erfaringsvidenskab}, Empirisme, men
som \emph{metaphysisk} Empirisme (jfr. \emph{Schelling}), der gjennem \emph{Anskuelsen} som Organ i
Virkningen seer den virkende Aarsag, i Factum \opage{101}Ideen, i det sig midt i Tiden Aabenbarende
det, som var før Verdens Grundvold blev lagt, og søger at udvikle Anskuelsen til Begreb.
Den begynder, idet den sætter Frihedens Gud som Sandheden, med en stor Bekræftelse, \emph{en stor Grundantagelse}, --- den gaaer herved ud over det Logiske og det Logiskes
Nødvendighed. Men den skal føre \emph{et fremskridende Beviis} for sin Grundantagelse, idet
den efterviser og erkjendende fordyber sig i det Beviis, som den villende og handlende Gud
selv vedbliver at føre for sin Existents i sit fortsatte Forhold til Menneskeheden.
Kun et \emph{uafbrudt fremskridende} Beviis og ikke et afsluttet skal nemlig egne sig for
den \emph{personlige} Guds Tilværelse. Ligesom vi kræve, at menneskelige Individer, der
ere os vigtige, vedvarende skulle give os nye Beviser paa deres Existents, saaledes kræve
vi det ogsaa med Hensyn til Gud. Vi kræve, at Guddommen altid kommer Menneskets Bevidsthed
nærmere, vi forlange, at den ikke blot i sine Virkninger, men \emph{selv}, bliver en Gjenstand
for Bevidstheden, hvortil der dog kun trinviis kan naaes hen. Det fyldige Beviis for den
personlige Guds Existents maa derfor føres gjennem \emph{hele} Virkeligheden; det gaaer tilbage
til vor Slægts Fortid og rækker ud i dens yderste Fremtid. \emph{I Frihedens Philosophie} kan der derfor \emph{ikke} være Tale om noget \emph{afsluttet System} fordi Virkelighedens System ikke er afsluttet, og Frihedens Gud, der har indført sig selv i Historien, ikke blot er
den, som er og som var, men ogsaa er den, som kommer\footnote{\emph{Martensen}, p.\ anf.\ St.,
S.\ 95 f.}.

Til denne høiere Videnskab: \emph{Aabenbaringsvidenskaben}, den ethiske Videnskab, det
Godes Videnskab i eminent Betydning, Videnskaben om Gud, den ene Gode, i hans frie
Selvmeddelelse til Mennesket, forholder \emph{Theologien} sig saaledes, at hiin bliver den \emph{universelle} Videnskab, Theologien den \emph{centrale} Videnskab, der udelukkende fordyber sig i
det Ene, i Guds Rige som saadant.

\subsection{Kritik af Forsøget paa at begrunde Troesvidenskabens Mulighed}
\label{ssec:troesvidenskab-kritik}

\opage{102}Kritiken af dette \emph{Videnskabsbegreb} maa først belyse Forsøget paa at godtgjøre \emph{Muligheden af en videnskabelig Viden af Troens Indhold}.

I hvilke Vendinger dette Forsøg er gjort, er i det Væsentlige allerede angivet paa et
tidligere Sted\footnote{See ovenfor, S.\ 44 ff.}. Det der Anførte bestemmes nærmere ved de
Udsagn, at \emph{Aabenbaringen} vel ikke vil give os et \emph{videnskabeligt Begrebssystem}, men
at dette ingenlunde vil sige, at den Intet har villet lære os om de høieste Spørgsmaal,
hvis Løsning Metaphysiken søger; thi om den end ikke giver os \emph{abstracte Begreber},
saa giver den os en omfattende \emph{Totalanskuelse}, der omfatter Himmel og Jord, Natur og Historie, det Forbigangne, Nærværende og Tilkommende\footnote{\emph{Martensen}: Om Tro og Viden, S.\ 26 f.}; og den „har \emph{sin egen Metaphysik}, af hvilken den ikke vil opgive det Ringeste“.

Det vil ved første Øiekast være indlysende, at det søgte Beviis for hiin videnskabelige
Viden kun naaes \emph{gjennem Subreptioner} ved Hjælp af \emph{Tvetydighed} i Brugen af
Begreberne. Det er Begreberne \emph{Viden, Sandhed, Tanke og Videnskab}, der bruges i
Dobbeltbetydning eller i udvisket Betydning. Det, at den Troende maa \emph{vide} hvad han troer ɔ: \emph{være sig bevidst}, hvad han troer, gjøres til, at han maa kunne vide det i
Betydning af at have en \emph{videnskabelig Erkjendelse} af det. Naar \emph{Troen} ikke selv var en \emph{Viden}, siges der, saa vilde den Troende ikke kunne bekjende sin Tro og ikke kunne
forkaste og bekæmpe de vildfarende Lærdomme. Her er atter den samme Begrebsombytning. Den Troende \emph{maa være sig Troesindholdet bevidst}, det maa være for hans Bevidsthed \emph{i udfoldet Form} forat han kan \emph{bekjende} det; men han kan være sig det bevidst som det, der
netop \emph{ikke kan vides videnskabeligt} ɔ: begribes ved et rationelt Princip, fordi
dets Metaphysik d.\ v.\ s. det metaphysiske Udtryk, til hvilket dets Bestemmelser kunne
reduceres, \opage{103}er et \emph{Modsigende}, en modsigende Forbindelse af de for Tanken gjældende
metaphysiske Bestemmelser. Den Troende kan, ud fra hiin Bevidsthed, \emph{bekæmpe} modsatte Lærdomme paa den Maade, paa hvilken det \emph{factisk} er skeet: \emph{enten} ved at paavise indre Modsigelser hos Modstanderne, \emph{eller} ved at argumentere ud fra de religiøse, af
Modstanderne ikke anerkjendte Forudsætninger. --- Der siges endvidere, at naar Troen ikke
allerede i sig var en Viden, \emph{saa var den ikke Sandheds Tro}. \emph{Troens Sandhed} betegner
for den religiøse Bevidsthed, at Troens Indhold er den \emph{sande} ɔ: \emph{virkelige} Gud, i Modsætning til de \emph{falske} ɔ: \emph{uvirkelige} Guder. Men Begrebet
\emph{Sandhed} overføres fra \emph{denne} Betydning til Betydningen af \emph{theoretisk Erkjenden
i Videnskabens Form}. Bevidstheden om den \emph{christelige Sandhed} i Betydningen af det,
\emph{der virkelig er} det Christelige, gjøres ved en Subreption til den \emph{christelige
Sandhedsidee}, og gjennem Bestemmelsen: \emph{Idee}, naaes der ved en \emph{ny} Subreption
til \emph{det Videnskabelige}. I de apokryphiske Skrifters fra den senere græske
Philosophie optagne Forestilling om \emph{Viisdommen}, „den guddommelige Tanke, der før
Verdens Skabelse legede for Guds Aasyn“, og som har udpræget sig i Aabenbaringen, søges en
Mulighed for, at den troende Bevidsthed \emph{tænkende} skal kunne ransage Aabenbaringen
og granske, ikke blot efter Sammenhæng, men ogsaa efter Grund; der tales om en \emph{guddommelig Aabenbarings- og Viisdomstanke}, der i den troende Bevidsthed begrunder en \emph{Enhed af Viden} og Samvittighed. Men naar Guds \emph{Væsen} bestemmes ved den \emph{al} Videns Nødvendighed \emph{transscenderende Villie}, saa forholder \emph{Viisdommen} sig til de „\emph{evige Raadslutninger}“,
der, som udgaaede fra en kun i sig selv grundende og derfor grundløs Villie ere \emph{uudgrundelige}. Derfor bliver den guddommelige Viisdom det, der overstiger al menneskelig
Viden; Aabenbaringen aabenbarer Raadslutningerne og Raadslutningernes Subject \emph{som det med Viden Incommensurable}, som det, om hvilket der kun kan vides \opage{104}et At, ikke et
Hvorledes. Troens Indhold kan med indre Consequents sammenknyttes, men \emph{Principet} bliver
det \emph{for Tanken Utilgængelige}, derfor bliver \emph{det Sammenhængende som Heelhed
utænkeligt}; kun \emph{den formelle Sammenhæng} i det \emph{Begrebne} bliver Gjenstand for Tanken.
At Troesindholdet ikke bliver Gjenstand for en \emph{egentlig Viden} ɔ: for en \emph{Begriben}, bliver ogsaa deels indirecte indrømmet gjennem de Bestemmelser, der gives af
Troesvidenskabens Væsen, deels fremgaaer det af Resultaterne af Forsøget paa at fremstille
en saadan Videnskab.

Bestemmelserne af \emph{Troesvidenskabens Eiendommelighed} ere saaledes beskafne, at de,
medens de vise en Usikkerhed i Opfattelsen af dens Forhold til \emph{den rent rationelle
Videnskab}: en Usikkerhed, der netop er begrundet i \emph{Umuligheden} af at begrændse dennes
Omraade, tillige, netop idet de skulle sikkre Troesvidenskaben en \emph{høiere} Rang, drage den ned \emph{under} Videns og Videnskabens Sphære, berøve den Retten til Videnskabens Navn.

\subsection{Kritik af Bestemmelsen af dens Forhold til den rationelle Videnskab}
\label{ssec:rationel-videnskab}

Hvad Opfattelsen af hint Forhold angaaer, saa bestemmes den \emph{negative} Philosophie først
saaledes, at den ikke blot søger, men \emph{finder}\footnote{Jfr.\ \emph{Martensen}: Om
Tro og Viden, S.\ 94: „det saaledes \emph{fundne} Princip“.} det sande Princip, høiere end
hvilket Intet kan \emph{tænkes} eller \emph{være}, altsaa naaer til \emph{Erkjendelsens Høieste}, det i Sandhed Absolute, og omfatter \emph{Metaphysikens hele Omraade}; medens den \emph{positive} Philosophie bestemmes saaledes, at den maatte falde sammen med \emph{Natur- og
Aandsphilosophien}, saa at altsaa begge, \emph{baade} den negative og den positive, vilde falde
\emph{indenfor} den \emph{rent rationelle Philosophies Sphære}, som \emph{Dele} af denne. Saa søges
imidlertid den \emph{rent rationelle} Videnskab, der \emph{identificeres} med den \emph{negative}, sondret fra den \emph{positive}; men Adskillelsen faaer kun Skin af Gyldighed naar man
vilkaarlig abstraherer fra det, der sættes som den negative Philosophies Slutning eller
bryder af hvor Maalet er berørt uden \opage{105}at udvikle det Absolutes Begreb. Derefter postuleres
som Betingelse for den \emph{positive} Philosophies Begyndelse „en stor Grundantagelse“. Men den \emph{negative} Philosophie skal jo allerede have „fundet Principet“; der behøves altsaa ingen
særskilt ny Grundantagelse. Ved disse Bestemmelser bliver derfor den positive Philosophie
ikke i Virkeligheden unddraget det rent Rationelles Omfang, og naar den „høiere“
Videnskabs Princip, der, i Modsætning til den blot logisk-physiske Sandhed, betegnes som den \emph{absolute} Sandhed, bestemmes som „det Gode selv i dets Aabenbaring \emph{for Tanken} som den
høieste Realitet“, saa viser dette hen i samme Retning. \emph{Men saaledes vilde altsaa
den positive Videnskab kun hente sin videnskabelige Charakteer fra Enheden med den
negative, rent rationelle.}

Den uvilkaarlige Indrømmelse af \emph{den rationelle Philosophies Universalitet} er med
Nødvendighed begrundet i Sagens Natur. \emph{Videnskaben er kun Videnskab som} System, \emph{som den efter sit Væsen universelle Tankes organiske Udfoldelse}. En Begrændsning af Videnskaben ved et \emph{Andet}, der dog skal være en \emph{Gjenstand for en Erkjendelse}, er en \emph{Umulighed}, fordi det er en \emph{Negation af Videnskabens Væsen}, og fordi dette er bestemt ved \emph{den udelelige Selvbevidstheds universelle Natur}.

Denne Conclusion vil man fra Martensens Standpunkt afvise ved at sige: det er \emph{ikke} Videnskaben \emph{som saadan}, der begrændses, men kun den \emph{lavere} Videnskab, og idet det, der
begrændser denne, er den \emph{høiere} Videnskab, saa bliver den lavere Videnskab netop i denne Begrændsning \emph{suppleret} til Videnskabens \emph{sande}, al \emph{given} Virkelighed \emph{omfattende System}.

Herved føres vi saa til at undersøge de Bestemmelser, ved hvilke den postulerede Videnskab
skulde hævde sig som den \emph{høiere}.

Det, der i Realiteten skal retfærdiggjøre Sondringen af den positive Videnskab fra den
rationelle, og hiins høiere Værd, er Postulatet af \emph{en ny Erkjendelsesart}, \emph{en fri}, \opage{106}\emph{ved Villien bestemt Erkjendelse}, der først skal gjøre sig gjældende i hiin „Grundantagelse“,
som altsaa i Virkeligheden ikke skal blive en Tilknytning til den rent rationelle
Philosophies Slutningspunkt, men en Sætten af et nyt, \emph{uensartet} Udgangspunkt, i Kraft af
en Underordning under Troens og Aabenbaringens Indhold, af et „frit Afhængighedsforhold“
til Aabenbaringen. Men denne høiere Erkjendelse, der postuleres, vil vise sig i Virkeligheden \emph{ikke at være Erkjendelse}, men kun \emph{en af Attraaen bestemt
Phantasievirksomhed}.

Den nye Erkjendelsens Form, der sættes som en høiere end den blot logiske Tænkning, skal
være den høiere og gjøre den derpaa byggede positive Videnskab til en høiere, fordi den er
\emph{en Erkjendelse} af den frie \emph{personlige Guds Aabenbaringer}. Men det Begreb om en
guddommelig Personlighed, til hvilket der vises hen, er allerede belyst, og det er viist,
at det kun har den endelige Enkeltpersonligheds Indhold, hvorved der bringes en
Endelighedsbestemthed ind, der gjør det af denne Kilde Fremgaaende til et Lavere end det
af Ideen Udledede. Man bliver stillet ligeoverfor det \emph{Dilemma}: naar Naturen og Historien skulle \emph{erkjendes} som Guds Aabenbaringer, saa maae de sees som begrundede i Guds Væsen, og dette kan da, som \emph{al} Tilværelses absolute Princip, ikke bestemmes ved Personlighed; --- fastholdes derimod \emph{Personlighedens} Bestemmelse, saa bliver Gud det \emph{absolut Ubegribelige}, og Naturen og Historien kunne \emph{ikke} af Tanken opfattes som Guds
Aabenbaringer. \emph{Personlighed} skal netop være det, \emph{der ikke gaaer op i Begrebet}, det,
der kun bliver Gjenstand for Anskuelse, Følelse, Tilbedelse. Den \emph{absolute} Personlighed vilde være den allerpersonligste Personlighed, det absolut singulære Væsen,
der aldeles ikke havde almindelige Prædicater, og derfor ingen Tilknytningspunkter for
Tænkningen. „Vel faaer Gud, naar han bestemmes som personlig, en Bestemmelse, der ogsaa
tilkommer mig, men dens Absoluthed er dens Singularitet, dens Distinction fra mig.
Angaaende en anden endelig Personlighed kan jeg, \opage{107}paa Grund af dens Beslægtethed med mig,
vel gjætte og formode Noget; men om den absolute Person kan jeg, da her endog Basis for
Forudsætningen om en mulig Lighed med mig i Tænken og Handlen falder bort, kun drømme og
phantasere.“ (Feuerbach.) --- Hvad der aabenbares af den \emph{frie}, personlige Gud, skal kun
kunne erkjendes af Mennesket \emph{ved Frihed}, og dette \emph{Friheds-Forhold}, der
modsættes Tankenødvendighedsforholdet, skal saa betegne den nye Erkjendelses og den
nye Videnskabs høiere Værd. Men denne \emph{frie} Erkjendelse, der skal være en \emph{høiere}
Erkjendelse, maa tænkes at blive til ved \emph{et Brud}, sat ved \emph{Villien}, --- en virkelig
Begrebsbestemmelse af Forholdet giver Martensen ikke ---; saaledes at Tankenødvendigheden
afbrydes ved Villiesacten, der sætter en ny Begyndelse. Villien, der sætter denne, bliver da \emph{en fra Erkjendelsen forskjellig Villie}, men for en saadan Villie have vi ingen
anden Betegnelse end \emph{Vilkaarlighed}, og den Erkjendelse, der skal hvile paa denne
Villiesact, bliver en saadan, der bæres af „Villiens Fjederham“ \emph{Phantasien}, og savner
ethvert Correctiv, enhver Regulator. Dens „Frihed“ er et \emph{Ringere} end Tankens Nødvendighed;
thi den Tankenødvendighed, der ved den rationelle Videnskab skal være Tegnet paa dennes
lavere Værd, er netop Udtrykket for Tankens Bestemthed ved sit eget Væsen, sin egen Lov;
den er derfor Særkjendet for \emph{al egentlig Videnskab}, og det er en
\emph{Misforstaaelse}, naar den indskrænkes til Videnskabens abstracteste og
indholdsløseste Former. Jo concretere og rigere Videnskabens Indhold er, desto \emph{vanskeligere} bliver den Assimilation, der giver Gjenstanden Tankenødvendighedens Form, men \emph{derfor bliver denne Assimilation ligefuldt Opgaven}, naar ikke det Concrete, kun successivt
Tilegnelige, forvexles med det ifølge sin irrationelle Natur med Tanken Incommensurable og
for den absolut Utilegnelige. Og paa disse Gebeter faae vi ikke en abstract
Tankenødvendighed, men \emph{den concrete Tankes Nødvendighed}, der er tvingende for Tanken fordi den er Tankens eget Væsen.

\opage{108}Den nye Erkjendelses og den nye Videnskabs høiere Værd motiveres endelig derved, at de
skulle give en \emph{Virkelighedserkjendelse}, hvad den negative Philosophie ikke skal
kunne. Her er imidlertid den samme Forvexling af Begreber tilstede, der paa et tidligere
Punkt blev paaviist. Den Tankenødvendighed, der er mit tænkende \emph{Væsens} Nødvendighed, har
\emph{existentiel} Vished for mig og er netop den høieste Forvisning, medens den ydre
Existentses Vished altid kan blive tvivlsom. Men den Existents, Martensen søger, er netop
en Existents i Analogie med den ydre, sandselige Virkeligheds, en saadan Existents, der
bliver Gjenstand \emph{for Anskuelse}. Til en saadan vises der ogsaa hen naar der tales om den
høiere Videnskabs personlige Princip og dette Princips Aabenbaring. Netop idet der fordres
et \emph{vedvarende} Beviis for denne personlige Guds Existents, drages han ned i Lighed med
Enkeltindividet, hvis Existents som endelig bestemt kan afbrydes. Den høiere Erkjendelse,
der betegnes ved det \emph{modsigende} Begreb: \emph{metaphysisk Empirisme}, bliver saaledes i
Virkeligheden kun \emph{en Empirisme, der forholder sig til en endelig, sandselig,
phantasiebestemt Gjenstand}, og Erkjendelsen af denne skal være \emph{Virkelighedserkjendelse},
medens Erkjendelsen af Ideen og det i Ideen Begrundede ikke skal være det. Og dog er det
vel indlysende, at naar der fordres af den høiere \emph{Videnskab}, at den skal forholde sig til
Existents, saa maa Existentsen, \emph{som sat i Tanken}, om end denne Tanke er en ved Villie
bestemt Tanke, altid blive \emph{Existentsens Tanke}, altsaa i sin Sandhed \emph{Ideens Væren i Tanken}. Skal Existents være noget Andet end den Existents, det Ideelle har i min Tanke,
saa bliver den \emph{sandselig} Existents, og Tanken, der skal optage den, bliver kun
\emph{formel} Bevidsthed \emph{om et sandseligt Forhold}, et \emph{Erfaringsforhold i almindelig} Betydning. Naar der lægges Eftertryk paa, at Gud aabenbarer sig i en Enkeltpersonlighed
for derigjennem at være \emph{umiddelbart} aabenbaret, saa glemmes der, at en saadan
Enkeltpersonlighed, i sin Adskillelse fra de andre menne\opage{109}skelige Individer og i sin Væren i
Menneskets Form, dog kun bliver en \emph{middelbar} Aabenbarelse og en Aabenbarelse
gjennem en \emph{Virkning} af Gud. Guds Aabenbarelse af \emph{sig selv} kan kun tænkes \emph{som en Væsensmeddelelse til Aanden}.

Naar saaledes Alt, hvad der skulde hævde den positive Videnskab en \emph{høiere} Rang, netop viser hen i den \emph{modsatte} Retning, og naar denne høiere Videnskab, idet den ikke skal kunne
give sin Gjenstand Nødvendighedens Form, idet den ikke skal naae til en virkelig Begriben
% PRINTED AS IS, p.109 [collation 2026-09-21]: page prints "Autorttet" (t for i); confirmed by a blind second reader.
af den, og idet den skal bøie sig under en Autorttet, der ikke er Tankens, og ikke kan
være høiere end Tankens, fordi, som det har viist sig, det Høieste i Aandens Verden finder
sit Udtryk i det Logiske, udenfor hvilket kun det Alogiske falder --- netop savner \emph{det}, der constituerer Videnskabens \emph{Væsen} og \emph{Form}, det, der charakteriserer Videnskaben \emph{som} Videnskab, og viser sig som hvilende paa et Fundament, der er lavere end Begrebet, saa \emph{taber den Retten til at bære Videnskabens Navn og endnu mere til at være den høiere
og herskende Videnskab}. Den bevarer af det for Videnskaben Betegnende kun det, at der fra
Principet drages sammenhængende Consequentser; men Principet, der bærer Consequentserne og
ikke kan omformes ved en Tilbagevirkning af Consequentserne, bliver et Ubegrebet.
\emph{Videnskab i sand Forstand og herskende Videnskab bliver den Videnskab, der begriber
det Absolute i dets sande Uendelighed som Tilværelsens Virkelighedsprincip}, og denne
Videnskab er \emph{Philosophien}.

\subsection{Troesvidenskabens Methode}
\label{ssec:methode}

Den Dom, der her er fældet med Hensyn til den „høiere“ Videnskab: Troesvidenskaben,
stadfæstes fuldstændig naar vi betragte \emph{dens Resultater}. Allerede ovenfor er det med
% PRINTED AS IS, p.109 [collation 2026-09-21]: the page prints "Uuderet" (turned n); both readers saw it.
Hensyn til væsentlige Punkter i Troeslæren: Guds Personlighed, Skabelsen og Uuderet,
paaviist, hvorledes Forsøget paa at gjøre dem til Gjenstand for en Erkjendelse mislykkedes,
og hvorledes den forsøgte Rationaliseren, consequent gjennemført, vilde sætte noget ganske
Andet i Stedet \opage{110}for de religiøse Lærdomme. Det Samme vilde kunne paavises med Hensyn til
hvilketsomhelst Dogme. Overalt maa Udviklingen føres frem ved Postulater, overalt maa
Begrebsbestemmelserne bruges dobbelttydigt, overalt maa Resultatet blive en umiddelbar
Forening af det Modsigende.

Methoden vil, naar alt det Skjulende fjernes, vise sig at være følgende. Forat give \emph{Tanke}bevægelsen et Tilknytningspunkt, gaaes der ud fra Bestemmelser, der hentes fra
Naturens og Aandens Virkelighed, men som, idet de skulle anvendes paa Troens Gjenstand,
vilkaarligt, tildeels phantastisk, omdannes forat den oprindelige Bestemthed kan glemmes og
% PRINTED AS IS, p.110 [collation 2026-09-21]: page prints "nn" (turned u) for "nu"; confirmed by a blind second reader.
% PRINTED AS IS, p.110 [collation 2026-09-21]: the page prints "Anrendelsen" (r for v, wrong sort); both readers saw it.
Anrendelsen blive mulig. Det saaledes Omformede bliver nn Grundlag for nye ved Postulater
vundne Bestemmelser, og idet Consequentserne af det Rationelle vilkaarlig afbrydes, eller
en Quantiteren, til hvilken intet Begreb kan knyttes, anbringes, bliver det Troesindhold,
der skal begrundes for Tanken, bragt frem gjennem en illusorisk Begrundelse, i
Virkeligheden kun gjennem Postulater, og der afsluttes dermed, at de to modsatte Elementer,
der mødes i Troesvidenskaben, begge finde deres Udtryk, men saaledes, at Udtrykkene
neutralisere hinanden i Ubestemthed, idet det rationelle Element paralyseres ved
Troeselementet, men dog ved sin Tilstedeværen afficerer dette. Subjectet i den Sætning, der
opstilles som Resultat, bestemmes ved Prædicater, der berøve det dets Indhold og kun lade
en ubestemt, vag Forestilling tilbage for Bevidstheden, som saa netop i Ubestemtheden, i
Mangelen paa den bestemte, klare Begrændsning, indbilder sig at have et Dybt, et
Speculativt, et Uendeligt.

\subsection{Dens Resultater}
\label{ssec:resultater}

Naar vi søge bestemte \emph{Exempler} paa denne Methode, hvis Nødvendighed ligger i den Opgave,
der er stillet, saa kunne vi f.\ Ex.\ finde dem i Udviklingen af Treenigheds\emph{dogmet} i
Dogmatiken og i Skriftet om Tro og Viden. Man vil der finde \emph{alle} Methodens Momenter
forenede: en Gaaen ud fra rationelle Bestemmelser: Selvbevidsthedens og Samfundets
Bestemmelser; en Omdannen af dem, ved hvilken de tabe deres specifiske Betydning
(Selvbevidsthedens \opage{111}ideelle Fordobbling bliver en flere Selvbevidstheder sættende
Mangfoldiggjørelse af Personligheden; Samfundet skal være Person); en vilkaarlig
Fremskriden gjennem Postulater (f.\ Ex.\ at Personen som absolut, guddommelig Person skal
være en Flerhed af Personer; at Personerne, trods deres Forskjel, dog skulle være Gud; at
Enheden, skjøndt Enhed af Modsætninger, ikke maa være en høiere Enhed); en vilkaarlig
Afbryden (ved Personernes \emph{Trehed}, medens Præmisserne føre til en ubestemmelig Flerhed); en
umiddelbar Sammenknytning af det hinanden Ophævende. Vi kunne finde dem i Udviklingen af \emph{Skabelsens Forhold til Kosmogonien}, hvor først Skabelsens Forestilling illusorisk
rationaliseres ved Kosmogoniens Begreb, men saa dette factisk ophæves ved at føres ind
under den ene skaberiske overnaturlige Begyndelse. Vi kunne finde dem i Udviklingen af \emph{Læren om Syndefaldet}, hvilket man paa een Gang skal tænke som en baade ubetinget og
betinget guddommelig Raadslutning, d.\ v.\ s. som en ikke blot fra Evighed bestemt
Raadslutning, men ogsaa som bestemmelig ved Skabningens Frihed, som en vordende
Raadslutning, der bestemmer sig til historisk Livsbevægelse; med andre Ord: ikke blot som
en logisk, men som en ethisk Raadslutning, den hellige Villies Raadslutning, i hvilken der
endnu er noget Uafgjort, men som dog ikke i sit Væsen ophører at være ubetinget. Vi kunne finde dem i Udviklingen af \emph{Forholdet mellem Naade og Frihed}, hvor der paa den ene
Side, idet Naaden, til Forskjel fra den blotte Almagt, ikke skal virke uimodstaaeligt, men
respectere den menneskelige Frihed, hævdes Villiens Frihed en Betydning, men denne saa
atter tages tilbage ved de tilføiede Bestemmelser, idet den væsentlige Frihed bestemmes som
den medskabte Naade, der kommer til Gjennembrud i den naturlige Villie, og er at betragte
som Naaden i Naturen, medens den blot naturlige, menneskelige Villie kun kan gjøre
Modstand. Vi kunne endelig, for ikke unødvendigt at mangfoldiggjøre disse Exempler, der kun
skulle supplere de allerede af Andre analyserede, til Slutning vise hen til Dogmati\opage{112}kens
sidste Afsnit, til \emph{Læren om de sidste Ting}, hvor den evige Fordømmelse fastholdes
ved Siden af Guds Raadslutning til Alles Salighed, hvor der tillægges de Fordømte en
ufortabelig Mulighed til det Gode, men saaledes, at denne Mulighed skal være absolut
afskaaren fra alle Virkelighedsbetingelser for sin Udvikling, og hvor Phantasien skaber
Forestillingen om en „Rumlighed“ indenfor Aandeligheden og Inderligheden, om en
Mellemtilstand efter Døden med en „Mellemlegemlighed“ o.\ s.\ v.

Saasnart det Tvetydige og Ubestemte i Troesvidenskabens Betegnelser fjernes, saa vise
Troesbestemmelserne og Vidensbestemmelserne sig her, istedenfor at være forbundne i
organisk Enhed, som forstyrrende og ophævende hverandre. \emph{En Troesvidenskab er altsaa
her factisk ikke given; om den overhovedet kan gives, afhænger af, om dette theologiske
Forsøg udtrykker noget for Theologien Almeengyldigt eller ikke.}

\subsection{Det for Theologien Typiske i den kriticerede Opfattelse
  af Troesvidenskaben}
\label{ssec:typiske-troesvidenskab}

Dette bliver let at afgjøre naar Forholdet mellem den \emph{Videnskabens Form}, hvilken
Theologien som Troesvidenskab tilstræber, og det \emph{Indhold}, der skal optages i denne Form,
betragtes. \emph{Tankens Former ere uadskillelige fra et Tankens, ved dens Væsen bestemt,
Indhold}; skal altsaa Troesindholdet \emph{virkelig} formes ved disse Former ɔ: gaae op \emph{som gjennemformet} i disse Former, ikke blot \emph{partielt og paa udvortes Maade} bestemmes ved dem,
maa det i egentlig Forstand kunne være et \emph{Vidensindhold}. Kan det ikke dette, vil en
saadan Anvendelse af Formerne, som den, ved hvilken Videnskaben betinges, være umulig.
Videnskaben, som begribende Erkjenden af et organisk udfoldet Hele, betinges ved, at
\emph{dens Indhold i sin Heelhed er et rationelt} ɔ: et ved Tanke bestemt, med Tanken ensartet Princips tankenødvendige Udformning, ikke blot et ved Tankebestemmelser formelt Sammenknyttet, med Tanken Uensartet.

Paa det Spørgsmaal: om Troesindholdet kan staae i et saadant Enhedsforhold til
Videns-Former og Indhold, er Svaret allerede ved dette Afsnits Begyndelse antydet, og \opage{113}ved
de sildigere Udviklinger ere nye Momenter givne til Besvarelsen. Der blev paa første Sted
fremhævet, at den religiøse Bevidstheds Sondring mellem den transscendente Gud og
Mennesket maatte føre til et Modsætningsforhold, idet Reflexionen bestemte Forholdet. Troen
med sin transscendente, overnaturlige Gjenstand og sit transscendente overnaturlige Princip
maa sondre sig fra Videns Indhold og Videns Princip; thi Viden, Tanken, er \emph{Menneskenaturens} Bestemmelse, dens Former altsaa \emph{naturlige} Former, dens Indhold \emph{bestemt ved disse Former}. Som Udtryk for \emph{Menneskenaturen} kan Tanken \emph{ikke} optage det \emph{Overnaturlige}; thi
hvad der kan modtages, er bestemt ved det modtagende Organ. En Modtagen af et \emph{qualitativt} Høiere vilde forudsætte en \emph{Omskaben} af Organet. Sondringen mellem den
overnaturlige Gjenstand og det naturlige Organ har saa, ved de sildigere Udviklinger, faaet
den bestemtere Form, at Troesindholdet bliver \emph{den frie, personlige, af ingen Væsenets
eller Videns Nødvendighed bundne Guds overnaturlige Aabenbarings Indhold}, et Indhold, der
angaaer Guds Tanken overstigende Væsen, Almagtens Gjerninger, de guddommelige
Raadslutninger. \emph{Den overnaturlige Aabenbaring er for den religiøse Bevidsthed et
Under: Modtagelsen af det ved Aabenbaringen Meddelte er et Under}, ved hvilket der sættes en
ny Erkjendelsens Begyndelse, der ikke har en Sammenhæng og en Væsensenhed med den
forudgaaende, i sig afsluttede Viden. \emph{Aabenbaringsindholdet bliver derfor et
Ubegribeligt, et Vidensindholdet Modsat}. At Modsætningen ligesaafuldt vilde komme frem,
naar der begyndtes med en \emph{quantitativt} lydende Differents mellem den menneskelige Viden og
det, der blivende overstiger den, er paaviist i det Foregaaende.

Den \emph{religiøse} Bevidsthed maa saaledes bestemme Forholdet som et \emph{Modsætningsforhold}, men
heri stemmer den overens med den \emph{videnskabelige} Bevidsthed, der ikke kan anerkjende noget
andet Princip for Videnskaben end den ved sit \opage{114}eget Væsens Lov bestemte Viden, og som i den
positive Religions Indhold seer det, der ved Forestillingsformen er bestemt som et
Endeligt, og som et saadant staaer i Modsætning til det ved Tankens Uendelighed bestemte
Vidensindhold.

Theologien viser sig altsaa, idet den vil grundlægge en Troesvidenskab, at faae en i
Principet modsigende Opgave. Dens Opgave er jo nemlig, idet den har hint, af den religiøse
Bevidsthed, fra hvis Standpunkt Theologien vil reflectere over Religionen, i Modsætning til
den menneskelige Viden og dens Indhold satte overnaturlige Aabenbaringsindhold til sin
Gjenstand, at gjøre det til et Videns-Indhold. Skal der gjøres et Forsøg paa at løse denne
Opgave, uden at Modsigelsen strax bliver iøinefaldende, saa maa der overalt tyes til hvad
der kan gjøre Forholdet uigjennemsigtigt: vilkaarlige Restrictioner og dobbeltlydige
Bestemmelser maae anvendes overalt hvor Grændser skulle angives, altsaa netop hvor
Modsætningen skulde fremtræde, og de to Modsætninger, der skulle forenes, maae begge
modtage en imod deres Natur stridende Omformning. Dette vil med andre Ord sige: \emph{paa
den ene Side maa Videns og Videnskabens Begreb forvanskes, medens et Skin af deres Bestaaen
bevares; paa den anden Side maa Begrebet om den Gud, hvis Væsen er Villie, hvis Virken er
Underet, faae en forvanskende Tilsætning af et tilsyneladende Rationelt:} \emph{det
Rationelle maa gjøres halvt irrationelt, det Irrationelle halvt rationelt.}

Hvad det \emph{Første} angaaer, saa bliver \emph{Videnskabens Begreb}, at være Erkjendelsen af det, der
med \emph{rationel} Consequents udledes af \emph{rationelle} Forudsætninger, \emph{udvidet} til at betegne en
Udledning af \emph{formelle} Consequentser af en \emph{antirationel} Forudsætning. Videnskaben faaer altsaa en \emph{Dobbeltbetydning}. \emph{Navnet} bevares som Fælledsnavn for det Væsensforskjellige, ja i \emph{Væsen Modsatte}, ved et modsat Princip Bestemte.
Troesvidenskabens Viden bliver en \emph{Bevidsthed} om hiin \opage{115}\emph{formelle Sammenhæng}, ikke en Begriben af Heelheden ud fra Principet. Og idet Theologien, i Bevidstheden om Forskjellen
mellem den saaledes forvandlede Videnskab og Videnskaben i dens almindelige Betydning, vil
bestemme hiins Eiendommelighed og istedenfor den Videnskabens Frihed, der er Udtrykket for
dens Bestemthed ved sit eget Princip, ved Tankens Nødvendighed, tillægger den en
Afhængighed af et med Tanken uensartet, og dog „Tanken corrigerende“ Princip, saa maa den,
ved et Postulat, forat bevare Skinnet af hvad der charakteriserer Videnskaben, bestemme
denne Afhængighed som \emph{fri} Afhængighed, og ved en Tvetydighed benytte den middelbare,
videnskabelige Erkjendelses Forhold til dens \emph{umiddelbare} Udgangspunkt forat tilveiebringe
et Skin af Lighed. Men ad ingen af disse Veie naaes det tilsigtede Maal. Hint Postulat af \emph{Frihed} kan ikke retfærdiggjøres. Videnskaben, der forholder sig til et \emph{realt Object},
kan stille sig i et \emph{frit} Afhængighedsforhold til dette, og maa det, for gjennem
Assimilationen af Objectet i dets Virkelighed at blive Herre over Objectiviteten; men den
gjør sig kun afhængig fordi den i Gjenstanden anerkjender et \emph{Rationelt}, altsaa et
Saadant, i hvilket \emph{dens eget Væsen} affirmeres. Men Videns Underordning under
Troesindholdets Autoritet er Underordningen under det \emph{Uensartede}, altsaa ikke Videnskabens
sande Afhængighed af Fornuftens almene Væsenslov, men en \emph{ufri}, Videnskabens Væsen
ophævende, Afhængighed. Videns Underordning under den fremmede Autoritet maa skee ved en Act af \emph{en fra Viden forskjellig Virksomhed, ved en ikke af Viden bestemt Villiesact};
derfor maa al Theologie consequent accentuere \emph{Villien} som den \emph{theologiske Videnskabs Organ}.
--- Tvetydigheden i Paaberaabelsen af det \emph{Umiddelbare} ligger deri, at idet det fremhæves,
at al Videnskab hviler paa et umiddelbart Udgangspunkt, og idet deraf sluttes, at der
altsaa ingen Modsigelse er i, at en Videnskab hviler paa Troesforholdets Umiddelbarhed, men
at der tværtimod maa fordres, at Troen maa være det Første, den primitive Sandhed,
Videnskaben det \opage{116}Afledede, Begrebet: det \emph{Umiddelbare} paa een Gang bruges om det, der er
Gjenstand for Erkjendelsens phænomenologisk første Form: Sandsningen, Erfaringen, og som
\emph{skal omsættes i et Begrebet}, deduceres og retfærdiggjøres ved at begribes, og om det, der
er Gjenstand for Tro i dette Ords egentlige --- ikke i den udviskede Jacobiske ---
Betydning, og som \emph{ikke kan eller skal omsættes i Viden}, eller retfærdiggjøres ved Tanken.

Det er saaledes kun ved Tvetydigheder og forvanskende Omformninger af de afgjørende
Begreber at Modsætningsforholdet kan skjules. I hvilke Former Troesvidenskaben end kan
fremstilles, saa maa den i Virkeligheden bestandig, fordi den hviler paa et med Tanken
absolut incommensurabelt, i Viden uomsætteligt Grundlag, mangle det, der constituerer
Videns Væsen og Form; den bevarer kun Videnskabens Navn, et Navn, der mister alt sit
Indhold ved at skulle sammenfatte det i Væsen Modsatte, og giver kun et Videnskabens Skin
ved partielt at bevare det blot Formelle. Det \emph{Rationelle} i Troesvidenskaben bliver den \emph{formelle Sammenhæng}; et kun \emph{tilsyneladende} Rationelt bliver \emph{Begrundelsen}, hvis
Nødvendighed bliver illusorisk, paa Grund af Principets Ubegrundethed for Tanken.
Videnskab efter Begrebets sande Betydning, sand Videnskab er kun een: den, der henter sit
Indhold fra Ideen, den universelle Videnskab, der ved sit Princip er i Enhed med det
personlige Livs Princip. Den er Videnskab uden particulariserende Prædicater, ikke
„christelig“ eller „troende“ Videnskab, men Videnskab alene.

Til Theologiens Forvanskning af Videnskabens Begreb svarer Forvanskningen af den positive
Religions Opfattelse af det for \emph{Guddommens} Væsen og Virken Eiendommelige: af \emph{Underet},
Udtrykket for den almægtigt villende Guds Virkeform og Aabenbaringsform. Underet, der
udtrykker den absolute Modsætning til Tanken og til den Tilværelses Lov, af hvilke Tanken
bæres, maa \emph{tilsyneladende} gjøres tilgængeligt for Erkjendelsen, naar en Viden om Troen skal
blive mulig. Det vil sige: \emph{Underets} \opage{117}\emph{Begreb maa bruges tvetydigt}, saa at det paa een Gang
betegner hvad der er et \textit{mirabile} og hvad der er et \textit{miraculum}, saa at der
anbringes en \emph{Quantiteren}, der drager det hen mod det blot Underlige, Forbausende, og
gjennem dette bringer en \emph{illusorisk Continuitet}, en „teleologisk“ Continuitet, tilveie med
det Almindelige og Naturlige. Men med Forvanskningen af Underet er \emph{Troens} Væsen
forvansket, thi i Underet er det positivt Religiøses Eiendommelighed concentreret og med
det staae alle Dogmernes eiendommelige Bestemmelser i uopløselig Sammenhæng. Til
\emph{Videnskabsbegrebets} Forvanskning svarer saaledes i Theologien \emph{Troesindholdets} Forvanskning.

\begin{center}---\end{center}

\section{Endeligt Resultat}
\label{sec:theologi-resultat}

Det endelige Resultat af disse Undersøgelser bliver altsaa, at de Mangler i videnskabelig
Henseende, der paavistes hos Martensen, maa blive typiske for Theologien overhovedet, og at
den Splid mellem Troes- og Videnskabsbestemmelserne, der factisk var tilstede hos
Martensen, maa, idet begge Bestemmelser sees fra Erkjendelsens Synspunkt, blive uundgaaelig
for hvilketsomhelst theologisk Forsøg.

Der bliver saa kun den Mulighed tilbage, at Forholdet kunde forandres naar Synspunktet
forandredes og de to Gebeter: Videns og Troens, saaes som principielt sondrede, og som ved
denne Sondring, der ligesom gjorde Indholdet incomparabelt, unddragne fra Sammenstødet. Hos
Martensen er, trods Polemiken mod denne Anskuelse, saaledes som den er udpræget hos Prof.
Nielsen, en saadan Sondring i Principet tilstede. Idet han lader Aabenbaringen „fra Først
til Sidst være praktisk“, og idet han accentuerer Villiens absolute Prioritet og dens
Differents fra Erkjendelsen, giver han i Virkeligheden Nielsen Ret med Hensyn til hans
Grundsondring. At Villien af Martensen sættes som høieste Enhed og constitutivt Princip,
fjerner ikke den principielle Sondring; thi det paavises ikke, hvorfra den Differents
kommer, der constituerer Erkjendelsen, hvorledes denne altsaa fremgaaer af det Constitutive.

\opage{118}Fra Kritiken af Theologiens uigjennemførlige Foreningsforsøg gaae vi da, forat prøve hiin
Mulighed, over til den kritiske Undersøgelse af den Theorie, der netop ved at statuere den
principielle Sondring af Troens og Videns Gebeter vil hævde begge i deres Gyldighed, og vi
beskæftige os først med Theorien i den Form, i hvilken den er fremstillet af sin Begrunder: \emph{S.\ Kierkegaard}.

\begin{center}---\end{center}


%% ============================================================
%% DEN PRINCIPIELLE SONDRING (pp. 118--224)
%% ============================================================
\chapter{Den principielle Sondring af Tro og Viden i Forsoningens
  Interesse}
\label{ch:sondring}

\section{Grundtrækkene af S.\ Kierkegaards Theorie}
\label{sec:kierkegaard}

\emph{S.\ Kierkegaard}. --- I den historiske Indledning ere Grundtrækkene af \emph{S.\ Kierkegaards}
Opfattelse af Forholdet mellem Tro og Viden allerede blevne angivne,
med særligt Hensyn til Sammenligningen med den Nielsenske Theorie. Paa
dette Punkt blive disse Grundtræk nærmere at udføre.

\emph{Kierkegaards} Opfattelse udvikles gjennem en Polemik mod Speculationen,
navnlig saaledes som den fremtræder i den Hegelske Philosophie; men
Polemiken rettes mod Speculationen fordi Kierkegaard i den seer det
stærkeste Udtryk for hele Tidens Misvisning: at have glemt at existere,
glemt hvad Inderlighed er. For den speculative Philosophie var det
Menneskeaandens høieste Opgave at hensætte sig i et Evighedens Medium,
leve et rent Evighedens Liv, og denne Opgave skulde realiseres ved den
rene Tanke, idet Individets Tænken af det Absolute sattes som Eet med
det Absolutes Tænken af sig selv, med den absolute, i Evighedens
Element værende, Tanke. Individets Liv, der fandt sit udtømmende Udtryk
i den logiske Tanke, gjennem dets Tænken af det Absolutes, al Sandhed
og Virkelighed omfattende, evige Livsproces, blev et Moment i denne
evige Proces, og frigjordes fra Endelighedsbestemtheden. Mod denne
\opage{119}Opfattelse af Menneskelivet og dets Opgave protesterer Kierkegaard.
Mod Speculationens Forflygtigelse af Existentsen, mod dets Glemmen af
Existentsens Krav til det det Ethiskes Fordringer underkastede Individ,
mod dets Overseen af Existentsens Betingelser, mod den illusoriske
Bevægelse i en abstract Evigheds Sphære accentuerer Kierkegaard
Existentsen, samler Alt om det Spørgsmaal: \emph{hvorledes udtrykker jeg
existerende Sandheden}? Som existerende er Mennesket ikke blot evigt,
hvad Speculationen vilde, men er en Synthese af det Timelige og Evige,
af Endeligt og Uendeligt; som existerende er den efter sit Væsen evige
Aand underkastet Endelighedsbetingelser, er i \emph{Vorden}; som existerende
skal Tænkeren i al sin Tanke tænke det med, at han er existerende; han
skal sætte al sin Tanke i Vorden. Men Vorden er \emph{Alternationen mellem Væren og Ikke-Væren}, ved Vorden bringes i hvert Moment \emph{Negationens} Bestemmelse ind i Aandens Væren og Tænken. Heraf følger saa med
Nødvendighed, at Mennesket ved Existentsen er forhindret fra at naae
den positive, betryggede Hvile i Evigheden, thi idet det, som
existerende, ikke kan fastholde det Uendelige og Evige, der er det
% PRINTED AS IS, p.119 [collation 2026-09-21]: page prints "man" (clear a-bowl, compared with the a in "fast-" on the same line) for sense-required "men"; confirmed by a blind second reader.
eneste Visse, i Uendelighedens Element, man skal fastholde det i
Endelighedens Element, i Existentsen, saa bliver Forholdet ikke den
betryggede Visheds, men det maa udtrykke den Negationens Bestemmelse,
der ligger i Vorden, og det Evige kan derfor af den Existerende, for
hvem det omsætter sig i et Tilkommende, kun haves for Erkjendelsen i \emph{Uvishedens} Form. Heri ligger saa, at den \emph{objective Viden i ingen} af
dens Former kan faae en saadan Sikkerhed, at den kan blive Grundlaget
for den uendelige Afgjørelse for Subjectet, for en Afgjørelse med
Hensyn til det, der angaaer Subjectets dybeste Trang. Den sandselige
Vished er et Bedrag, den historiske Viden er en Approximation, den
mathematiske Erkjendelse er en Abstraction, den speculative Tænkning
ligeledes en Abstraction og i sidste Instants en Hypothese i Forhold
til Tilværelsen; paa ingen af disse Videns Former kan derfor den
Individet uendeligt \opage{120}interesserende Afgjørelse bygges. Men kan Maalet \emph{ikke} naaes ad \emph{Objectivitetens} Vei, saa maa det naaes ad
\emph{Subjectivitetens}: i et Subjectivitetsforhold maa det Evige gribes, og
dette Subjectivitetsforhold er det, i hvilket Visheden om det Evige med
Troens Lidenskab gribes ud af den objective Uvished. \emph{Troen er} „\emph{den
objective Uvished, fastholdt i den meest lidenskabelige Inderligheds
Tilegnelse}“.

Opgaven for Individet er altsaa: \emph{at blive subjectiv}, og idet det, der
sættes som en Opgave, er det, ethvert Menneske kunde synes, uden videre
at være, saa faaer det at blive subjectiv, den bestemtere Betydning:
at blive det, hvad Mennesket efter sit Væsen er, at lade
Subjectiviteten udvikles, omskabes i sig selv ligeoverfor det Evige.
Først for en saadan Subjectivitet, der er den udviklede Mulighed af
Subjectivitetens første Mulighed, bliver Sandheden til. Subjectet, der
er henviist til sig selv, faaer den Opgave, uendelig at bekymre sig om
sig selv; men denne uendelige Bekymring om sig selv er Retningen mod
det Evige i Forhold til Existents; den fører hen mod det Ethiske og det
Ethisk-Religiøse, den væsentlige Sandhed og den eneste væsentlige
Erkjenden. Existentsens Krav forhindrer Individet fra at blive i det
Metaphysiske, stiller det den Opgave: at forstaae sig selv i Existents.
Idet Subjectet henvises til et Subjectivitetsforhold, henvises det
saaledes til det, at det kun naaer det gjennem et \emph{Valg} mellem dette og
Objectivitetsforholdet, og dette Valg er Valget mellem det \emph{personlige} og det \emph{upersonlige} Forhold til Sandheden. Der gives for den Existerende
ingen Mediationens Enhed af det subjective og det objective Forhold.
Mediationen, der vil forbinde dem, fører os kun tilbage til
Abstractionen, den forklarer ikke, hvorledes den evige Sandhed er at
forstaae i Tidens Bestemmelse af den, der ved at existere, selv er i
Tiden.

I det \emph{subjective} Forhold bliver \emph{Forholdets Hvorledes} det Væsentlige. I
det \emph{objective} Forhold gaaer Retningen \emph{bort fra} Individet, stræber hen
imod det objec\opage{121}tive, færdige Resultat, og dets Maximum er den
Selvmodsigelse, at den existerende Subjectivitet forsvinder, medens den
Objectivitet, der bliver til, subjectivt seet, enten er en Hypothese
eller en Approximation, fordi al evig Afgjørelse ligger i
Subjectiviteten. I det subjective Forhold gaaer Objectiviteten ud, og
Forholdets Hvorledes bliver det Afgjørende. I sit Maximum er dette
Hvorledes \emph{Uendelighedens Lidenskab}, og Uendelighedens Lidenskab er \emph{selve Sandheden}. \emph{Men Uendelighedens Lidenskab er netop Subjectiviteten
og saaledes er Subjectiviteten Sandheden}.

\emph{Troesforholdet}, Udtrykket for Subjectivitetsforholdet, for den
uendelige lidenskabelige Interesses Forhold, er saaledes det \emph{høieste},
fordi det i Existentsen bliver det eneste mulige Forhold til Sandheden,
Sandhedens Form i Existentsens Virkelighed. Sandheden i den objective
Tænknings Opfattelse, som Enheden af Tanke og Væren, bliver kun en
Abstractionens Chimære og i Sandhed kun en Skabningens Forlængsel, ikke
fordi Sandheden ikke er dette, men fordi den Erkjendende er en
Existerende, og saaledes Sandheden ikke kan være dette for ham saalænge
han existerer, fordi Existents spatierer Tanke og Væren og for den
existerende Aand sætter Sandheden selv i Vorden. Troens Lidenskab, den
uendelige lidenskabelige Interesse, bliver saaledes den Form, i hvilken
det Evige gribes; i den er det Evige paa \emph{negativ} Maade tilstede i
Subjectiviteten, og den \emph{negative Uendelighed} bliver i \emph{Existentsen det ene sande Udtryk} for det Uendelige, \emph{den uendelige negative Beslutning} bliver \emph{Individualitetens Form for Guds Væren i den}. Og \emph{Troen i strængeste Forstand} er \emph{Troen paa det absolute Paradox}, der ikke som det
sokratiske constitueres ved, at den evige Sandhed forholder sig til den
Existerende, men derved, at den evige Sandhed selv bliver Paradoxet ved at \emph{blive til i Tiden}. Som Tro paa Paradoxet er Troen \emph{den uendelig lidenskabelige Fastholden af det for Tanken Absurde}, og netop ved det Absurdes Frastød er Lidenskaben potentseret \opage{122}til sit Høieste. Det \emph{absolute Paradox} er \emph{Christendommens} Centralpunkt, og Spørgsmaalet om,
\emph{hvorledes jeg existerende udtrykker Sandheden}, falder da sammen med det
Spørgsmaal, \emph{hvorledes bliver jeg Christen}? Svaret er altsaa: ved Troen
paa Paradoxet, og det, der i sidste Instants bærer Paradoxets Byrde, og
tvinger Individet til at fastholde det i Troen trods dets absolute
Utænkelighed, er \emph{Syndsbevidsthedens Fortvivlelse}, der ikke finder
Forsoning i det Ethiske, ikke i den Religiøsitet, der endnu hviler i
det ubrudte Immanentsforhold; men først i det Christelige, i Troen paa
„Guden i Tiden“.

\begin{center}---\end{center}

I denne Theorie sættes saaledes en \emph{Videns Selvbegrændsning}, begrundet i
den existerende Aands Grundforhold. Viden faaer sin, ved
Existentsmodsigelsen betingede, af den selv anerkjendte Begrændsning,
der bliver en Begrændsning i en dobbelt Henseende: i Henseende til
Vished og i Henseende til Concretion, og ved denne dobbelte
Begrændsning statueres et Incommensurabilitetsforhold mellem den
objective Viden og Personlighedens Stræben og Trang. Men under
Begrændsningen hævdes der Viden en \emph{berettiget} Plads, om end en \emph{underordnet} Plads, og en \emph{Gyldighed}, om end kun en \emph{relativ} Gyldighed.
Den faaer sin Gyldighed og Sikkerhed indenfor det \emph{Abstractes} og \emph{Upersonliges} Sphære, indenfor det Mathematiske og Logiske; den finder
sin Skranke hvor den naaer til det \emph{Concrete} og \emph{Existentielle}; den
bliver, paa Grund af denne Skranke, ikke afgjørende for Personlighedens
dybeste og høieste Interesser, for Forholdet til den evige Salighed;
den afgjørende Afslutning bliver ikke mulig paa dens Gebet. Dens
Sphære strækker sig indtil det Punkt, hvor Troens Sphære begynder,
hvor Troesforholdet sættes gjennem et Brud med den objective Viden, og
indenfor Troens Sphære faaer Tanken, idet den træder i et tjenende
Forhold til Troen, atter en Betydning ved den Opgave, at \emph{udfolde de} \emph{Be\opage{123}stemmelser}, i hvilke der skal existeres, og at \emph{hævde Paradoxet som Paradox}, at værne om Paradoxet mod ethvert Forsøg
paa at begribe det. En Begriben af Paradoxet vilde netop være en
Tilintetgjørelse af dets eiendommelige Betydning, af Troens specifiske
Gjenstand; derfor hviler \emph{Theologien}, der sætter en saadan Begriben
som sit Maal, paa en Grundmodsigelse.

Idet \emph{Troen} bliver til \emph{gjennem et Brud} med den \emph{objective} Viden, hævder den sig sit \emph{eiendommelige Gebet} og sin \emph{principielle} Gyldighed; den
bliver --- medens den er utilegnelig for Viden --- utilgængelig for
Videns Angreb. Viden og Tro faae hver sin Sphære, til hvilken de
forholde sig som eiendommelige Principer. Troen som religiøs Tro
hviler paa et Valg, der i en høiere Potents udtrykker den samme
Subjectivitetens Afgjørelse som det Valg, ved hvilket Individets
ethiske Liv constitueres. Forsoningen mellem Tro og Viden begrundes i
Sondringen af Sphærerne og Principerne, og i Underordningen af Viden
under Troen. Troesvisheden, der seierrig løfter sig op over Tvivlen,
er den høiere Vished end Erkjendelsens, der, i Forholdet til det
Concrete og Virkelige, vedblivende er behæftet med Tvivlen.

\section{Dens Forhold til andre samtidige Opfattelser}
\label{sec:samtidige}

Gjennem det Udviklede bliver Forholdet mellem \emph{Kierkegaards} Opfattelse
af Religionen og Troen og andre samtidige Opfattelser tydelig.

Den \emph{speculative Religionsphilosophie} havde ladet Religionen være et
Trin i Aandens Erkjendelsesudvikling, der fandt sin Fuldkommelse i
Philosophien, i den speculative Tanke, der med logisk Nødvendighed
udfoldede sig af den lavere Form. Først i Begrebet lod den Troens
Indhold finde dets sande Form; ved at omsættes i den philosophiske
Erkjendelses Form fik Troesindholdet en høiere Sandhed; thi Tankens,
Begrebets Sandhed, er den høieste Sandhed. Hos \emph{Kierkegaard} sees
Sphærernes Forhold fra et andet Synspunkt. Videnskaben bliver Et,
Livet et Andet; Udfoldelsen af Tanken i Objectivitetens Form, i
Mulighedens Medium, efter Nødvendighedens Lov, skilles fra Sandhedens
\opage{124}Tilegnelse i Individet i Subjectivitetens Form, i Virkelighedens
Medium, efter Frihedens Lov. Sandhedens personlige Existents bliver
accentueret som det Høieste, Subjectivitetsforholdet bliver denne
Existentses Form, og indenfor de personlige Existentsformer bliver den
religiøse den høieste, og indenfor den atter Troesforholdet til det
for Tanken absolut Utilgængelige: det Paradoxe, Christendommens
specifiske Bestemmelse. Tankens Opgave bliver her ikke at assimilere
Troesgjenstanden, at opløse den i Begrebet, men at opdage Tankens
Modsætning: det absolut Utænkelige, Paradoxet, og at fastholde det
som Paradox forat der saa kan troes i Kraft af det Absurde.

Den \emph{almindelige Theologie} fordyber sig \emph{objectivt} i lærde
Undersøgelser; den vil gjennem disse bringe frem hvad Christendommen
er som \emph{objectiv} Lære; ad Historiens og Kritikens Vei vil den sikkre
sig Visheden herom, og under denne Stræben, der, efter Sagens Natur,
taber sig i en uendelig Approximation, lader den det være
Forudsætningen, at den Undersøgende er Christen, og at det, det
kommer an paa, er det videnskabelige Resultat, hvis Tilegnelse skal
følge af sig selv. For \emph{Kierkegaard} er Existentsens uendelige
Afgjørelse Hovedsagen: Approximationen, ved hvilken netop Afgjørelsen
forhindres, er Bedraget; Tilegnelsen er Opgaven, og det Væsentlige
ikke et Hvad, men et Hvorledes. Der spørges ikke om, hvad nu
Christendommen historisk er; men der spørges: hvorledes bliver jeg
Christen? Sandheden er Tilegnelsens Inderlighed, Sandheden er
Subjectiviteten.

Hos \emph{Feuerbach} bliver, ligesom hos Kierkegaard, Videns og Troens
Sphærer sondrede, men idet Feuerbach bestemmer Religionen som det
Irrationelle, hævder han Tankens Ret imod Troen og sætter Religionen
som det Usande, der skal forlades. Modsætningen mellem Feuerbach og
Kierkegaard i Opfattelsen af Religionen og særligt af Christendommen
er derfor den størst mulige, og dog er der ingen Samtidig, med hvem
Kierkegaard har flere Berøringspunkter i Bestemmelsen af hvad
Christendommen \opage{125}er. Dette har heller ikke undgaaet ham, og han har,
netop paa Grund af sin Opfattelse af Christendommen, ikke kunnet
forundre sig derover. Er Individets Forhold til Christendommen et
Subjectivitetsforhold, og bestemtere et saadant, i hvilket
Subjectiviteten er i sin høieste Potentsation, i sin meest energiske
Lidenskab, saa kan, foruden den Troende, der i Lidenskab fastholder den
som Sandhed, kun den opfatte dens Væsen, der i Lidenskab forkaster
den. Den Sidste vil, staaende udenfor Christendommen, afgjørende og
skarpt kunne sige hvad Christendommen er, han vil kunne skildre den
saaledes, at den Troende i Et og Alt kan vedkjende sig hans Skildring.
Det er upaatvivlelig Feuerbach, paa hvem Kierkegaard har tænkt, naar
han (i Afsl.\ Efterskrift S.\ 473), efterat have omtalt Theologiens og
Opvakthedens Forvirren af det Christelige, siger: „paa den anden Side:
en Spotter angriber Christendommen og foredrager den til samme Tid saa
forsvarligt, at det er en Lyst at læse ham, at den, der er i
Forlegenhed for at faae den bestemt fremsat, næsten maa tye til
ham.“ Vi finde ogsaa en Samstemmen i Opfattelsen i en Række af de
væsentligste Punkter. Feuerbach og Kierkegaard stemme overens i at
bestemme Religionens Forhold som et \emph{subjectivt}, om end Subjectivitetens
Art opfattes forskjelligt; de stemme overens i Opfattelsen af det
religiøse Standpunkt som et \emph{praktisk}, og af Troens \emph{Modsætning} til
Tanken; i Benægtelsen af Muligheden af en \emph{Viden} om et Væsen, der er \emph{høiere end det menneskelige}; i at hævde Kategorien: \emph{den Enkelte} for
den Religiøse; i Opfattelsen af \emph{Religionens absolute Forhold} som \emph{sondrende sig fra} og derfor i Virkeligheden \emph{negerende de relative} Forhold; i at fastholde \emph{Incarnationens ligelige Forhold til alle efterfølgende} Tider; i at opfatte Christendommens Forhold til Verden
som et \emph{uforanderligt}, som et \emph{blivende Modsætningsforhold}; i
Opfattelsen af \emph{Lidelsen} som Christendommens Kjendemærke, af \emph{Afdøen}
som den Troendes Ønske; i Bestemmelsen af den \emph{Christnes Forhold til det Naturlige}, i \emph{Begrændsningen af de Christ\opage{126}nes} \emph{Sympathie til de Christne}, og i den nøie Forbindelse, i hvilken de sætte Begreberne: \emph{Gud} som det absolute τέλος, og den \emph{evige Salighed}. Disse Paralleler
vilde endnu yderligere kunne forøges, men de anførte ere de meest
fremtrædende.

Forholdet mellem \emph{Kierkegaards} og \emph{Nielsens} Opfattelse er i det
Foregaaende omhandlet og saaledes bestemt, at de to Former af Theorien
i alle for Problemets Udvikling betydningsfulde Hovedpunkter stemme
overens, medens det af Nielsen Tilføiede ikke fører Problemet videre
og er et Selvmodsigende og Uholdbart. Paa Grund af dette Forhold bliver
den sammenhængende Fremstilling af Theoriens secundære Form
umiddelbart føiet til Fremstillingen af den oprindelige, og den
udførligere Kritik af Standpunktet knyttet til hiin, for saa at
fuldstændiggjøres ved en Kritik af det for Theorien i dens primitive
Form Eiendommelige.

\section{Grundtrækkene af Prof.\ R.\ Nielsens Theorie}
\label{sec:nielsen-theorie}

\emph{R.\ Nielsen}. --- Idet \emph{Nielsen}\footnote{Ved Fremstillingen af Nielsens Theorie er der, foruden til Nielsens
egne Skrifter, kun taget Hensyn til den authentiske Fremstilling af Theorien i \textit{Dr.}\ \emph{Heegaards} Indledning til den rationelle Ethik.}, ligesom \emph{Kierkegaard}, bestemmer
Grundforholdet mellem Tro og Viden som et Sondringens, og nærmere som en principiel Sondrings
Forhold, og netop gjennem Sondringen søger Muligheden af at hævde
begge Principer og begge Sphærer i deres Gyldighed og lade dem bestaae
som forsonede i den samme Bevidsthed; saa concentrerer han sin Theorie
i den Sætning: de \emph{absolut uensartede} Principer: Tro og Viden kunne, netop \emph{paa Grund af deres absolute Uensartethed}, forenes i samme Bevidsthed. Den \emph{principielle} Modsætning, der her er bestemt som en
Modsætning mellem \emph{Tro og Viden}, bringes ogsaa ind under andre
Bestemmelser, der allerede findes hos Kierkegaard: under Bestemmelserne
Liv (Existents) og \emph{Viden}, \emph{Subjectivitet og Objectivitet}, \emph{personlig og
upersonlig Sandhed}, eller under Bestemmelserne: \emph{Villie og Viden}, og denne
\opage{127}sidste Modsætning, og gjennem den de andre, føres saa af
Nielsen tilbage til den \emph{metaphysiske} Modsætning, Modsætningen mellem
Grundideerne: \emph{Magt og Viden}, eller \emph{det Gode og det Sande}.

Modsætningen mellem \emph{Villie og Viden}, mellem \emph{Villiesprincipet og
Tænkningsprincipet}, bliver den fremherskende i de sildigere, mere
udførlige Meddelelser fra Nielsens egen og fra hans Referenters Haand.
Begge de modsatte Aandsvirksomheder søges opfattede som Totaliteter,
saaledes, at der i Erkjendelsesvirksomheden er en Villie tilstede, men
overalt som tjenende Tanken, som en Erkjendelsesvillie, medens der i
Villiesvirksomheden skal være en Tanke tilstede, men som underordnet
under og tjenende Villien, som en Villiestanke. Til \emph{Villien} føres \emph{Troen} tilbage; det Religiøse, der opfattes som \emph{udelukkende praktisk}, knyttes ved \emph{Villiesprincipet} uadskilleligt sammen med det \emph{Ethiske}, der
har sin høieste Form i det \emph{antirationelle Religiøs-Ethiske}.

Fra det \emph{psychologiske} Udgangspunkt accentueres da Villiens af Viden
uafledelige Eiendommelighed. Villien skal vel forudsætte Tanken, men
ikke den theoretiske, men kun den praktiske Tænkning, og i denne
Forudsætten skal den være ligesaa oprindelig som Tanken. I sin
Eiendommelighed har Villien sin principielle Betydning og sit Ideal
ligesom Viden, og idet Villiesbestemmelserne og Vidensbestemmelserne
udvikle sig i deres Selvstændighed, beherskede af det eiendommelige
Princip, danne de to i sig afsluttede Totaliteter, der ikke komme i
Conflict med hinanden. Dette ligger netop i den absolute Uensartethed.
Modsætningens Form er Udtrykket for den factiske Dualisme i den
menneskelige Aand. I det Absolute er Videns og Villiens Princip: det
Godes Idee og det Sandes Idee, i concret, absolut gjennemsigtig Enhed:
i den endelige Aand, den psychologiske Bevidsthed, er en Dualisme
tilstede, der lader det i Ideen Forenede træde ud i \emph{absolut} Modsætning. Opgaven for Individet er netop at begribe Modsætningen
som saadan og derigjennem at vinde Forsoningen. Philosophiens
Hoved\opage{128}opgave er at opdage den menneskelige Videns, altsaa sine egne,
Grændser, og, idet den opdager dem, at anerkjende det andet Princip,
der, tilligemed den hele Sphære af Bestemmelser, der knyttes til det,
afspeiles af Viden, men ikke assimileres, ikke begribes. Idet de to
Sphærer rummes i samme Bevidsthed, maae de vel berøre hinanden, og
forsaavidt have en Grændse mod hinanden, men Grændsen er ikke en
hæmmende, indskrænkende Grændse; thi den er sat indenfra ved hvert
Princips eiendommelige Bestemthed, og „det absolute Mysterium“
forhindrer Sammenstødet. Medens Viden saaledes, afspeilende sin
Modsætning, opklarer Principernes Forhold, udfolder den sig uhæmmet
indenfor sin egen Sphære, kun følgende sine egne Love: den udfolder en
Videnskab paa Ideens Grundlag, der omfatter alle de humane Livsformer
og giver et sandt Udtryk for den concrete, reale Tilværelse. Og
samtidig udfolder Troen sit Indhold for den troende Bevidsthed, hvis
Erkjendelse er praktisk, ikke theoretisk. Kløften mellem Religion og
Philosophie skal Ethiken udfylde i sin dobbelte Skikkelse som rationel
og religiøs Ethik. Grændsevidenskab er Religionsphilosophien, der som
Videnskab hører under Videns Kreds, medens den, som havende det
Religiøse til sit Indhold, skal orientere i Troens Sphære. Den skal
vise, hvorledes eet og samme Begreb kan blive uensartet med sig selv
eftersom det udvikles til „Videns-“ eller til „Troesbegreb“. Den
skal paavise en alsidig articuleret Vexelbestemmelse, en nøiagtig
gjensidig Tilsvaren, mellem de uensartede Begreber, hvilket skeer
saaledes, at Vidensbegreberne tjene til intellectuelt Underlag,
Troesbegreberne derimod træde i Forgrunden og udvikles systematisk.

Idet Villien, der har sin Autonomie i Kraft af sit eiendommelige
Princip, og i Valget gjør sin Selvstændighed gjældende, autonomisk
potentserer sig gjennem bestandig intensivere Concentrationer,
fremtræder den fremskridende Række af Villiesstadier, det Ethiskes
Stadier, hvis Fremadskriden betegnes ved, at Erkjendelsesmomentet mere
og mere absorberes og bliver usynligt, at Henblikket til Theorien mere
\opage{129}og mere fortrænges af Villiens Energie. Det \emph{rationelt} Ethiskes Maximum
er tilstede hvor Subjectiviteten, i Kraft af en Villiesbeslutning,
vender Theorien Ryggen. Men selv over dette Standpunkt skrides der ud
i det \emph{paradoxe, antirationelt} Ethiske. „Concentrationen er her saa
afgjørende, at alle Hensyn til Theorien svinde“. Villiens Energie er
saa intensiv, at den ikke alene ikke tager Hensyn til Theorien, men
endogsaa bryder med den. Ja end ikke den rationelle Ethik paa sit
Maximum kan her komme i Betragtning, fordi dens Opgaver ere immanente,
dens Tro en Fornufttro. Paa den paradoxe Ethiks Høide tales der til
Mennesket i lutter Imperativer; det handler i Kraft af det Absurde.

Det \emph{Ethiske} bliver her sat i nøieste Forhold til det \emph{paradox
Religiøse}, ligesom overhovedet dets Princip er sat i Enhed med
Religionens Princip; og det Ethiske opfattes som \emph{ved sig selv} visende
hen til det Religiøse, fordi \emph{det} Gode kun skal begribes ved \emph{den} Gode.
Der bliver derfor allerede paa det Ethiskes Gebet Tale om Tro, men
denne Tro er en „rationel Tro“ inden den, gjennem den fremskridende
Villiesconcentration, bliver en Tro i Forhold til Paradoxet. --- \emph{Religionen} har sine Stadier ligesom det Ethiske, og de betegnes ligesom dettes ved
Villiens og Troens Fordybelse, en Fordybelse, der bliver tydelig naar
f.~Ex.\ den græske Religion med dens Myther sammenlignes med Jødedommen
med dens Forestilling om Jehovah som hellig Villie, og med
Christendommen.

Idet de afsluttende Stadier af det Ethiske og det Religiøse naaes,
forvandler det oprindelige tilsyneladende Coordinationsforhold mellem
Tro og Viden sig til et Subordinationsforhold. Troesforholdet og
Villiesidealet bliver det Høieste, Viden det Underordnede. Der tales
om „Religionens overvidenskabelige Sandheder“, om „et Høiere end
Videnskaben“, om en „overvidenskabelig Teleologie“. Hertil svarer
saa nøiagtigt Opfattelsen af Villien som det egentlige Constitutive i
det menneskelige Væsen. Paa det Ethiskes afsluttende Stadium, hvor
Bruddet er skeet med Theorien, \opage{130}gaaer hele Livet saaledes op i
Villiesacter, i Handling, i en Arbeiden paa den Opgave at blive en
ethisk Charakteer, at der ingen Plads bliver for en selvstændig
Erkjendelse, for et Liv i Viden. ---

Ved den givne Fremstilling af Hovedpunkterne i Kierkegaards og Nielsens
Theorie vil den tidligere Paavisning af Overensstemmelsen mellem dem
være bleven suppleret, medens de virkelige og tilsyneladende Afvigelser
tillige bestemtere træde frem. Gjennem de efterfølgende Udviklinger vil
Forholdet mellem den oprindelige Theorie og den afledede endnu
fuldstændigere blive tydeliggjort.

\section{Kritik af Nielsens Theorie}
\label{sec:nielsen-kritik}

Kritiken af Theorien i den Skikkelse, som \emph{Nielsen} har givet den, maa
først holde sig til Bestemmelsen af \emph{den principielle Modsætnings} Form.
Men i flere af de Udtryk, der ere givne for Modsætningen, navnlig i
dem, der ere tagne fra \emph{Kierkegaard}, viser der sig en Tvetydighed og en
Vilkaarlighed i Opfattelsen, der gjøre dem uskikkede til at tjene til
Udgangspunkter for Kritiken; de maae kun belyses for at gjøre
Forholdet mellem Nielsen og Kierkegaard tydeligere.

Det her Sagte gjælder Modsætningen mellem \emph{Existents} (Liv) og \emph{Viden} (Theorie), og mellem \emph{Subjectivt} og \emph{Objectivt}. Hvad den første
Modsætning angaaer, saa tages \emph{Existentsens} Begreb i en vilkaarlig
Dobbeltbetydning, paa hvilken en sophistisk Slutning bygges.
\emph{Existentsen} tages nemlig \emph{først}, forat hævdes som almeen Forudsætning
og nødvendig Betingelse for enhver concret Bestemmelse, altsaa ogsaa
for Viden, i Betydningen af \emph{individuel Væren i Almindelighed}, og \emph{derefter}, forat der gjennem Begrebet Existents kan hævdes det Praktiske
en Prioritet og en Overvægt, i Betydningen af \emph{Villiesværen}, („Existentsen er
praktisk“). Uden denne Tilsnigelse vilde Argumentationen, der skal føre til, at det er
muligt at sætte Existentsen øverst og dog beholde Theorien, men umuligt at sætte Theorien
øverst og dog beholde Existentsen, og \opage{131}videre til at godtgjøre Troens Superioritet over
Viden, slet ikke kunne skride fremad. Og i Beviset for hiin Sætning indeslutter den første
Præmis: en Bevidsthed kan existere uden at være theoretisk, men ikke være theoretisk uden at existere, tillige en Usandhed. Al Væren i Individualitetens Form er Existents; der
kan altsaa gives en saadan Existents, der ikke er Bevidsthed, ikke theoretisk. Men naar
\emph{Bevidsthed} bliver Existentsens Subject, saa \emph{indeslutter Existents Theorien, thi
Bevidstheden er kun existerende som værende theoretisk}.

Uklarheden opstaaer derved, at Nielsen optager Kierkegaards Bestemmelser, men forandrer
Forudsætningerne. Kierkegaard accentuerede i Bestemmelsen: Existents, \emph{Alternationen
mellem Væren og Ikke-Væren}, der satte den Existerendes Tanke i \emph{Vorden}, og afskjar den
objective Videns Vished fra den. Han fik da Modsætningen mellem \emph{existentiel} (subjectiv) Tænkning, og \emph{objectiv} Tænkning, og den \emph{subjective} Tænkning blev for ham
den Tænkning, hvis Maal var \emph{Existentsinderlighed, Existeren i det Tænkte}. Nielsen
lader Viden udfolde sig til et System, med en anden Vished og et andet Forhold til
Tilværelsens System end hos Kierkegaard. Existentsen opfattes da ikke som det, der, som
Bestemmelse for den Tænkende, forhindrer den objective visse Viden; men Existentsen faaer
\emph{Enkelthedens} Bestemmelse i Modsætning til \emph{Almindelighedens}, og da Almindeligheden er
\emph{Videns} Bestemmelse, sættes \emph{Existentsen} som \emph{praktisk}. Det, at Existentsen, efter
Kierkegaards Opfattelse, bliver det, hvortil der maa tages Hensyn, hvorfra der ikke uden
Strid med det Ethiske og uden indre Modsigelse kan abstraheres, bliver saa for Nielsen til,
at \emph{Existentsen} ɔ: det \emph{Praktiske}, er Forudsætningen, men, forat hævdes som saadan, maa den
tages i den ovenfor paaviste \emph{Dobbeltbetydning}.

En lignende Uklarhed og Vilkaarlighed træder frem hvor Modsætningen bestemmes ved det \emph{Subjective} og det \emph{Objective}, af hvilke hint forholder sig til \emph{Livet}, denne til \emph{Videnskaben}. Hvor Modsætningen bestemmes ved \opage{132}disse Former (f.\ Ex.\ i \emph{phil.\ Propædeutik}
§ 15) tages „Subjectiviteten“ og „det Subjective“ i \emph{Dobbeltbetydning}, saaledes at
det snart kommer til at udtrykke det tilfældigt Individuelle og det Egoistiske, der falder
udenfor det objectivt Gyldige, snart Subjectivitetens almene Form. Ogsaa her er Forholdet
til \emph{Kierkegaard} bestemmende. For \emph{Kierkegaard} betegnede den subjective Sandhed den \emph{tilegnede} Sandhed. Denne Sandhed er paa det Ethiskes Gebet \emph{det Almene, der realiseres i Existents, og som maa tænkes forat der kan existeres deri}. Paa det \emph{paradox Religiøses}
Gebet har den Sandhed, der skal tilegnes, Paradoxets Bestemmelse. I alle Sphærer sættes
Sandhedens Tilegnelse, fordi Existentsen forhindrer den objective Vished, som betinget ved
en Act, der overalt kan betegnes som en \emph{Troens} Act. Hos \emph{Nielsen}, hvor det ikke er
Existentsens Usikkerhed, men den abstracte Modsætning mellem det Almene og det Enkelte, der
betinger Disjunctionen mellem Viden og Existents, bliver det Subjective, i Forhold til
Videns Objectivitet, sat under en ensidig Modsætnings Bestemthed, som Væren i
Individualitetens Modsætning til det Almene, og derfra kommer Dobbeltbetydningen, idet
deels den Kierkegaardske Bestemmelse eftervirker, deels den nye Betydning udelukkende
træder frem. Men idet Subjectivitetens Begreb opfattes i denne Betydning gjennem en ensidig
og abstract Modsætning, oversees det, at Subjectivitetens Form ogsaa rummer det \emph{objectivt Gyldige}, at Subjectet optager det som Norm i sig uden at Indholdet forandres, at den subjective Handling kan have objective Principer, at der er et objectivt Sandt og Gyldigt i den fornuftbestemte Subjectivitet. Subjectivitetens, Livets, Standpunkt kommer i hiin Opfattelse til at falde sammen med Egoismens Standpunkt.

Paa samme Maade som Subjectivitetens Begreb, faaer ogsaa \emph{Objectivitetens} Begreb en \emph{Dobbeltbetydning}. Hvor Subjectivitetens Hovedstadier udvikles\footnote{\emph{Nielsen}: Philosophisk Propædeutik, p.\ anf.\ St.} betyder det \opage{133}Objective snart det objective Ydre
eller det blot abstracte Almene, snart det det Subjectives Totalitet omfattende Almene og
Almeengyldige. Og idet den abstracte Modsætning definitivt fastholdes, bliver det
vilkaarligt ignoreret, at paa de forskjellige Stadier Subjectets Bestemthed --- hvad der
især fremtræder tydeligt paa det Stadium, der betegnes som den uendelige Subjectivitets,
--- netop er i Kraft af hvad der er bestemt som Tilværelsens objective Princip og at den
udtrykker dette. Det bliver overseet, at den subjective Sandhed netop er den tilegnede og i Tilegnelsen med sig væsentlig lige objective Sandhed.

Modsætningen mellem Liv og Videnskab, Subjectivt og Objectivt, føres af Nielsen tilbage til
Modsætningen mellem \emph{Villie og Erkjendelse}, der bliver det \emph{psychologiske} Grundlag for Modsætningen mellem \emph{Tro} og \emph{Viden}, ligesom denne sidste Modsætning har sit \emph{metaphysiske} Grundlag i Grundideernes Modsætning, i den absolute Subjectivitets ontologiske
Grundbestemmelser. Modsætningen mellem \emph{Villie og Erkjendelse}, den \emph{psychologiske} Grundmodsætning, maa blive Kritikens egentlige Udgangspunkt, Undersøgelsen af Bestemmelsen
af deres Forhold den fundamentale, og fra det \emph{Psychologiske} ville vi da, gjennem det \emph{Ethiske}, føres til det \emph{Metaphysiske}, hvor den definitive Afgjørelse skal finde Sted, det
endelige Resultat vindes.

\subsection{Den psychologiske Modsætning mellem Villie og Viden}
\label{ssec:psych-modsaetning}

Hvad der med Rette urgeres af Nielsen, er det, at Villien, som \emph{primitiv} Aandsfunction, har en Eiendommelighed, der gjør det umuligt at reducere
den til Tanken, i den Forstand, at den skulde blive en blot
Modification af denne. Deraf, at den, ligesom Erkjendelsen, har sin
Eiendommelighed, maa saa følge en \emph{total} Aandens Virken, en
Tilstedeværen af alle Aandens Functioner, indenfor hver af de to
Aandsvirksomheder, og en eiendommelig Grundform for disse, \opage{134}og de maae i
denne eiendommelige Form objectiveres for Bevidstheden, der omfatter det
Theoretiske og det Praktiske.

Men naar Nielsen med Rette hævder Eiendommeligheden, og med Rette
polemiserer mod en Ensidighed, der, i den Form, i hvilken han bekæmper
den, ikke har anden Repræsentant end Spinoza, saa taber han sin Ret
naar han ved en Subreption forvandler Eiendommeligheden til en
Modsætning, der fra Formen overføres paa Indholdet, og som nærmere
bestemmes som en \emph{absolut} Uensartethed mellem de to Aandsvirksomheder med Hensyn til Indhold, Maal og Ideal; og han bliver inconsequent baade
naar han statuerer den absolute Uensartethed medens han fordrer den
dermed uforenelige Enhed i det absolute Princip, og naar han fastholder
hiin absolute Uensartethed, skjøndt han ved ensidigt at sætte Villien
som det egentlig Constitutive i det menneskelige Væsen, lader den
blive et Høiere end Viden og begrunde denne.

Det Usande i disse videregaaende Bestemmelser vil blive tydeligt ved en
nærmere Undersøgelse af den Maade, paa hvilken Nielsen opfatter
Forholdet.

Naar vi gaae ud fra den Indrømmelse, at Erkjendelse og Villie maae
sættes som eiendommelige, selvgyldige Yttringsformer af Væsenet, der
ikke kunne reduceres til Modificationer af hinanden, saa bliver det
strax Opgaven at bestemme, Hvormeget der ligger i denne Indrømmelse og
om den maa føre til de postulerede Consequentser:
Subordinationsforholdet og den absolute Uensartethed. Nielsen bestemmer
Villiens Forhold til Selvet i den Sætning: „det i Subjectiviteten
Allersubjectiveste er Selvet; men det Selviske i Selvet er
Selvbestemmen, er Villie“. Ved at bruge Begrebet Selvbestemmen
saaledes, at det paa een Gang bliver det væsentlige Udtryk for
Selvhedsenergien og identificeres med Villien, falder Eftertrykket paa
denne sidste; den tilegner sig --- om end Tænkningens, ligesom ogsaa
den kunstneriske Phantasies, oprindelige Eiendommelighed udtales, --- i
Virkeligheden Principatet, og den nærmeste og umiddelbare Consequents
bliver, at den Sidestillen af Villie og Viden, til \opage{135}hvilken den
principielle Modsætning fører, bliver til Vidensprincipets
Subordination under Villiesprincipet: en Subordination, der under Theoriens
Udvikling træder frem i mange Former. Til denne Opfattelse bidrager ogsaa den ovenfor
paaviste ensidige
Opfattelse af Begrebet Existents (Liv), der identificerer det exclusivt
med det Praktiske, med Villien („Livet er i Villien“). Anderledes
stiller Forholdet sig naar hiin Vilkaarlighed undgaaes, naar vi, i
Overensstemmelse med hvad der paa et tidligere Punkt blev antydet, ved
Opfattelsen af det Menneskelige gaae ud fra den Bestemmelse, der
kategorisk sondrer Menneskelivet fra Livets lavere Stadier: fra \emph{Væsensuniversalitetens} Bestemmelse, der udtrykker sig i den
Forsigværen i den frigjorte Almindeligheds Form, der er
Selvbevidsthedens, og naar vi opfatte Væsenets Grundbestemmelse som
\emph{Energie}, nærmere som Selvhedsenergie. Begge Former: Erkjendelse og
Villie, blive da primitive Udtryk for det i Selvbevidstheden som Enhed
aabenbarede Væsens Energie, for den Selvets Selvbestemmen, gjennem
hvilken det, bestemmende sig, realiserer sig: i begge Functioner hævder
Selvheden og Energien sig under eiendommelige Grundformer. Opgaven
bliver da, nærmere at bestemme Yttringsformernes Eiendommelighed, og
Bestemmelsen af den vil afgjøre Spørgsmaalet om Modsætningsforholdet
og nærmere belyse det postulerede Subordinationsforhold.
Eiendommeligheden bliver at bestemme saaledes, at i Erkjendelsen
Selvhedsenergien nærmest har Almindelighedens og Idealitetens, i
Villien Enkelthedens og Realitetens Form. Men disse som Udgangspunkter
givne Bestemmelser vise hen paa hinanden og bestemme sig henimod
hinanden. Realiteten, der constitueres ved Villien, er en Realisation
af det ideelle Væsens Bestemmelser i Selvet eller udenfor Selvet.
Idealiteten, der virkeliggjøres gjennem Erkjendelsen, er en ved
Selvets Realitetsside betinget Idealiseren af et Realt og er en saadan
selv i det ideelle Selvs Selverkjendelse. Og paa samme Maade er
Almindeligheden, der charakteriserer Erkjendelsen, en Almindelighed, der
stræber henimod Concretionen, Punktualiteten, og hvis \opage{136}sidste Maal er
den concrete Erkjendelse af Tilværelsen, i hvilken Erkjendelsen af Selvet
efter dets individualiserede Væsensbestemmelser og væsentlige Forhold
er indbefattet; med andre Ord: det individuelt-almene Væsens
Gjennemsigtighed for sig selv som Led i en Tilværelsens Totalitet.
Omvendt er Enkelthedsbestemmelsen i Villien en saadan, der med
Nødvendighed udfolder det Almene i sig, fordi Villien kun er Villie som
det selvbevidste, ɔ: efter Begrebet almene og universelle, Væsens
Villie. Villien maa udfolde sig fra den ved det Enkelte bestemte, det
Enkelte sættende, Villie til den ved det, Væsenet udtrykkende,
Almeenfornuftige bestemte, det Almene realiserende Villie.
Villiesyttringens Enkelthedsform er her ikke ophævet, men
Enkelthedsbestemtheden er her bleven gjennemsigtig i det Almene. Paa
lignende Maade som med Almindelighedens og Enkelthedens Bestemmelser
forholder det sig med de dermed saa nøie sammenhængende:
Objectivitetens og Subjectivitetens. Erkjendelsen har nærmest
Objectivitetens Præg; gjennem Bestemtheden ved Andetheden, gjennem
% PRINTED AS IS, p.136 [collation 2026-09-21]: wrong sort r for d (print "Indivirer"); confirmed by a blind second reader.
Subjectets Forhold til de for alle erkjendende Indivirer fælleds
almeengyldige Tankeformer og Tankelove, der bestemme det med
Nødvendighed. Men Erkjendelsen fuldender sig kun i den Subjectet
omfattende ligesom om det for sig selv klare Selv concentrerede
Erkjendelse, og den er, som Virkeliggjørelse af en Subjectets
Væsensbestemmelse, som ideel Udvidelse af Selvet, Gjenstand for en
Interesse, der kan have det intensiveste Subjectivitetens Præg, som en
Erkjendelsens Lidenskab, og i den er overhovedet, fordi den er
\emph{Væsens}yttring, det \emph{subjective} Væsen virksomt \emph{efter sin Heelhed}, sin udeelte Personlighed. Paa den anden Side har Villiens Subjectivitet, jo
mere Villiesrealisationen fjerner sig fra det Instinctmæssige,
Driftagtige og Egoistiske, jo mere den bliver en Virkeliggjørelse af
det Almeenfornuftige, af Væsenets Lov, en Væsensvillen, Objectivitetens,
det Almeengyldiges, Bestemthed i sig.

Paa Villiens og Erkjendelsens ufuldkomne Udviklingstrin divergere de og
kunne fremtræde i Isolation; paa deres \opage{137}Culminationstrin ɔ: paa den
fornuftige Bestemtheds Trin, bliver der, medens Eiendommeligheden
bevares, den uopløseligste Sammenhæng og den fuldstændigste Harmonie.
Den til sit Maal naaede Erkjendelse vilde, som Udtryk for det i sig
gjennemsigtige Væsens Energie, ifølge det ene Væsens nødvendige Væren
i begge Former, umiddelbart omsætte sig i Villiens, ved
Væsenserkjendelsen, Væsenets Selvaabenbarethed, bestemte Energie, og
den Villie, der med den intensiveste Energie, i den meest omfattende
Betydning, var Væsensvillen: Realisation af Væsenet efter dets
Concretion, vilde være en Villie, der umiddelbart omsatte sig i
Bevidsthed, i Erkjendelse, og i en Erkjendelse, der var bestemt ved
Videns almene, objectivgyldige Love.

Ved det her Udviklede ophæves saaledes den af Nielsen statuerede skarpe
Modsætning mellem Villie og Viden, og mellem den til disse svarende
„praktiske“ og „theoretiske Bevidsthed“. Den praktiske Bevidsthed
skal være personlig, interesseret, i Conflict med Existentsen, bevæget
af Frygt og Haab, fuld af Selvhensyn; den theoretiske upersonlig,
contemplativ, uinteresseret, fortabt i det Almene, bevæget af
Nødvendighed, glemmende sig selv i Begrundelsen af det Objective. Den
relative Sandhed i disse Modsætninger er anerkjendt gjennem det ovenfor
Udviklede, men den abstracte og absolute Modsætning er benægtet.

Og naar det Spørgsmaal opkastes: om der, uden at de to
Væsensyttringers Eiendommelighed og Primitivitet opgives, kan
tilkjendes den ene af dem et ideelt Principat, saa maa dette
Spørgsmaal besvares paa en anden Maade end Nielsen har gjort det ved
at tillægge Villien Principatet; hvilket han gjør deels af den ovenfor
anførte Grund, og deels fordi Erkjendelsen skal standse ved det Almene
og ved Væsensnødvendigheden og ikke naae til det Enkelte som saadant og
til Friheden. Ud fra det ovenfor angivne Udgangspunkt maa Principatet
tillægges Erkjendelsen. Væsenets Universalitet udtrykker sig
\emph{umiddelbart} i den frigjorte Almindeligheds Form, der er Erkjendelsens
Grundform; Enkelthedsbestemtheden, der er Villiens Form, hviler paa
\opage{138}Almindelighedens Grund, og Villien er kun Villie som det erkjendende
Væsens Villie. Den ethiske Villen som primitiv Beslutning, der skaber
Forholdet til det Gode, i hvilket Villien constituerer sig som ethisk,
og som concret Villen af det Gode i dets Bestemthed, er betinget ved
Viden og er kun mulig ved denne. Bevidstheden om det \emph{Forpligtende} forudsætter i sin Klarhed og Sikkerhed Tanken.

Men naar saaledes Principatet tillægges Erkjendelsen, der i
Menneskevæsenet derigjennem faaer samme Forhold til Villien som i det
Absolute Idealitetsbestemmelsen til Realitetsbestemmelsen; saa bliver,
ifølge det ovenfor Udviklede, begge Væsensyttringers Eiendommelighed at
fastholde forat Villien ikke skal reduceres til en Modification af
Tænkningen. Ved Overgangen fra den ene Væsensyttring til den anden
skeer der en \emph{Omsætning, en Formforvandling}, hvis Art definitivt
bestemmer Forholdet. Idet Erkjendelsesindholdet omsættes i et
Villiesindhold, bliver Omsætningen nærmest at charakterisere som en
Omsætning fra Reflexionens og det Almenes Sphære til den umiddelbare
Værens og Enkelthedens. Idet Villiesindholdet omsættes i et
Erkjendelsesindhold --- og dette vil ikke blot sige: formelt gjengives
af Bevidstheden --- skeer Omsætningen i modsat Retning. I Beslutningens
Øieblik bliver det ved Reflexionen Adskilte sluttet sammen til Enhed,
Reflexionen gaaer til Hvile i Beslutningens Villiesact, og der handles
i denne Forstand reflexionsløst. Omvendt, idet Villiesacten med dens
Indhold omsættes i Erkjendelse, kaldes Reflexionen atter frem, den
umiddelbare Enhed opløses i Reflexionens Sondring, Handlingens Enkelte
oversættes i Erkjendelsens Almene.

Hvilken bliver nu altsaa Forholdets kategoriske Bestemthed ifølge denne
Omsætning? Væsensyttringernes Eiendommelighed bringer en
\textit{differentia specifica}, en Artforskjel, ind mellem Villie og
Erkjendelse, men hvorledes bliver Uensartetheden nærmere at bestemmes?
Vi ville her see Sagen alene fra det psychologiske Standpunkt, og see
bort fra det Postulat, at en absolut Uensartethed skulde være be\opage{139}grundet
i Ideeforholdene. Spørgsmaalet faaer da den Form: om Uensartetheden
bliver en saadan, at den, efter Nielsens Udtryk, kun tillader en
Villiens \emph{Afspeiling} i Viden som et Fremmedt, kun ved sig selv Bestemt,
der ikke bliver gjennemsigtigt for Viden, eller om Differentsen er
behersket af Væsensenheden. Svaret afhænger væsentlig af Opfattelsen af
\emph{Enkelthedsbestemmelsen} i Villien. Oversees det ved Opfattelsen af
denne, at Villien kun er Villie som det selvbevidste ɔ: efter sin Natur
universelle og almene, Væsens Villie, saa bliver Villiens Enkelthed og
Umiddelbarhedsmomentet i den ført tilbage til den umiddelbare
Naturbestemtheds Virksomhedsyttring i Driften, den sandselig-egoistiske
Attraa, og idet denne nærmest sees fra den Side, gjennem hvilken den
hænger sammen med den individuelle Naturværens uoplukkede Enkelthed,
der staaer i \emph{umiddelbar} Modsætning til Videns Almene, saa bliver den
Virksomhed, der udgaaer derfra, opfattet som overalt bestemt ved denne
Enkelthedens Form, og, hvor den ikke overhovedet unddrager sig
Bevidstheden, som den, der idethøieste gjengives af Bevidstheden som
det i sin Umiddelbarhed med Erkjendelsens Almene absolut
Incommensurable. Forlades derimod denne Opfattelse af Enkeltheden ---
der lidet vilde stemme med, at Friheden tillægges Villien --- og
fastholdes det, at Villiens Enkelthed er det universelle Væsens
Fremtrædelsesform, er \emph{Aands}bestemthed, bliver Forholdet et væsentligt
andet. I det saaledes bestemte Enkelte gaaer det Almene kun paa den
Maade til Hvile, at det Enkelte bliver gjennemsigtigt: Omsætningen er
en Omsætning fra Tænkningens, og udtrykkelig fra den „theoretiske“
Tænknings, Almindelighed til den personlige Levens Enkelthed; men
Individet som fornuftigt villende udtrykker i Enkeltheden det Almene,
giver det Væren i dens Form, der gjennemsigtig aabenbarer det Almene.
Formdifferentsen beherskes altsaa af den væsentlige Enhed;
Bevidsthedsindholdet bliver under begge Former: som Erkjendelsesindhold
og Villiesindhold, væsentlig det samme. Jeg erkjender Pligten --- og
det er atter: jeg erkjender den theoretisk --- f.~Ex.\ den Pligt, \opage{140}at
kæmpe for Fædrelandet; jeg beslutter mig til at existere i det
Erkjendte, og idet Beslutningen er virkeliggjort, er Indholdet, som
realiseret i den enkelte Villie, væsentlig det samme. Bestemtest
fremtræder Enheden i det Forhold, hvor selve Erkjendelsen bliver
Villiens Gjenstand.

Imod den fremsatte Opfattelse af det psychologiske Forhold mellem
Villie og Erkjendelse vil der fra det af Nielsen indtagne Standpunkt
kunne gjøres forskjellige Indvendinger gjældende, der nærmere maae
belyses, idet de berøre væsentlige Grundforhold. Vi begynde med de
Indvendinger, der kunne hentes fra virkelige eller tilsyneladende
Kjendsgjerninger, og gaae derfra over til dem, der hentes fra
Begrebsanalogier og fra Begrebernes Grundforhold.

\medskip

\noindent 1.~Mod det ovenfor opstillede Forhold mellem Erkjendelse og
Villie kan der hentes en Indvending fra det \emph{forskjellige} Forhold, i
hvilket Erkjendelse og Villie staae til \emph{Indvirkningen} paa det \emph{Reale} og til \emph{Fremkaldelsen} af \emph{Handling}. Hvad det Første angaaer, saa synes der
at være en Modsætning tilstede, der viser hen til en lignende
Modsætning mellem Villie og Viden som den Modsætning, der statueres
mellem det Absolutes Bestemmelser: en Modsætning, der, reflecteret i
den psychologiske Bevidsthed, skal vise sig som en absolut
Uensartethed. Villien indvirker nemlig, idetmindste indenfor visse
Grændser, umiddelbart bestemmende paa Organismen: Mediet for
Indvirkningen paa det reale Andet; Erkjendelsen synes derimod ikke at
have en Virkning i denne Henseende. Dette er dog kun tilsyneladende,
idet Erkjendelsens Energie, gjennem Udviklingen af det denne Energie
tjenende Nerveapparat, medfører Virkninger, der ere blivende i
Organismen. Differentsen bliver en relativ Differents med Hensyn til
det Bevidste og det Ubevidste i Virkningen. Erkjendelsesvirksomhedens
bestemmende Indvirkning paa Organismen er ubevidst, Villiens er, med
Hensyn til det sidste Maal, dette være nu Bevægelse eller Uddannelse,
bevidst, men de \opage{141}mellemliggende organiske Functioner, gjennem hvilke
Maalet virkeliggjøres, foregaae ubevidst og blive kun indenfor visse
Grændser bevidste gjennem en Erkjendelsesvirksomhed, ligesom de
Virkninger, ved hvilke Erkjendelsesenergien uddanner Organerne. ---
Hvad det andet Punkt: Indvirkningen paa Handlingen, angaaer, saa
urgeres der, at selv den fuldstændigste, meest gjennemførte Erkjendelse
som saadan skal være magtesløs til at skabe Handling, og dette oplyses
ved, at f.~Ex.\ den fuldstændigste Viden med Hensyn til Modet er
magtesløs til at skabe den modige Handlen, til at give den af Naturen
Feige det virkelige Mod\footnote{Andensteds udtrykker Nielsen sig ubestemtere: „den
virkelige Praxis kan ikke \emph{ligefrem} afledes af Theorien“. Om theor.\ og pr.\ Erkj.\ S.\ 66
(i Nordisk Universitetstidskrift, Supplementsheftet.)}. I disse Sætninger, der lide af den
væsentlige
Tvetydighed, at eet og samme Udtryk: Mod, snart betegner det
umiddelbare Naturanlæg, snart en ethisk Bestemmelse, blandes det Sande
med det Falske. Det er for det Første sandt, og en Sandhed, som Ingen
benægter, at der gives et naturbestemt Villiens Grundlag, en
naturbestemt Villiens Forudsætning, der ikke kan gives eller umiddelbart
erstattes gjennem Erkjendelsen; men hvad der gjælder om dette, gjælder
ligesaafuldt om Erkjendelsens, idet dennes, af det enkelte Individs
Selvbestemmen uafhængige, naturbestemte Forudsætninger ikke kunne gives
af Erkjendelsen; og, hvad der gjælder om Erkjendelsen, at den ikke kan
skabe Naturanlægget i begge Retninger, gjælder ligesaa fuldt om
Villien; saa at det følgelig bliver usandt, at der af Erkjendelsens
Forhold til Naturgrundlaget skal kunne drages nogen Consequents med
Hensyn til dens Modsætning til Villien. Der er for det Andet i hine
Sætninger den Sandhed, at Erkjendelsen, forat blive praktisk virksom,
maa undergaae den ovenfor berørte Omsætning, ved Beslutningen til at
existere i det Erkjendte. Men usandt er det, at Erkjendelsen skulde
være magtesløs i Forhold til Handlingens, navnlig den ethiske Handlings
Virkeliggjørelse --- en Consequents, der i sin strængeste Opfattelse
følger af den absolute Modsætning, der udelukker \opage{142}al Indvirkning af
Erkjendelse paa Villie. Den sande Erkjendelse af det ethisk Rette, af
Pligten og Loven som \emph{Væsens}bestemmelse, \emph{Væsens}fordring, forenet med den
sande Erkjendelse af de det Ethiske hæmmende Drifters og Lidenskabers
Væsen, vil saaledes give det i Menneskenaturen Væsentlige Overvægten
over det abnorme Hæmmende, at dettes Magt vil være brudt, og
Handlingens og Tilstandens Bestemthed resulterer nu, under hiin ovenfor
angivne Omsætning, ved Erkjendelsens Hjælp af hint Overmægtige. Dette
kan ogsaa udtrykkes saaledes: idet gjennem Erkjendelsen det efter
Anlægget sig selv i sin Concretion erkjendende Væsen actualiserer
Bevidstheden om en Væsenslov og et Væsensforhold, bliver det i
Erkjendelsen sig selv i concretere Bestemthed aabenbare Væsen
bestemmende i Forhold til det Væsenets Virken Fordunklende og
Hæmmende. Hvor Erkjendelsen factisk bliver magtesløs, vil, --- bortseet
fra de reent physiske Hæmmelser --- Magtesløsheden bestandig hænge
sammen med Erkjendelsens Ufuldstændighed. Den Betydning, der her er
tillagt Erkjendelsen, indeslutter ikke den Consequents, at Villien
bliver en blot Modification af den, og saaledes taber sin Primitivitet.
Forholdet bliver det, at Væsenet, der actualiserer sig som forsigværende
i Erkjendelsens Form efter sin universelle Bestemthed, ifølge
Personlighedens Enhed med indre Nødvendighed kræver den Fordobbling,
der finder Sted, idet det i Universalitetens Form Actualiserede, ved
Idealiteten Bestemte, actualiseres i Enkelthedens, i Realitetens Form i
Villien, men i denne Fordobbling bliver i gjennemsigtig Enhed med sig
selv. Viden opfattes altsaa ikke som det, der frembringer Villie, men
som det, der, som Væsenserkjendelse, som Form for Væsenets
Selvaabenbarethed, gjennem Væsenets Enhed giver den Væsenet udtrykkende
Villie Indhold, altsaa gjennem Væsensenheden bestemmer Villien.

\medskip

\noindent 2.~En anden Indvending mod den ovenfor fremsatte Opfattelse
kan hentes fra den Kjendsgjerning, at Villiesacter, Handlinger i deres
Afsluttethed, deels kunne have en ufuldstændig Erkjendelse til
Forudsætning, deels virkeliggjø\opage{143}res uden Hensyn til eller i Strid mod
Theorien. Paa Grundlag af denne Kjendsgjerning, der altsaa deels viser
hen til en factisk Incongruents, deels til en Incommensurabilitet, kan
der nemlig argumenteres saaledes, at der siges, at da Villien i begge
Tilfælde maa hente Bestemmelser i sig selv, saa viser dette tilbage
til et eiendommeligt \emph{Villiesprincip}, hvoraf saa følger Villiens og
Erkjendelsens absolute Uensartethed og gjensidige Uafhængighed. Hvilken
Gyldighed denne Argumentation har, vil blive tydeligt af Følgende. Hvad det \emph{første} af de ovenfor paapegede Forhold angaaer, saa er den
Kjendsgjerning ubestridelig, endog ubestridelig som almeen
Kjendsgjerning, at Handlingen, bestemtere: den ethiske Handling, til
Erkjendelsesforudsætning kan have en ufuldstændig Reflexion, en blot
partiel Erkjenden, saa at altsaa Beslutningen, idet den forvandler
dette Ufuldstændige, Partielle, til et Fuldstændigt, Afsluttet, idet
den forvandler det til Handlingens \textit{ratio sufficiens}, tilføier
et ved Erkjendelsesforudsætningen ikke Givet, og altsaa ialfald med
Hensyn til dette viser sig som ikke bestemt ved den, men som hentende
sin Bestemmelse i sig selv. Den ufuldstændige
% PRINTED AS IS, p.143 [collation 2026-09-21]: print has transposed "Tlifælde" (the earlier transcription had silently corrected it); confirmed by a blind second reader.
Erkjendelsesforudsætning suppleres i dette Tlifælde ved hvad man kunde
betegne som en praktisk Conclusion, for den ethiske Villies Vedkommende
som en paa en ethisk Tro hvilende Conclusion. Men analysere vi dette
Forhold nærmere, saa vil det vise sig, ikke at indeslutte hvad der
søges lagt deri. Det, der endeligt bestemmer Handlingens
Virkeliggjørelse, kan være det rent Driftagtige, Instinctmæssige, der
afgiver en „praktisk Grund“, som kan falde udenfor de theoretiske og
være dem modsat. Men dette Factiske viser foreløbig kun, at
Menneskenaturens \emph{sporadiske} Yttringer kunne divergere; thi
Drifttilskyndelsen er netop det Isolerede, Enkelte og Tilfældige, der
ikke er ført tilbage til Væsenet, ikke gjort gjennemsigtigt som
Væsensbestemmelse. Det beviser ikke definitivt Noget med Hensyn til
Villiens Forhold til Erkjendelsen, fordi Driften kun er Villiens
rudimentære Form, der med indre Nødvendighed udvikler sig til Villiens
sande Form, ud fra \opage{144}hvilken hint Forhold maa bestemmes. Vi komme
derfor den definitive Afgjørelse nærmere ved at betragte det Tilfælde,
hvor det, der bestemmer Beslutningen, har det Bevidstes, i høiere
Instants det Ethiskes, Charakteer. I dette Tilfælde er det Udtrykket
for den i sig concentrerede Personligheds ikke i Reflexionens Form
udfoldede intensive Virken; men i denne Personlighed er en forudgaaende
Erkjendelse virksom, der ligesom er nedfældet i umiddelbar Væren og i
denne virker som Totalitet, paa samme Maade som en
Erkjendelsesudvikling, der er gaaet tilbage til umiddelbar Væren i
Aanden, kan bestemme Aandens theoretiske Virken. Hvad Beslutningen
føier til den forudgaaende ufuldstændige Erkjendelse, er derfor ikke et
med denne Uensartet: det er et, den ufuldstændige, og som saadan
relativ abstracte, Erkjendelse til Concretion Supplerende, og dette
Supplerende viser gjennem den Væsenstotalitet (Personlighedstotalitet),
der er dets umiddelbare Kilde, tilbage til en forudgaaende, Villien, og
navnlig den ethiske Villen, bestemmende Erkjendelse. --- Hvad det \emph{andet} Punkt angaaer: at Villien kan bestemme sig ligegyldig mod eller endog
tiltrods for Theorien, saa er dette ligeledes en ubestridelig
Kjendsgjerning; men den Fortolkning, der lægges ind i
Kjendsgjerningen: at den skulde godtgjøre Villiens Utilgængelighed for
Erkjendelsen, dens antirationelle Charakteer, er det Uberettigede. Den
„Theorie“, i Modsætning til hvilken der skal handles, maa, ifølge det
ovenfor Udviklede, være en abstract og uigjennemført eller --- hvad
der ogsaa bestandig indeslutter en Abstraction --- en ensidig Theorie.
Med Hensyn til den første ligger Divergentsens Mulighed i det
Sporadiske, det endnu ikke til Enhed Sammenarbeidede i Menneskenaturen,
der ogsaa gjør sin Indvirkning gjældende indenfor det Praktiskes Sphære, \emph{tagen for sig alene}. Med Hensyn til den sidste gjælder det, at
naar Individet, i Kraft af den sunde Forstands praktiske Sands, i sin
Handling corrigerer en saadan Theorie, --- ligesom det, vel at mærke,
\emph{indenfor det Theoretiske} corrigerer den ifølge en theoretisk Sands, --- saa handler det, indenfor den praktiske Sphære, ligesom \opage{145}Theoretikeren
selv gjør indenfor den theoretiske naar han gjennem ubevidste, men
uundgaaelige, Inconsequentser, ud fra et Apriorisk, der virker med
Nødvendighed, corrigerer sin egen Theories Ensidighed; saaledes som
f.~Ex.\ \emph{Hume}, idet han benægter Causalitetsbegrebets Gyldighed for
Erkjendelsen, affirmerer den ved at bevare Kategorien: Indtryk, i
hvilken et Causalitetsforhold er indesluttet. Og den praktiske Sands
vil, naar man ikke \emph{overhovedet} vil benægte en Indvirken af Erkjendelse
paa Villie, fra en Side være at opfatte som Anbringelsen af et \emph{theoretisk} Correctiv, i Forhold til \emph{den Erkjendelse, der skal anvendes}.
Naar det \emph{ethiske} Individ handler ligegyldig imod eller endog i
Modsætning til ensidige Theorier, saa skeer dette, ligesom ved den
ovenfor berørte Anbringen af Correctiver, ud fra et Apriorisk, og dette
er nærmere at bestemme som et Apriorisk i det Ethiske. Begreber som det
Gode, Forpligtelsen o.~desl.\ ere saadanne aprioriske Bestemmelser, der
kunne fordunkles ligesom Forstandens aprioriske Bestemmelser, men som
aldrig kunne trænges saaledes tilbage, at de ikke udøve en ubevidst
Virkning. Men disse ethiske Grundbegreber, dette Aprioriske i det
Ethiske, ud fra hvilket Handlingen bestemmes, maae, som
\emph{Væsens}bestemmelser, ɔ: som Udtryk for det udeelte Væsen, for hvilket
ogsaa \emph{Viden} er Udtryk, ogsaa være Vidensbestemmelser, og det ikke blot
i den Betydning, at de „afspeiles“ i Viden, og saaledes maa der blive
en Væsensenhed mellem „det Gode, der danner Grundlaget for den
praktiske Existents, og det Gode, der afhandles under det Sandes
Idee“. Det Aprioriske i det Ethiske vil være et oprindelig Virkende,
men i sin Virken er det ifølge Væsensforholdet til Viden
gjennemskueligt for Erkjendelsen, og naar det skal faae Concretion,
naar det Gode skal udfolde sig til dette bestemte Gode, naar
Pligtbegrebet skal faae et bestemt Indhold, saa bliver Erkjendelsen
nødvendig Betingelse for Handlingens Sandhed. --- Naar der, forat
hævde det Praktiskes Uafhængighed af det Theoretiske, vises hen til
Theoriernes Mangfoldighed, kan der med samme Ret vises hen til de
mangfoldige Differentser i den prak\opage{146}tiske Bevidsthed: Differentser, der
kun kunne bringes til Enhed ved den sande Erkjendelse.

Fra det kriticerede Standpunkt vilde det her Udviklede vel blive
indrømmet i den Forstand, at Tænkningens Nødvendighed for Handlingen
anerkjendtes, men denne for Handlingen nødvendige Tænkning vilde saa
blive bestemt som \emph{praktisk Tænkning}, og mellem denne og den theoretiske
Tænkning vilde der blive sat en diametral Modsætning. Nielsen lægger et
Eftertryk paa denne Bestemmelse: \emph{praktisk Tænkning}, men den faaer
imidlertid hos ham ingensteds nogen virkelig Begrebsforklaring,
ligesaalidt som det tilsvarende Begreb: \emph{praktisk Erkjendelse}, og det
maa overhovedet være paafaldende, paa hvilken Maade han, under
Udviklingen af sin Theorie, anvender sine psychologiske Kategorier, uden
skarp Begrebsbestemmelse og uden Sondring af det væsentlig
Forskjellige. Saaledes bruges Bestemmelserne: Bevidsthed, Viden,
Tænkning, Erkjendelse, ganske iflæng. --- \emph{Praktisk Tænkning} maatte
nærmest betyde en Reflexion, der bestemmer Sammenhængen mellem Princip
eller Forudsætninger og Consequentser, mellem Midler og Øiemed, paa det
Praktiskes Gebet, altsaa, efter Forudsætningen, en formel Udviklen af
Consequentser af det, der ikke opløses i Reflexionen, og denne
Reflexion, der dog overhovedet ikke vilde faae nogen Plads paa det
paradox Ethiskes Gebet, maatte, idet Handlingen indtræder, tænkes som
absorberet, usynliggjort i Handlingens Enkelthed. En saadan praktisk
Tænken vilde imidlertid i sin \emph{Form} ikke adskille sig fra den
theoretiske Tænken, der uddrager Consequentser af et endnu ikke af
Erkjendelsen Assimileret, f.~Ex.\ af et blot empirisk Bestemt, eller
som danner Constructionsbegreber. Bestemmelsen af Udgangspunktet
(Forudsætningen) og af Øiemedet kan tænkes som sat ved en
Driftbestemthed eller ud fra den ethiske Subjectivitets Villiesbestemmen.
I første Tilfælde er den det, der skal gjennemklares ud fra
Væsenserkjendelsen, i sidste Tilfælde har den en saadan
Væsenserkjendelse til Forudsætning, men Væsenserkjendelsen er paa een
Gang theoretisk og praktisk. I den udprægede ethiske \opage{147}Charakteer kan den
praktiske Tænkning antage en saadan Umiddelbarhedens Form, at der
handles tilsyneladende reflexionsløst, ligesom ifølge en
Naturnødvendighed, der har Analogie med Geniets Virken. Men denne
reflexionsløse Handlen er ensartet med den Virken af den udviklede
Erkjendelse, der tilsyneladende overspringer Mellemleddene og
umiddelbart naaer Resultatet, men i Virkeligheden gjør det ved, ifølge
Reflexionens Uddannethed, uendelig hurtigt at gjennemløbe
Mellemleddene. Naar saaledes den praktiske Tænkning i sit Væsen ---
thi Væsenet bestemmes ved Formen --- bliver liig med den theoretiske,
der forudsættes gjennem og i den praktiske, saa gjælder det Samme om
den „\emph{praktiske Erkjendelses}“ Forhold til den theoretiske. De kunne
ikke adskilles ved den \emph{Tænkningens Form}, ved hvilken de komme tilveie;
thi den praktiske Erkjendelse, der individualiserer den almene Lov i
dens Anvendelse paa Subjectet, giver den existentiel Anvendelse, følger
netop de selvsamme Love som den theoretiske \emph{combinerende} Tænkning.
Differentsen skulde da ligge i, at den praktiske Erkjendelse
(Erkjendelsen „til praktisk Brug“) tages i \emph{Kantisk} Betydning om det \emph{ud fra en Subjectets Trang Postulerede}, der \emph{ikke} retfærdiggjøres ved
den theoretiske Erkjendelse\footnote{Af \emph{Troens} Gud skal der saaledes kun gives en \emph{praktisk}
Erkjendelse, jfr.\ Afhandl.\ om theor.\ og prakt.\ Erkjendelse. S.\ 72.}. Men et saadant
Postulat, der, forat have
Gyldighed, maa have sin Grund i \emph{Subjectets Væsensbestemmelse}, i det
Aprioriske i det Ethiske, har kun sin Sandhed ved at kunne erkjendes
gjennem Væsenets Selvaabenbarethed i Erkjendelsen. Postulatet er en
Anticipation i Kraft af en i Bevidstheden virkende apriorisk
Bestemmelse, svarende til de Anticipationer, der kunne finde Sted paa
den theoretiske Erkjendelses Gebet, og paa hvilke den almindelige
psychologiske Iagttagelse og Erfaring, f.~Ex.\ af Barnets Anvendelse af
Begreber, og Philosophiens Historie frembyder utallige Exempler. Hos
Nielsen sættes der imidlertid en \emph{absolut Modsætning} mellem \emph{theoretisk}
\opage{148}og \emph{praktisk} Erkjendelse. Der siges, at Valgfrihed, Valg, Beslutning,
kun kunne være Gjenstand for den sidste, ikke for den første. De skulle
vel falde ind under Tankens Belysning, men i Principet ligefuldt være
Villieskategorier og deres Erkjendelse derfor væsentlig praktisk. Det,
der erkjendes theoretisk, skal ikke kunne erkjendes praktisk, fordi i
enhver afgjørende Erkjendelse Indholdet bliver bestemt ved Maaden og
Maaden ved Indholdet. --- I disse Udsagn er \emph{praktisk Erkjendelse} brugt som
enstydig med \emph{praktisk Bevidsthed}, d.~v.~s.\ med den \emph{umiddelbare Bevidsthed} om \emph{Subjectets praktiske Bestemmethed}: en Bevidsthed, der
naturligviis i sin Umiddelbarhed sondrer sig fra den paa Reflexion hvilende Erkjenden, men i hvilken Forholdet er det \emph{samme som overhovedet i Erfaringsforholdet til Gjenstanden}. Indeslutter den praktiske
Bevidsthed ikke blot Bevidstheden om sin Gjenstands \emph{Væren}, men om sin Gjenstands \emph{Gyldighed}, saa viser dette kun, at den i Gjenstanden
forholder sig til et Apriorisk, men ikke, at dette Aprioriske er et for
den theoretiske Erkjendelse Utilgængeligt. Er den praktiske Bevidstheds
Indhold \emph{Væsensbestemmelse} i Subjectet, maa det kunne blive
gjennemsigtigt i \emph{theoretisk} Erkjendelse, begrebet \emph{ud fra Væsenserkjendelsen}. I dette Forhold indrømmer Nielsen dog kun en „\emph{Afspeiling}“, ikke en virkelig Begriben, en tilegnende Erkjenden af
Gjenstanden efter dens Virkelighedsbetydning. Det Gode, som praktisk
Livsidee, siges at indeholde et Umiddelbarhedens Overskud, en Villiens
Oprindelighed, der er principielt utilgængelig for Videnskaben, og
derfor bliver Forholdet kun Afspeilingens. Men denne Betegnelse:
\emph{Afspeiling}, er ikke en \emph{Kategorie}, men kun en uheldig \emph{Metaphor}, der,
hvorledes vi end see paa Sagen, fører ind i Modsigelser. Ved dette
Billede skal det Forhold udtrykkes, i hvilket det med Viden \emph{absolut Uensartede} er værende for Viden. Idet Intet absolut skal falde udenfor
Viden, men der dog statueres et med Viden absolut Uensartet, bliver det
nødvendigt at bestemme Forholdet saaledes, at Viden optager ligesom et
Billede af Tingen, \opage{149}men ikke det Afbildede selv efter dets
Virkelighedsbestemthed; den skal saaledes i Forholdet til det af et
andet Princip bestemte Ethiske, i „theoretisk Afspeiling for sig selv
kunne eftergjøre de Bevægelser, der existentielt, punktuelt, ja
antirationelt foregaae i Villien“; men om en assimilerende Begriben
skal der ikke være Tale. Men denne Opfattelse af Forholdet viser sig
let som uholdbar. \emph{Fastholdes} det \emph{Afspeiledes absolute
Uensartethed}, vil \emph{end ikke en Afspeiling kunne tænkes som mulig}, thi mellem det absolut Uensartede kan \emph{intet} Vexelvirkningsforhold finde Sted. Hvert Princip vilde udfolde sine
Bestemmelser af sig, ligegyldige mod hinanden og uden nogen gjensidig
Bestemmen, og man vilde, paa Grund af den absolute Uensartethed, i den
uafhængige Parallelisme ikke kunne faae Noget, der svarede til \emph{Spinozas}
\textit{idea corporis}: en Bestemmelse, der hos ham, idet den gjør det
med Tanken absolut uensartede Legemlige til Tankens Indhold, kun er et \emph{Postulat}, da Tankens Begreb Intet indeholder, der viser hen til Legeme.
Afspeilingens Mulighed er altsaa, under den antagne Forudsætning, et
uberettiget Postulat. Men selv om man \textit{per impossibile}
supponerede, at det med Viden \emph{absolut Uensartede} kunde være tilstede i
Viden, vilde Forestillingen om en Afspeiling dog kun føre til
allehaande Modsigelser. Den Afspeiling, fra hvilken Udtrykket hentes,
er, som det fremgaaer af den Maade, paa hvilken Betegnelsen bruges, den
normale Speilflades troe Afbilden af det, der afspeiles ɔ: af en
legemlig Gjenstands \emph{Overflade}. Ved en saadan er der en virkelig \emph{Lighed}, en virkelig \emph{Tilsvaren} mellem Speilbilledet og det umiddelbart for Synet
Værende. Men naar Afspeilingens Bestemmelse anvendes paa Forholdet
mellem Viden og Villie, paa hiins ideelle Gjengiven af dennes
Bestemmelser, saa er der \emph{ingen Tilsvaren}, idet i den ideelle Gjengiven
af det ideelle Væsentlige dette Væsentlige ikke gjengives efter sin
Væsensbestemthed, og saa maa den Lighed, der skulde berettige
Anvendelsen af Begrebet, forsvinde paa Grund af det optagende Mediums
\emph{absolute Formforskjellighed} fra det Optagne. Sætter \opage{150}man nemlig den
afspeilende Viden som under Afspeilingen virkende \emph{efter sin Eiendommelighed}, saa bliver det „Afspeilede“ i Realiteten \emph{ikke afspeilet}, men \emph{absolut forvandlet} ved det optagende Mediums med det
uensartede Natur. Vilde man undgaae dette ved at lade Viden, trods
Selvvirksomheden, \emph{ikke} forvandle det Afspeilede, saa vilde det
indeslutte, at den, imod Forudsætningen, var \emph{ensartet} med det, og da blev Forholdet \emph{Mere end en Afspeiling}. Vilde man begrunde
Ikke-Forvandlingen ved, at Viden forholdt sig \emph{absolut passiv}, saa vilde
en saadan Passivitet være uforenelig med Videns Væsen, og Afspeilingen,
der fra det Afspeiledes Side maatte sees som en Indvirken paa det
Afspeilede, vilde blive utænkelig, da en Indvirkning forudsætter
Tilbagevirkning, altsaa Activitet. \emph{Som en Indvirkning fra det Afspeilede maatte} i ethvert Tilfælde \emph{Afspeilingens Tilbliven forklares};
thi der kan \emph{ikke} ved en \emph{immanent Udvikling} ud fra Viden selv med logisk
Nødvendighed naaes til det Afspeilede, da dette er det \emph{absolut Fremmede}, der ikke kan komme frem i \emph{Sammenhæng med Videns Bestemmelser}.
Men \emph{Indvirkningen} viser jo gjennem det \emph{Vexelvirkningsforhold}, den forudsætter, hen til \emph{Ligheden i Væsen}, til Væsensenhed; den \emph{ophæver} altsaa den Forudsætning, hvorpaa Forestillingen om en blot Afspeiling
hviler. \emph{Hvorledes man altsaa end søger at opfatte Forholdet, bliver Betegnelsen Afspeiling uheldig og ubrugelig til at afgive en kategorisk Bestemmelse}.

\medskip

\noindent 3.~Ved Siden af de Beviser for den postulerede absolute
Uensartethed mellem Viden og Villie, der kunne søges i
Kjendsgjerninger, men, som det har viist sig, kun gjennem en uberettiget
Fortolkning af disse, vilde der saa kunne stilles et Beviis, bygget paa en \emph{Analogieslutning} fra Forholdet mellem den \emph{kunstneriske Phantasies} Virken og \emph{Erkjendelsen} af dens Productioner, idet den supponerede
absolute Uensartethed i dette, som analogt med Villiens og Videns
forudsatte, Forhold benyttedes til at støtte Paastanden om hiin
postulerede Uensartethed. Der \opage{151}kunde saa siges: det Skjønne og det Sande ere absolut
uensartede, hvad der sees deraf, at en æsthetisk Analyse aldrig kan give Kunstværket: det
Skjønne kan ikke tænkes saaledes ind i Tanken, at de to Ideer blive identiske. Men hvad der
gjælder om det, gjælder ogsaa om det Gode. Villien, Tanken og Anskuelsen afspeile alle den
hele Idee. --- Med Hensyn til denne Parallel vil det først være at undersøge, hvorvidt
Analogieslutningen vilde være berettiget, hvilket afhænger af, om Phantasiens og Villiens
Differents fra Erkjendelsen har en \emph{væsentlig} Lighed, og dernæst, hvorledes den kunstneriske
Frembringelses Forhold til Erkjendelsen er at opfatte. --- Hvad det \emph{første} Punkt angaaer,
saa viser der sig i Phantasiens og Villiens Forhold til Viden en Forskjel, der maa faae en
afgjørende Betydning. I den kunstnerisk frembringende Phantasie er
Umiddelbarhedsbestemmelsen en anden end i Villien: det er en \emph{Natur}bestemthedens
Umiddelbarhed, i hvilken Aandsvirksomheden træder frem i den genialt skabende Phantasie; det
er en \emph{Aands}bestemthedens Umiddelbarhed, der gjør sig gjældende i Villien \emph{som} Villie ɔ: i den
over den blotte Drift hævede bevidste Villie. Hvad der i den kunstneriske Genialitets
umiddelbare Forbindelse og Sammensmeltning af Natur og Aand gjør sig gjældende, kan ikke
uden Videre overføres paa Villiesvirksomheden, og vilde derfor, selv om det viste hen til en
absolut Uensartethed mellem Phantasie og Viden, ikke bevise Noget med Hensyn til Villiens og
Videns Forhold. Men at det ikke viser hen til \emph{en saadan} Uensartethed, bliver tydeligt ved
Betragtningen af det \emph{andet} ovenfor udhævede Forhold. Med Hensyn til dette betragte vi
særskilt den ved Geniets Naturgave betingede \emph{active Frembringen} af det Skjønne og den
\emph{receptive Tilegnelse} af det genialt Frembragte i en ved Anskuelse ledet Erkjenden.
Hvad den geniale Frembringelse angaaer saa forholder det sig med den paa Phantasiens Gebet
som paa Videns: den er betinget ved et af Individets frie Beslutning uafhængigt, oprindeligt
ɔ: forud for Individets Selvbestemmen givet, Anlæg, i hvilket det Aandelige gaaer
\opage{152}i umiddelbar Enhed med Naturbestemtheden, ligesom under Geniets Udvikling den
kunstneriske Reflexion gaaer i umiddelbar Enhed med, ligesom smelter sammen med, den
geniale Frembringen. At der bliver en \emph{Nuance} i Forholdet mellem Naturbestemtheden og
Aandsbestemtheden i den Genialitet, der forholder sig til Erkjendelse, og i den, der
forholder sig til kunstnerisk Phantasie, faaer i denne Sammenhæng ingen Betydning: her
falder Eftertrykket kun paa \emph{Naturanlæggenes Uafhængighed af den bevidste
Virksomhed og den frie Selvbestemmen}. Naar der tænkes paa denne Uafhængighed, kan det
naturligviis med rette siges, at „den æsthetiske Analyse aldrig kan give Kunstværket“,
ɔ: at man ikke ved Erkjendelsen kan naae til at give sig Betingelserne for en genial
Produceren paa det kunstneriske Gebet; at man ikke ud fra Erkjendelsen kan naae til
Frembringelsen af det originale og som saadant individuelt betingede Kunstværk. Men dette
beviser Intet med Hensyn til det foreliggende Spørgsmaal. --- Hvad Erkjendelsens Forhold
som \emph{receptivt} opfattende det kunstnerisk Frembragte angaaer, saa bliver Tilegnelsen vel,
som overhovedet Tilegnelsen af det concrete Givne, en, om end ikke ved faste Grændser
bestemt, \emph{approximerende} Tilegnen, men under denne Approximation bevares Harmonien mellem
Erkjendelsens Fordringer og Kunstværkets Skjønhedslov, og Erkjendelsen stiller sine
Fordringer i Overensstemmelse med en, ikke blot i Erkjendelsen afspeilet, men i
Erkjendelsen indoptagen Skjønhedsidee. Den analyserende Erkjendelses Optagen af
Kunstværket er vel ledet af en Anskuelse, der i Receptivitetens Form danner Parallelen til
den kunstnerisk frembringende Phantasies Productivitet; men denne Anskuelses Umiddelbarhed
omsættes i Reflexion ligesom Erkjendelsen omsættes i Umiddelbarhed, og det ved
Genialitetens aandelige Naturmagt Frembragte bliver, efter at være frembragt, den
erkjendende Aands Eiendom i en endnu høiere Betydning end den concrete Naturfrembringelse bliver det. --- Ikke heller ad denne Vei kan derfor den absolute Uensartethed hævdes.
\opage{153}4.~Det \emph{afgjørende} Beviis mod den ovenfor udviklede psychologiske Theorie og den
afgjørende Begrundelse af den statuerede absolute Uensartethed søges endelig i det
\emph{Metaphysiske}: i \emph{Grundforholdet i Ideerne}, der bestemmer det psychologiske
Forhold, og i Begrebet om det \emph{Antilogiske}, der anvendes paa Villien, samt i den paa
metaphysiske Forudsætninger hvilende Bestemmelse af \emph{Videns Grændse} som en \emph{absolut} Grændse.

Idet Villie og Viden sættes som primitive, ved en \emph{principiel} Adskillelse sondrede,
Aandsvirksomheder, saa faae de deres særskilte Idealer i forskjellige \emph{Grundideer}:
\emph{Villien i det Godes Idee, Viden i det Sandes}. I det Absolute, siges der, maa det
Godes og det Sandes Idee tænkes i concret Enhed; men denne høieste, concrete Enhed skal
ikke kunne være for vor psychologiske Bevidsthed; for den er Dualismen: den absolute
Sandhed sondrer sig for os i en dobbelt Sandhed, en upersonlig, der ensidig svarer til
Tænkningen, og en personlig, der (ensidig?) svarer til Villien, og disse to Sandheder
blive for os absolut uensartede\footnote{jfr. Nielsen: Om theor. og praktisk Erkj. S. 68.}.
Men Idealets og Principets Bestemmelse maa bestemme det mod Idealet efter sit Væsens
Nødvendighed Stræbende, ved Principet Beherskede; Viden og Villie maa derfor sættes som
absolut uensartede, og consequent kan der aldeles ingen gjensidig Indvirkning indrømmes at
finde Sted. --- Hvad nu denne Deduction angaaer, saa gjælder først i Almindelighed
Følgende. Det maa, med Hensyn til Forholdet mellem den ponerede Dualisme i vor Bevidsthed
og Enheden i det Absolute, urgeres, at det er \emph{en Modsigelse, at statuere denne Enhed
og dog at fastholde Dualismen i Existentsen}. Tankens Udgangspunkt, naar det Absolute skal
bestemmes, er det \emph{psychologiske} Forhold; fra dette gaaer Tankebevægelsen til det
\emph{Ontologiske}, og derfra, begrundende, \emph{tilbage} til \emph{det Psychologiske}. Men er Dualismen
tilstede i Udgangspunktet, vil der ikke kunne tilveiebringes en Enhed i \opage{154}det Ontologiske,
thi Dualismen maa, naar den i det Psychologiske skal have en \emph{blivende} Charakteer,
være \emph{Væsens}bestemmelse i Modsætningerne, og kan da ikke hæves, naar der gaaes
tilbage til Principet for Modsætningerne. Og, omvendt, naar der statueres en \emph{sand
Enhed} i Principet, vil der ikke ud fra dette kunne naaes til en psychologisk Dualisme: en
saadan vil altid forudsætte en i det Ontologiske skjult tilstedeværende om end tilhyllet
Dualisme. En Deduction af den absolute Modsætning i det Psychologiske ud fra Harmonien i
det Absolute vilde være en Deduction af et Forhold ud fra det absolut modsatte,
\textit{qvod absurdum est}. Nielsen har paa dette Punkt ogsaa kun kunnet hjælpe sig med
Postulater og Subreptioner.

I det Udviklede ligger Umuligheden af overhovedet ud fra Ideens Grundforhold at komme til
de Conseqventser, Nielsen søger at naae; men ogsaa bortseet fra det vilde det være
umuligt, paa Basis af hvad „\emph{Grundideernes Logik}“ giver, at kunne bygge en paalidelig
Slutning. Vel er den allerede foreliggende Deel af Grundideernes Logik kun en næsten
forsvindende Deel af det hele Arbeide, der er anlagt efter en saadan Maalestok, at det
uvilkaarligt minder om Aristoteles's Udsagn om det Maleri, der ikke mere vilde være et
Kunstværk: Maleriet af 10,000 Stadiers Størrelse; men den bliver dog af Forfatteren
opfattet som et relativt Hele, af hvilket Principet for det hele Arbeide tydeligt skal
kunne erkjendes, og til hvilket derfor Kritiken har Ret til at holde sig. I denne
foreliggende Deel møder os nu den største Forvirring, de meest iøinefaldende Modsigelser i
Bestemmelsen af Grundideerne. I Grundideernes Logik, 1ste Bind, angives ved Indledningens
Slutning --- uden Begrundelse af disse Bestemmelsers Berettigelse --- \emph{Videns}, \emph{Magtens} og \emph{Sandhedens} Idee som \emph{Grundideer}, og disse Grundideer fungere saa i de Udviklinger, der
udfylde Størstedelen af dette Bind. Men i samme Bind (S. 434) opkastes det Spørgsmaal:
\emph{hvilke og hvormange Grundideer gives der}? og vi faae som Svar en Angivelse af en ny
Trehed af Grundideer: \opage{155}den \emph{kosmologiske}, \emph{geologiske} og \emph{biologiske} Idee, der som
\emph{reale} Grundideer stilles i Modsætning til hine --- ogsaa til Magtens Idee ---
som \emph{ideelle} Grundideer (S. 449). Og \emph{over} Grundideerne staaer saa den \emph{absolute} Idee, medens før \emph{Sandhedsideen} som den \emph{concrete Enhed} af Videns og Magtens Idee
maatte blive den absolute Idee (jfr. S. 432). Sandhedsideen, der i Indledningen sattes som
svarende til Aanden, medens Videns Idee svarede til det Logiske, Magtens Idee til Naturen,
bliver nu dog, skjøndt den opfattes som en enkelt af de ideale Ideer, sat som den, der
sammenslutter Naturens Idee, der udtrykkes ved Realsystemet, og Aandens Idee, der udtrykkes
ved Idealsystemet. Sandhedens Idee bliver, netop som sammensluttende Enhed, Modsætningen
til Aandens Idee og Naturens Idee (S. 451 f.), medens det dog atter siges (S. 455), at „de
to Riger, formedelst Aandens overgribende Magt, i Sandhedens Idee ere eet Rige“. Efterat
der saaledes allerede er bragt et begyndende Anarchie ind i Grundideernes Rige, optræder
nu, i \emph{andet} Bind af Grundideernes Logik, pludselig ligesom et nyt Dynastie af Grundideer:
det \emph{Sandes}, det \emph{Godes} og det \emph{Skjønnes} Idee, der ogsaa (II. 437) faae
Prædicatet: \emph{absolute Ideer}. Her bliver det \emph{Sandes} Idee bevaret fra den
foregaaende Trehed, det Skjønnes og det Godes Ideer ere nye og deres Forhold
til Sandhedens er et andet end Videns og Magtens. Og \emph{Sandheden} har \emph{nu skiftet Function}. I
første Bind svarede den til \emph{Særegenhedens} Moment i Begrebet, som den concrete Enhed
af det Enkelte og det Almene. Nu svarer den til det \emph{Almenes} Moment og faaer til sin
Modsætning den det \emph{Særegnes} Moment repræsenterende \emph{Skjønhedsidee} (II. 376).
De nye Ideer fremtræde ikke ifølge nogen indre Nødvendighed, men bringes kun frem gjennem
Postulater og Subreptioner. Medens vi bevæge os i en Udvikling, der angaaer Subjectivitetens
Selvbestemmen i Almindelighedens, Særegenhedens og Enkelthedens Form som \emph{tænkende},
\emph{reflecterende-anskuende} og \emph{exact forestillende}, forvandler saaledes pludselig
--- \opage{156}medens rigtignok for længe siden \emph{Villiens} Bestemmelse er brugt, men uden at være
deduceret --- den sidste Bestemmelse sig til Villiens, ifølge det \emph{Postulat}, at kun den
villende Subjectivitet kan siges i egentlig Forstand at være det absolute Enkelte.

Til disse forskjellige Generationer af Grundideer, der leilighedsviis forøges med
\emph{Verdensideen} og \emph{Guddommen}, af hvilke hiin er absolut, men denne endnu absolutere,
sættes nu \emph{Villien} saaledes i Forhold, at den snart, som ovenfor angivet, sættes i
Forbindelse med det \emph{Godes} Idee, snart med \emph{Magtens} Idee; men idet
Grundideernes Forhold opløser sig i Forvirring, synker Grundvolden sammen, paa hvilken
gyldige Consequentser skulde hvile.

Ligesaa usikker bliver den Støtte, som søges for Adskillelsen mellem Viden og Villie i
Bestemmelsen af det \emph{Antilogiske}, der, anvendt paa Villien og det Gode, betegner det
% PRINTED AS IS, p.156 [collation 2026-09-21]: stray comma after "den" printed (file silently dropped it); confirmed by a blind second reader.
i Villien, og bestemtere: i enhver Form af den, Antirationelle, det med Tanken
Incommensurable i det praktisk Gode. Bestemmelsen: det \emph{Antilogiske}, der staaer i
Forbindelse med \emph{Magtens Idee}, spiller en stor Rolle i Grundideernes Logik, men er
deels ingensteds deduceret eller retfærdiggjort, deels ikke til at forene med Logikens
øvrige Bestemmelser. Der ligger ingen Retfærdiggjørelse for Begrebet i en aphoristisk
Bemærkning, der (I. 259 f.), længe efter at der er gjort Brug af Begrebet, vil begrunde
dets Nødvendighed ved en Distinction mellem Videns, ɔ: den ontologiske Videns, Indhold og
Gjenstand: en Distinction, der uberettiget overføres fra den menneskelige Videns Sphære og
i Anvendelsen modsiger de andensteds af Nielsen selv givne Bestemmelser om den
guddommelige Videns Indhold og Form. Det Antilogiskes Bestemmelse naaes kun ved en
Subreption, idet strax ved Logikens Begyndelse Væren som almeen Bestemmelse, der tilkommer
Tanken fordi den \emph{er} Tanke, forvandles til den objective Væren, Tankens og Værens Enhed
altsaa til Viden, og idet senere, ved Udviklingen af Reflexionens Former (I. 66),
Andethedens Bestemmelse, i Betydning af den reale Værens Andethed, \opage{157}bringes ind gjennem en
lignende Subreption. Det Begreb, der vilde blive afgjørende for Opfattelsen af Begrebet det
Antilogiske: Begrebet \emph{Materie}, er ingensteds forklaret og retfærdiggjort ud fra
Grundideerne; det omskrives ialfald billedligt som „Magtens Udslag“, og medens Materiens
Realitet paa dogmatisk Maade uden videre forudsættes, naaes dens eiendommelige Prædicater
kun ved Postulater, men benyttes dog, som om de vare deducerede, som Grundlag for
afgjørende Consequentser. --- En Retfærdiggjørelse af Begrebet: det Antilogiske, vilde i
ethvert Tilfælde strande paa den Modsigelse, det indeslutter. Det kan anvendes paa dette
Begreb, hvad Nielsen siger, med Henblik paa Kants Undersøgelser: at det er umuligt at
antage Gjenstandsformer, der ikke oprindelig ere Subjectivitetsformer. Og det er en
umiddelbar Modsigelse --- hvad Nielsen selv fremhæver --- at vide om Noget, at det absolut
falder udenfor Viden. Men det Antilogiske \emph{som saadant} bliver netop kun Betegnelsen for hvad
der skulde staae i dette Forhold til Viden, og det ophører ikke at være det fordi der skal
være en „punktlig Tilsvaren mellem det Antilogiske og det Logiske“ (I. 190 f.), eller fordi der \emph{paastaaes}, at hver Gjenstandsbestemmelse, ved hvilken Viden indskrænkes, er identisk
med en tilsvarende Indholdsbestemmelse, hvorved Viden udvides (I. 261); thi deels er en
saadan Parallelisme mellem det i Form absolut Uensartede --- og kun den absolute
Uensartethed vilde retfærdiggjøre Begrebet: det Antilogiske --- utænkelig, deels vilde det
Antilogiske \emph{som saadant} ligefuldt falde udenfor Viden fordi et \emph{med det parallelt Uensartet} blev et Videns Indhold. Men dette Forhold fører os til Spørgsmaalet om Videns Grændse.

Bestemmelsen af en \emph{Videns Grændse}, ved hvilken der tilsigtes en
metaphysisk-psychologisk Begrundelse af den absolute Uensartethed mellem Villie og Viden,
staaer i umiddelbar Sammenhæng med Begrebet om det Antilogiske, og Videns Grændse bliver
kun et almenere Synspunkt. Videns Grændse bliver opfattet ikke blot som en relativ, men som
en absolut Grændse, som Tanken selv skal aner\opage{158}kjende, saaledes at den, ved at anerkjende den
som saadan, anerkjender det, der ligger udenfor Grændsen, som det absolut Uensartede, der
har Gyldighed i sin absolute Uensartethed. Denne Bestemmelse af Videns Grændse har Nielsen
bragt frem paa flere Steder, og draget vidtgaaende Consequentser deraf. I Grundideernes
Logik er den \emph{menneskelige Videns} Grændse accentueret, deels ved Fremhævelsen af det \emph{Antilogiske}, deels ved Sammenligningen med \emph{den guddommelige Videns Ideal}. Der vises
hen til det Antilogiske som det Utænkelige, til den rene ɔ: abstracte Tankes Afmagt til at
frembringe et Existerende: en Afmagt, der begrundes derved, at paa det psychologiske
Standpunkt Magten er skilt fra Viden, saa at der gives et System af Magtbestemmelser uden
Viden, et System af Vidensbestemmelser uden Magt, og i forskjellige Vendinger bestemmes den
blivende Forskjel mellem den guddommelige, absolut concrete Viden og den menneskelige
(f.\ Ex. II. 432 f., jfr. I. 263). I den philosophiske Propædeutik (fra 1857) kommer
Bestemmelsen af Grændse frem som \emph{Tankens Begrændsning ved Sandsningen og det
Sandselige}. Til de der givne Udviklinger vises der senere tilbage i Stridsskriftet mod
Martensen, men Begrændsningen faaer nu, skjøndt der udtrykkeligt bygges paa hine
Udviklinger, pludselig Betydningen af \emph{Tankens Begrændsning ved Villiesindholdet og
Troesgjenstanden}, og der drages saa Consequentser som om denne Betydning af Grændsen var
deduceret ved hiin.

Men naar det undersøges, hvorvidt Begrebet om en absolut Videns Grændse, ved hvilket den
absolute Uensartethed mellem Viden og Villie kunde hævdes, lader sig retfærdiggjøre, saa
vise der sig afgjørende Vanskeligheder. Tanken for sig \emph{som Tanke} kan ikke erkjende sin egen
absolute Begrændsning, thi idet den erkjendte den, vilde den paa een Gang sætte og ophæve
Grændsen, da denne Erkjenden vilde hvile paa en Erkjenden, d.\ v.\ s. paa en ideel
Assimilation, af det Begrændsendes Natur. Og naar man, forat undgaae denne Modsigelse,
opfatter Forholdet saaledes, \opage{159}at det Subject, der har Tanken til sin Function, ud fra sin
ogsaa andre Functioner omfattende Enhed erkjendte Tankens absolute Begrændsning ved det til
de andre Functioner sig Forholdende, og derved for Bevidstheden Værende, saa vilde en
saadan Erkjenden af Videns Grændse netop \emph{forudsætte} hvad der her søgtes beviist gjennem
den: den absolute Uensartethed mellem Viden og de andre Functioner, idet Functionernes
absolute Uensartethed med Viden er Betingelse for, at det, hvortil Functionen forholder
sig, virkelig falder absolut udenfor Viden; Beviset vilde altsaa blive en Cirkelslutning.

Naar den menneskelige Videns absolute Grændse begrundes ved den Kjendsgjerning, at denne
Viden er afmægtig til at frembringe det Reale, og at der altsaa er sat en
Endelighedssondring mellem Viden og Magt, og naar der saa heraf drages Slutninger med
Hensyn til Uensartetheden, saa anbringes der et kun halvt sandt Synspunkt og drages en
uberettiget Conseqvents. Det er en ubestridelig --- men ogsaa ubestridt --- Kjendsgjerning,
at vor Viden, ligesom vor Villie, ikke har Magt til en umiddelbar Frembringen i det Ydre;
det følger ligefrem deraf, at vor Viden, ligesom alt Værende, er sat, afledet, af det
guddommelige Princip, og at det Ydre ikke først behøver at sættes. Men Synspunktet er
forrykket idet det er overseet, at der i os er en bestandig Omsætten af det Ideelle i det
Reale og omvendt; og Slutningen til den absolute Grændse er falsk, idet det reale Ydres
Uafhængighed af vor Tanke netop kun beviser den \emph{fælleds Afhængighed} af Principet,
men ikke en Uensartethed, der sætter Viden en absolut Grændse, og idet hint Forhold mellem
det Ideale og Reale i os modsiger en saadan Uensartethed. Vi kunne altsaa ikke ad denne Vei
naae til at hævde Grændsen som absolut, og idet intet absolut Uensartethedsforhold skiller
Tanken i os fra det Reale, der gjennem Magtens Idee skulde bringes i Begrebsforbindelse med
den Villie, hvis Afmagt til \emph{umiddelbar} Indvirken paa det Ydre netop er den selvsamme som
Tankens, saa kunne vi altsaa ikke gjennem dette Reale hævde Villiens absolute Uensartethed
med Viden.

\opage{160}Hvad endelig det Absolute og vor Videns Forhold til det angaaer, saa vilde Antagelsen af en
absolut Grændse for vor Viden enten, naar det indrømmedes, at Dualismen i vort Væsen maatte
forudsætte en tilsvarende Dualisme i vort Væsens Princip, føre til, at der statueredes en
Deling i det Absolute; thi Dualismen er her ikke blot et rhetorisk Udtryk for „Spænding“,
„energisk Modsætning“; eller den vilde føre til, at den endelige Tanke ved en uoverstigelig
Kløft maatte sættes som adskilt fra den uendelige, der statueres i det Guddommelige, medens
den dog erkjender dennes Væsen.

Benægte vi altsaa en Videns Grændse? Nei; vi sætte netop en saadan i Kraft af
Grundforholdene; men \emph{som en bestandig bevægelig Grændse}. See vi Tilværelsen som en
organiseret Enhed, behersket af det absolute Princip, saa ligger deri den Consequents, at i
det enkelte, ved Væsenets Universalitet bestemte Led i Totaliteten Heelheden og
Heelhedsprincipet kan være tilstede for en Viden; men det Forhold mellem Totaliteten og de
enkelte Led, ifølge hvilket Totaliteten og Totalitetsleddene, idet de ere tilstede i det
enkelte organiske Led, ere tilstede paa Idealitetens Maade, ikke i den udfoldede Realitets
Form, indeslutter, at hiin Viden maa have Abstractionens Charakteer, men tillige, at
Abstractionen er uendelig integrabel. En absolut Videns Grændse, i Betydning af en fast,
uoverstigelig Grændse, indrømme vi derfor ikke, og forkaste derfor den Slutning, der bygges
paa denne Forudsætning.

\begin{center}---\end{center}

Inden vi gaae over til at paavise \emph{Consequentserne} af den af Nielsen opstillede
psychologiske Theorie, og navnlig til at godtgjøre, at den statuerede absolute Uensartethed
mellem Principerne indeslutter \emph{et Brud paa Bevidsthedens Enhed}, betragte vi først
den psychologiske Theories Reflex i Opfattelsen af det Ethiskes og det Religiøses Forhold
til Erkjendelsen, og anvende Resultatet af Kritiken af hiin paa den Bestemmelse, der gives
af dette Forhold.

\begin{center}---\end{center}

\subsection{Det Ethiske og det Religiøse i Forhold til Erkjendelsen} \opage{161} 
\label{ssec:ethisk-religioes}

Idet Tanke og Villie psychologisk ere bestemte paa den ovenfor angivne Maade, følger heraf
som Consequents det Ethiskes ubetingede Selvstændighed i Forhold til Viden og det
Videnskabelige. Dette udtrykkes naar der siges: den ethiske Existents har sin Grund i det
Praktiske, medens Videnskaben aldrig kommer ud over det Theoretiske. Livet og
Virkelighedsethiken gaaer sin egen Gang, ubekymret om Alverdens Theorier. De ethiske
Principer og de ethiske Virkelighedsbestemmelser kunne vel gaae ind under Videnskabens
theoretiske (rationelle) Belysning, men høre hjemme i en anden Sphære og afspeile sig kun i
hiin. Den theoretiske Ethik kan derfor aldrig blive Grundlag for den praktiske. Dens
Betydning ligger i, at den, ved den theoretiske Erkjendelse af de ethiske Problemer, søger
at tilfredsstille de Fordringer, der stilles af Videns Idee, men indenfor denne Idee har
den theoretiske Realitet sin kategoriske Grændse. I Livet faae Ethikens Problemer en
praktisk Realitet. I den praktiske Virkelighed er der et Element, der aldrig vil lade sig
omsætte i Viden, men idet Ethiken dog giver en Afspeiling af det, faaer den en
eiendommelig Stilling mellem den umiddelbare Livsexistents og den objective Videnskab, og
bliver, selv som rationel Ethik, ikke afhængig af Videnskaben alene. ---
Idet Villien er sat i sin principielle Selvstændighed, fremtræder der, som Former for
Villiesprincipets Udvikling, en Række af bestandig intensivere Villiesconcentrationer, hvis
Stigen betegnes derved, at der mere og mere handles uden Hensyn til Theorien, ja endog i
Strid med den, og disses theoretiske Conseqvents opdager den rationelle Videnskab netop i
Kraft af Erkjendelsen af det Existentielles Ubegribelighed. Af det Rationelle kan paa det
Ethiskes Standpunkt et Antirationelt fremkomme som Resultat og Consequents fordi
Existentsprincipet er et Villiesprincip, der ikke søger Maalestokken for sin Betydning og
Værdi i nogen theoretisk Viden, men ene og alene i sin Concentrations Inderlighed. Villien
concentrerer sig ud over det Rationelle, og Theorien \opage{162}indseer denne nye Concentration som en
existentiel Consequents af de foregaaende, men kan ikke paa existentiel Maade gjennemskue
den, endmindre frembringe den. Paa den antirationelle, paradoxe Ethiks Standpunkt tales der
til Villien i lutter Imperativer, og Mennesket opfatter disse i Troen som lutter
Imperativer. Al Erkjendelse forvandler sig her til ubetingede Postulater; thi den i absolut
Villie concentrerede guddommelige Subjectivitet, til hvilken Mennesket forholder sig, kan
ikke udtrykke sig i nogen almeen, objectiv Form, og Subjectet kan ikke faae Anvendelse for
almene Theorier eller theoretisk Viden. Det Religiøse og Troesforholdet er allerede tilstede
paa det Ethiskes første Trin, i det rationelt Ethiske, men Troen er paa dette ethiske
Stadium en rationel Tro. Den praktisk-rationelle Ethik bliver religiøs idet det Godes Idee
sættes, ikke i relativ, men i absolut Forstand, som Villiesidee, hvoraf følger, at Villien
maa blive baade guddommelig og menneskelig, saa at der bliver et Gudsforhold. Paa den
antirationel-religiøse Ethiks Trin bliver Gud den hellige, almægtige Gud, til hvem den
% PRINTED AS IS, p.162 [collation 2026-09-21]: wrong sort: third letter after "Vil" is a full-height l with no dot (positive signal, checked at 3x against "Villie" elsewhere on the page); ABBYY also reads "V illle"; needs % comment
syndige Villle staaer i Forhold. Guds Væsen er ikke længere conformt med Naturlovene, men
Guds Tanke er paa een Gang Eet med Loven og hævet over Loven. Sit Forhold til Gud udtrykker
Subjectet, der i Syndsbevidstheden har sin Villies høieste Concentration, alene existentielt
i Tro og ubetinget Lydighed. Kun ud fra Syndsbevidstheden naaes der til den Livsforudsætning,
at Gud er den hellige, almægtige Villie, der er hævet over Naturens Love. Her faae
Bestemmelserne: aabenbaret Religion og det Overnaturlige deres Betydning: „den
aabenbarede Religion“, siges der, „er vitterligen en Aandsmagt“. --- Idet det Absolute, seet
fra \emph{Religionens} Synspunkt, bestemmes som \emph{Villiesideal}, bliver \emph{Troen udelukkende} refereret til \emph{Villien} og Menneskets \emph{Troesforhold} bestemt som et \emph{Villiesforhold}.
Villien bliver altsaa sat i et eiendommeligt Forhold til det Aabenbarede, Supranaturale,
Overfornuftige.

\begin{center}---\end{center}

\opage{163}For denne Opfattelse af det Ethiske og det Religiøse ligger der et Correctiv i de i det
Foregaaende givne Udviklinger. Naar Villien, der kun er Villie som menneskelig Villie, som
det bevidste, tænkende Væsens Villie, i Væsenets Enhed er i Harmonie med Erkjendelsen; naar
Indholdets væsentlige Enhed bevares under Omsætningen fra Tankens til Villiens og fra
Villiens til Tankens Form; naar Villien efter sin sande Form som Væsensvillen maa udtrykke
Væsenets Heelhed, have Væsenet efter dets Totalitet til Indhold; og naar det metaphysiske
Princips Enhed indeslutter Umuligheden af at sætte Villiesprincipet og Vidensprincipet som
absolute Modsætninger, at statuere en factisk Dualisme i det Menneskelige, saa følger deraf
for det Første med Hensyn til \emph{Ethiken}, at Erkjendelsen, altsaa ogsaa Erkjendelsens
videnskabelige Udvikling, maa faae en anden Betydning for det Ethiske end i en Theorie, der
vel tillægger den videnskabelige Ethik en væsentlig Betydning for Erkjendelsen; men
statuerer, at den hverken kan blive Grundlag for en ethisk Existents, eller faae
nogensomhelst Gyldighed paa det Ethiskes høieste Trin, --- der altsaa, med eet Ord, tillægger
den en \emph{væsentlig} Betydning --- \emph{kun ikke for det Ethiske}. Erkjendelsen maa faae en
væsentlig og afgjørende Betydning for den ethiske Handlen, som Betingelse for Bevidstheden
om dennes Væsen, Maal, Indhold og Betingelser. --- Der følger dernæst deraf, at \emph{kun} en \emph{rationel} Ethik --- en saadan, om hvilken det vil vise sig, at den i Virkeligheden
ikke kan fremstilles fra Nielsens Standpunkt --- kan faae Sandhed og Gyldighed, og at en
\emph{antirationel} Ethik, som hvilende paa en Villie, der forholder sig ligegyldig til det
Theoretiske, ikke kan være en sand Ethik, ikke kan have Gyldighed, men bliver en
„Absurditetsethik“. Hvad en \emph{religiøs} Ethik angaaer, saa vil det i Fortsættelsen af dette
Skrift blive paaviist, hvorledes den rationelle Ethik alene i sandeste Betydning kan hævde
sig denne Bestemmelse, medens det, der pleier at kaldes en religiøs Ethik ɔ: en ved det
positivt Religiøse bestemt Ethik, enten bliver indholdsløs, ifølge et abstract supranaturalt
Udgangspunkt, eller udfyldt \opage{164}med et vilkaarligt statutarisk Indhold, ifølge Forholdet til den
transscendente guddommelige Enkeltpersonlighed; i intet Tilfælde altsaa en sand
Virkelighedsethik, der kun kan fremkomme hvor den ethiske Handlen sees som udtrykkende Menneskets \emph{Væsen} efter dets concrete ideal-reale Bestemthed.

Med Hensyn til det \emph{Religiøse} indeslutte Resultaterne af det tidligere Udviklede
Consequentser, der svare til dem, vi have draget med Hensyn til Forholdet mellem Ethikens og
Videnskabens Princip og til Ethikens sande Form.

Sammenhængen mellem det \emph{Ethiske} og det \emph{Religiøse} tilveiebringes hos Nielsen derved, at
\emph{Villiens Ideal} bestemmes som \emph{Idealet af den absolute Villen af det Gode}, og som saadant som ligt med den \emph{guddommelige} Villie, den \emph{Godes Villie}, og det Religiøses Forhold til
Viden bliver det samme som det Ethiskes. Idet Gud, seet fra det Ethiskes Synspunkt,
aabenbarer sig for Bevidstheden \emph{exclusivt} under Bestemmelsen: den Gode, bliver der en
\emph{absolut Uensartethed} mellem \emph{Villiesloven}, den guddommelige Villies Fordring, og \emph{Erkjendelsens Væsensbestemmelser}, og idet den guddommelige ubetingede Villie bliver den
guddommelige Subjectivitets absolute Udtryk, og Villien er Udtrykket for den guddommelige
absolute \emph{Frihed}, hvilken det Godes Idee fordrer, saa bliver den ved Villie bestemte Gud
\emph{Underets} Gud ɔ: den Gud, hvis Villie har ubetinget Magt over alle sine Betingelser;
han bliver den \emph{skabende} Gud, og netop ved den ved Skabelsen og Underet hævdede
Ubetingethed skal Gud være den \emph{hellige} Gud, ethisk Villiesideal. „Gud vilde ophøre at
være den Hellige, den i Sandhed Almægtige, dersom han i sit Forhold til Verden ikke var den
absolut frie Gud, den Gud, der skaber af Intet“. \emph{Aabenbaringens} Gud skal være en
\emph{Vilkaarlighedens} Gud og som saadan ikke kunne være Gjenstand for \emph{theoretisk} Erkjendelse, men kun for \emph{Tro}, og bestemtere ikke for en rationel Tro, men for en Tro, der
forholder sig til \emph{Paradoxet}. Vilkaarlighedens Gud træder i Forhold til Mennesket gjennem
ubetingede \opage{165}Bud, og Menneskets Forhold til disse er ubetinget Lydighed, et
\emph{Villiesforhold}, ligesom Troen er en \emph{Villiesact}. Det \emph{Gudsbegreb}, Philosophien
har udviklet, skal, fordi det er theoretisk bestemt, ikke kunne være Gjenstand for Tro; det
skal være tankeløst at sige, at der kan troes paa Philosophiens Gud.

I de foregaaende Udviklinger ere Momenterne givne til en Kritik af disse Bestemmelser. Det
er \emph{Gudsbegrebet} og \emph{Troesbegrebet}, der ud fra dem blive at corrigere.

Nødvendigheden af at lade det Ethiske ligesom udmunde i det Religiøse, eller rettere, at
lade det vise tilbage til dette som til sin Forudsætning, er allerede fremhævet i de
foreløbige Begrebsudviklinger af det Religiøses Væsen; i dette Punkt kunne vi derfor kun
samstemme med Nielsen. Ogsaa naar Forholdet udtrykkes i den Sætning: Villiens Ideal er et
Villiesideal, have vi ingen Indvending at gjøre mod dette Udtryk. Men et Usandt kommer frem,
idet Villien indenfor det Guddommeliges Sphære bestemmes ud fra det samme falske Begreb af
Villie, der statueredes indenfor det Ethisk-Menneskeliges Sphære. Den \emph{isolerede} Villie, der
forholder sig til Erkjendelsen som udelukkende denne --- og det bliver kun en Talemaade, at
den indeslutter det absolut Uensartede: Tanken som Moment --- bliver overført paa det
Guddommelige, der er den ethiske Bevidstheds Gud, Troens Gud, idet Forestillingen om denne
Villie paa phantastisk Maade udvides til Absoluthed. \emph{I Virkeligheden er saaledes ikke
Villiens absolute Væsen hævdet, men den menneskelige illusoriske Vilkaarligheds og selviske
Drifts Væsen er paa modsigende Maade uendeliggjort}. Det Spørgsmaal er omgaaet: paa hvilken
Maade Villiens Begreb maa metamorphoseres idet Villien faaer Absoluthedens Bestemmelse og
overføres paa det Guddommelige. Med Hensyn til dette er det indlysende, at paa den ene Side
alle de Bestemmelser maae fjernes, der hæfte ved den menneskelige Villie som relativ
ubevidst, som forholdende sig til et Ydre, og som føiet ind i Aarsagernes og Virkningernes
Kjede, og derved, bevidst eller ubevidst, i hvert \opage{166}Moment bestemt ved et Andet; og at, paa
den anden Side, Villien i det Absolute kun kan være Væsensvillen, i Harmonie med de i Viden
aabenbare Væsensbestemmelser af det Guddommelige, og bestemmende sig ved de Betingelser, der
ere de i det for sig selv absolut gjennemsigtige Væsen begrundede. Den Nødvendighed, der
ligger i den guddommelige Villies \textit{libera necessitas}, er kun Frihedens Lov,
Selvbestemmelsens Bestemmethed ved det i Erkjendelse og Villen for sig aabenbare ideal-reale
Væsen. En Transscenderen af de ved Væsenet bestemte Betingelser er en phantastisk
Sublimation, ved hvilken \emph{Væsensyttringen} Villie transscenderer \emph{Væsenet} ɔ: sig selv.
Lader man Forestillings- og Phantasieopfattelsen af den guddommelige Villie falde, saa faaer
man i det Absolute, som svarende til Villiesvirksomheden i det Psychologiske, den
Subjectivitetens Energie, ved hvilken det Absolute evigt sætter Realiteten som bestemt ved
det Ideelle, --- og i denne Bestemmelse af den guddommelige Villie forenes det Synspunkt, fra
hvilket den sees som Udtryk for den Energie, ved hvilken et Realt sættes og bestemmes, og
det, fra hvilket den sees som Norm for de endelige Subjecters Villie, som Udtryk for det
Godes Idee. I denne Opfattelse viser den guddommelige Villie sig i Harmonie med den til
Erkjendelsesvirksomheden i det Psychologiske svarende Energie, ifølge hvilken den absolute
Subjectivitet evigt forholder sig til sig selv i sin Idealitet gjennem det Reale. Den sande
Opfattelse af Villiesbestemmelsen i dens Overførelse paa Gud vilde derfor see den
guddommelige Villie, det Godes Idee, i uopløselig Enhed med det Sandes Idee, og ikke blot
postulere denne Enhed i Guds Væsen i og for sig, men statuere den i Gud saaledes som han er
for vor ethiske og religiøse Bevidsthed. Modsætningen, hvis Nødvendighed for os hos Nielsen
bestandig kun postuleres, vilde føre til, at \emph{begge} de Former, i hvilke Gud skal være
for den menneskelige Bevidsthed, for Tro og Viden, vilde blive \emph{usande}, netop fordi de
absolut adskilte hvad der kun er sandt som værende i Enhed.

\opage{167}Den Opfattelse af Villien, som Nielsen gjør gjældende, fører ham til, uden noget Beviis for
Sammenknytningens Berettigelse, at sætte Guds \emph{Hellighed} som betinget ved Guds \emph{undergjørende Magt}, til altsaa at lægge et afgjørende Eftertryk paa Forestillingen om \emph{Underet}. Videnskaben skal vel ikke kunne føre et Beviis for dette Under, som Troen sætter
ved et Postulat, men den skal ikke heller kunne bevise Underets Umulighed. Af praktisk
Interesse for Underet kommer Nielsen imidlertid til at føre en Art videnskabeligt Beviis for
dets Mulighed ud fra Begrebet om en Subjectivitet med ubetinget Frihed og ubetinget Valg af
Midler og Betingelser, og han kommer til at glemme, at Videnskaben ikke kan give Plads for
Underets Mulighed --- naar da ikke Underet tages i uegentlig Betydning --- uden at fornægte
sig selv. Den maa statuere Underets Umulighed, i Kraft af det Forhold mellem Frihed og
Nødvendighed, som Nielsen (Gr. L. II, 442) selv har hævdet, og fordi Underet, naar det ikke
tages i den uegentlige, aandrigt eller tankeløst forflygtigende Betydning, men i den
religiøse Forestillings Opfattelse, bliver en guddommelig Vilkaarlighedshandling, der
negerer \emph{alle} reale Betingelser, og altsaa maa benægtes ud fra ethvert Gudsbegreb, der sætter
Gud i et Væsensforhold til det Reale, altsaa ogsaa ud fra det Gudsbegreb, Nielsen selv
opstiller. De Modsigelser, til hvilke Forestillingen om Underet fører den Bevidsthed, for
hvilken Viden, under sin absolute Uensartethed med Troen, har absolut Gyldighed, ville i det
Efterfølgende blive paaviste.

Naar Villien i det Guddommelige opfattes i sand Enhed med Viden, det Godes Idee i Enhed med
det Sandes Idee, bliver den menneskelige Villies Ideal et saadant, der ikke bringer
Mennesket i indre Splid med sig selv ved at lade Villien blive ude af Stand til at udtrykke
Væsenet efter dets Heelhed. Hvad der hos Nielsen tilstræbes ved Bestemmelsen af Guds Villie
som Villiens Ideal: at give den ethiske Villiesfordring ubetinget og evig Gyldighed ved at
see den som udgaaende fra den det Gode sættende guddommelige Villie, vil ligefuldt blive
naaet selv om den af \opage{168}ham fastholdte Opfattelse af den guddommelige Villie forlades og Villien
sees som i \emph{Harmonie} med Viden udtrykkende det guddommelige Væsen, og som erkjendt i denne
Harmonie af den menneskelige Tanke.

Den Opfattelse, Nielsen gjør gjældende af Gudsbegrebet, fører, hvad \emph{Troesbegrebet} angaaer,
som allerede paaviist, til, at Troen bliver \emph{Tro paa det Paradoxe}, og til at den sættes
i et særligt Forhold til Villien. Det, at Gudsforestillingen concentrerer sig i Bestemmelsen
af den absolute Villie, der sættes som incongruent med Erkjendelsens almene
Fornuftbestemmelser, gjør Gud til Gjenstand for Tro i den positive Religions specifiske
Betydning; og at denne Tro bringes i særligt Forhold til Villien, er forstaaeligt, deels paa
Grund af Troesgjenstandens Bestemthed (Gud = absolut Villie), deels paa Grund af Troens
psychologiske Betingelser (fra et psychologisk Synspunkt kan den, som stridende mod
Erkjendelsen, kun sees som fremgaaende af en Erkjendelsen modsat Aandsvirksomhed). Troen som
Tro paa det Paradoxe bliver til ved et Brud med Erkjendelsen, ved en af Villie fremgaaende
Beslutning. Idet Religionen forholder sig til Tro, faaer den Sætning Gyldighed, at
„Religionen fra Først til Sidst er praktisk“. Naar Religionen af de gamle Dogmatikere
bestemtes som en \textit{modus} \emph{\textit{cognoscendi colendique}} \textit{Deum}, saa indrømmes af Nielsen vel en
\textit{cognitio}, men ikke som en theoretisk, men som en praktisk Erkjendelse, idet
Mennesket i Religionen underordner alle Sjæleevner under Villien og lader dem repræsenteres
ved denne. Erkjendelsen som praktisk Erkjendelse bliver derfor Et med Bevidstheden om
Nødvendigheden af at statuere Noget ifølge dets Uundværlighed for Subjectets Trang. Denne
praktiske Erkjendelse begriber ikke, men fordrer; den drager Consequentser, ikke af et
Rationelt, men af et Antirationelt.

Mod disse Bestemmelser maa først det gjøres gjældende, der paa et tidligere Punkt blev
udviklet angaaende Videns absolute Grændse. Det viste sig, at Viden ikke kunde anerkjende en
Grændse, der som en urokkelig Grændse \opage{169}skilte den fra et med den absolut Uensartet. Idet
nemlig Viden, forat anerkjende et fra sig forskjelligt Princip, ikke kunde tænkes som
isoleret, men maatte føres tilbage til et mere omfattende Subject, der ligesom nøder Viden
til Anerkjendelsen, saa vil dette omfattende Subject, der idet det adskiller Viden fra det
fra Viden forskjellige Princip, \textit{eo ipso} sammenslutter begge i sin Enhed, ophæve
Muligheden af en absolut Uensartethed, af en blivende Modsigelse mellem dem. Naar der som
Grund til at antage det mod Viden absolut Stridende vises hen til Subjectets Trang, saa maa,
som allerede ovenfor paaviist, denne Trang retfærdiggjøres, den maa godtgjøres som berettiget.
Trangen som Aandsbestemthed er ikke en umiddelbar, instinctmæssig Følelse, er ikke umiddelbar
sikker paa Tilstanden og paa det Trangen Tilfredsstillende, men hviler paa et
Reflexionsforhold, og den Dom, den indeslutter med Hensyn til Tilstandens Bestemmelse og til
det Trangen Tilfredsstillende, staaer i et andet Forhold til Bevidstheden end den Dom, den
sandselige Trang indeslutter og har ikke den samme Art af Sikkerhed. Hvad der kaldes
„existentiel Erfaring“ paa Aandens Gebet, forholder sig derfor til en Dom, der maa
retfærdiggjøres. En saadan Retfærdiggjørelse kan ikke finde Sted ad anden Vei end ved at
paavise, at i Trangen en oprindelig gyldig Væsensbestemmelse finder sit Udtryk. Men
Væsensbestemmelsen maa, som Væsensbestemmelse, ifølge Væsensenheden, kunne finde sit Udtryk i
Viden, om end ikke i en ensidig Form af Viden. Hvad der modsiger Videns Væsen, Viden som
Viden, kan, som det Viden Fornægtende, ikke heller have praktisk Gyldighed.

Hvad dernæst det angaaer, at der vises hen til Villien som Organet for Troen paa Paradoxet,
saa har det, som ovenfor udviklet, sine forstaaelige Motiver; men tilfredsstiller ligesaalidt
den religiøse Bevidsthed som den videnskabelige Erkjendelse. Det tilfredsstiller ikke den
religiøse Bevidsthed; thi denne seer Troens Tilbliven ikke som sat ved en Villiesact, men som
fremgaaende af en overnaturlig Virkning, der paa een Gang omdanner alle Aandens Evner, \opage{170}altsaa
Erkjendelsen saavelsom Villien, idet den, medens den bestemmer Villien til at modtage
Aabenbaringens Indhold --- en Bestemmen, i hvilken Villien er et Stof, \emph{med} hvilket Noget
gjøres --- luttrer Erkjendelsen til at kunne opfatte det. Men ligesaalidt tilfredsstiller hiin
Opfattelse den videnskabelige Bevidsthed. Fra dennes Standpunkt maa det fastholdes, at
Villien som Væsensyttring af den menneskelige \emph{Natur} ligesaalidt væsentligt kan fastholde sig
til det „\emph{Overnaturlige}“ som hvilkensomhelst anden af denne Naturs Væsensyttringer, og et
saadant Forhold er ogsaa kun postuleret. Hiin Sammenknytning af Villie og Tro har praktisk
viist sig i sine Consequentser i alle Kjætterforfølgelser og i Inqvisitionens Rasen mod de
Ikke-Troende, der, ifølge hiin Sammenknytning, dømtes som dem, der ikke \emph{vilde} troe.

Idet Villien er gjort til Organet for det Religiøse, der udelukkende drages ind under
Bestemmelsen af det Praktiske, men Villien her opfattes i sin Modsætning til Erkjendelsen og
det Almeenfornuftige, saa bliver det endelig at fremhæve, at den Villie, hvorom der her
tales, kun er det Villiens Naturgrundlag, der træder frem som Drift, Attraa, Ønske,
overhovedet som det Alogiske i Menneskenaturen.

Fjernes Bestemmelsen af Villiens og Videns absolute Modsætning, saa beholder Troen og
Troesforholdet ligefuldt sin Betydning baade paa det Ethiskes og paa det Religiøses Gebet,
men Troen bliver \emph{rationel} Tro, Yttring af den hele ɔ: fornuftbestemte, i Erkjendelsen
sig aabenbare, Menneskenaturs Drift. Religionens Absolute, i hvilket ogsaa det Ethiske
grunder, bliver det, der ligeligt tilfredsstiller Erkjendelsens og Villiens Fordring, det, i
hvilket \emph{hele} Menneskenaturen kommer til Hvile. Ogsaa saaledes opfattet faaer Troen
bestandig et Forhold til Villie; men dette Forhold bliver et væsentligt andet end i hiin
Opfattelse. Villiesacten i Troen bliver en ved Erkjendelsens Fordring bestemt Villiesact, en
Act, i hvilken den dobbelte Form af Væsensenergien gaaer i Enhed, ligesom i den Act, i
hvilken \opage{171}Erkjendelsen villes, og i den, i hvilken det Erkjendte fastholdes som blivende,
sikkert Resultat.

Hvorledes det fra et rationelt Standpunkt, --- fra den rationelle Ethiks Standpunkt og ud fra
en Tro, der staaer i Harmonie med Erkjendelsen --- er muligt, at hævde de ethiske og
ethisk-religiøse Grundbegreber, der paastodes kun at kunne hævdes ud fra en med Viden absolut
uensartet Tro, i deres væsentlige Betydning, vil blive paaviist i det Skrift, der danner
Supplementet til dette.

\subsection{De psychologiske Bestemmelsers Forudsætninger i det
  Metaphysiske}
\label{ssec:metaphysisk}

Modsætningen i det \emph{Psychologiske} mellem Erkjendelse og Villie viste, efter hvad der allerede
tidligere blev udviklet, tilbage til en Modsætning paa det \emph{Metaphysiskes} Gebet, ved hvilken
den skulde begrundes. Det var en Modsætning i Grundideerne, der af Nielsen bestemtes i
forskjellige Former. Deels opfattedes nemlig Modsætningen mellem Vidensprincipet og
Villiesprincipet som constitueret ved Modsætningen mellem det \emph{Sandes} Idee og det
\emph{Godes} Idee: en Modsætning, der vel evigt skulde være hævet i det absolute
Subjectivitetsprincips concrete Forsigværen, men som i den endelige Bevidsthed nødvendigen
skulde forme sig til en absolut Modsætning. Deels førtes den psychologiske Modsætning tilbage
til Modsætningen mellem \emph{Videns} Idee og \emph{Magtens} Idee.

I denne Tilbageføren og Begrunden er der hos Nielsen deels en Uklarhed, deels en
Utilstrækkelighed. --- Uklarheden viser sig i det Usikkre i Forholdet mellem de to
Sondringer, idet Modsætningen mellem det Godes og det Sandes Idee allerede falder indenfor
% PRINTED AS IS, p.171 [collation 2026-09-21]: turned n: page prints "Modsætniugen" (u positively open at top, bowl joined at foot, against closed n in "-gen" of the same word); already silently corrected, restore with comment; confirmed by a blind second reader.
Videns Kreds, og derfor ikke kan opfattes under samme Synspunkt som Modsætniugen mellem Magtens og Videns Idee. Ogsaa faae det Sandes og det Godes Idee Skjønhedens Idee ved Siden
af sig, medens Videns og Magtens Idee ikke have noget tilsvarende Sideordnet, da Sandhedens
Idee, naar \opage{172}den ikke fungerer som Bestemmelse indenfor Videns Kreds, bestemmes som deres
høiere Enhed og altsaa indtager en anden Plads i Forhold til dem. Utilstrækkeligheden viser
sig deels i det Ufuldstændige i de psychologiske Grundfunctioners Tilbageførelse til et
Metaphysisk, deels i Uoverensstemmelsen mellem det, der føres tilbage, og det, hvortil det
føres tilbage, deels i Mangelen paa en virkelig Begrundelse af den absolute Modsætning.

Naar Functionerne tilbageføres til Videns og Magtens Idee, bliver der en Ufuldstændighed,
idet den kunstneriske Phantasie, i sin relative Autonomie, ikke faaer noget Tilsvarende i de
metaphysiske Principer. Denne Ufuldstændighed undgaaes vel naar det Sandes og det Godes Idee
sættes i Stedet for Videns og Magtens, men der opstaaer saa den ovenfor berørte
Vanskelighed med Hensyn til Forholdet mellem de to Bestemmelser af Modsætningen.

Sees der nu bort fra det Kunstneriske og fra Phantasien, og tages der udelukkende Hensyn til
Villie og Erkjendelse, saa viser der sig en væsentlig Incongruents, der umiddelbart træder
frem ved det ene af Leddene. Villien er opfattet som det i eminent Forstand \emph{Subjective}
i Selvet; i den tilsvarende Idee: Magtens Idee, er derimod, i dens Modsætning til Videns
Idee, det \emph{Objectives} Moment fremtrædende.

Men hvorledes end den til den psychologiske Modsætning svarende Bestemmelse i det Absolute
formuleredes, saa vilde det, som det allerede ovenfor blev godtgjort, blive umuligt, ud fra
en sand Enhed i det Absolute at naae til den absolute Modsætning i det Psychologiske, eller
omvendt. Ved det Forhold, Nielsen statuerer, vilde der kun blive en \emph{Navne-Enhed} mellem
den menneskelige og guddommelige Viden; og hiin vilde, idet Grændsen sattes som absolut
Grændse, ved en uudfyldelig Kløft være adskilt fra denne. Et egentligt Forsøg paa at
godtgjøre Nødvendigheden af den absolute Modsætning i det Psychologiske, har Nielsen, som
% PRINTED AS IS, p.172 [collation 2026-09-21]: the page sets "ikkeheller" solid, as one word; two readers saw it.
tidligere fremhævet, ikkeheller gjort. Der gives overalt kun Postulater, og det Bestemmende
er for Nielsen ingensteds \opage{173}en Begrebsnødvendighed, men kun den individuelle Trang til at faae
det Religiøse i hans Opfattelse reddet ved en principiel Sondring fra Videns Gebet.

Den i det Psychologiske statuerede Dualisme faaer nu ogsaa, som sit Forbillede, en
\emph{Dualisme i det Metaphysiske}, og Begrebet: Dualisme er atter her at tage i den
alleregentligste Betydning som en uforsonet og uforsonlig Modsætning, der ikke kan
sammenfattes i en Modsætningsenhed. Dualismen træder \emph{først} frem som en \emph{Dualisme mellem
Viden og Magt}, ifølge det Antilogiske i Magten, der er Videns, ogsaa den ontologiske Videns,
blivende Skranke. Det nytter ikke at vise hen til, at der til hver antilogisk
Magtbestemmelse \emph{svarer} en logisk Vidensbestemmelse\footnote{jfr. Gr. L. I, 190: det
Antilogiske „afgiver en Reflex“, der beriger Reflexionen.}; thi dette bliver kun en
Parallelisme, der ovenikjøbet bliver uforklarlig naar det Antilogiskes og det Logiskes
Bestemmelser, der hver følge sine Bevægelseslove, hver for sig sees i deres Sammenhæng
indenfor deres Sphære; --- det bliver ikke en virkelig Enhed; thi det Antilogiske \emph{som} Antilogisk kan ikke assimileres af Viden, og idet det ikke blot er en ligesom overfladisk
Formbestemmelse, men gjennemtrængende Væsensbestemmelse, bliver dets Ikke-Optagethed
afgjørende for Dualismens Vedvaren, der ikke heller hæves derved, at der rhetorisk tales om
en „vidende Magt“ og en „mægtig Viden“.

\emph{For nu at fjerne denne Dualisme gives der Magtbegrebet en Fordobbling}. Magten
bliver bestemt som den, der, ifølge sin oprindelige indre Enhed med Viden, paa een Gang er
ideel og reel; der tales saavel om logiske som om antilogiske Magtbestemmelser. Men denne
dobbelte Bestemthed i Magten, der kun er sat \emph{ved et Postulat}, lægger kun den \emph{dualistiske}
Modsætning over i Magten uden at paavise dens Ophævelse i Enheden (jfr. Gr. L. I, 189 f.,
191, 392. II, 142, jfr. II, 278). I Magten, hvis \opage{174}Stilling til Viden og til Subjectiviteten
overhovedet bliver uklar ved denne Fordobbling, er saaledes Dualismen blivende i den \emph{uforklarede} Dobbelthed. Dualismen i Magten skjules kun ved Assertioner, som saa danne
Forudsætninger for alle de efterfølgende Udviklinger, der hvile paa den Subreption, at den
principielle Enhed er hævdet. Og \emph{idet Dualismen i Magten ikke hæves, udøver den en Tilbagevirkning paa Viden}, hvor den viser sig deri, at Viden skal omfatte baade hvad
den assimilerer og hvad den ikke assimilerer; denne Adskillelse gjælder nemlig ikke blot om
den menneskelige Viden, men om den ontologiske, der forholder sig til det uopløselige
Antilogiske og i det har en Skranke.

I den sidste Form for Dualismen reflecterer sig den dualistiske Modsætning mellem det
Antirationelle og det Rationelle (jfr. Gr. L. II, 278 Anm.), den Modsætning, hvis Bevaring
for Nielsen er det oprindelige Diemed, og idet nu Dualismen i Videns Idee reflecterer sig i
den menneskelige Viden, saa viser den sig her deels directe gjennem Dobbeltheden i Viden som
tilegnende og som blot afspeilende: en modsigende Dobbelthed, der maa medføre en Spaltning i
Vidensindholdet\footnote{Var „Afspeilingen“ en \emph{væsentlig} Bestemmelse i Viden, en
væsensnødvendig Videns Form, maatte der ved en \emph{immanent} Udvikling naaes til det
„Afspeilede“; men dette er en Umulighed paa Grund af dettes \emph{absolute Uensartethed} med det
begrebne Vidensindhold. Viden maa som Viden udfolde sine Muligheder systematisk, men
Mulighederne vilde, naar de skulde omfatte det absolut Uensartede, umiddelbart modsige den
Enhed, den systematiske Udfoldning forudsætter.}; deels indirecte derigjennem, at Viden, idet
den ifølge sit Princip skal udfolde sig i Selvstændighed, dog skal tillæmpe sine Bestemmelser
efter det Antirationelle. Fordringen af en saadan Tillæmpning ligger i det tidligere anførte
Udsagn, at Videnskaben skal lade Miraklets Mulighed staae som det, der ikke kan modbevises:
et Udsagn, der trodser de Modsigelser, til hvilke Indrømmelsen af denne Mulighed
\opage{175}fører\footnote{I Forestillingen om Miraklet træder netop den Modsigelsen indesluttende
Dualisme frem; thi naar Miraklet tages i sin \emph{egentlige} Betydning, betegner det det abstract
% PRINTED AS IS, p.175 [collation 2026-09-21]: dittography "er er" in footnote *) (mid-line, not a line-break repeat); confirmed by a blind second reader.
Ideelles vilkaarlige Virken med Negation af \emph{ethvert} Realt, af \emph{alle} reale Betingelser: den samme Virken, der er er sat i Skabelsesforestillingen. Muligheden af Miraklet i denne Betydning
søger Nielsen at hævde ved sin efter den religiøse Forestilling tillæmpede Bestemmelse af den \emph{absolute Subjectivitet}, der frit skal kunne transscendere alle Betingelser. Men derved faaer
man enten umiddelbart den angivne Dualisme mellem et Ideelt og det Reale, eller, naar den
ideelle Subjectivitet dog fra Begyndelsen af tillægges Realitetsbestemmelser, bliver dets
Transscenderen af disse, ved hvilken de negeres, idet nye sættes, kun tænkelig som et Brud,
hvorved der altsaa sættes en Dualisme i selve det absolute Princip og i den reale
Tilværelse. Og Dualismen i Principet bliver saa ikke blot den dualistiske Sondring mellem det
Absolute, bestemt ved de oprindelige og ved de senere satte Realitetsbestemmelser; men mellem
det Absolutets Idealitetsside og dets Realitetsside, hvilken sidste sættes som et Intet, idet
den transscenderes. --- Vilde man begrunde Umuligheden af at modbevise Miraklets Mulighed
derved, at Viden, anerkjendende sin Grændse, maa anerkjende Personlighedens Fordringer, og at
Personligheden gjennem Fordringen af en hellig Gud fordrer en almægtig, skabende Gud, saa er
denne Vending allerede belyst i det Foregaaende. Personlighedens Krav er i sin Sandhed
Væsenets Krav, der ikke kan modsige den Væsenet udtrykkende Viden.}, og staaer i Strid med
andre bestemte Udsagn\footnote{F.\ Ex. Gr. L. II, 190. „En Almagt, der ubetinget skaber ---
--- af Intet, er ikke blot for os ufattelig, men i Ideen ubegribelig“.}. En saadan
Tillæmpning bliver synlig i Opfattelsen af Magtbegrebet, idet \emph{Magtens Begreb} hos
Nielsen vakler over mod \emph{Almagtens Forestilling}, og den bliver synlig idet en
\emph{Forestillingssondring} fastholdes mellem den absolute, i Villie concentrerede mægtige
Subjectivitet og den reale Gjenstand; thi gjennem denne Forestillingssondring trænger sig
Forestillingen ind om en almægtig ɔ: absolut vilkaarlig, \emph{skabende} Villie, der i ubetinget
Frihed hæver sig over alle Betingelser: den Forestilling, paa hvilken i sidste Instants
Paastanden om det for Videnskaben Umulige i at modbevise Miraklets Mulighed hviler.

\opage{176}Det her Udviklede fører os til at tilføie nogle Bemærkninger med Hensyn til den
videnskabelige Charakteer af det Værk af Nielsen, i hvilket de anførte metaphysiske
Bestemmelser ere givne.

Naar \emph{Dualismen} er uovervunden i det Metaphysiske, maa dette influere paa
\emph{Dialektikens Form}. Den dialektiske Udviklings indre immanente Nødvendighed maa blive
umulig; dogmatiske Paastande og Subreptioner maae træde i Stedet. I Grundideernes Logik er
Dogmatismen og Manglen paa sand Dialektik iøinefaldende.

Det er allerede tidligere berørt, at \emph{Materiens Realitet} hos Nielsen overalt dogmatisk
forudsættes og at dens Forhold til Magtens Idealitetsside kun bestemmes ved en Metaphor, idet
den betegnes som „Magtens Udslag“. Paa samme Maade bliver \emph{Videns Mulighed} (Løsningen
af Subjectivismens Problem) og Begrebet om den \emph{absolute Viden}, ligesom ogsaa
\emph{Magtens Begreb} i Virkeligheden dogmatisk forudsat, idet Deductionen er illusorisk.

Dialektiken mangler i Værket som Heelhed; den dialektiske Bevægelse finder kun Sted indenfor
de enkelte Partier paa Basis af det ved Subreptioner eller Postulater Naaede.

Vi fremhæve, til Belysning af den videnskabelige Fremgangsmaade, følgende Punkter:

1.~I Indledningen til Gr. L. gaaes der ud fra et historisk Udgangspunkt: Hegels Bestemmelse
af Væren som Philosophiens Begyndelse. Det viser saa --- hvad Hegel ikke er uvidende om ---
at Væren forudsætter Tænken. Tænken skal altsaa være det Første, men Tænken kan kun tænkes
som værende ɔ: som havende Tænkens Væren. Her er altsaa Enhed af Tænken og Væren i denne
Abstractionens Betydning. Men denne Tænkens og Værens Enhed gjøres saa ved en Subreption til
Tænkens Enhed med den \emph{objective} Væren: til \emph{Viden}, og idet den \emph{objective, reale Værens} Bestemmelse er bragt ind, faae vi Viden som \emph{Modsætningsenhed af Tænken og real
Væren}, og paa Basis af dette ved et Ordspil Naaede gives saa de efterfølgende Udviklinger.

\opage{177}2.~Hvor Begrebsmomenternes Forhold bestemmes, paa hvilket hele Læren om Domme og Slutninger
hviler, bliver først \emph{Almindelighedens} og \emph{Enkelthedens} Momenter saaledes bestemte, at
deres Bestemmelser fuldstændig congruere. Den uendelige Negativitet, der er Subjectivitetens
egen Form, bliver uden sondrende Bestemmelser Udtrykket baade for Almindeligheden og
Enkeltheden. Men nu optages fra \emph{Forestillingen om det Enkelte} Individualitetens Bestemmelse, og denne individuelle Væren, der i sin \emph{Umiddelbarhed} er det Almenes \emph{Modsætning},
skydes ind under Enkelthedens Begreb, saaledes at nu dette, der oprindelig var identisk med
det Almenes, paa een Gang er i Enheds- og i Modsætningsforhold til det. Naar saaledes
Modsætningen er bragt ind i Enheden, saa er det naturligt, at det \emph{Særegnes} Moment som
det mæglende maa udtrykke den \emph{concrete Enhed} af Enheden og Modsætningen, og naar først de tre
Momenter ere givne i denne Bestemthed, kan der combineres i det Uendelige, og man kan, ved
tilsyneladende immanent Dialektik, naae til alle Slutninger og Slutningssystemer o.\ s.\ v.
Men Grundvolden, paa hvilken hele denne dialektiske Bygning hviler, er upaalidelig.

Hvad der betegner Gr. L. som Heelhed, er, at \emph{Forestillingssondringen er overført ved et
Postulat paa Begrebet og bevaret i det Absolute}. Idet saa det Absolutes Begreb indeslutter Fordringen af \emph{Enhed}, bliver det postuleret, at Forestillingsmodsætningens Dualisme gaaer ind
under Enheden: der postuleres, at Viden og Magt gaae sammen til Enhed idet Magten bliver
Enhed af Idealitet og Realitet. Men den postulerede Enhed bliver uholdbar, og det ensidige
Udgangspunkt: Forestillingsmodsætningen i dens dogmatiske Fastslaaethed, bliver det, der
forhindrer den sande Enhed.

\begin{center}---\end{center}

Paa Grundlaget af de foregaaende Udviklinger kan der nu gives en Paavisning af de
Consequentser, til hvilke \opage{178}den Nielsenske Theorie fører i psychologisk, ethisk og religiøs
Henseende.

\section{Theoriens Consequentser}
\label{sec:consequentser}

\subsection{Med Hensyn til det Psychologiske}
\label{ssec:conseq-psych}

Af den abstracte Modsætning mellem det Praktiske og det Theoretiske, mellem Tro og Viden, og
af den derved betingede uforenelige Dobbelthed i Opfattelsen af Tilværelsen og dens Princip,
følger som Consequents \emph{en blivende, uforsonlig Spaltning i den menneskelige
Bevidsthed}, Bevidsthedens indre uløselige Modsigelse; og af hiin absolute Modsætning følger
endvidere, \emph{at de som absolut uensartede opfattede Grundfunctioner, der kun have Sandhed
i deres Enhed og Harmonie, forvanskes i deres Væsen}.

Den første af disse Consequentser, der i sig indeslutter den anden, søges fra Nielsens
Standpunkt afbødet, og, i Modsætning til den, den Sætning hævdet, at \emph{Tro og Viden, netop paa Grund af den absolute Sphæreforskjel, den absolute Uensartethed, skulle kunne forenes,
uden Modsigelse, i samme Bevidsthed}. Men Begrundelsen af denne Sætning bliver aldeles
utilstrækkelig og Consequentsen lader sig ikke afværge.

Vi berøre kun i Forbigaaende det Argument for de absolut uensartede Opfattelsers
modsigelsesfrie Væren i samme Bevidsthed, der af Referenter af Nielsens Theorie er blevet
søgt i Phantasieopfattelsens Forhold til Virkelighedserkjendelsen. Bortseet nemlig fra de
øvrige allerede af Andre paaviste Mangler ved denne med ringe Eftertanke opstillede Parallel,
bliver der en Ulighed mellem dette Forhold og Forholdet mellem Tro og Viden netop i det
afgjørende Punkt, og denne Ulighed gjør Analogien ubrugelig. I Troens og Videns Forhold til
deres Gjenstand er Forholdet paa \emph{begge} Sider et \emph{Virkelighedsforhold}. Gjenstanden bliver ikke
i Troesopfattelsen, som i Phantasieanskuelsen, sat som Skin, i Illusionens Medium, men den
sættes, af Troen som af Viden, som den sande, ved et absolut \opage{179}gyldigt Princip bestemte,
totale, fuldstændige Realitet, og Troens og Videns \emph{modsatte} Bestemmelse af det som den sande \emph{Virkelighed} Satte er det, paa hvilken vor Paastand om Conflictens Uundgaaelighed
hviler.

Af større Interesse end denne forfeilede Analogie er det Forsøg, \emph{Nielsen} selv\footnote{\emph{Nielsen}: Om den gode Villie. S. 28 ff.} gjør paa at hævde Foreneligheden i samme Bevidsthed ifølge den
absolute Uensartethed ved at bestemme Charakteren af Videns Grændse, og ved at drage
Consequentser af denne Bestemmelse.

Besindelsen over vor Videns Betingelser skal føre til den Indsigt, at al vor Viden er
begrændset (d.\ v.\ s. „\emph{ifølge Sandsningens og Tænkningens Uensartethed}“\footnote{I
Stridsskriftet mod Martensen S. 133 f. udledes Bestemmelsen af Videns \emph{absolute} Grændse af Forskjellen mellem menneskelig og guddommelig Viden, mellem vor Viden og vor (\emph{vor}!) Videns
Ideal.}), og at Videns Grændse ikke blot er en relativ, men en \emph{absolut}, og Forholdet
mellem det Absolute og det Relative i vor Videns Begrændsning skal opklare Forholdet mellem
Tro og Viden. „Havde den menneskelige Viden kun \emph{relative} Grændser, saa maatte der --- siges
der --- naturligviis altid kunne gives «\textit{plus ultra}», da maatte alle vore Erkjendelser uden
Forskjel omsider blive at begrunde i Videnskaben, og hvad Videnskaben ikke kunde begrunde,
maatte være i og for sig usandt. Har den menneskelige Viden derimod, foruden de relative
Grændser, en \emph{absolut} Grændse, saa er her altsaa Noget, vi i denne Tilværelse aldrig
faae at vide: da er her et \emph{Mysterium}, en Tilværelsens Gaade, hvis Opløsning i Videnskaben maa
blive usand. Til dette Mysterium maae henføres det personlige Livs absolut ideale
Forudsætninger“. --- Begrebet om Videns \emph{absolute} Grændse skal saa benyttes til at
fjerne Modsigelsen fra den Bevidsthed, der i sig skal rumme baade Viden og Tro. „Det er vor
Videns absolute Grændse, det er Tilværelsens Gaade, hvorpaa saavel den Videndes som den
Troendes Opmærksomhed her maa fæste sig. Har al menneskelig Vi\opage{180}den en absolut Grændse, er her
virkelig et Aandens Mysterium, en absolut Gaade, da er her ogsaa \emph{Noget, der maa troes,
fordi det ikke kan vides}. Naar nu det i Gaaden Skjulte, det, der ene og alene er bestemt til
at være Gjenstand for Tro, ikke desmindre drages ind under Viden og Videnskab, hvad saa? Saa
overskrider Viden sin absolute Grændse --- --- og giver sig frisk væk til at udgrunde,
begrunde og begribe det, der efter sin Natur hverken kan eller skal videnskabelig begribes
--- og den kommer saa nødvendigviis til \emph{modsigende} Begreber. Men hvad der kun kan udtrykkes i
modsigende Begreber, maa være \emph{usandt}: Troens Mysterier lade sig kun udtrykke i modsigende
Begreber, altsaa ere de usande. At Noget kan være \emph{sandt} i Troen og \emph{usandt} i Videnskaben, følger nu ligefrem deraf, at Videnskaben, ved at ville begribe hvad der ligger
udenfor dens Grændse, \emph{selv} bliver usand“. Paa tilsvarende Maade kan saa Noget, der
virkelig er \emph{sandt} i Videnskaben, være \emph{usandt} i Troen. „Det er nemlig ikke blot Viden,
men Troen, der har en absolut Grændse: at Viden har Troen til Grændse, forudsætter nødvendigt,
at Troen har Viden til Grændse. Men at Troen har Viden til Grændse, vil sige, at et
Troesbegreb umuligt kan udvikles til et Vidensbegreb. Idet Troen støtter sig til Miraklet,
forudsætter den en virkelig Natur med virkelige Love. Men hvorledes forholder Naturen sig til
Miraklet? Ja, det er jo netop dette, der er Tilværelsens Gaade. Hvis nu Troen, forglemmende
sin absolute Forudsætning, nemlig Gaaden, overskrider sin Grændse og begynder at ville
omkalfatre det virkelige Naturbegreb til Bedste for Miraklet: da bliver Naturlærens
Naturbegreb ikke blot en Usandhed i Troen; men Troen, der har indladt sig paa Noget, den ikke
magter, bliver \emph{selv} usand“. --- „At Tro og Viden lade sig forene i een Bevidsthed, ikke \emph{endskjøndt}, men \emph{fordi} de ere absolut uensartede Principer, vil da med andre Ord sige, \emph{at de kun lade sig forene i et absolut Mysterium}. Det absolute Mysterium er paa een Gang
\emph{exclusiv Grændse og negativ Midte}. Som exclusiv \opage{181}\emph{Grændse} sætter Mysteriet
Skilsmisse imellem Tro og Viden som imellem to Regioner. Hver Region har sit Ubetingede, sit
Midtpunkt: Troens Midtpunkt er Friheden, den almægtige, hellige Villie; Videns Midtpunkt er
Nødvendigheden, den rationelle, i uforanderlige Love sig aabenbarende Nødvendighed. Naar nu
den Troende seer mod Troens, den Vidende mod Videns Midtpunkt, ere Extremerne spændte til det
Yderste: hvad der er sandt i Troen: at nemlig det Ubetingede er Miraklet, er usandt i Viden;
hvad der er sandt i Viden: at nemlig det Ubetingede maa søges i Fornuftnødvendigheden, er
usandt i Troen. Som \emph{negativ Midte} forener Mysteriet derimod de adskilte Extremer, og
viser sig som \emph{Regionernes Gjennemgangsmedium}: gjennem Mysteriet gaaer Veien fra Tro til
Viden, fra Viden til Tro. Idet den Troende og den Vidende begge (!) paa een Gang rette deres
Blik mod Mysteriet, udlignes Striden og Extremerne enes. \emph{Seet fra Videns Standpunkt er
nu Miraklet muligt}. Thi vel er det umuligt, at der nogensinde skulde kunne skee et
vilkaarligt Brud paa de evige, med Fornuftnødvendigheden identiske Naturlove, men med Lovenes
Uforanderlighed kan jo, \emph{naar de i Mysteriet skjulte Mellembestemmelser underforstaaes},
den Aabenbarelse af Villiens, Viisdommens og Kjærlighedens Magt, \emph{hvori den Troende seer Miraklet}, meget vel tænkes at bestaae. Seet fra \emph{Troens} Standpunkt kan
Fornuftnødvendigheden med de uforanderlige Naturlove forenes med Forudsætningen af Miraklet.
Thi vel er det umuligt, at Nødvendigheden skulde kunne hæmme Friheden, og Naturens Love
indskrænke Guds ubetingede Villie, men \emph{naar de i Mysteriet skjulte Mellembestemmelser underforstaaes}, lader det sig jo \emph{tænke}, at Naturlovene, langt fra at brydes, tværtimod i
allerhøieste Forstand opfyldes, hvergang en guddommelig Villiesact fuldbyrdes“\footnote{Begrebet
om „det \emph{absolute Mysterium}“ bliver i Skriftet: Om den gode Villie, af Nielsen paa flere
Punkter anvendt forat hæve Conflicten. Saaledes (S. 89 f.) med Hensyn til \emph{Incarnationen}.
„Enheden [af guddommelig og menneskelig Natur] er og bliver et absolut Mysterium. Det Første,
% PRINTED AS IS, p.182 [collation 2026-09-21]: footnote: page prints "dogmatkske" (k for i, positive ascender; text layer agrees); confirmed by a blind second reader.
den religiøse Tænkning har at gjøre, er at stille sig af med alle dogmatkske Incarnationsproblemer. For den troende Bevidsthed er Grundmiraklets Modsigelse ganske vist
løst, men ikke derved, at den selv tænker Løsningen, men derved, at den \emph{troer}, at Løsningen er \emph{tænkt} i Mysteriet.“}.

\opage{182}Af det ved det absolute Mysterium satte Forhold mellem Tro og Viden udleder Nielsen nu den
Consequents, at der mellem disse, der ikke kunne indskrænke hinanden, ikke kan tænkes en
endelig vilkaarlig Overenskomst, og at ligesaalidt deres Adskillelse kan bringe Adsplittelse
ind i den menneskelige Bevidsthed. Utænkeligheden af en ensidig, vilkaarlig Overenskomst, af
et Concordat, skal fremgaae deraf, at den kun kunde tænkes mellem \emph{relativt ensartede} Magter, der paa udvortes og endelig Maade \emph{indskrænkede} hinanden, hvilket de kun kunde gjøre
under den Forudsætning, at deres Begrændsning var en udenfra kommende. „Men Tro og Viden kunne --- \emph{paa Grund af Mysteriet} --- umulig indskrænke hinanden. Vel er der i Troen en Grændse for
Viden, forsaavidt det, der kan troes, ikke kan vides; men denne Grændse er Eet med Videns \emph{egen, iboende} Grændse. Med sin absolute Grændse, sin indre Begrændsning, er Viden ligeoverfor
Troen \emph{aldeles uindskrænket}, saa uindskrænket som om der ingen Tro var; ligesaa er Troen med
sin \emph{iboende} Vidensbegrændsning uindskrænket ligeoverfor Viden, saa \emph{uindskrænket}, som om der
slet ingen Viden var. Saalidt som Viden nogensinde er istand til at angribe Troen, saalidt vil
Troen nogensinde være istand til at angribe Viden: saalidt som Viden er istand til at rette
sig efter Troen, ligesaalidt er Troen istand til at rette sig efter Viden. Kunne de nu, hvad
det absolute Mysterium absolut forhindrer, hverken stride mod hinanden eller understøtte
hinanden, saa er en endelig Overenskomst imellem dem utænkelig.“ (S. 35; jfr. S. 139.)

At Adskillelsen mellem Tro og Viden ikke kan \emph{halvere} den menneskelige Bevidsthed, skal ligge
deri, „at Be\opage{183}vidstheden halveres ved \emph{relativt ensartede}, men ikke ved \emph{absolut
uensartede} Factorer. Naar en Bevidsthed halveres i Tanken, faaer den en halv Tanke istedenfor
en heel: det, der mangler, er Tanke; halveres den i Villie, faaer den en halv Villie istedenfor
en heel: det, der mangler er Villie. Derimod kan man ikke i strængere Forstand (!) sige, at en
Bevidsthed halveres af Tanken og Villien saaledes, at Tankens Halvdeel bliver Villie, Villiens
Halvdeel Tanke; thi Tanke kan kun i høist uegentlig Mening forvandles til Villie og Villie til
Tanke. Staaer det da fast, at Troen paa ingen Maade kan forvandles til Viden eller Viden til
Tro, saa staaer det ogsaa fast, at Troen ligesaa umuligt kan halveres ved Viden som Viden ved
Tro. Kan én Bevidsthed ikke halveres af Venskab og Mathematik, af Kjærlighed og Chemie, saa
kan den endnu mindre halveres af Ethik og Astronomie. Der gives mathematiske Venner, men der
gives ingen venskabelig Mathematik; der gives forelskede Chemikere, men ingen chemisk
Forelskelse, ingen forelsket Chemie. Der kan gives ethiske og uethiske Astronomer, men
Astronomien er hverken ethisk eller uethisk. Videnskaben er med andre Ord hverken ethisk eller
uethisk, hverken christelig eller uchristelig, hverken troende eller vantro; Kjærlighed er
praktisk, men Videnskab er theoretisk“. (S. 35 f.)

Vi have gjengivet disse Udviklinger, ogsaa de mod Slutningen anbragte «\textit{lazzi}», saa
fuldstændig, fordi kun en saadan Fuldstændighed kan give et tydeligt Billede af den Maade, paa
hvilken Theorien maa forsvares. Vi belyse nu Udviklingen i det Enkelte.

Hvad der ved disse Udviklinger først maa være paafaldende, er den Maade, paa hvilken det, der
i Propædeutiken er fremsat med Hensyn til Uensartetheden mellem \emph{Tænkningen} og \emph{Sandsningen}, det Tænkelige og det Sandselige, ved et Spring overføres paa
Forholdet mellem \emph{Tanke} (Viden) og \emph{Tro}, skjøndt dette Forhold i ethvert Tilfælde
først gjennem mange Mellemled kunde bringes sammen med hint, og skjøndt der mellem begge strax
\opage{184}viser sig den tidligere berørte væsentlige Uoverensstemmelse, der bestemmes ved Differentsen
mellem Naturexistentsens Enkelthed og Aandsbestemthedens Enkelthed.

Idet saa \emph{Videns Grændse} er sat som en \emph{Grændse mod Troen} og Tro og Viden i
Grændsen gjensidig skulle begrændse hinanden, --- hvilket ikke kan undgaaes idet de sættes i \emph{samme} Bevidsthed --- saa synes dette at maatte ophæve Bestemmelsen af deres \emph{absolute
Uensartethed}, idet kun det, der falder indenfor en \emph{væsentlig Enhed}, kan begrændse hinanden.
Men naar den absolute Uensartethed hæves, saa er Værnet mod en \emph{forstyrrende Conflict}, hvilket
netop skulde findes i Uensartethedens Absoluthed, tilintetgjort.

Denne Consequents søges nu undgaaet ved den nye Bestemmelse, der drages frem: Bestemmelsen af
\emph{det absolute Mysterium}, der for Tro og Viden skal være det, der forhindrer en udvortes
Begrændsning. Men denne Bestemmelse, der \emph{forudsætter} den \emph{absolute} Uensartethed, hjælper
paa ingen Maade ud over Vanskelighederne. Først bliver der nemlig en væsentlig Uklarhed i
Anvendelsen af Mysteriets Begreb. Idet Viden anerkjender sin Grændse, anerkjender den i det,
der er Troens Sphære, et Mysterium for Viden. Idet Troen har en Grændse mod Viden, bliver det,
der er Vidensgjenstand, et Mysterium for Troen. \emph{Troen er altsaa Videns Mysterium, Viden
Troens}, og Mysteriet, der, saaledes opfattet, ikke har noget Mysteriøst, forhindrer paa denne
Maade ikke Indskrænkningen, men udtrykker den netop. Men nu bliver pludselig, forat
Sammenstødet i Grændsen kan forhindres, Mysteriet til et \emph{Tredie}, et \emph{Gjennemgangsmedium},
gjennem hvilket Veien gaaer fra Tro til Viden, fra Viden til Tro. Dog bevirker ikke heller
denne Metamorphose hvad der tilsigtedes. Idet det søges paaviist, hvorledes Mysteriet, medens
det ubetinget adskiller, er det Forsonende, løber der ikke blot en Tvetydighed ind, men
\emph{Ensartetheden} afløser den \emph{absolute Uensartethed}. Tvetydigheden løber ind idet det
ubetinget adskillende Mysteriums \opage{185}forsonende Betydning skal paavises med Hensyn til \emph{Miraklet}.
Viden skal, ved at underforstaae \emph{de i Mysteriet skjulte Mellembestemmelser}, kunne tænke sig
den Magtaabenbarelse, „i hvilken \emph{den Troende} seer Miraklet“, som bestaaende med Lovenes
Uforanderlighed, men dette kan den kun, naar den, medens Miraklets \emph{Navn} beholdes, lader
\emph{Miraklet som Mirakel forsvinde i et supponeret Rationelt}. Og Ensartetheden kommer saa
ind med det Samme, idet \emph{Troen} ved en lignende Forudsætning af Mellembestemmelser skal forene
Forudsætningen af Miraklet med Forudsætningen af de uforanderlige Naturlove, d.\ v.\ s. idet
den \emph{ligeledes skal rationalisere Miraklet}. Men Ensartetheden kommer ogsaa ind i den
Form, at Viden og Tro, i Forhold til Mysteriet og de deri skjulte Mellembestemmelser, ligesom
\emph{vexle Rolle}, saa at \emph{Viden} bliver \emph{troende, Troen tænkende}. Viden, der bygger Muligheder
paa skjulte Mellembestemmelser --- Mellembestemmelser \emph{mellem Noget og Intet}; thi
saaledes maa Mellembestemmelserne med Hensyn til Miraklet, den partielle \emph{Skabelse}, opfattes
--- og paa Mellembestemmelser i det \emph{absolute Mysterium}, om hvilket \emph{Intet} kan \emph{vides, Intet supponeres}, forholder sig i Virkeligheden \emph{troende} til det Supponerede, og \emph{Troen} skal jo „kunne \emph{tænke}“ Muligheden af Naturlovenes Fuldbyrdelse trods Miraklet. Saaledes bringer altsaa
Berøringen i Grændsen, trods Mysteriets Negativitet, en Ensartethed ind og kun i denne sees
det Forsonende.

At en saadan uvilkaarlig Ophæven af den \emph{absolute} Uensartethed maa være uundgaaelig,
ligger med Nødvendighed i det menneskelige \emph{Selvs}, i den menneskelige \emph{Bevidstheds} Natur. Væsenets og Bevidsthedens Enhed gjør sig uimodstaaelig gjældende ligeoverfor den
absolute Sondring af Yttringsformerne. Muligheden eller Nødvendigheden af den paastaaede
absolute Uensartethed i det Psychologiske er der, som ovenfor paaviist, ikke heller nogensinde
af Nielsen gjort et Forsøg paa at godtgjøre.

Og ligesaalidt som Mysteriet kunde hævde den \emph{abso\opage{186}lute Uensartethed}, ligesaalidt kan den, selv
naar den ovenfor paaviste Metamorphose foretages, forhindre den \emph{ydre Begrændsning}. Da
det tidligere har viist sig, at en Selvbegrændsning af Viden, i den Betydning, som Nielsen
opfatter den, er umulig, saa kunde Mysteriet kun tænkes at hjælpe til at Viden i sin
Begrændsning kunde være ubegrændset, naar Viden, idet den statuerede Mysteriet, statuerede det
\emph{som en Videns Mulighed}, og det ikke som en med vor Viden absolut incommensurabel Videns
Mulighed, men som en med vor Viden væsenslig Videns. Ellers vilde nemlig Modsigelsen
blive i den menneskelige Bevidsthed saalænge den er som saadan. Men naar Troens
Eiendommelighed fastholdes, saa bliver hiin Videns Mulighed hævet, og Begrændsningen er atter,
trods Mysteriet, sat som en virkelig \emph{ydre} Indskrænkning, der ikke omsættes i en indenfra
given Grændse.

En saadan Omsætten af den ydre Begrændsning i en indenfra kommende kunde ogsaa ialfald blive
forestillelig med Hensyn til \emph{Viden}, der, idet den, ifølge sin Natur, afspeilede Modsætningen,
ved sig selv skulde blive sig sin Grændse bevidst og blive sig selv bevidst i sin nødvendige
Selvbegrændsning. Troen derimod kan ikke forestilles, ved sig selv at blive sig sin Grændse
bevidst og at forvandle den ydre Grændse til en indre; den Troende faaer kun Grændsen at vide
gjennem Viden, og kan netop derfor ikke foretage hiin Omsætning, idet Omsætningens
Forudsætning: at Viden som Viden vedvarende sætter den ydre Grændse, ikke kan slippes uden at
al Bevidsthed om en Grændse forsvinder, hiin Omsætning altsaa bliver selvmodsigende. Forsvandt
al Bevidsthed om Grændsen, saa vilde Troen ikke være forsonet med Viden, men vilde være
uvidende om Viden, hvilket den ikke kan, naar de, efter Forudsætningen, ere forenede i samme
Bevidsthed, og hvilket den, efter Forudsætningen, ikke maa være, da Tro og Viden netop skulle
forsones med gjensidig Bevidsthed om Modsætningen.

Begrændsningens Ophævelse som ydre Begrændsning vilde kun kunne blive tænkelig under
Forudsætning af, at \opage{187}Viden og Tro i gjensidig Gjennemsigtighed hver for sig udfoldede et
uendeligt Indhold. Men et saadant Forhold er udelukket ved den statuerede absolute
Uensartethed.

Naar saaledes det absolute Mysterium viser sig uskikket til det Hverv, der betroedes det, saa
staaer den Consequents urokket, at \emph{Bevidstheden adsplittes ved Troens og Videns
Modsætning}. Tro og Viden gjøre hver for sig Krav paa \emph{hele} Aanden; de forholde sig
begge til \emph{Virkeligheden} som Heelhed og til \emph{denne Virkeligheds Princip}, altsaa til
\emph{samme Object}, og idet de skulle gjøre det som \emph{absolut uensartede}, saa ligger
heri Modsigelsens og Adsplittelsens \emph{Uundgaaelighed}. Nielsens Udviklinger, der skulle
godtgjøre det Modsatte, ere i høi Grad forvirrede. Der er ikke Tale om Tankens Halvering ved
Villie eller om Villiens ved Tanke, men om Halveringen af den Bevidsthed, som skal rumme det til det samme Object sig forholdende Uensartede og derved kommer i Splid i sig selv. De „livlige“
Analogier, Nielsen tilføier til Oplysning, ere nærmest kun et lidet heldigt Forsøg i det
Burleske, og bevise, med Hensyn til den omhandlede Sag, om muligt, mindre end
Intet\footnote{Nielsens Exempler med de „mathematiske Venner“ og den „venskabelige
Mathematik“ oplyse kun det lidet Interessante, at man ikke uden videre, fordi et Adjectiv kan
prædiceres om et Substantiv, derfor, omvendt, kan prædicere et til Substantivet svarende
Adjectiv om et til Adjectivet svarende Substantiv, at man f.\ Ex. ikke, fordi man kan tale om
en flau Vittighed, kan tale om en vittig Flauhed. Det, at der kan gives ethiske og uethiske
Astronomer, men ikke en ethisk eller uethisk Astronomie, oplyser kun, at forskjellige
Aandsvirksomheder kunne have samme Subject uden at kunne gjøres til Prædicat for hinanden
(saaledes som f.\ Ex., indenfor Videns Sphære, Astronomie og Lægevidenskab, idet der vel gives
lægekyndige Astronomer, men ingen lægevidenskabelig Astronomie); men det beviser Intet med
Hensyn til en absolut Uensartethed, eller til, at disse Aandsyttringer kunde prædiceres om det
samme Subject uden at de væsentlig gik tilbage til det samme Princip, eller uden at det samme
udelte Væsen gjorde sig gjældende gjennem dem.}.

Til et ligesaa negativt Resultat som Benyttelsen af Forestillingen om Mysteriet fører Nielsens
Forsøg paa, ved \opage{188}Anvendelsen af en \emph{Troeskritik} og en \emph{Videnskritik} og ved
Adskillelsen mellem dem\footnote{\emph{Nielsen}: Om den gode Villie, S. 58 ff.}, at forhindre
Conflicten mellem Sphærerne. Han sætter det som nødvendigt, naar Tro og Viden forenes i een
Bevidsthed, ved en Kritik at forebygge, at de blandes, skjøndt dog Uensartethedens Absoluthed
maatte synes at forhindre Blandingen uden Anvendelsen af særegne Midler. Kritiken skal saa
være en dobbelt: en Kritik, der sørger for, at de to Erkjendelseskredse under alle
Vexelbestemmelser, Grændsesætninger og Reflexbelysninger holdes ude fra hinanden, og en
Kritik, der paaseer, at der i enhver af de adskilte Kredse gjøres Forskjel mellem Væsentligt
og Uvæsentligt, Nødvendigt og Tilfældigt, Gyldigt og Ugyldigt. Under den dobbeltsidige (eller
egentlig tredobbelte!) Anvendelse af Kritiken fremkommer da en Troeskritik, der staaer i samme
Forhold til Videnskritiken, som Troen til Viden: er Troen absolut uensartet med Viden, saa er
Troeskritiken absolut uensartet med Videnskritiken. Troeskritiken, der udelukkende seer paa
den religiøse Ægthed, lader sig slet ikke anfægte af Videnskritikens Betænkeligheder.
Historiske Urigtigheder i Evangelieberetningerne for Exempel kunne allerede af den Grund ikke
anfægte den, at \emph{den aldrig faaer dem at vide}, „thi i det Øieblik, den faaer dem at vide, er den
ikke Troeskritik, men Videnskritik“\footnote{Man anvende dette paa Troens Forhold til Viden,
og paa den Kritik, der sondrer Gebeterne!}. Tvivl om det Overnaturlige i Historien eller
Forsøg paa at udskille det som en mistænkelig Tilsætning, berøre ikke Troeskritiken; den vilde
derved ligeledes give Plads for en Videnskritik. Saaledes skal hver Sphære blive uforstyrret i
sig selv.

Det er let indlysende, at man, ved disse Bestemmelser, ikke undgaaer Conflicten. Naar først
den Kritik, der skal forhindre Sammenblandingen, skal være baade Videns og Troens Sag, saa
bliver ikke blot Viden vidende om Troen, men Troen vidende om Viden og Conflicten er, ifølge
det \opage{189}Foregaaende, uundgaaelig. Og selv om man saa, hvor Troeskritiken og Videnskritiken vende
sig hver mod sin „Erkjendelseskreds“, ved et Ordspil fører Tvivlen fra hins Kreds over paa
dennes, saa slipper man dog ikke for det dobbelte Bevidsthedsforhold til det indenfor samme
Gebet Modsigende, og derved er Conflicten given. Den subtile Distinction, der gjøres med
% PRINTED AS IS, p.189 [collation 2026-09-21]: page prints "Nyttte" (triple t; text layer agrees); confirmed by a blind second reader.
Hensyn til Tvivlens Subject, er saaledes til ingen Nyttte.

I Følelsen af, at Muligheden af et Sammenstød ikke er afværget, søges der saa endelig forat
forhindre Bruddet et ligesom \emph{neutralt} Gebet, paa hvilket en Forhandling og en Udligning
af Modsigelsen kan finde Sted. Dette neutrale Gebet, der bliver det psychologiske Medium, i
hvilket en Bevidsthed, der skal kjende og erkjende baade Troens og Videns Princip, de to
absolut uensartede Principer kunne skilles og mødes og paa en Maade træde i Underhandling med
hinanden, bestemmes som \emph{Meningen}. Viden skal det ikke kunne være, fordi Viden hører til
Videnskaben og ikke er neutralt Gebet; Tro ikke heller, eftersom Troeserkjendelsen udelukker
Videnskaben og altsaa heller ikke frembyder noget Neutralt. „Den vældige Modsætning mellem
Oldtidens og Nutidens Anskuelse af Himmel og Jord --- der fra Culturbevidsthedens Standpunkt
maa synes at undergrave det gamle og nye Testamentes, og derved Aabenbaringens, guddommelige
Autoritet, vil ikke forsvinde, men lade sig udligne (!) ved fortsatte Forhandlinger paa
Meningens neutrale Grund. --- Den væsentlige Forskjel mellem Viden og Mening blev allerede
Grækerne klar; men den principielle Forskjel mellem Tro og Mening synes endnu ikke at være
bleven Theologerne klar. Ligesom Forskjellen mellem Viden og Menen kun kan klares ved en
Videnskritik, saaledes kan Forskjellen mellem Tro og Mening kun klares i en Troeskritik. Et
Exempel: Hvad der i første Mosebog fortælles om Paradiset, om Livets Træ, Kundskabens Træ
o.\ s.\ v. er ganske vist fortalt i den Mening, at det altsammen hører under det udvortes
Virkelige, at Livets Træ og Kundskabens Træ ligesaavist have havt deres Rødder i Havens
sandselige \opage{190}materielle Jordbund, som enten Viintræet eller Figentræet. I den Mening have vore
luthersk-orthodoxe Fædre endnu for 100 Aar siden ganske vist opfattet Fortællingen om Paradiis.
Men i denne Mening kunne vi (som Troende?) paa Culturbevidsthedens nuværende Trin umulig opfatte
den. Dersom der nu ingen væsentlig Forskjel var mellem Tro og Mening, saa vilde en Forsoning af
Religion og Cultur ganske vist være umulig, men er her en Forskjel, en principiel Forskjel, saa
er Forsoningen mulig. Men at der virkelig er en principiel Forskjel mellem Mening og Tro, skal
kjendes paa, at vi kunne bevare Fædrenes Tro og forkaste deres Mening. Til Troens ufravigelige
% PRINTED AS IS, p.190 [collation 2026-09-21]: dittography across line break: "Forudsæt-" / "sætninger"; confirmed by a blind second reader.
Forudsætsætninger hører, at Forestillingen om Paradiset med alle deri indeholdte Charakteertræk
tillægges væsentlig og sand Realitet, væsentlig og sand Virkelighed. Men hvilken Virkelighed?
En ydre eller en indre? Ja, derom er der forskjellige Meninger. Men saa er det jo kun en
Mythe; eller rettere: hvorvidt det er Mythe eller Aabenbaring, vi have for os i dette Sagn,
derom kan der dog vel være forskjellig Mening? Nei, Sagnet om Paradiis kan, saa sandt det i
rene, klare Træk for Troens Blik afslører et af de dybeste Villiesmysterier, umulig være
Mythe. Men hvorledes skal man da forestille sig Forholdet mellem det Indre og det Ydre, mellem
det, der skeer legemligt, og det, som foregaaer i Aandens Verden? Ja, derom er der
forskjellige Meninger.“\footnote{\emph{Nielsen}: „Den gode Villie“, S. 62--64.}

Ved Hjælp af „Meningen“ skal altsaa en Forsoning tilveiebringes; men Meningen, der skal udøve
en saa væsentlig Function, faaer, skjøndt den kaldes „kategorisk Betegnelse“, en noget
proteusagtig Charakteer. Mening bruges snart som Synonym med „Betydning“, snart som Betegnelse
for den blot subjective Opfattelse, der skal gjøre Plads for det i Viden og Tro Almeengyldige,
og snart som Betegnelse for det, der bliver en Grændse, over hvilken der ikke kan naaes ud
fordi der endes i et \textit{non liquet}. At \opage{191}anvende dette flydende Begreb i en fast Function
synes derfor ugjørligt. Men idet man saa, ved at bruge det, statuerer, at Forsoningen mellem
Religion og Cultur vilde være umulig, naar ikke Conflicten udjævnedes derved, at Noget
udsondredes fra Troen, saa synes det at være en Modsigelse, at Modsætningen mellem det
\emph{absolut Uensartede} skulde blive mindre skarp ved Meningens Udsondring. Hvad der sættes
som et Forsoningsmiddel, maatte tværtimod blive et Middel til at skærpe Striden; thi naar
Meningen, det „Neutrale“, udsondres, og altsaa det, fra hvilket den udsondres, kommer frem i
sin principielle Eiendommelighed, maa Modsætningen netop blive skarpere end før. En Tilsætning
af den neutrale Mening kunde kun formilde Modsætningen.
En Forsoning mellem Tro og Viden ved Hjælp af Meningen --- der kun i den Forstand bliver
neutral, at den ikke skaanes af nogen af de kæmpende Parter --- vilde kun kunne tænkes paa to
Betingelser. Enten maatte man, gjennem Udskilningen af hvad der tildeles Meningen --- d.\ v.\ s.
Meningen om Troen, thi kun den bliver der Tale om, --- \emph{rationalisere Troen}, „bringe Troen
til Fornuft“, ved en vilkaarlig Omfortolkning, der ikke kan normeres; men dette strider jo mod
Forudsætningerne. Eller man maatte gjennem Sondringen mellem Tro og Mening tilveiebringe en
\emph{Fordobbling af Virkeligheden}, ved hvilken det forhindredes, at Tro og Viden, ved at
mødes indenfor \emph{samme} Virkeligheds Gebet, skulde komme i Strid, og ved at søge denne
Fordobbling vilde man da indrømme, at Troens og Videns modsatte Opfattelse af det \emph{samme}
Object dog maatte føre til \emph{Modsigelse}, skjøndt Tro og Viden skulle være afsluttede i sig selv.
Og en saadan Fordobbling søges nu ogsaa virkelig. Men den ved Virkelighedsfordobblingen
statuerede „høiere“ Virkelighed, „en Virkelighed, som Troen levende anerkjender, en
Virkelighed, hvorom Viden Intet veed“ --- skjøndt Grundideernes Logik dog idetmindste
deducerer den --- vakler imellem at betegne Syn, Symbol og en phantastisk Idealitets
Virkelighed, og bliver i denne usikkre Betydning en saadan, som hverken Viden \opage{192}eller Tro ville
anerkjende. Naar f.\ Ex. Fortællingen om Paradiset „med \emph{alle} deri indeholdte Charakteertræk“
siges at have væsentlig og \emph{sand} Virkelighed, saa er for Viden en „indre“ Virkelighed, en
„høiere Virkelighed“, i hvilken \emph{alle} i hiin Fortælling indeholdte Charakteertræk ere bevarede,
kun en Phantasievirkelighed, og Troen vil ikke nøies med en sublimeret Virkelighed, der mangler
den Facticitet, paa hvilken Troen som Tro maa lægge afgjørende Vægt. Troen fordrer, at det
Factiske congruerer med det Ideelle, med Troesidealet, at Virkeligheden i denne Forstand bliver
en høiere Virkelighed, men den protesterer paa det Bestemteste mod hiin Fordobbling af
Virkeligheden, der berøver Troens Gjenstand det Historiskes reale Betydning. Vilde man saaledes
lade Christi Fødsel, Død og Opstandelse foregaae ikke i den ydre historiske Virkelighed, i
hvilken Menneskelivet har sin Skueplads, men kun i hiin høiere Virkelighed, til hvilken de, som
gaaende ind under Underets Begreb, efter Theorien maae høre hen, saa maatte Troen heri see en
Tilintetgjørelse af hele Christendommens Grundlag. Og Videnskaben vilde kun kunne give den Ret
deri, ligesom den paa egne Vegne overhovedet, medens den protesterer imod at unddrages den
sande Idealitetens Virkelighed, ligesaa bestemt maa protestere mod en Fordobbling, der
forfalsker Idealitetens Virkelighed ved at sætte en Bastardvirkelighed, født af Phantasien,
istedet, og som ovenikjøbet maa blive vilkaarlig og inconseqvent. Ethvert Under maatte faae
hiin høiere Virkelighed til sit Element, men af de Undere, der have et Ydre, tør man kun føre
dem, der kunne reduceres til Visioner, til „Udslag“ af den Troendes indre Forhold til det
Guddommelige, tilbage til den, idetmindste tør man kun gjøre det uforbeholdent med Hensyn til
dem. Saaledes kan f.\ Ex. Miraklet med den brændende Tornebusk blive en indre Virkelighed, men
nu saadanne Mirakler som f.\ Ex. Overgangen over det røde Hav, eller Vandets Fremspringen af
Klippen ved Moses's Slag? Disse tør man ikke berøve den ydre Realitets Virkelighed, og hvorfor
ikke? Mon af en Troesgrund? Nei, af et inconsequent, fra Troens Syns\opage{193}punkt vilkaarligt, rationelt
Hensyn til de historiske Kjendsgjerninger, der jo tale for, at Jøderne i den ydre Virkeligheds
Forstand kom over det røde Hav, og at de tørstende Vandrere i Ørkenen ikke maatte nøies med et
„indre“, „subjectivt“ Vand, eller et „tilkommende“ Vand, men reddedes fra at omkomme ved et i
den ydre Virkeligheds Forstand virkeligt Vand\footnote{Jfr. 4. Moseb. XX, 11: og de drak
deraf, \emph{og deres} Qvæg.}, som Troen endnu søger i en virkeligt existerende Kilde.

De Forsøg, Nielsen har gjort paa at undgaae den Conseqvents, at Troens og Videns Gyldighed i
deres Modsætning maatte bringe en blivende, uforsonlig Splid ind i den menneskelige
Bevidsthed, vise sig saaledes som mislykkede. Og hvad der gjælder om dem, maa gjælde om
\emph{ethvert} andet Forsøg, der maatte gjøres paa at afværge Consequentsen. Det er en Udvei,
der maa forsøges naar man vil undgaae et forstyrrende Sammenstød mellem det Modsatte, der i
Modsætningen beholder Gyldighed: at postulere principielle Modsætninger, der i \emph{absolut} Uensartethed hver udfolder sig i en i sig afsluttet Heelhed\footnote{At Troens og Videns
Forhold hos Nielsen ikke i Virkeligheden svarer til denne Fordring, er paaviist i det
Foregaaende.}; men en absolut Uensartethed mellem Aandens Grundfunctioner har viist sig, ikke
at kunne begrundes ved deres factiske Eiendommelighed, og hvis den skulde finde Sted, saa
maatte, hvad der er en Umulighed, Bevidsthedens Enhed, der udtrykker Selvets Enhed, forvandles
til to med hinanden absolut uensartede og dog det samme Selv tilhørende Bevidstheder, og det
ligesaa Umulige maatte forudsættes: Tilværelseshelhedens Delthed i \emph{to} absolut uensartede
Tilværelses\emph{helheder}.

\emph{Den anden} psychologiske Consequents af Theorien, der ovenfor blev fremhævet, var den, at de to
Aandsfunctioner: Viden og Villie, der kun i deres Enhed og Harmonie vare sande, i den abstracte
Modsætning forvanskedes. Dette viser sig med Hensyn til \emph{Viden} i den Fordobbling af dens Begreb,
der ovenfor er paaviist, og hvis Virkning \opage{194}vil blive udviklet i det nærmest Efterfølgende. Med
Hensyn til \emph{Villien} bliver Consequentsen, at den reduceres til \emph{blind, egoistisk Drift},
fordi Villien i dens \emph{Modsætning} til Erkjendelsen, det Almeenfornuftige, kun er det \emph{Villiens Naturgrundlag}, der træder frem i \emph{Driften}, i det \emph{Alogiske} i Menneskenaturen. Den uigjennemsigtige
Naturenkelheds Bestemmelse, der er Driftens, overføres paa Villien overhovedet. I de
Villiesformer, der sættes som de høieste, er i Virkeligheden Villiens Fornuftbestemthed forsvunden
og det Alogiske blevet eneherskende. Dette faaer saa \emph{Frihedens} Navn, medens det i
Virkeligheden er \emph{den blinde Drifts Naturnødvendighed}. Om den sande Frihed, det universelle
Væsens \textit{libera necessitas}, eller rettere: \textit{necessaria libertas}, veed den
„paradoxe Subjectivitet“ Intet.

Den med Hensyn til \emph{Viden} og \emph{Villie} dragne Consequents kan ogsaa faae et andet Udtryk, der
allerede bestemtere berører det Ethiske og det Religiøse, nemlig, at \emph{Viden og Tro}, idet hver af
dem faaer en Modsætning, der i sin principielle Oprindelighed er absolut, har absolut
Selvgyldighed, maae forstyrres i deres Væsen. Det er den factisk stedfindende gjensidige udvortes
Begrændsning, som forstyrrer begge de Begrændsede i deres Væsen, der kun kunde være uforstyrret i
den frie, fuldstændige Udfoldelse. Hvad \emph{Viden} angaaer, saa maa den Indskrænkning ved Troen, der
finder Sted trods alle Forsøg paa at benægte den, angribe den i dens Rod og gjøre den usikker i
sig selv. Videns Væsen indeslutter i sig \emph{Universalitetens} Bestemmelse; uden denne
Bestemmelse er Viden ikke i Sandhed Viden\footnote{At dette gjælder fra et psychologisk
Synspunkt ligesaa fuldt som fra et metaphysisk, fremgaaer af det tidligere Udviklede angaaende
Videns Grændse.}. Den illusoriske Universalitet, der gives Viden ved at lade den „afspeile“ sin
absolute Modsætning, er utilstrækkelig; thi Viden kan ikke afslutte sig til en Heelhed med dette
dobbelte Indhold, den kan ikke udfolde sig i et System, i hvilket det Afspeilede faaer Plads.
Idet Viden, \opage{195}i sin Begrændsning ved Troen, skal anerkjende Umuligheden af at modbevise Miraklets
Mulighed, altsaa med Hensyn til denne Mulighed skal slaae sig til Ro i et \textit{non liquet},
saa er en saadan Videns Indskrænkning statueret, der ligesom unddrager Viden dens hele Fundament
og gjør den værgeløs imod en Skepsis. Og naar det tidligere paaviste Subordinationsforhold
statueres, og Troens Indhold sættes som den absolute Sandhed, saa bliver det af Videns Væsen og
Love Fremgaaende, der i Underordningen dog opfattes som Troesindholdets Modsætning, sat som et
Usandt.

Med Hensyn til \emph{Troen} faaer Begrændsningen den samme Følge som med Hensyn til Viden: Troen taber
ved Sammenstødet sin Hvile i sig selv, sin uforstyrrede Vished. Og Sammenstødet er uundgaaeligt;
thi medens det lader sig tænke, at Viden --- opfattet paa en anden Maade end i Nielsens Theorie
--- ifølge sit Væsen og sit Forhold til den menneskelige Natur, kan afslutte sig i sig selv til
en sand Heelhed, erkjendende den Troes psychologiske Oprindelse, der forholder sig til det
positivt Religiøse, saa har Troen ingen saadan Mulighed i sig. Den Troende, der søger at afslutte
sig i Troen, kan kun gjøre det ved at afvise Viden; men i al Afvisen af Viden, i al Hævden af
Troessphærens Afsluttethed, bruges der Vidensbestemmelser, der indeslutte en Vidensheelhed; den
Troende kan ikke komme til Hvile i et \textit{credo quia absurdum}, i Bevidstheden om det
Paradoxe, fordi Viden er menneskelig Væsensbestemmelse, og gjør sig gjældende i al Bevidsthed.

\subsection{Med Hensyn til det Ethiske}
\label{ssec:conseq-ethisk}

Den ensidige psychologiske Opfattelse af Villien, der stillede den i abstract Modsætning til
Erkjendelsen, medfører umiddelbart Consequentser for det \emph{Ethiske}, hvis sande Bestemmelse ikke
kan naaes ud fra det usande Villiesbegreb. Mangelen ved Opfattelsen af det Ethiske bliver en
\emph{dobbelt}: at det ikke bliver renset fra det Egoistiske, og at det ikke \opage{196}bliver et
Væsenet efter dets Fuldstændighed Udtrykkende, men kun et Fragmentarisk, og principielt
sættes som et saadant. Det Første ligger deri, at Villien, i dens Modsætning til Erkjendelsen,
kun bliver den blinde Drift, der, med hvor høie Navne man end udsmykker den, altid kun bliver
et Selvisk. Det Sidste ligger deri, at idet Villien, ifølge Modsætnings-Forholdet til
Erkjendelsen, kun kan udtrykke Væsenet partielt, saa kan det Ethiske ogsaa kun udtrykke et
Partielt og derved Usandt, ikke Realisationen af Væsenet efter dets Fuldstændighed. Og idet
Grundbestemmelsen af det Ethiske er ensidig og usand, saa maa det, der betegnes som høiere og
høiere Stadier i det Ethiskes Udvikling, fjerne sig mere og mere fra det i Sandhed Ethiske.
Idet „Villiesconcentrationerne“ skulle lade det Theoretiske mere og mere forsvinde, saa at
Villien paa sit høieste Stadium rent taber det af Sigte, saa taber den saakaldte ethiske
Handling enhver Norm; den bliver afhængig af en uklar Drift i Subjectet, der ingen Control kan
underkastes, ikke bindes ved noget Almeengyldigt. Paa det Antirationelles Trin bliver Villien,
der kun lyder den erkjendelsesløse Tilskyndelses Bud, og som opfatter det fra denne Udgaaende
som den guddommelige Villies Bud, fanatisk og hovmodig i sin Egoisme. Det Ethiske slaaer her,
hvor der brydes med Theorien, om i sin absolute Modsætning\footnote{Det nytter ikke at beraabe
sig paa, at Subjectet, idet det har brudt med alle objective Grunde, ikke skal have lukket sit
Øie for Indsigt, Erkjendelse, Reflexion; men at dets Erkjendelse skal være ethisk Erkjendelse,
Selverkjendelse, Livserkjendelse, Gudserkjendelse; thi Erkjendelsens Bestemmelse bliver her
kun en Tvetydighed: „den religiøse Bevidsthed danner sig Erkjendelser, men de ere egentlig kun
Postulater“. --- Naar der hos \emph{Kierkegaard} ved Valget og Troesforholdet lægges et Eftertryk paa
Villiesmomentet som det i Valget Afsluttende og Afgjørende i Forhold til Theoriens
uinteresserede Neutralitet og uendelige Bevægelse, saa oversees ikke den af Valget forudsatte
theoretiske Erkjendelse: --- „jeg maa existere forat kunne tænke, og jeg maa kunne tænke
(f.\ Ex. det Gode) forat existere deri,“ --- Erkjendelse og Villie stilles ikke i principiel,
metaphysisk begrundet, Modsætning til hinanden, og det theoretisk Erkjendtes Omsætten i
Villiesbestemmelse opfattes ikke som forandrende Indholdet.}; det Egoistiske og Absurde bliver
\opage{197}Handlingens Princip. \emph{Den antirationelle Ethik er det Ethiskes Forvanskning; al
sand Ethik maa være rationel Ethik, thi kun i Harmonie med Erkjendelsen har Villien Sandhed.}

Som den \emph{rationelle} Ethiks Opgave bestemme vi Virkeliggjørelsen af det udelte menneskelige
Væsen ved den frie Handlen. Som dens Grundlag sætte vi den Villie, der i Harmonie med
Erkjendelsen og i sin Udvikling og Fuldendelse bestemt ved Erkjendelsen virker som
Væsensvillen ɔ: som Villen af Væsenstotaliteten. Som det, hvortil det Ethiske med Nødvendighed
stræber hen, og hvori den først finder sin Fuldbyrdelse, sit Begrebs Udfyldelse, have vi
allerede tidligere bestemt det Religiøse, men det Religiøse efter dets evige og uendelige
Væsensbestemmelse. Den rationelle Ethik bliver ikke, ved at accentuere Erkjendelsens
Betydning, en Handlingens Naturlære, slipper ikke, fordi Friheden sees i Enhed med
Væsensnødvendigheden, den ethiske Fordrings Ubetingethed. Som \emph{ubetinget} sees den ethiske
Fordring allerede naar den sees som \emph{Væsens}fordring, som Udtryk for
\emph{Menneskevæsenets Lov}; thi Væsenets Fordring til Individet, i hvis Form det skal
virkeliggjøres, er ubetinget, fordi Individet kun er ved Væsenet; definitivt sees Fordringen
som ubetinget ved Forholdet til \emph{Væsenets guddommelige Grund}.

En \emph{rationel} Ethik anerkjendes ogsaa \emph{efter Navnet} i den Nielsenske Theorie, men de nærmere
Bestemmelser, der gives med Hensyn til den, føre, naar Uklarhederne og Tvetydighederne i
Udtrykkene fjernes, til at \emph{benægte} dens Mulighed paa dette Standpunkt. Den rationelle Ethik
skal „staae under Videnskabens Indflydelse“, medens den religiøse Ethik staaer under Troens.
Herved synes, da ogsaa den religiøse Ethik som Ethik maa have Videnskabens \emph{Form}, at maatte
forstaaes, at Videns \emph{Lov} bliver bestemmende med Hensyn til den rationelle Ethiks \emph{Gjenstand}.
Men i Virkeligheden forholder det sig ikke saaledes. Villien skal have sin Udviklingslov i sig
selv, og, ligesom den „aldeles ingen Indflydelse“ indrømmer Viden paa sig, saaledes skal \opage{198}den
kun afspeile sig som et absolut Uensartet i Viden. „Rationel“ viser altsaa ikke hen til en
begribende Erkjenden af den ethiske Villiesvirksomhed, begrundet i Villiens Overensstemmelse
med Viden og Bestemmelighed ved Viden, men kun til det Negative: at Villien endnu ikke er
bestemt som den paradoxe, syndige Villie, og Villiens Ideal endnu ikke som paradox
Subjectivitet. Men i \emph{enhver} Form af Villien, ogsaa i den, der fremtræder i den
rationelle Ethik, skal der være et \emph{antirationelt} Element: Bestemmelsen „rationel“ i
saadanne Sammensætninger som „rationel Selvbestemmelse“, „rationel Ethik“\footnote{Jfr. Gr. L.
I. XXVII: „personlige Sandheder“ indenfor det \emph{Rationelle}.}, faaer derfor ingen anden Betydning
end i Sammensætningen „rationelt Paradox“ d.\ v.\ s. den faaer \emph{ingen} Betydning. Naar
det Rationelle „principielt angiver Videns Grændse“, kan der heller ikke blive anden
Bestemmelse for det med Viden absolut Uensartede end det Antirationelle: \emph{alt} Ethisk maa, som
Udtryk for Villiesprincipet, blive et \emph{Antirationelt}\footnote{Paa den „rationelle Ethiks
Maximum“ skal Subjectiviteten „vende Theorien Ryggen“. \emph{Heegaard}: Indledn. til den rat. Ethik.
S. 444.}, og til det antirationelle Ethiske kan der ikke svare en rationel Ethik. Forskjellen
mellem et rationelt og et antirationelt Ethisk, mellem en rationel og antirationel
Ethik\footnote{Begreberne: et rationelt \emph{Ethisk} og en rationel \emph{Ethik} blive i Fremstillingerne
% PRINTED AS IS, p.198 [collation 2026-09-21]: turned sort: page prints "nundgaaelig" (first letter n, arched top, against the u that follows it); confirmed by a blind second reader.
af Nielsens Theorie paa en uklar Maade brugte iflæng.} svinder derfor nundgaaelig
bort\footnote{Dette viser sig ogsaa i Usikkerheden i Bestemmelsen af Gudsbegrebet i den
rationelle Ethik. Jfr. \emph{Heegaard}: Indledn. til den rat. Ethik. S. 442 ff.}. Og naar der
paastaaes: „den vide Kloft mellem Philosophie og Religion \emph{udfyldes} af Ethiken“, i hvilken
„Livstankerne \emph{mødes} med de videnskabelige Grundtanker“, saa er dette „mødes“ kun en tvetydig
Vaghed, hiint „udfyldes“ en Usandhed. Var den rationelle Ethik i Sandhed rationel, saa vilde
den, ifølge \opage{199}Bestemmelsen af Grundforholdet mellem Erkjendelse og Tro, falde over paa
Philosophiens Side, og den samme Kloft, der adskiller Viden fra Tro, vilde adskille den
rationelle Ethik fra den antirationelle religiøse Ethik. Men idet den rationelle Ethik kun er
rationel af Navn, bliver den, trods Navnet og trods det, at en Adskillelse sættes ved den uden
% PRINTED AS IS, p.199 [collation 2026-09-21]: turned sort n→u: page prints "Overlegeuhed" (seen at 3x zoom); confirmed by a blind second reader.
videre postulerede Overlegeuhed, der tillægges den religiøse Ethik i Forhold til den (Gr. L.
I. XXVII), draget over til den modsatte Side, og Kloften falder mellem den og Videnskaben.
„Ethikens Sandheder“ maae, efter Nielsens Opfattelse, blive „overvidenskabelige“ paa samme
Maade som Religionens; det er den selvsamme Modsigelse at tale om en ethisk Videnskab i
strængere Forstand som om en „Troesvidenskab“, en Theologie.

\subsection{Med Hensyn til det Religiøse}
\label{ssec:conseq-religioes}

Af den ensidige Opfattelse af Villien og af det exclusive Forhold, i hvilket den saaledes
bestemte Villie sættes til Troen, følger \emph{først} som Consequents \emph{det religiøse Forholds
Endeliggjørelse}. Det Religiøse kommer ikke i Forhold til hele Menneskevæsenet, men kun til en
enkelt, isoleret Side deraf. Idet Troens religiøse Forhold udelukkende henføres til Villien,
urgeres det vel, at alle Sjælekræfter her ere tjenende Villien, ligesom implicite virkende
gjennem dens Form. Men den abstracte Modsætning mellem Viden og Villie som absolut uensartede
modsiger dette postulerede Forhold\footnote{Bestemmelserne: praktisk Tænkning, praktisk
Erkjendelse, blive, i den Betydning, de tages af Nielsen, kun postulerede. Den absolute
Modsætning udelukker Tænkning og Erkjendelse overhovedet fra Villien fordi Grundformen i Villie
og Tænkning bliver absolut forskjellig.}, og gjør den hele menneskelige Naturs Tilfredsstillelse
i det religiøse Forhold umulig. Det religiøse Forhold bestemmes som heelt igjennem værende et
praktisk \opage{200}Forhold; men det Praktiske kommer ikke til at udtrykke den praktiske Aands Heelhed, men
den ensidige og usande Form deraf, der udelukker Erkjendelsen. Den postulerede Totalitet faaer
ved den \emph{absolute} Modsætning til den postulerede Erkjendelsens (den theoretiske
Virksomheds) Totalitet, i hvilken de samme Sjælekræfter skulle være virksomme under Videns
Suprematie, en nærmere Bestemmelse, der ophæver Totalitetsbestemmelsen.

En anden Consequents af den absolute Modsætning mellem Villien og Troen og Viden bliver den, \emph{at der, skjøndt der postuleres et Regulerende med Hensyn til hvad der hører Religionen
væsentlig til, dog ikke kan gives nogen Norm for dette eller noget Værn mod det rent Meningsløse
og Absurde}. Der vises i denne Henseende hen til en Troeskritik, der skal have den Opgave, ikke
ud fra en ubestemt Følelse, men fra et Princip, at sondre det \emph{ægte Religiøse} d.\ v.\ s.
det, der \emph{har alle til Aabenbaringens} Væsen \emph{svarende Kjendemærker}, fra det uægte. Denne Opgave
skal Troeskritiken løse ved at gjøre Skilsmisse mellem \emph{Phantasiens Ideer og den absolute
Villiesidee}, en Adskillelse, ved hvilken Modsætningen mellem Mythe og Aabenbaring skal
bestemmes. Det \emph{Regulerende} med Hensyn til det Religiøse skal altsaa søges i den \emph{absolute
Villiesidee}, Troens Ideal. Men den „\emph{absolute Villiesidee}“ bliver, \emph{naar Erkjendelse og
Villie sættes som absolute Modsætninger, ikke i Sandhed ethisk Idee, det Godes Idee, og kan ikke
for Bevidstheden udfolde sine Bestemmelser i Concretion, saaledes at der kan blive et
Regulerende}. Opfattes Villien paa denne Maade, bliver det, der, som sat ved den guddommelige
Villie, kaldes det Gode (»bonum est quia Deus vult«), det \emph{ubegribelige, ved ingen
Consequents beregnelige Udtryk for den absolute Vilkaarlighed}, og et Kriterium paa det, der i
Religionernes concrete Indhold er det ægte Religiøse, kan ikke angives, fordi „Vilkaarlighedens
Gud“ kan gjøre \emph{Alt} til sit Udtryk. \opage{201}Og denne „Vilkaarlighedens Gud“\footnote{Jfr. \emph{Nielsen}: Om theor. og praktisk Erkj. S. 74.}, til hvilken Villiesideen er forvandlet ved den abstracte
Modsætning mellem Viden og Villie, bliver lig med den \emph{almægtige, undergjørende} Gud i den religiøse \emph{Forestillings} Betydning. Troeskritiken, der faaer „den almægtige Villie til Princip og
Underet til absolut Forudsætning“, og ikke udskiller „det Overnaturlige, som den hellige Villie
til virkelig og virksom Fremstilling af den religiøse Idee har sammenføiet med det Naturlige“,
maa uden Forskjel anerkjende \emph{ethvert} Under; thi ethvert Under er \emph{som} Under denne
Almagts Aabenbarelse. Idet den guddommelige vilkaarlige Villie bestemmes som transscenderende
Lovene og sine reale Betingelser, aabnes der et Mulighedernes Panorama, der viser \emph{Alt} som
muligt, idet Alt kan blive et Udtryk for den uberegnelige Villie. Henvisningen til den absolute
„Villiesidee“, der kun ved Navnet staaer i Forhold til Idealet for den ethiske Handling, hjælper
altsaa ikke til at finde et Kriterium med Hensyn til det concrete Religiøse: denne bestemte Lære,
dette bestemte Bud, denne bestemte Historie, dette bestemte Mirakel\footnote{At lade vor
subjective Trang være Sandhedskriterium og Erkjendelsesgrund paa dette Gebet, vilde fra dette
Standpunkt ikke heller kunne gjennemføres, da det „Subjective“ er blevet berøvet det
Almeengyldiges Norm (jfr. „den gode Villie“ S. 137).}. Derfor maa Problemet om Videns Forhold til
Troen paa dette Standpunkt altid holdes i Almindelighedens Form; Hensynet til de bestemte
Religionsformer, de bestemte Dogmer, de bestemte Mirakler, maa saavidt muligt undgaaes, forat
ikke det Alogiske, og derigjennem den uundgaaelige Conflict med Viden, skal træde for tydeligt
frem. Man holder sig i Regelen til den Religion, der har det meest aandelige Præg: til
Christendommen, og taler om den som om den udtrykte Religionen overhovedet, og ved den taler man
atter i al Ubestemthed om det Concrete, netop fordi en Indgaaen paa dette vilde vise Conflictens
Uundgaaelighed.

\opage{202}Og ligesaalidt som der kan findes en Norm for Sondringen mellem det ægte og det uægte Religiøse,
ligesaalidt kan der bestemmes en \emph{Udviklingslov for det Religiøse}, en Maalestok for hvad
der er at sætte som de høiere og som de lavere Former af de historisk fremtraadte Religioner. Der
postuleres vel en Maalestok i den dybere og dybere Villiesconcentration; men idet
Grundopfattelsen af Villien lader det, der bestemmes som den dybeste Villiesconcentration, ligge
fjernest fra den fornuftige Erkjendelse, fra Bevidstheden om Fornuftens almene Lov, saa maatte
consequent det, der skulde sættes som det Høieste, blive det Høieste i Retning af det
Superstitiøse og Absurde, og Rangforholdet vilde ligefrem være sat paa Hovedet. Kun ud fra det
rationelt Ethiske, fra Idee-Erkjendelsen, og fra Opfattelsen af Loven for den usondrede
Menneskeaands historiske Udvikling kan i Virkeligheden en Norm og et Kriterium findes; men denne
Norm og dette Kriterium falder da udenfor det Religiøse, taget i den Forstand, i hvilken Nielsen
opfatter det\footnote{Med Hensyn til de \emph{hedenske} Religioner er der hos Nielsen en Uklarhed i
Opfattelsen. Saaledes naar der siges (Om theor. og pr. Erkj. S. 74): „det er i strængere
Forstand ikke Villiens, men Tankens Gud, der har givet sig tilkjende for Hedningene“. Herved er
nemlig, efter Forudsætningerne, egentlig Religionens Væsen ophævet, og dog kan det vel ikke
nægtes, at denne Tankens Gud virkelig er bleven \emph{troet} af Hedningene.}.

Saavel Regulativet for det i Sandhed Religiøse og Loven for Religionernes Udviklingsgang, som den
indre Sammenhæng i „Troesbegreberne“ skulde \emph{Religionsphilosophien}\footnote{Jfr. \emph{Nielsen}: Om den gode Villie, S. 134 ff.} paavise; men Aspecterne for Muligheden af en saadan Videnskab og
for Løsningen af dens Opgave have allerede viist sig saa mislige som vel tænkeligt.

Naar Religionsphilosophien skal bestemme det væsentlige Religiøse i dets indre Sammenhæng gjennem
en „religiøs Motivering“, saa bliver, efter det ovenfor Udviklede, denne Opgave umulig at løse,
fordi det, at Erkjen\opage{203}delsen er sat som absolut uensartet med Villien (Troen), indeslutter, at
Villiesprincipets (Troesprincipets) Aabenbaringsformer ikke kunne reguleres\footnote{At det kun
er en Tvetydighed, naar der siges (Om theor. og pr. Erkj. S. 75): „At Aabenbaringens Gud
fremstilles som Vilkaarlighedens Gud, vil jo ingenlunde sige, at Guds ubetingede Villie skulde
drive Spil enten med Nødvendigheden eller med Sandheden eller med Godheden: tværtimod! Al Grund,
al Sandhed og Godhed har netop (Rom. XI, 32--33) sit Midtpunkt, sit Udgangspunkt og sit
Hvilepunkt i den guddommelige Villie (\textit{bonum est quia Deus vult})“, --- behøver ikke nærmere at
udvikles.}. Beraaber man sig paa det enkelte Individs Erfaring af, at dets Villie er bragt i
indre Harmonie og Forsoning ved Forholdet til en Kreds af givne ethisk-religiøse Bestemmelser,
saa bliver denne Paaberaabelse unyttig; thi ethvert Individ kan gjøre sin modsatte individuelle
Erfaring gjældende med lige Ret, og skal denne Erfarings Gyldighed benægtes, saa maa der gaaes
tilbage til \emph{Almeenbestemmelser}, der kun kunne være givne ved den \emph{theoretiske} Erkjendelse af Villiens \emph{Væsen} og væsentlige Grundforhold; men derigjennem kommer man, da
Almeenbestemmelserne som theoretisk bestemte vise hen paa en anden Opfattelse af Villiens Væsen,
tilbage til Villiens og Erkjendelsens principielle Enhed i det menneskelige Væsens Enhed, og
forlader altsaa det Standpunkt, den Nielsenske Theorie indtager.

Faaer Religionsphilosophiens Opgave: at udvikle det væsentlig Religiøse i indre Sammenhæng, den \emph{bestemtere} Form, at det skal udvikles \emph{systematisk paa Basis af Vidensbegreberne}, saaledes at under den \emph{absolute Uensartethed} en \emph{alsidig articuleret Vexelbestemmelse} mellem Troesbegreberne og Vidensbegreberne paavises, saa vilde Løsningen af denne Opgave, selv om det væsentlige
Religiøse \emph{var} sikret som saadant, ligeledes strande paa Bestemmelsen af den \emph{absolute
Uensartethed}. Der tilsigtes --- for, tilbagevirkende, at hævde Logikens Theisme --- en
Parallelisme mellem Vidensbegreberne og „Troesbegreberne“ og en „systematisk“ Udvikling af disse
sidste, der \opage{204}skal svare til Vidensbegrebernes Systematik. Men \emph{Parallelismen} maa blive \emph{umulig}, paa Grund af \emph{Grundformens Differents}, der end ikke lader en „\emph{Afspeiling}“ være tænkelig,
og hvad der tales om Begrebernes „\emph{formelle Ensartethed}“, er kun et tomt Ord. De Exempler, der
gives paa Begreber, der \emph{blive} absolut uensartede i Videns og i Troens Sphære, ere valgte
saaledes, at Betydningen kan blive elastisk; andre Exempler vilde strax røbe det Tomme eller
Uigjennemførlige i Paastanden om Enheden.

Hvad der her er sagt angaaende Religionsphilosophien bliver nu nærmere belyst idet vi undersøge
Maaden, paa hvilken dens \emph{Opgave} begrundes, og Midlerne, der antydes til Opgavens Løsning.

Forat bestemme \emph{Religionsphilosophiens Væsen} og \emph{Opgave} gaaes der ud fra Begrebet om en
\emph{absolut Grændse for vor Viden}, hvilken Grændse nu søges begrundet ved Forskjellen mellem
vor menneskelige Viden og den guddommelige Viden, \emph{vor} Videns Ideal. Med den absolute
Vidensgrændse skal saa \emph{Mysteriet} være sat og med det et med \emph{Vidensprincipet absolut uensartet Troesprincip}, og naar Viden og Tro udgjøre to absolut uensartede Regioner i vort høiere
Bevidsthedsliv, saa skal den ene og samme Gud vise sig for os med en anden Skikkelse i Videns end
i Troens Lys.

I denne Udledning af Consequentser fra Begrebet om vor Videns Begrændsning i Sammenligning med
den guddommelige er Forholdet mellem den abstracte og den concrete Erkjendelse af det
Guddommelige gjort til et Forhold mellem to, i deres Dobbelthed ikke deducerede, \emph{absolut
uensartede}, Erkjendelser, Videns og Troens, og det bliver, idet Tro og Villie bringes sammen,
uden videre forudsat, at et andet menneskeligt Organ skal kunne gribe hvad der ligger udenfor
Erkjendelsens Grændse, uden at en saadan Forudsætning staaer i indre Forbindelse eller i Harmonie
med Bestemmelsen af vor Videns Begrændsning.

Der fortsættes saa: „naar Gud som Gud kan erkjendes baade i Viden og i Troen, og naar Guds Væsen
\opage{205}aabenbarer sig i to Skikkelser: Nødvendighedens (Videns) og Frihedens (Troens), saa maa der
kunne gives \emph{to Slags Theisme, en videnskabelig og en religiøs}, en Videns og en Troens
Theisme, der indbyrdes forholde sig som de Erkjendelsessphærer, de tilhøre, altsaa i Principet
ere \emph{absolut uensartede}“. Men denne absolute Uensartethed skal ingenlunde udelukke
Muligheden af, at Theismen i den ene Skikkelse kan have Charakteermærker, der \emph{nøiagtig svare} til
Charakteermærkerne i den anden. Tværtimod skal det ligge i Sagens Natur, at Charakteermærkerne i
Videnstheismen maae \emph{speile} Characteermærkerne i Troestheismen som Nødvendighedsbestemmelser
Frihedsbestemmelser, som Vidensbegreber Troesbegreber.

Idet der her tales om en „Speiling“, bliver dette Udtryk efter Sammenhængen at forstaae som
enstydigt med \emph{Parallelisme}, og til denne Betydning af Ordet knytte der sig saa videre
Consequentser.

Af de givne Bestemmelser skal nemlig følge, „at en christelig Verdensbetragtning, ikke
umiddelbart, men middelbart, skal kunne faae Indflydelse paa Philosophien som autonomisk
Videnskab, at et \emph{Henblik til Troestheismen} kan være \emph{veiledende} ved en stræng \emph{videnskabelig} Udvikling af \emph{Videnstheismen}; at f.\ Ex. den i Grundideernes Logik udprægede Theisme kan,
uden i noget Punkt at overskride Grændsen for det Logisk-Physiske, være bestemt ved og
bestemmende for en tilsvarende Religionsphilosophie, der altsaa faaer Logiken til Grundlag.“

Umuligheden af det her postulatsviis Fremsatte fremgaaer af det ovenfor Udviklede. Vagheden og
Tvetydigheden i de brugte Bestemmelser er ved sig selv iøinefaldende.

Religionsphilosophien bestemmes saa som „\emph{Grændsevidenskab}“, idet den som Videnskab hører under
Videns Sphære, men, som havende Religionen til Indhold \emph{maa} orientere os i Troens, det positivt
Religiøses Sphære. „Den \emph{philosopherende Bevidsthed} er, hvad der forstaaer sig af sig selv, det \emph{subjective Medium}, hvori de \emph{uensartede} Begreber mødes; heraf følger, \emph{at Be\opage{206}vidsthedsenheden maa forudsætte en formel til Grændsebestemmelsen svarende Begrebsenhed}, at, med andre Ord,
de \emph{absolut uensartede} Begreber maae i Grændsen selv blive \emph{formelt ensartede}. Hvad en
saadan formel Ensartethed har at betyde, kan let paavises. Vi tale om \emph{absolut uensartede
Principer}, men \emph{Princip} er jo \emph{\textit{in abstracto}} altid \emph{Princip}. Vi tale om \emph{uensartede} (hvorfor ikke: \emph{absolut} uensartede?) Erkjendelser, uensartede Begreber o.\ s.\ v., men idet vi abstrahere
fra det Eiendommelige, bliver Erkjendelse = Erkjendelse, Begreb = Begreb“; --- det vil sige: idet
vi abstrahere fra det for de absolut Uensartede Eiendommelige, saa faae vi et \emph{absolut
indholdsløst Navn} tilbage som deres Fælledsbestemmelse, altsaa i Virkeligheden et \emph{Intet}, og paa
Basis af dette Intet skal saa „Religionsphilosophien“, den „omvendte Videnskab“, operere.

Som Grændsevidenskab skal nu Religionsphilosophien faae sine \emph{Troesbegreber} fra \emph{Troeslivet} og sine \emph{Vidensbegreber} fra \emph{Videnskaben}. I sidste Henseende skal det vise sig hvad „en
logisk Theisme“ har at betyde; i første Henseende skal det komme an paa religiøs Psychologie, paa
en Bedømmen af Erfaringernes psychologiske Værdi i Forhold til Troesprincipet. --- Den
„tilsyneladende Modsigelse“, at Religionsphilosophien paa de opstillede Betingelser maa blive en
Videnskab, der \emph{handler} baade om det, der falder udenfor, og det, der falder indenfor Videnskaben,
skal være let at hæve. „Vil Nogen paastaae, at der umulig kan gives en \emph{Viden} om en Troeserkjendelse, naar Troeserkjendelsen efter sin Natur skal falde \emph{udenfor} Viden, lad ham
da begynde med Begyndelsen og paastaae Umuligheden af, at der kan gives en \emph{Bevidsthed} (!) om
Jorden naar den virkelige Jord efter sin Natur skal falde \emph{udenfor} Bevidstheden\footnote{Med
Hensyn til denne lidet behændige Sophisme henviser jeg til Dr. \emph{Heegaards} klare og taalmodige
Analyse: Prof. Nielsens Lære om Tro og Viden. S. 36 ff.}. En Indvending som den, at
\opage{207}Troesbegreberne, ved at reflectere sig i Viden, absolut maae blive Vidensbegreber, har Intet at
betyde; Videnskabens Opgave er netop at vide det Uensartede som uensartet“.

„Som Grændsevidenskab maa Religionsphilosophien vise os, \emph{hvorledes eet og samme Begreb kan blive uensartet} (absolut uensartet?) \emph{med sig selv, eftersom det udvikles til Videns- eller Troesbegreb}. Til Troesbegrebet \emph{Skabelse} svarer f.\ Ex. et Vidensbegreb (?) Skabelse, hint
udviklet af Friheden, dette af Nødvendigheden; til Troesbegrebet \emph{Gudmenneske} svarer
Vidensbegrebet (?) Gudmenneske o.\ s.\ v. Hvor forskjelligartede disse Begreber ere, kan sees af
Strauß, Feuerbach og Kierkegaard“. (Hvad vil \emph{Videns}begrebet: Gudmenneske --- ikke Gudmenneskehed
--- sige? og hvad skulde et \emph{Videns}begreb: Skabelse, eller, hvad vi kunne substituere: et
\emph{Videns}begreb: \emph{Miraklet}, betyde? Det er tomme Tvetydigheder eller ligefremme Absurditeter.)
--- Religionsphilosophien skal da paavise en „\emph{alsidig articuleret Vexelbestemmelse}“ \emph{mellem de uensartede} (absolut uensartede?) \emph{Begreber}; hvilket skal skee saaledes, at \emph{Vidensbegreberne} tjene som \emph{intellectuelt Underlag}, \emph{Troesbegreberne} derimod træde i Forgrunden og \emph{udvikles systematisk}. Af Vexelforholdet mellem Begreberne fra de to Sphærer skal saa
Vexelforholdet fremlyse mellem de to Arter \emph{Theisme}, mellem den Theisme, hvis Gudsbegreb kan
udvikles med logisk Nødvendighed i en Grundideernes Logik, og den Theisme, hvis aabenbarede Gud er
Gjenstand for Tro.

Idet disse Fordringer stilles, bliver det aldeles uklart, i \emph{hvilket Forhold Troens „Begrebssystem“} staaer til \emph{Videns Begrebssystem}. Den principielle Bestemmelse af Troesbegrebene maa, ligesom den
udelukker den „alsidig articulerede Vexelsbestemmelse“ mellem de uensartede Begreber, udelukke en
Optagelse af Troesbegreberne i den immanente Udvikling af de til intellectuelt Underlag tjenende
Vidensbegreber, og overhovedet en \emph{systematisk} Udvikling ved Viden, og naar denne Opgaves
Løsning betinger „Religionsphilosophiens“ Existents, bliver denne \opage{208}umulig. Religionsphilosophiens
Opgave er af Nielsen stillet saaledes, at Løsningen af den bliver \emph{selvmodsigende}, idet den
maatte ophæve Forudsætningen af den absolute Uensartethed, paa hvilken den hele Theorie og selve
denne Religionsphilosophie hviler. Uklarheden i hint Forhold bliver iøinefaldende ved
Bestemmelsen af Forholdet mellem det ethiske og det logiske \emph{Guds}begreb. „En Udvikling af det
logiske Gudsbegreb“, siges der, „er saa langt fra at beroe paa nogen Fornægtelse af det ethiske,
at den tværtimod indeholder \emph{alle} for det ethiske Gudsbegreb nødvendige Momenter. Vel er Logik
ikke Ethik og de logiske Begreber ikke ethiske; men derfor maa Ethiken ligefuldt have sin Logik
og de ethiske Begreber \emph{reflectere sig i} den logiske Idee“.\footnote{Nielsen: paa anf. St. S. 118.} Det Spørgsmaal fremtrænger sig her: Menes der, at denne Reflecteren af Begreberne om det
med Videns Indhold absolut Uensartede skal give hine Begreber en Plads i det logiske System,
eller menes der kun, at der i Logiken skal findes logiske Begreber, der i deres absolute
Uensartethed med de udenfor Logikens Sphære faldende ethiske „svare til“ disse? Det Første viser sig \emph{umiddelbart} som \emph{umuligt}, og dets Umulighed er maaskee ogsaa anerkjendt naar det siges:
„Logiken kan, som Logik, kun fremhæve Selvbestemmelsens logiske, ikke dens ethiske Momenter“. Det
andet er ovenfor blevet paaviist som en tom Forestilling. Med Hensyn til Forholdet mellem de to
Arter Begreber bliver det Dilemma afgjørende, der maa stilles den Nielsenske Theorie: vides
Nødvendigheden af det Absolutes Spaltning i Troens og Videns Absolute i Kraft af en logisk
Udvikling, eller kun i Kraft af Troesindholdets factiske Væren for den Bevidsthed, der ogsaa
rummer Videns Indhold? I første Tilfælde bliver Modsætningen logisk deduceret, der naaes i \emph{Deductionen} fra den ene af Modsætningerne til den anden, og Modsætningens Absoluthed er derved
hævet; i sidste Tilfælde kan Dobbeltheden ikke deduceres i Logiken, og den ene af Modsætningerne
\opage{209}ikke faae nogen Plads der eller i nogen Forstand reflectere sig i et Logisk.

Den Opgave, der stilles \emph{Religionsphilosophien}, viser sig da som \emph{uløselig}, og den \emph{videnskabelige} Charakteer, der skal gives den udenfor det Videnskabelige „i videre Forstand“ kan
kun hævdes den ved et Behændighedskunststykke. Forat faae Troesbegreberne ind i en videnskabelig
Forms systematiske Udvikling, maae de først paa skjult Maade være bragte ind i Vidensbegrebernes
Kreds, for, efter der ligesom at have modtaget en logisk Daab, atter at tages ud med en
videnskabelig Farvning. Ligesom det lykkedes Guldmagerne, af den „philosophiske Jord“, rørt
sammen med „det philosophiske Vand“, at uddrage netop saa meget Guld som de iforveien havde
blandet deri\footnote{Et Udsagn af Cornelius Agrippa.}, saaledes vil det lykkes den Nielsenske
Religionsphilosophie at skaffe Troesbegreberne netop det Qvantum videnskabelig systematisk Form,
som de „troestheistiske“ Elementer, der ere indbragte i Logiken, have faaet hængende ved sig fra
Blandingen, --- men ikke heller mere, og der vil kun kunne skaffes dem denne Form derved, at det
Videnskabelige forvanskes og Vidensbegrebernes System forstyrres.

Den Forsoning, der skulde naaes ved Religionsphilosophien, derved at den, netop idet den paaviste
Begrebernes nødvendige Sondring i de absolut uensartede Troes- og Vidensbegreber, bragte
\emph{Overensstemmelse} tilveie mellem disse uensartede Begreber, maa efter det tidligere
Udviklede blive uopnaaet. Overensstemmelsen skulde tilveiebringes ved den fra Theologiens helt
forskjellige Forklaring af de fra Troesbegreberne charakteristiske Modsigelser, hvilken
Religionsphilosophien som Grændsevidenskab skal have. „At Troesbegreberne, naar de udvikles af
deres eget Princip, indeholde Modsigelser, kan ikke, som fra et scholastisk-orthodox Standpunkt,
forklares deraf, at den syndige Fornufts oprindelige og nødvendige Kategorier skulde staae i
Strid med den aabenbarede Sandhed; heller ikke, som fra et kri\opage{210}tisk-videnskabeligt Standpunkt
deraf, at Troesbegreberne selv skulde være usande; nei! Modsigelsen \emph{forklares} (!?) i Consequents
af Principet, udtrykkeligt deraf, at \emph{de manglende Mellembestemmelser ligge skjulte} i det \emph{absolute Mysterium}. Da her altsaa ligesaalidt kan være Tale om at presse Troesbegreberne ind
under en videnskabelig som Vidensbegreberne ind under en religiøs Maalestok, saa vil det heraf
kunne skjønnes, at der i Religionsphilosophien end ikke vil kunne være Tanke om et Concordat
mellem Tro og Viden“: et saadant er ikke blot overflødigt, men umuligt. --- Men \emph{Nødvendigheden} af
Bevidsthedsindholdets Spaltning i absolut uensartede Modsætninger bringer ikke, selv om den var
factisk, en Harmonie tilveie, og den uklare Forestilling om et absolut Mysterium, der som et
uudtømmeligt Forraadskammer gjemmer modsigende „Mellembestemmelser“, er allerede bleven
tilstrækkelig belyst og viist i sin Tomhed. Dens Værd er netop den samme som den
Religionsphilosophies, der skal benytte den.

\begin{center}---\end{center}

Vi kunne nu uddrage det \emph{positive Resultat} af Kritiken af den Nielsenske Theorie. Det kan
sammenfattes i følgende Bestemmelser:

\section{Resultat af Kritiken af Nielsens Theorie}
\label{sec:nielsen-resultat}

1.~Med Hensyn til det \emph{Psychologiske}: Det menneskelige Væsens i Selvbevidstheden
aabenbare Enhed udelukker enhver dualistisk Sondring i Bevidstheden, enhver absolut
Uensartethed mellem Væsenets Yttringer, de eiendommelige Aandsfunctioner. Under Functionernes
Eiendommelighed maa deres væsensbegrundede Enhed, deres Harmonie, deres Gjennemsigtighed for
Erkjendelsen bevares, og Functionernes Enhed maa, ifølge den psychologiske Erfaring og ifølge
Analogien med Metaphysikens Grundforhold, blive en Enhed under Erkjendelsens Principat. --- I
den menneskelige Bevidsthed kan da Intet constituere sig i den absolute Uensartetheds Sondring
fra Erkjendelsen, og Intet kan med Sandhed hævde sig i Modsætning til denne, der \opage{211}som Væsenets
Udtryk staaer i Harmonie med Væsenets andre Functioner i deres normale Udfoldelse. Naar derfor
Troen, opfattet som religiøs Tro, som Aabenbaringstro, sættes som forholdende sig til det
Antirationelle, og skal faae Plads uden Modsigelse i samme Bevidsthed som den over sig klare
Viden, bliver dette en Umulighed: det forstyrrende Sammenstød er uundgaaeligt, og naar under
Videns og Troens Modsætning Viden viser sig at kunne afslutte sig til Heelhed og blive et
Udtryk for Aanden efter dens Totalitet, medens Troen, der ikke kan frigjøre sig fra
Vidensbestemmelserne, ikke kan det, saa maa Dommen over Troen i sidste Instants tilkomme Viden.

2.~Med Hensyn til \emph{det Ethiske}: Ingen sand Ethik kan gives uden som rationel Ethik. Naar
det Ethiskes Sphære omfatter den totale menneskelige Væsensbestemmelses Virkeliggjørelse ved
den frie menneskelige, ved det Godes Idee bestemte, Villie; naar det sande Ethiske fordrer, at
alle Menneskeaandens væsentlige Yttringsformer ere i Harmonie, at Menneskeaanden er tilstede
efter sin ubrudte Heelhed i hver af dem; og naar i Menneskevæsenet Erkjendelsen har det ideelle
Principat, --- saa bliver den sande og concrete Erkjendelse Betingelsen for den sande ethiske
Handlen: for den, ved hvilken Forholdet til det Godes Idee sættes i det Valg, i hvilket den
ethiske Villie constituerer sig som saadan, og for den paa Valget hvilende enkelte ethiske
Handling, i hvilken den sædelige Villie giver sig Concretion. Livsheelhedens Ethiceren
forudsætter ligesaafuldt som den første Griben af Ideen i dens Sandhed den luttrede og sande
Erkjendelse af det Forpligtende, og er kun mulig ved denne.

3.~Med Hensyn til \emph{det Religiøse}: Det sande religiøse Forhold: det absolute Forhold til
det Absolute, maa være et saadant, der kan omfatte hele Aandens Liv og give Aanden som Heelhed
Forsoning og Uendeliggjørelse. Men dette bliver kun muligt hvor Spaltningen er hævet i det
menneskelige Væsen og den menneskelige Bevidsthed; hvor Enheden er tilstede som
Aandsfunctionernes bevidste Harmonie og Gjennemsigtighed i Erkjendelsen. Den kriticerede
\opage{212}Theorie gjør Forholdet til det Absolute relativt og endeligt, gjør Forsoningen ufuldstændig og
aabner Døren for alt det Superstitiøse og Uethiske, og som saadant Irreligiøse.

\begin{center}---\end{center}

Efterat saaledes den Nielsenske Theorie er kritisk undersøgt og dens Consequentser
belyste\footnote{Efterat Manuscriptet til denne Afhandling var gjort færdigt til Trykken, ere
Prof. Nielsens Forelæsninger fra Christiania udkomne. De have imidlertid ikke gjort
Tilføininger eller Forandringer nødvendige.}, bliver det, der er fælleds for den og
Kierkegaards, og det, der er eiendommeligt for hver af dem, at sammenfatte, forat det
bestemtere kan vise sig, i hvilken Udstrækning Kritiken af den ene Theorie faaer Gyldighed for
begge.

\section{Det for Kierkegaard og Nielsen Fælles og det for hver af dem
  Ejendommelige}
\label{sec:faelles-ejendommelige}

1.~\emph{Det for begge Theorier Fælleds} er i det Foregaaende paaviist. Det var først og
fremmest Bestræbelsen for at hævde Tro og Viden i deres af hinanden uafhængige Gyldighed, og
at sikkre Troens Ret og Gyldighed gjennem Videns Selvbegrændsning. Men tillige viste der sig at
være en factisk Overensstemmelse deri, at Viden, hvis Betydning for Subjectet kun opfattes som
relativ, subordineres under Troen, der faaer absolut Betydning, at den videnskabelige Sandhed
underordnes den „overvidenskabelige“; deri, at Videns indre Uendelighed ogsaa hos Nielsen
factisk giver Plads for en ydre Indskrænkning, og endelig deri, at Aandens Forhold til det
paradox Religiøse absorberer dens Forhold til det rationelt Ethiske.

2.~\emph{Det for Nielsen Eiendommelige} i Udformningen af Theorien er Bestemmelsen af
Uensartetheden som \emph{absolut} og Forsøget paa en \emph{psychologisk og en metaphysisk
Begrundelse} af denne absolute Uensartethed. Med Hensyn til dette har nu Kritiken viist,
hvorledes den absolute Uensartethed indeslutter en Selvmodsigelse, og hvorledes den factisk
negeres ved de Bestemmelser, der gives med Hensyn til Forholdet mellem Modsætningerne, og
lige\opage{213}ledes, hvorledes saavel den psychologiske som den metaphysiske Begrundelse er uholdbar. Af
positiv Betydning for Theoriens Udvikling bliver altsaa dette for Nielsen Eiendommelige ikke:
han giver den ikke et fastere Grundlag eller en fuldkomnere Form, men omformer den kun saaledes,
at den bliver fuld af Inconsequentser og Modsigelser\footnote{Det er en Selvfølge, at der ved
Modsigelser her ikke tænkes paa saadanne, der ere Udtryk for Modsætningernes nødvendige
Fremtræden under den Begrebsudvikling, der fører hen mod Modsætningernes Forening, men om de
Modsigelser, der ikke kunne opløses i en høiere Enheds Harmonie.}. Disse ere paa forskjellige
Steder i det Foregaaende paaviste; vi minde her om de væsentligste:

Der statueres en absolut Uensartethed indenfor Menneskeaanden, skjøndt en saadan vilde ophæve
Menneskevæsenets, Selvbevidsthedens og Tilværelsens Enhed.

Der statueres en absolut Uensartethed mellem Viden og Villie, mellem Viden og Tro, og dog
indrømmes der, med Hensyn til disse, Paavirkninger, Begrændsninger, Blandinger,
Gjennemgangsmedier, Afspeilinger.

Der vises ved den oprindelige Bestemmelse af Forholdet mellem Villie og Viden hen til deres
Sidestillethed, og dog endes der med at sætte et Subordinationsforhold.

Viden skal, i sin principielle Eiendommelighed, have absolut Gyldighed, og dog skal den bøie sig
under Troen, idet den skal anerkjende Miraklets Mulighed.

Troen (i den positive Religions Betydning) skal have absolut Gyldighed, men dog bliver Troen
læmpet efter Viden ved den forsøgte Rationaliseren af „det, hvori den Troende seer Miraklet“.

Der sættes absolut Uensartethed mellem Troesbegreber og Vidensbegreber, og dog postuleres en
„alsidig articuleret“ Vexelbestemmelse og Muligheden af en systematisk Udvikling af
Troesbegreberne paa Basis af Vidensbegreberne.

Der tales, trods den absolute Sondring mellem Viden og Villie, om en „praktisk Tænkning“, en
„praktisk Erkjen\opage{214}delse“, der statueres altsaa en „tænkende Villie“, en „erkjendende Villie“, ---
et psychologisk „Amphibium“.

Villien faaer Friheden til sin Grundbestemmelse, men Villiens Enkelthedsform opfattes saaledes,
at den gaaer ind under den blinde Naturnødvendighed.

Villien skal repræsentere det meest Subjective i Mennesket, og dog føres den tilbage til
Magtens Idee, Objectivitetens Repræsentant.

Der statueres en rationel Ethik; men dens Mulighed ophæves ved Bestemmelsen af det Ethiskes og
af Villiens Væsen.

„Kløften“ mellem det absolut Uensartede, mellem Philosophie og Religion, skal \emph{udfyldes} af
Ethiken: Ethiken skal altsaa ved sin Existents modbevise Uensartethedens Absoluthed; og naar
denne Modsigelse skal undgaaes, skeer det kun ved den nye, at Ethiken saaledes identificeres med
det Religiøse, at den, skjøndt Videnskab, lader en uendelig Kløft aaben mellem sig og
Philosophien.

Disse Inconsequentser og Modsigelser vilde med Lethed kunne forøges yderligere; men det Anførte
maa være tilstrækkeligt.

Den Betydning, de faae for Theorien, er ikke blot eller nærmest en theoretisk, men tillige og
væsentlig en praktisk. Theorien, i den Skikkelse, den har faaet hos Nielsen, faaer gjennem dem
Præget af at være et blot Tanke-Experiment, der af en bevægelig Dialektik kan holdes i Svæven,
men som der aldrig kan have været gjort et Forsøg paa at udtrykke i Existents; thi isaafald
maatte Modsigelserne være blevne bevidste. At Theorien, trods det, at dens oprindelige Motiv kan
have været en Interesse for det positivt Religiøse, dog definitivt kun faaer Tanke-Experimentets
Charakteer, viser sig ogsaa derigjennem, at den, der fremstiller den, uagtet han i Theorien
tilkjender det paradox Ethiske, der sættes som det høieste Udtryk for Villien, et Suprematie,
der maatte føre til, at al Theoriseren lagdes tilside, dog vedbliver at være theoretiserende.
Paa det paradox Ethiskes Standpunkt har Individet „ingen \opage{215}Tid til at beskæftige sig med den
videnskabelige Theorie, thi Livet giver det nok at tænke paa“. Opgaven kan da kun blive den,
idet „Subjectiviteten i Kraft af en Villiesbeslutning vender Theorien Ryggen, at fordybe sig i
den Opgave, at blive en ethisk Charakteer“. Der bliver da ingen Plads mere for en Syslen med det
Theoretiske. Det kan ikke vælges som ethisk Livsopgave, fordi det, som absolut uensartet med det
Ethiske, ingen Betydning kan faae for dette, for den høieste, personlige Sandhed. Det kan ikke
dyrkes for Andres Skyld, for den Opklaring f.\ Ex., der ved Opstillingen af den omhandlede
Theorie kan ydes dem; thi „Livet gaaer sin Gang, ubekymret om Alverdens Theorier“, altsaa vel
ogsaa om den efter Prof. Nielsen opkaldte. Det eneste Consequente vilde være, at den, der
opstillede denne Theorie, saasnart den var opstillet, strax slap Theorien \emph{som} Theorie for
blot at existere, og at den saa optoges af en Anden for strax at passere videre ligesom i den
Selskabsleeg, der kaldes: „Lad Fuglen ikke døe“, indtil den --- rimeligviis snart --- døde. Naar
for den, der opstiller denne Theorie, Theoriens Opstilling ikke bliver hans Theoretiserens
absolute Grændse, saa er der derved givet en Maalestok for den „Alvor“, med hvilken han omfatter
Theorien. Vel vilde det: trods Theorien at vedblive som Theoretiker, i en vis Forstand stemme
med Læren om den absolute Forskjel mellem Theorie og Praxis, og i samme Forstand være consequent,
men idet Consequentsen kun er som Inconsequents, og Inconsequentsen er den eneste Consequents,
bliver den dialektiske Bevægelse saa livlig, at den i sin Sammenføining ikke meget faste Theorie
neppe længe vil kunne udholde den.

3.~\emph{Det for Kierkegaard Eiendommelige} fremtræder først i Theoriens hele Physiognomie og i
dens Forhold til Personligheden. Fri for hine Inconsequentser udfolder hos ham Theorien sig
klart og gjennemskueligt, og idet den uden Misvisning sigter hen mod et Maal, mod hvilket Aanden
drages med en uimodstaaelig Kraft, som i det gamle Sagn Skibet mod Magnetbjerget, faaer den sin
Stadfæstelse i Personlighedens Stræben, der praktisk ud\opage{216}trykker hvad Theorien udsiger, giver det
Virkelighed i et Liv, der „kun har villet Eet“.

Det Eiendommelige fremtræder dernæst, som allerede tidligere angivet, i \emph{Begrundelsen}, i
Paavisningen af det Grundforhold, ved hvilket Videns Selvbegrændsning skal blive en
Nødvendighed. I \emph{Existentsens} Grundforhold, i den \emph{Existerendes} Forhold til det \emph{Evige}, er
Nødvendigheden søgt for Videns Erkjendelse af sin Mangel paa Evne til at gribe dette Evige, det
eneste Visse, i Vishedens Form, i den Form, i hvilken det kun skal kunne være, gjennem den
objective Uvisheds Ophævelse, for Troen. Med Hensyn til dette Støttepunkt for Theorien maa der
derfor til hvad der, angaaende det for Kierkegaard og Nielsen Fælleds, allerede er givet i
Kritiken af den Nielsenske Theorie, føies en supplerende Kritik, der først skal undersøge
Kierkegaards Opfattelse af „det Eviges“ Begreb og derefter Consequentserne af det ved denne
Opfattelse Betingede: af Bestemmelsen af Videns Begrændsning ved Paradoxet, og af det Ethiske og
af dets Forhold til det Religiøse, særligt til det Christelige.

\section{Kritik af de for Kierkegaard ejendommelige Bestemmelser}
\label{sec:kierkegaard-kritik}

1)~\emph{Opfattelsen af det Eviges Begreb.} --- I Kierkegaards Opfattelse af dette Begreb røber
der sig en dobbelt Indvirkning: paa den ene Side en Indvirkning fra den nyere speculative
Idealisme, paa den anden Side en Indvirkning fra den græske Philosophie. Den første Indvirkning,
der gjør sig gjældende gjennem selve Polemiken mod Speculationens Forflygtigelse af Existentsen,
viser sig i hans Opfattelse af det sandseligt Reale, der kun sees som en uforklarlig, i sig
bedragerisk, Tilføining til det Evige. Naar saa, hvortil Speculationen selv giver Berettigelse,
Endelighedens Bestemmelse accentueres, bliver Consequentsen den, at det Evige, der dog altid skal
være tilstede i Existentsen forat give Continuitet, som det med hint Fremmede Sammenknyttede, for
den Existerende bestemmer sig som et Saadant, der ikke rolig kan besiddes af ham i Vishedens Form,
men kun i Lidenskabens Subjectivitetsforhold momentviis kan faae Vished og momentviis give
\opage{217}Continuerlighed\footnote{Jfr. \emph{Kierkegaard}: Afsluttende Efterskrift.}. Og idet det Evige i den
Existerende saaledes faaer momentan Existents og Tilbliven, maa det opfattes som det, der dog
ikke bestemmes for Tanken ved dette Forhold. Den saaledes opfattede Evighed bliver netop den
speculative Philosophies abstracte, realitetsløse Evighed; kun saaledes opfattet, at den først
for Menneskeaanden faaer Værens Nærværenhed naar Existentsen med dens Vorden er hævet. Men at det
Reale med dets Vorden skal forsvinde forat det Evige kan komme tilsyne, bestemmer dette Reale med
dets Væreform som et forstyrrende Accessorisk, der ikke skal være: hiin realitetsløse Evighed som
det egentlige Sande for de i Evighedens ene sande Existents hjemmehørende Aander.

Den græske Distinctionsdialektiks Opfattelse, der skarpt sondrer det evige Værende fra det
Ikke-Værende, træder frem gjennem Bestemmelsen af det Eviges Forhold til Vorden. Idet det Evige
berører det i Vorden givne μὴ ὄν, kan det ikke hævde sig for den existerende Bevidsthed i
Vishedens Form, men det positive Forhold til det bliver uadskilleligt fra et negativt, og
Bevidstheden skal, som værende i Vorden, ikke ved Tanken kunne gribe det Evige, der er Aandens
eget Væsen.

Begge de paaviste Indvirkninger føre altsaa til en abstract Grundopfattelse af det Evige, der
enten seer det under den evige Værens Bestemmelse eller idethøieste lader dets Bevægelse blive
en evig ideel Bevægelse, ikke den Bevægelse, gjennem hvilken oprindelig den reale Virkelighed i
dens Væsensgyldighed sættes og det Evige, i Idealiseringen af den, er i Concretion som
idealt-realt. Hiin abstracte Opfattelse af det Evige og Absolute medfører, at det Evige, skjøndt
dets Forhold til Vorden ikke kan fjernes, ikke kan faae Bestemthed gjennem dette Forhold, saa at
en uendelig integrabel Erkjendelse af det Absolute ikke bliver mulig, og, paa Grund af Forholdet
til den i Vorden indesluttede Negation, end ikke en Vished om det i dets Abstraction gjennem den
objective Viden. Medens i den concrete Opfattelse \opage{218}af det Evige og dets Bevægelse Negationen vilde
være optagen i og behersket af det Eviges Selvbekræftelse, berører, i den abstracte Opfattelse,
det Evige Negationen i det Reale som et μὴ ὄν, der gjør det Evige usynligt.
Ud fra hiin concrete Opfattelse viser sig Correctivet for Consequentsen af den abstracte. I
Opfattelsen af Individets Forhold til Tilværelsens ideelt-reale Heelhed og til Tilværelsens
ideelt-reale Princip som det enkelte organiske Leds til Organismens Heelhed og til Organismens
Princip ligger Muligheden til, at det Evige, det Absolute, kan være for den individuelle
Bevidsthed i Vishedens Form. Gjennem det i Bevidstheden Aprioriske, der er Menneskeaandens
Væsensudtryk, medens det udtrykker det Absolutes, det Aanden Begrundendes, Væsensbestemmelser,
bliver hiin Vished mulig, og idet det Aprioriske kan assimilere det Empiriske, og gjennem denne
Assimilation udfolde sig i Concretion, er den relative Abstraction, i hvilken det Evige er
tilstede for den Bevidsthed, der har objectiv Vished om det, en bestandig integrabel. Det Eviges
Væsen bliver ikke for vor Tanke blot i den bestemmelsesløse evige Værens Form, i hvilken det
staaer i uforenelig Modsætning til det Ikke-Værende; det udfolder sig for Bevidstheden, saaledes,
at det i hvert Moment er et, trods Abstractionen, i sig ideelt Afsluttet, der dog viser hen mod
den voxende Concretions Mulighed. Den Existerende, der bliver sig det Aprioriske bevidst som sin
Væsensbestemmelse, har gjennem det en Garantie for den objective Sandhed af sin Viden om det
Evige. Hvad han ikke kan undlade at tænke uden at hans Bevidsthed forstyrres, uden at den kommer
i Splid med sig selv, er for ham det objectiv Sande, og er dette i Kraft af den subjective,
individuelle Bevidstheds Fordring.

Her oversees altsaa ikke Forskjellen mellem Erkjendelsesforholdet i det Absolute selv og i den
enkelte menneskelige Bevidsthed; men Forskjellen sees som en saadan, der ikke ophæver det Eviges
Vished ogsaa for den endelige Erkjendelse. Dennes Viden om det Absolute faaer, idet den bliver en
Viden i den relative Abstractions bestandig integrable \opage{219}Form, det Aprioriskes Vished, og naar, paa
det Ethiskes Gebet, en concretere Bestemmelse fordres end Erkjendelsesforudsætningen giver; naar
der altsaa handles ud fra hvad der kunde kaldes en „ethisk Tro“, saa bliver denne Tro ikke en
Modsætning til hiin Viden, men den bliver i Harmonie med Erkjendelsen; den bliver en Stadfæstelse
af Erkjendelsens Principer, i hvilke Aanden gjenkjender sit eget Væsens Udtryk.

2)~\emph{Consequentserne af Videns Begrændsning ved Paradoxet.} --- Idet det Evige ikke skal
kunne være i Vishedens Form for den Existerendes Tanke, men kun gribes i Troen gjennem Uvisheden,
saa bliver Troens Subjectivitetsforhold Formen for Sandheden i Existentsen, og
Subjectivitetsforholdets høieste Form bliver det troende Forhold til Paradoxet: det sig som
syndigt anerkjendende Individs Tro paa Guden i Tiden. Men dette subjective Forhold maa, som et
Bevidsthedsforhold til et Aandeligt, for det efter sit Væsen tænkende Individs Bevidsthed faae et
Tankeudtryk, og idet dette kommer i Modsigelse med Tankens Love og Former, deels directe, deels
gjennem de Consequentser, der ligge i det, saa maa der i Bevidstheden opstaae en Strid mellem
Troens Tankebestemmelser og Videns Tankebestemmelser, og Conflicten, der maa komme frem saasnart
en Troessætning er udtalt, maa skærpes for hver Consequents, der udledes af denne. Men ved denne
Modsigelse maa Bevidstheden komme i Splid med sig selv. Troessætningen sættes som Udtryk for den
absolute Sandhed: Tankemodsigelsen i den maa da virke tilbage paa Bevidsthedens øvrige
Tankeindhold og tilintetgjøre dets Sandhed, skjøndt det, som Udtryk for Viden, er Udtryk for en
Væsensbestemmelse i Mennesket. Og naar dette er skeet, saa maa den Viden oprindelig indrømmede
relative Gyldighed ogsaa i den Henseende tages tilbage, at Forstandens Anvendelse til Værn for
Paradoxet, dens Distingueren mellem Paradoxets høiere Sandhed og det Meningsløse, bliver umulig.
Naar Paradoxet indeslutter det Utænkelige: Foreningen af Gud og det enkelte Menneske i Christi
Person, det Eviges Tilbliven i et Moment af Tiden; naar det, ved \opage{220}Bestemmelsen Synd, udtaler
Forvandlingen af Menneskenaturen, det i sin Selvbevidsthed dog Identiteten med sig bevarende
Individs paradoxe Forvandling, dets Taben af Evigheden, af det Evige, som Mennesket dog ikke kan
blive af med, og som i Syndsforladelsen, i Troen paa Christus, atter vindes, og vindes saaledes,
at det i samme Øieblik atter forudsætter sig selv som det Utabte, saa indeholde alle disse
Bestemmelser hvad der, omsat i Tankeudtryk, kommer i Strid med alle Tænkningens Love og Former,
og idet det i sin Modsætning til Vidensbestemmelserne udtrykker Sandheden, gjør det det, der
strider derimod, til det Usande.
\emph{Paradoxet er Underet i Tankens Verden}, og ligesom Underet i Naturens Verden ophæver
Naturens og Naturlovens Begreb, saaledes maa Underet i Tankens Verden omstyrte Tankens Love, og
lade al Videns Vished opløse sig i Tvivlen. Ligesom der fordres af det troende Individ, at det i
sit Liv skal afdøe fra alt det Naturlige og kun leve i Religionens abstracte Aandelighed,
saaledes maatte der ogsaa, consequent, fordres, at det i sit Tankeliv skal kaste al Videns
Virkelighedsindhold bort for blot at være med Syndsbevidsthedens Paradox i det paradoxe
Gudsforhold. Men denne Opgave for Aanden: at bortkaste hvad der hører dens Væsen til, indeslutter
en Selvmodsigelse og bringer en uforsonlig Splid ind i Bevidstheden.

3)~\emph{Consequentserne af Opfattelsen af det Evige med Hensyn til det Ethiske og særlig til
dets Forhold til det Religiøse.} --- Den abstracte Opfattelse af det Evige indeslutter den
Consequents, at det Ethiske maa blive bestemmelsesløst og derved uvirkeligt. Idet Kierkegaard
lader Individet, der ethisk skal vælge sig selv efter sin Uendelighed, i Valget drage sig ud af
al Endelighedsbestemthed, bliver det Opgaven for ham, at paavise Muligheden af, at den ethiske
Villie atter kan faae Indhold og Concretion; men den Dialektik, ved hvilken han vil naae dertil,
forudsætter en anden Grundopfattelse af det Evige og af Individets Forhold til det end hiin
abstracte. Er det Ethiske kun ved Forholdet til det Evige, og beholder \opage{221}dette Evige bestandig den
samme Abstractionens Form, saa maa det Ethiske blive et Uudfoldeligt; men som saadant bliver det
et Uvirkeligt; thi den ethiske Handlen fordrer Bestemthed, Concretion.

Det, at det Ethiske ikke kan faae Virkelighed, fremgaaer, paa middelbar Maade, af den ved
Opfattelsen af det Evige betingede Bestemmelse af det Religiøses Forhold til det Ethiske. Idet
Kierkegaard lader det ethiske Forhold ved sig selv vise ud over sig til det religiøse, bestemmer
han dette som den absolute Retning mod det absolute τέλος, saaledes at dette Forhold, netop fordi
det Evige, til hvilket det er et Forhold, aldrig kan udfolde sig for Viden, bliver et særligt
Forhold, der ikke udtømmer sig eller bliver concret i de relative ethiske Forhold, og, som den
absolute Pligt, kan fordre Forsagelsen af dem alle. Men denne Bestemmelse af Forholdet maa føre
til, at de ethiske Forhold, skjøndt de her supponeres som relativt gyldige i deres Bestemthed,
dog maae forsvinde for det religiøse absolute Forhold; thi idet det, at forholde sig absolut til
det absolute τέλος, er, \emph{altid} at forholde sig saaledes til det, saa maa dette absolute Forhold,
saaledes opfattet som Kierkegaard gjør det, optage den hele Menneskeaands Kraft, og ikke lade
nogen blive tilovers til de relative Forhold. Forholdt Mennesket sig til disse med en Kraft, en
Aandsvirksomhed, forskjellig fra den, med hvilken det forholder sig til det absolute, saa vilde
denne sidste blive begrændset og Forholdet relativt. Naar der opstilles den Fordring, at
Mennesket paa een Gang skal forholde sig absolut til det absolute τέλος og relativt til de
relative, saa bliver denne Fordring uopfyldelig. Den kunde kun opfyldes, naar den Kraft, med
hvilken Individet forholdt sig til de relative Formaal, tænktes som en Deel af den Kraft, med
hvilken det forholdt sig til det absolute. Men saa vilde det altsaa i Forholdet til hine forholde
sig til dette, og dette vilde sige: det absolute Formaal vilde blive concret i de relative. Og
til et saadant Forhold vises vi ogsaa hen ved Umuligheden af den Fordring, der stilles med Hensyn
til Forsagelsen. Mennesket \opage{222}skal, til Fordeel for det absolute τέλος, kunne forsage \emph{alle} de
relative. Men en saadan Forsagelse er det Utænkelige. Mennesket kan, i et Uendelighedens
Øieblik, idet det sammenfatter Existentsen i en Abstraction ligeoverfor det Eviges Abstraction,
forsage den hele „Endelighed“, alle de relative Formaal, ligesom Selvbevidstheden, idet den
slutter sig sammen i et rent Jeg = Jeg, kan løsne sig fra al Bestemthed. Men i den tidsbestemte
Virkelighed, i hvilken Individet existerer i en Concretion, forsager det kun den ene
„Endelighed“, det ene Forhold, ved den Kraft, ved hvilken det existerer i det andet, og selve den
Abstractionens Act, ved hvilken i hiint Øieblik den totale Forsagelse skete, er, ligesom
Selvbevidsthedens yderste Abstraction, betinget ved en Concretion, som den forudsætter og af
hvilken den bæres i samme Øieblik, den hæver sig derover.

Der naaes saaledes aldrig i Virkeligheden ud over det Ethiskes Totalitet, men saasnart det
Ethiske faaer Gyldighed, om end det enkelte ethiske Forholds Gyldighed bliver relativ, saa bliver
det Religiøse at opfatte som det, der ikke sondrer sig fra det Ethiske, men bliver Sjælen i det
Ethiske, hvilket sidste ikke kan gjennemtænkes uden hint, uden at blive til det Ethisk-Religiøse.
Dertil vilde man ogsaa komme gjennem den Kierkegaardske Bestemmelse, at det Ethiske kun bliver
til idet Uendelighedens Bevægelse gjøres. Det Evige, til hvilket Forholdet er i det Religiøse,
kommer altsaa ikke først frem for Bevidstheden paa det Religiøses Gebet; og den Abstraction, i
hvilken det forbliver paa dette (som det uendeligt Negative), gjør det ikke muligt at adskille
Religionens Evige fra Ethikens Evige, saa at der altsaa ingen Grund bliver til at fastholde hiin
Sondring.

Umuligheden af at hævde det Ethiskes Virkelighed, viser sig atter, naar Christendommens Fordring
stilles og naar det Christeliges afgjørende Bestemmelser komme frem, og disse Bestemmelser ophæve
ikke blot det Ethiskes Virkelighed, men ogsaa Virkeligheden af den Religiøsitetens Form, hvis
Enhed med det Ethiske røbede sig gjennem de ovenfor paaviste Antydninger.

\opage{223}Naar Christendommens Fordring til Mennesket, i Overensstemmelse med Opfattelsen af det abstracte
Modsætningsforhold mellem Christendommen og Verden, bestemmes som Fordringen til at afdøe, til at
hade sit naturlige Liv, saa indeslutter denne Fordring Umuligheden af at faae et ethisk Indhold
til at bestaae sammen med det Christelige. Den naturlige Tilværelses Indhold forsvinder som et
Intet i den absolute Modsætning til den evige Sandhed, og med den forsvinder Betingelsen for det
Ethiskes Concretion, og det vil sige: for dets Virkelighed. Denne Consequents er dragen i
Kierkegaards Polemik mod Ægteskabet (i „Øieblikket“) og det i Afsl. Efterskrift (S. 374) Antydede
kan ikke føre til et andet Resultat. Der maa blive den samme Umulighed i at lade det Ethiske
beholde sin Plads trods hiin Fordring, som at lade Forstanden beholde en Betydning trods
Paradoxet.

Naar den afgjørende christelige Bestemmelse \emph{Arvesynd} bringes frem, og naar den opfattes som
indesluttende Tabet af det evige Væsen og Uvidenheden om dette Tab, saa berøver den det Ethiske
Virkeligheden ved at berøve Mennesket det Ethiskes Betingelse: Muligheden af Forholdet til det
Evige. Naar ved Arvesynden det evige Væsen sættes som tabt og som dobbelt tabt ved Uvidenheden om
Tabet, saa sættes Muligheden til at vinde det tilbage som betinget ved Troen i Christendommens
Betydning; udenfor denne Tro bliver der intet Evighedsforhold. Naar altsaa det Ethiske ikke kunde
faae en Plads indenfor det Christelige, saa kan det ikke heller existere udenfor dette, fordi
dets Betingelse der mangler. Og hvad der gjælder om det Ethiske, kommer ogsaa til at gjælde om
den Religiøsitetens Form, som Kierkegaard i Afsl. Efterskrift betegner som Religiøsiteten \textit{A}, og som, uden at svare til en bestemt, historisk fremtraadt Form af Religionen, skal omfatte den
Religiøsitet, der endnu ikke er den christelige, og i sin idealiserede Form vel nærmest viser
tilbage til græske Motiver. Om denne Religiøsitetens Form, der som Religion maa indeslutte et
Forhold til det Evige, siges der, at den idetmindste \emph{kunde} \opage{224}findes i Hedenskabet, da den kun har
Menneskenaturen som saadan til Forudsætning; men fra det Synspunkt, der er givet ved den
christelige Bestemmelse af Arvesynden, viser hiin Religiøsitetens Form sig kun som en Mulighed i
den ufordærvede Menneskenatur, der aldrig kan have havt Virkelighed. Kun i Christendommen kan
Evighedsforholdet virkeliggjøres, men Christendommens ideale Fordring er løftet saa høit, at det
sorgfulde Spørgsmaal trænger sig frem: „mon Christendommen nogensinde er kommen ind i Verden?“.
Naar Theoriens Consequentser med Hensyn til det Ethiske blive disse, saa vil der deri for den,
for hvem det Ethiske har absolut Betydning, ligge en uoverstigelig Hindring for at anerkjende
Theorien som sand og som løsende Bevidsthedens Splid.

\section{Endeligt Resultat}
\label{sec:sondring-resultat}

Det endelige Resultat med Hensyn til den Kierkegaardske Theorie er anticiperet i de
foregaaende Udviklinger. Ogsaa den har viist sig, trods sine væsentlige Fortrin for den
Theorie, der er afledet af den, og trods den Berettigelse, den har ligeoverfor den ensidige
Speculation, ikke at kunne tilfredsstille. Forsøget paa en Principsondring, anerkjendt af den
sin egen Begrændsning erkjendende Viden, fører ogsaa i den, ikke til Bevidsthedens Forsoning,
men til dens uundgaaelige Adsplittelse, til en uundgaaelig Splid i Aanden, der ikke til
Fordeel for Troen kan fornægte Viden, der er dens egen Væsensbestemmelse, ikke til Fordeel
for det Religiøse opgive det Ethiske, der er Aandens sande Virkeliggjørelsesform, og derfor
for Aanden et ubetinget Gyldigt.


%% ============================================================
%% SLUTNING (pp. 225--226)
%% ============================================================
\chapter*{Slutning}
\addcontentsline{toc}{chapter}{Slutning}
\label{ch:slutning}

\opage{225}Saaledes har den kritiske Undersøgelse af den Theorie, der ved Sondringen mellem Principerne
vilde tilveiebringe Aandens Forsoning, og af Theologiens Forsøg paa, i en Troesvidenskab
organisk at forene Tro og Viden, ført til Stadfæstelsen af det Resultat, til hvilket de
foreløbige Begrebsbestemmelser af de positive Religioners Væsen førte, og som stemmede overens
med det, til hvilket den umiddelbare religiøse Bevidsthed naaede fra det modsatte Udgangspunkt.

Dette Resultat var, at der mellem Troen i den positive Religions Betydning ɔ: mellem Troen som
„Aabenbaringstro“, og Viden, ifølge deres forskjellige Indhold og forskjellige Kilde, bliver en
uophævelig Incongruents, en uforenelig Modsætning, og at det derfor bliver umuligt, at de uden
Modsigelse kunne forenes i samme Bevidsthed, bevare deres Gyldighed for den samme Personlighed.

Deri ligger saa som umiddelbar Consequents Nødvendigheden af et Valg mellem Modsætningerne, af
hvilke en maa være Udtrykket for Menneskeaandens høieste Energie og i sig indeslutte Maalet for
dens høieste Stræben. Og dette Valg maa, ifølge den i det menneskelige Væsens Enhed liggende
Fordring, være et saadant, i hvilket Aandens Grundvirksomheder ere i Harmonie i sig og med
hverandre, i hvilket altsaa Væsenet er i Enhed med sig i sin Erkjenden og Villen.

I denne Enhed med sig selv kan Aanden være idet \emph{Viden} vælges. Idet Viden vælges, vælges nemlig
det, der, som Væsensudtryk for det universelle Menneskevæsen, i sig bærer Muligheden til at
afslutte sig til en Totalitet med indre Uendelighed; det, der, idet det har det ideelle
Principat i den udelelige menneskelige Aand, kan blive Normen for en ethisk Livsheelhed; det,
der, medens det \opage{226}gjennemskuer de endelige Religionsformers Væsen, begriber dem i deres Udspring
af den sig udfoldende, hen mod sit Princip stræbende Menneskenatur, og derfor ikke i dem
beholder en forstyrrende Bevidsthedsskranke, og som, idet det i de historiske Religionsformer
udskiller det Forgængelige og Endelige fra det Blivende, Væsentlige og Evige, det Ikke-Religiøse
fra det Religiøse, kan staae i Harmonie med og i sig optage dette Væsentlige, det Religiøse efter
dets sande og evige Betydning.

Medens altsaa Videns Uforenelighed med den Tro, hvis Gjenstand er det positivt Religiøse,
sættes, sættes tillige ikke blot Muligheden, men Nødvendigheden af, at under Valget af Viden det
i Religionerne Væsentlige, „det Religiøse i Religionerne“, bevares. Modsætningsforholdet til de
positive Religioner i deres Endelighed forvandler sig saaledes til et Enhedsforhold til det
Religiøse efter dets Uendelighed.

Hvorledes Problemet om Troens og Videns Forhold fra dette Synspunkt bliver at løse, vil blive
udviklet i det Skrift, der skal slutte sig til dette. Idet det hævder Enhedsforholdet mellem den
sande Erkjendelse og den sande Handlen, den ethiske Handlen, og idet det paaviser, hvorledes det
i sand Betydning Religiøses væsentlige Kjendemærker ere tilstede i denne Erkjendelse og i denne
Handlen, hvorledes disse altsaa ere i Harmonie med hint, vil det forsøge, ud fra Philosophiens
Begreb af det concret bestemte Absolute, den høieste \emph{Videns}bestemmelse, at udvikle en
Livsanskuelse, der tilfredsstiller ethisk-religiøst, og at naae til de væsentlige Bestemmelser
af det ethiske og religiøse Forhold, der paastodes, kun at kunne udledes af et med
Vidensprincipet absolut uensartet Villies- og Troesprincip.

I dette Skrift vil altsaa det
negative Resultat af det foreliggende finde sit positive Supplement.

\end{document}
